% Auto-generated from data/segments.json + layout.json (+ transforms) — do not edit by hand.
% volume: 31Abhi03
% mode: printing
% segments: 1261
% transform_rules: 8
% toc_marks: 160

% Volume layout defaults (layout.json ``layout``).
\csromanlayoutapply{1}{1.5}{21.6pt}{8pt}{8pt}{65pt}{2.5em}%

\pitaka{อภิธมฺม\-ปิฏก}
\csromanmatikaanchor{mtk.0}
\csromantocmarknopage{chapter}{ธาตุกถาปาฬิ}{mtk.0}
\bookmark[dest={mtk.0},level=0]{ธาตุกถาปาฬิ}
\gambhira{\textbf{ธาตุกถาปาฬิ}}
\namakkaram{นโม ตสฺส ภควโต อรหโต สมฺมา\-สมฺพุทฺธสฺส}
\csromanmatikaanchor{mtk.1}
\csromantocmarknopage{section}{อุทฺเทส}{mtk.1}
\bookmark[dest={mtk.1},level=1]{อุทฺเทส}
\csromanheader{1}{\textbf{อุทฺเทส}}
\csromanmatikaanchor{mtk.2}
\csromantocmarkref{subsection}{1. นยมาติกา}{mtk.2}
\bookmark[dest={mtk.2},level=2]{1. นยมาติกา}
\csromanheader{2}{1. \textbf{นยมาติกา}}
\proseitem{1}{\csromanfolio{1}สงฺคโห อสงฺคโห. สงฺคหิเตน อสงฺคหิตํ. อสงฺคหิเตน สงฺคหิตํ. สงฺคหิเตน สงฺคหิตํ. อสงฺคหิเตน อสงฺคหิตํ. สมฺปโยโค วิปฺปโยโค. สมฺปยุตฺเตน วิปฺปยุตฺตํ. วิปฺปยุตฺเตน สมฺปยุตฺตํ. สมฺปยุตฺเตน สมฺปยุตฺตํ. วิปฺปยุตฺเตน วิปฺปยุตฺตํ. สงฺคหิเตน สมฺปยุตฺตํ วิปฺปยุตฺตํ. สมฺปยุตฺเตน สงฺคหิตํ อสงฺคหิตํ. อสงฺคหิเตน สมฺปยุตฺตํ วิปฺปยุตฺตํ. วิปฺปยุตฺเตน สงฺคหิตํ อสงฺคหิตํ.}

\csromanmatikaanchor{mtk.3}
\csromantocmarkref{subsection}{2. อพฺภนฺตรมาติกา}{mtk.3}
\bookmark[dest={mtk.3},level=2]{2. อพฺภนฺตรมาติกา}
\csromanheader{2}{2. อพฺภนฺตร\-มาติกา}
\proseitem{2}{ปญฺจกฺขนฺธา. ทฺวาทสา\-ยตนานิ. อฏฺฐารส ธาตุโย. จตฺตาริ สจฺจานิ. พาวีสติ\-นฺทฺริยานิ. ปฏิจฺจ\-สมุปฺปาโท. จตฺตาโร สติปฏฺฐานา. จตฺตาโร สมฺมปฺปธานา. จตฺตาโร อิทฺธิปาทา. จตฺตาริ ฌานานิ. จตสฺโส อปฺปมญฺญาโย. ปญฺจินฺทฺริยานิ. ปญฺจ พลานิ. สตฺต โพชฺฌงฺคา. อริโย อฏฺฐงฺคิโก มคฺโค. ผสฺโส เวทนา สญฺญา เจตนา จิตฺตํ อธิโมกฺโข มนสิกาโร.}

\csromanmatikaanchor{mtk.4}
\csromantocmarkref{subsection}{3. นยมุขมาติกา}{mtk.4}
\bookmark[dest={mtk.4},level=2]{3. นยมุขมาติกา}
\csromanheader{2}{3. นยมุข\-มาติกา}
\proseitem{3}{\csromanfolio{2}ตีหิ สงฺคโห. ตีหิ อสงฺคโห. จตูหิ สมฺปโยโค. จตูหิ วิปฺปโยโค.}

\csromanmatikaanchor{mtk.5}
\csromantocmarkref{subsection}{4. ลกฺขณมาติกา}{mtk.5}
\bookmark[dest={mtk.5},level=2]{4. ลกฺขณมาติกา}
\csromanheader{2}{4. \textbf{ลกฺขณมาติกา}}
\vspace{\tipitakaparskip}
\csromancenter{สภาโค. วิสภาโค.}
\csromanmatikaanchor{mtk.6}
\csromantocmarkref{subsection}{5. พาหิรมาติกา}{mtk.6}
\bookmark[dest={mtk.6},level=2]{5. พาหิรมาติกา}
\csromanheader{2}{5. \textbf{พาหิรมาติกา}}
\vspace{\tipitakaparskip}
\csromancenter{สพฺพาปิ ธมฺมสงฺคณี ธาตุกถาย มาติกาติ.}
\csromanmatikaanchor{mtk.7}
\csromanmatikaanchor{mtk.8}
\csromantocmarknopage{chapter}{1. ปฐมนย}{mtk.7}
\bookmark[dest={mtk.7},level=0]{1. ปฐมนย}
\csromantocmarknopage{section}{1. สงฺคหาสงฺคหปทนิทฺเทส}{mtk.8}
\bookmark[dest={mtk.8},level=1]{1. สงฺคหาสงฺคหปทนิทฺเทส}
\csromanheader{1}{1. สงฺคหา\-สงฺคห\-ปท\-นิทฺเทส}
\csromanmatikaanchor{mtk.9}
\csromantocmarkref{subsection}{1. ขนฺธ}{mtk.9}
\bookmark[dest={mtk.9},level=2]{1. ขนฺธ}
\csromanheader{2}{1. \textbf{ขนฺธ}}
\proseitem{6}{\csromanfolio{3}รูปกฺขนฺโธ กติหิ ขนฺเธหิ กติหายตเนหิ กติหิ ธาตูหิ สงฺคหิโต, รูปกฺขนฺโธ เอเกน ขนฺเธน เอกาทสหา\-ยตเนหิ เอกาทสหิ ธาตูหิ สงฺคหิโต. กติหิ อสงฺคหิโต, จตูหิ ขนฺเธหิ เอเกนายตเนน สตฺตหิ ธาตูหิ อสงฺคหิโต. (1)}

\proseitem{7}{เวทนากฺขนฺโธ กติหิ ขนฺเธหิ กติหายตเนหิ กติหิ ธาตูหิ สงฺคหิโต, เวทนากฺขนฺโธ เอเกน ขนฺเธน เอเกนายตเนน เอกาย ธาตุยา สงฺคหิโต. กติหิ อสงฺคหิโต, จตูหิ ขนฺเธหิ เอกาทสหา\-ยตเนหิ สตฺตรสหิ ธาตูหิ อสงฺคหิโต. (2)}

\proseitem{8}{สญฺญากฺขนฺโธ กติหิ ขนฺเธหิ กติหายตเนหิ กติหิ ธาตูหิ สงฺคหิโต, สญฺญากฺขนฺโธ เอเกน ขนฺเธน เอเกนายตเนน เอกาย ธาตุยา สงฺคหิโต. กติหิ อสงฺคหิโต, จตูหิ ขนฺเธหิ เอกาทสหา\-ยตเนหิ สตฺตรสหิ ธาตูหิ อสงฺคหิโต. (3)}

\proseitem{9}{สงฺขาร\-กฺขนฺโธ กติหิ ขนฺเธหิ กติหายตเนหิ กติหิ ธาตูหิ สงฺคหิโต, สงฺขาร\-กฺขนฺโธ เอเกน ขนฺเธน เอเกนายตเนน เอกาย ธาตุยา สงฺคหิโต. กติหิ อสงฺคหิโต, จตูหิ ขนฺเธหิ เอกาทสหา\-ยตเนหิ สตฺตรสหิ ธาตูหิ อสงฺคหิโต. (4)}

\proseitem{10}{วิญฺญาณ\-กฺขนฺโธ กติหิ ขนฺเธหิ กติหายตเนหิ กติหิ ธาตูหิ สงฺคหิโต, วิญฺญาณ\-กฺขนฺโธ เอเกน ขนฺเธน เอเกนายตเนน สตฺตหิ ธาตูหิ สงฺคหิโต. กติหิ อสงฺคหิโต, จตูหิ ขนฺเธหิ เอกาทสหา\-ยตเนหิ เอกาทสหิ ธาตูหิ อสงฺคหิโต. (5)}

\proseitem{11}{รูปกฺขนฺโธ จ เวทนากฺขนฺโธ จ กติหิ ขนฺเธหิ กติหายตเนหิ กติหิ ธาตูหิ สงฺคหิตา, รูปกฺขนฺโธ จ เวทนากฺขนฺโธ จ ทฺวีหิ ขนฺเธหิ \csromanfolio{4}เอกาทสหา\-ยตเนหิ เอกาทสหิ ธาตูหิ สงฺคหิตา. กติหิ อสงฺคหิตา, ตีหิ ขนฺเธหิ เอเกนายตเนน สตฺตหิ ธาตูหิ อสงฺคหิตา. (1)}

\prose{รูปกฺขนฺโธ จ สญฺญากฺขนฺโธ จ ฯเปฯ ทฺวีหิ ขนฺเธหิ เอกาทสหา\-ยตเนหิ เอกาทสหิ ธาตูหิ สงฺคหิตา. กติหิ อสงฺคหิตา, ตีหิ ขนฺเธหิ เอเกนายตเนน สตฺตหิ ธาตูหิ อสงฺคหิตา. (2)}

\prose{รูปกฺขนฺโธ จ สงฺขาร\-กฺขนฺโธ จ ฯเปฯ ทฺวีหิ ขนฺเธหิ เอกาทสหา\-ยตเนหิ เอกาทสหิ ธาตูหิ สงฺคหิตา. กติหิ อสงฺคหิตา, ตีหิ ขนฺเธหิ เอเกนายตเนน สตฺตหิ ธาตูหิ อสงฺคหิตา. (3)}

\prose{รูปกฺขนฺโธ จ วิญฺญาณ\-กฺขนฺโธ จ ฯเปฯ ทฺวีหิ ขนฺเธหิ ทฺวาทสหา\-ยตเนหิ อฏฺฐารสหิ ธาตูหิ สงฺคหิตา. กติหิ อสงฺคหิตา, ตีหิ ขนฺเธหิ น เกหิจิ อายตเนหิ น กาหิจิ ธาตูหิ อสงฺคหิตา. (4)}

\proseitem{15}{รูปกฺขนฺโธ จ เวทนากฺขนฺโธ จ สญฺญากฺขนฺโธ จ กติหิ ขนฺเธหิ กติหายตเนหิ กติหิ ธาตูหิ สงฺคหิตา, รูปกฺขนฺโธ จ เวทนากฺขนฺโธ จ สญฺญากฺขนฺโธ จ ตีหิ ขนฺเธหิ เอกาทสหา\-ยตเนหิ เอกาทสหิ ธาตูหิ สงฺคหิตา. กติหิ อสงฺคหิตา, ทฺวีหิ ขนฺเธหิ เอเกนายตเนน สตฺตหิ ธาตูหิ อสงฺคหิตา. (1)}

\proseitem{16}{รูปกฺขนฺโธ จ เวทนากฺขนฺโธ จ สงฺขาร\-กฺขนฺโธ จ ฯเปฯ ตีหิ ขนฺเธหิ เอกาทสหา\-ยตเนหิ เอกาทสหิ ธาตูหิ สงฺคหิตา. กติหิ อสงฺคหิตา, ทฺวีหิ ขนฺเธหิ เอเกนายตเนน สตฺตหิ ธาตูหิ อสงฺคหิตา. (2)}

\proseitem{17}{รูปกฺขนฺโธ จ เวทนากฺขนฺโธ จ วิญฺญาณ\-กฺขนฺโธ จ ฯเปฯ ตีหิ ขนฺเธหิ ทฺวาทสหา\-ยตเนหิ อฏฺฐารสหิ ธาตูหิ สงฺคหิตา. กติหิ อสงฺคหิตา, ทฺวีหิ ขนฺเธหิ น เกหิจิ อายตเนหิ น กาหิจิ ธาตูหิ อสงฺคหิตา. (3) (ติกมูลกํ.)}

\proseitem{18}{\csromanfolio{5}รูปกฺขนฺโธ จ เวทนากฺขนฺโธ จ สญฺญากฺขนฺโธ จ สงฺขาร\-กฺขนฺโธ จ กติหิ ขนฺเธหิ กติหายตเนหิ กติหิ ธาตูหิ สงฺคหิตา. รูปกฺขนฺโธ จ เวทนากฺขนฺโธ จ สญฺญากฺขนฺโธ จ สงฺขาร\-กฺขนฺโธ จ จตูหิ ขนฺเธหิ เอกาทสหา\-ยตเนหิ เอกาทสหิ ธาตูหิ สงฺคหิตา. กติหิ อสงฺคหิตา, เอเกน ขนฺเธน เอเกนายตเนน สตฺตหิ ธาตูหิ อสงฺคหิตา.}

\proseitem{19}{รูปกฺขนฺโธ จ เวทนากฺขนฺโธ จ สญฺญากฺขนฺโธ จ วิญฺญาณ\-กฺขนฺโธ จ ฯเปฯ จตูหิ ขนฺเธหิ ทฺวาทสหา\-ยตเนหิ อฏฺฐารสหิ ธาตูหิ สงฺคหิตา. กติหิ อสงฺคหิตา, เอเกน ขนฺเธน น เกหิจิ อายตเนหิ น กาหิจิ ธาตูหิ อสงฺคหิตา. (จตุกฺกมูลกํ.)}

\proseitem{20}{รูปกฺขนฺโธ จ เวทนากฺขนฺโธ จ สญฺญากฺขนฺโธ จ สงฺขาร\-กฺขนฺโธ จ วิญฺญาณ\-กฺขนฺโธ จ กติหิ ขนฺเธหิ กติหายตเนติ กติหิ ธาตูหิ สงฺคหิตา, รูปกฺขนฺโธ จ เวทนากฺขนฺโธ จ สญฺญากฺขนฺโธ จ สงฺขาร\-กฺขนฺโธ จ วิญฺญาณ\-กฺขนฺโธ จ ปญฺจหิ ขนฺเธหิ ทฺวาทสหา\-ยตเนหิ อฏฺฐารสหิ ธาตูหิ สงฺคหิตา. กติหิ อสงฺคหิตา, น เกหิจิ ขนฺเธหิ น เกหิจิ อายตเนหิ น กาหิจิ ธาตูหิ}

\proseitem{21}{ปญฺจกฺขนฺธา กติหิ ขนฺเธหิ กติหายตเนหิ กติหิ ธาตูหิ สงฺคหิตา, ปญฺจกฺขนฺธา ปญฺจหิ ขนฺเธหิ ทฺวาทสหา\-ยตเนหิ อฏฺฐารสหิ ธาตูหิ สงฺคหิตา. กติหิ อสงฺคหิตา, น เกหิจิ ขนฺเธหิ น เกหิจิ อายตเนหิ น กาหิจิ ธาตูหิ อสงฺคหิตา.}

\csromanheader{2}{(ปญฺจกํ.)}
\csromanmatikaanchor{mtk.10}
\csromantocmarkref{subsection}{2. อายตน}{mtk.10}
\bookmark[dest={mtk.10},level=2]{2. อายตน}
\csromanheader{2}{2. \textbf{อายตน}}
\proseitem{22}{จกฺขายตนํ กติหิ ขนฺเธหิ กติหายตเนหิ กติหิ ธาตูหิ สงฺคหิตํ, จกฺขายตนํ เอเกน ขนฺเธน เอเกนายตเนน เอกาย ธาตุยา สงฺคหิตํ. กติหิ อสงฺคหิตํ, จตูหิ ขนฺเธหิ เอกาทสหา\-ยตเนหิ}

\proseitem{23}{\csromanfolio{6}โสตายตนํ. ฆานายตนํ. ชิวฺหายตนํ. กายายตนํ. รูปายตนํ. สทฺทายตนํ. คนฺธายตนํ. รสายตนํ. โผฏฺฐพฺพา\-ยตนํ ฯเปฯ เอเกน ขนฺเธน เอเกนายตเนน เอกาย ธาตุยา สงฺคหิตํ. กติหิ อสงฺคหิตํ, จตูหิ ขนฺเธหิ เอกาทสหา\-ยตเนหิ}

\proseitem{24}{มนายตนํ เอเกน ขนฺเธน เอเกนายตเนน สตฺตหิ ธาตูหิ สงฺคหิตํ. กติหิ อสงฺคหิตํ, จตูหิ ขนฺเธหิ เอกาทสหา\-ยตเนหิ เอกาทสหิ ธาตูหิ อสงฺคหิตํ.}

\proseitem{25}{ธมฺมายตนํ อสงฺขตํ ขนฺธโต ฐเปตฺวา จตูหิ ขนฺเธหิ เอเกนายตเนน เอกาย ธาตุยา สงฺคหิตํ. กติหิ อสงฺคหิตํ, เอเกน ขนฺเธน เอกาทสหา\-ยตเนหิ สตฺตรสหิ ธาตูหิ อสงฺคหิตํ.}

\proseitem{26}{จกฺขายตนญฺจ โสตายตนญฺจ เอเกน ขนฺเธน ทฺวีหายตเนหิ ทฺวีหิ ธาตูหิ สงฺคหิตา. กติหิ อสงฺคหิตา, จตูหิ ขนฺเธหิ ทสหายตเนหิ}

\proseitem{27}{จกฺขายตนญฺจ ฆานายตนญฺจ. จกฺขายตนญฺจ ชิวฺหายตนญฺจ. จกฺขายตนญฺจ กายายตนญฺจ. จกฺขายตนญฺจ รูปายตนญฺจ. จกฺขายตนญฺจ สทฺทายตนญฺจ. จกฺขายตนญฺจ คนฺธายตนญฺจ. จกฺขายตนญฺจ รสายตนญฺจ. จกฺขายตนญฺจ โผฏฺฐพฺพายตนญฺจ เอเกน ขนฺเธน ทฺวีหายตเนหิ ทฺวีหิ ธาตูหิ สงฺคหิตา. กติหิ อสงฺคหิตา, จตูหิ ขนฺเธหิ ทสหายตเนหิ โสฬสหิ ธาตูหิ อสงฺคหิตา.}

\proseitem{28}{จกฺขายตนญฺจ มนายตนญฺจ ทฺวีหิ ขนฺเธหิ ทฺวีหายตเนหิ อฏฺฐหิ ธาตูหิ สงฺคหิตา. กติหิ อสงฺคหิตา, ตีหิ ขนฺเธหิ ทสหายตเนหิ ทสหิ ธาตูหิ อสงฺคหิตา.}

\proseitem{29}{จกฺขายตนญฺจ ธมฺมายตนญฺจ อสงฺขตํ ขนฺธโต ฐเปตฺวา จตูหิ ขนฺเธหิ ทฺวีหายตเนหิ ทฺวีหิ ธาตูหิ สงฺคหิตา. กติหิ อสงฺคหิตา, เอเกน ขนฺเธน ทสหายตเนหิ โสฬสหิ ธาตูหิ อสงฺคหิตา.}

\proseitem{30}{\csromanfolio{7}ทฺวาทสา\-ยตนานิ กติหิ ขนฺเธหิ กติหายตเนหิ กติหิ ธาตูหิ สงฺคหิตา, ทฺวาทสา\-ยตนานิ อสงฺขตํ ขนฺธโต ฐเปตฺวา ปญฺจหิ ขนฺเธหิ ทฺวาทสหา\-ยตเนหิ อฏฺฐารสหิ ธาตูหิ สงฺคหิตา. กติหิ อสงฺคหิตา, น เกหิจิ ขนฺเธหิ น เกหิจิ อายตเนหิ น กาหิจิ ธาตูหิ}

\csromanheader{2}{(ทฺวาทสกํ.)}
\csromanmatikaanchor{mtk.11}
\csromantocmarkref{subsection}{3. ธาตุ}{mtk.11}
\bookmark[dest={mtk.11},level=2]{3. ธาตุ}
\csromanheader{2}{3. \textbf{ธาตุ}}
\proseitem{31}{จกฺขุธาตุ กติหิ ขนฺเธหิ กติหายตเนหิ กติหิ ธาตูหิ สงฺคหิตา, จกฺขุธาตุ เอเกน ขนฺเธน เอเกนายตเนน เอกาย ธาตุยา สงฺคหิตา. กติหิ อสงฺคหิตา, จตูหิ ขนฺเธหิ เอกาทสหา\-ยตเนหิ สตฺตรสหิ ธาตูหิ อสงฺคหิตา.}

\proseitem{32}{โสตธาตุ. ฆานธาตุ. ชิวฺหาธาตุ. กายธาตุ. รูปธาตุ. สทฺทธาตุ. คนฺธธาตุ. รสธาตุ. โผฏฺฐพฺพ\-ธาตุ. จกฺขุวิญฺญาณ\-ธาตุ. โสตวิญฺญาณ\-ธาตุ. ฆานวิญฺญาณ\-ธาตุ. ชิวฺหาวิญฺญาณ\-ธาตุ. กายวิญฺญาณ\-ธาตุ. มโนธาตุ. มโนวิญฺญาณ\-ธาตุ เอเกน ขนฺเธน เอเกนายตเนน เอกาย ธาตุยา สงฺคหิตา. กติหิ อสงฺคหิตา, จตูหิ ขนฺเธหิ เอกาทสหา\-ยตเนหิ สตฺตรสหิ ธาตูหิ อสงฺคหิตา.}

\proseitem{33}{ธมฺมธาตุ อสงฺขตํ ขนฺธโต ฐเปตฺวา จตูหิ ขนฺเธหิ เอเกนายตเนน เอกาย ธาตุยา สงฺคหิตา. กติหิ อสงฺคหิตา, เอเกน ขนฺเธน เอกาทสหา\-ยตเนหิ สตฺตรสหิ ธาตูหิ อสงฺคหิตา.}

\proseitem{34}{จกฺขุธาตุ จ โสตธาตุ จ เอเกน ขนฺเธน ทฺวีหายตเนหิ ทฺวีหิ ธาตูหิ สงฺคหิตา. กติหิ อสงฺคหิตา, จตูหิ ขนฺเธหิ ทสหายตเนหิ}

\proseitem{35}{จกฺขุธาตุ จ ฆานธาตุ จ. จกฺขุธาตุ จ ชิวฺหาธาตุ จ. จกฺขุธาตุ จ กายธาตุ จ. จกฺขุธาตุ จ รูปธาตุ จ. จกฺขุธาตุ จ สทฺทธาตุ จ. \csromanfolio{8}จกฺขุธาตุ จ คนฺธธาตุ จ. จกฺขุธาตุ จ รสธาตุ จ. จกฺขุธาตุ จ โผฏฺฐพฺพ\-ธาตุ จ เอเกน ขนฺเธน ทฺวีหายตเนหิ ทฺวีหิ ธาตูหิ สงฺคหิตา. กติหิ อสงฺคหิตา, จตูหิ ขนฺเธหิ ทสหายตเนหิ}

\proseitem{36}{จกฺขุธาตุ จ จกฺขุวิญฺญาณ\-ธาตุ จ ทฺวีหิ ขนฺเธหิ ทฺวีหายตเนหิ ทฺวีหิ ธาตูหิ สงฺคหิตา. กติหิ อสงฺคหิตา, ตีหิ ขนฺเธหิ ทสหายตเนหิ โสฬสหิ ธาตูหิ อสงฺคหิตา.}

\proseitem{37}{จกฺขุธาตุ จ โสตวิญฺญาณ\-ธาตุ จ. จกฺขุธาตุ จ ฆานวิญฺญาณ\-ธาตุ จ. จกฺขุธาตุ จ ชิวฺหาวิญฺญาณ\-ธาตุ จ. จกฺขุธาตุ จ กายวิญฺญาณ\-ธาตุ จ. จกฺขุธาตุ จ มโนธาตุ จ. จกฺขุธาตุ จ มโนวิญฺญาณ\-ธาตุ จ ทฺวีหิ ขนฺเธหิ ทฺวีหายตเนหิ ทฺวีหิ ธาตูหิ สงฺคหิตา. กติหิ อสงฺคหิตา, ตีหิ ขนฺเธหิ ทสหายตเนหิ โสฬสหิ ธาตูหิ}

\proseitem{38}{จกฺขุธาตุ จ ธมฺมธาตุ จ อสงฺขตํ ขนฺธโต ฐเปตฺวา จตูหิ ขนฺเธหิ ทฺวีหายตเนหิ ทฺวีหิ ธาตูหิ สงฺคหิตา. กติหิ อสงฺคหิตา, เอเกน ขนฺเธน ทสหายตเนหิ โสฬสหิ ธาตูหิ อสงฺคหิตา.}

\proseitem{39}{อฏฺฐารส ธาตุโย กติหิ ขนฺเธหิ กติหายตเนหิ กติหิ ธาตูหิ สงฺคหิตา, อฏฺฐารส ธาตุโย อสงฺขตํ ขนฺธโต ฐเปตฺวา ปญฺจหิ ขนฺเธหิ ทฺวาทสหา\-ยตเนหิ อฏฺฐารสหิ ธาตูหิ สงฺคหิตา. กติหิ อสงฺคหิตา, น เกหิจิ ขนฺเธหิ น เกหิจิ อายตเนหิ น กาหิจิ ธาตูหิ}

\csromanheader{2}{(อฏฺฐารสกํ.)}
\csromanmatikaanchor{mtk.12}
\csromantocmarkref{subsection}{4. สจฺจ}{mtk.12}
\bookmark[dest={mtk.12},level=2]{4. สจฺจ}
\csromanheader{2}{4. \textbf{สจฺจ}}
\proseitem{40}{ทุกฺขสจฺจํ กติหิ ขนฺเธหิ กติหายตเนหิ กติหิ ธาตูหิ สงฺคหิตํ, ทุกฺขสจฺจํ ปญฺจหิ ขนฺเธหิ ทฺวาทสหา\-ยตเนหิ อฏฺฐารสหิ ธาตูหิ สงฺคหิตํ. กติหิ อสงฺคหิตํ, น เกหิจิ ขนฺเธหิ น เกหิจิ อายตเนหิ น กาหิจิ ธาตูหิ อสงฺคหิตํ.}

\proseitem{41}{\csromanfolio{9}สมุทยสจฺจํ. มคฺคสจฺจํ เอเกน ขนฺเธน เอเกนายตเนน เอกาย ธาตุยา สงฺคหิตํ. กติหิ อสงฺคหิตํ, จตูหิ ขนฺเธหิ เอกาทสหา\-ยตเนหิ สตฺตรสหิ ธาตูหิ อสงฺคหิตํ.}

\proseitem{42}{นิโรธสจฺจํ น เกหิจิ ขนฺเธหิ เอเกนายตเนน เอกาย ธาตุยา สงฺคหิตํ. กติหิ อสงฺคหิตํ, ปญฺจหิ ขนฺเธหิ เอกาทสหา\-ยตเนหิ}

\proseitem{43}{ทุกฺขสจฺจญฺจ สมุทยสจฺจญฺจ ปญฺจหิ ขนฺเธหิ ทฺวาทสหา\-ยตเนหิ อฏฺฐารสหิ ธาตูหิ สงฺคหิตา. กติหิ อสงฺคหิตา, น เกหิจิ ขนฺเธหิ น เกหิจิ อายตเนหิ น กาหิจิ ธาตูหิ อสงฺคหิตา.}

\proseitem{44}{ทุกฺขสจฺจญฺจ มคฺคสจฺจญฺจ ปญฺจหิ ขนฺเธหิ ทฺวาทสหา\-ยตเนหิ อฏฺฐารสหิ ธาตูหิ สงฺคหิตา. กติหิ อสงฺคหิตา, น เกหิจิ ขนฺเธหิ น เกหิจิ อายตเนหิ น กาหิจิ ธาตูหิ อสงฺคหิตา.}

\proseitem{45}{ทุกฺขสจฺจญฺจ นิโรธสจฺจญฺจ อสงฺขตํ ขนฺธโต ฐเปตฺวา ปญฺจหิ ขนฺเธหิ ทฺวาทสหา\-ยตเนหิ อฏฺฐารสหิ ธาตูหิ สงฺคหิตา. กติหิ อสงฺคหิตา, น เกหิจิ ขนฺเธหิ น เกหิจิ อายตเนหิ น กาหิจิ ธาตูหิ}

\proseitem{46}{ทุกฺขสจฺจญฺจ สมุทยสจฺจญฺจ มคฺคสจฺจญฺจ ปญฺจหิ ขนฺเธหิ ทฺวาทสหา\-ยตเนหิ อฏฺฐารสหิ ธาตูหิ สงฺคหิตา. กติหิ อสงฺคหิตา, น เกหิจิ ขนฺเธหิ น กาหิจิ อายตเนหิ น กาหิจิ ธาตูหิ}

\proseitem{47}{ทุกฺขสจฺจญฺจ สมุทยสจฺจญฺจ นิโรธสจฺจญฺจ อสงฺขตํ ขนฺธโต ฐเปตฺวา ปญฺจหิ ขนฺเธหิ ทฺวาทสหา\-ยตเนหิ อฏฺฐารสหิ ธาตูหิ สงฺคหิตา. กติหิ อสงฺคหิตา, น เกหิจิ ขนฺเธหิ น เกหิจิ อายตเนหิ น กาหิจิ ธาตูหิ อสงฺคหิตา.}

\csromanheader{2}{(ติกมูลกํ.)}
\proseitem{48}{\csromanfolio{10}ทุกฺขสจฺจญฺจ สมุทยสจฺจญฺจ มคฺคสจฺจญฺจ นิโรธสจฺจญฺจ อสงฺขตํ ขนฺธโต ฐเปตฺวา ปญฺจหิ ขนฺเธหิ ทฺวาทสหา\-ยตเนหิ อฏฺฐารสหิ ธาตูหิ สงฺคหิตา. กติหิ อสงฺคหิตา, น เกหิจิ ขนฺเธหิ น เกหิจิ อายตเนหิ น กาหิจิ ธาตูหิ อสงฺคหิตา.}

\proseitem{49}{จตฺตาริ สจฺจานิ กติหิ ขนฺเธหิ กติหายตเนหิ กติหิ ธาตูหิ สงฺคหิตานิ, จตฺตาริ สจฺจานิ อสงฺขตํ ขนฺธโต ฐเปตฺวา ปญฺจหิ ขนฺเธหิ ทฺวาทสหา\-ยตเนหิ อฏฺฐารสหิ ธาตูหิ สงฺคหิตานิ. กติหิ อสงฺคหิตานิ, น เกหิจิ ขนฺเธหิ น เกหิจิ อายตเนหิ น กาหิจิ ธาตูหิ อสงฺคหิตานิ.}

\csromanheader{2}{(จตุกฺกํ.)}
\csromanmatikaanchor{mtk.13}
\csromantocmarkref{subsection}{5. อินฺทฺริย}{mtk.13}
\bookmark[dest={mtk.13},level=2]{5. อินฺทฺริย}
\csromanheader{2}{5. \textbf{อินฺทฺริย}}
\proseitem{50}{จกฺขุนฺทฺริยํ กติหิ ขนฺเธหิ กติหายตเนหิ กติหิ ธาตูหิ สงฺคหิตํ, จกฺขุนฺทฺริยํ เอเกน ขนฺเธน เอเกนายตเนน เอกาย ธาตุยา สงฺคหิตํ. กติหิ อสงฺคหิตํ, จตูหิ ขนฺเธหิ เอกาทสหา\-ยตเนหิ}

\proseitem{51}{โสตินฺทฺริยํ. ฆานินฺทฺริยํ. ชิวฺหินฺทฺริยํ. กายินฺทฺริยํ. อิตฺถินฺทฺริยํ. ปุริสินฺทฺริยํ เอเกน ขนฺเธน เอเกนายตเนน เอกาย ธาตุยา สงฺคหิตํ. กติหิ อสงฺคหิตํ, จตูหิ ขนฺเธหิ เอกาทสหา\-ยตเนหิ}

\proseitem{52}{มนินฺทฺริยํ เอเกน ขนฺเธน เอเกนายตเนน สตฺตหิ ธาตูหิ สงฺคหิตํ. กติหิ อสงฺคหิตํ, จตูหิ ขนฺเธหิ เอกาทสหา\-ยตเนหิ เอกาทสหิ ธาตูหิ อสงฺคหิตํ.}

\proseitem{53}{ชีวิตินฺทฺริยํ ทฺวีหิ ขนฺเธหิ เอเกนายตเนน เอกาย ธาตุยา สงฺคหิตํ. กติหิ อสงฺคหิตํ, ตีหิ ขนฺเธหิ เอกาทสหา\-ยตเนหิ}

\proseitem{54}{\csromanfolio{11}สุขินฺทฺริยํ. ทุกฺขินฺทฺริยํ. โสมนสฺสิ\-นฺทฺริยํ. โทมนสฺสิ\-นฺทฺริยํ. อุเปกฺขินฺทฺริยํ. สทฺธินฺทฺริยํ. วีริยินฺทฺริยํ. สตินฺทฺริยํ. สมาธินฺทฺริยํ. ปญฺญินฺทฺริยํ. อนญฺญาต\-ญฺญสฺสามี\-ตินฺทฺริยํ. อญฺญินฺทฺริยํ. อญฺญาตาวิ\-นฺทฺริยํ เอเกน ขนฺเธน เอเกนายตเนน เอกาย ธาตุยา สงฺคหิตํ. กติหิ อสงฺคหิตํ, จตูหิ ขนฺเธหิ เอกาทสหา\-ยตเนหิ สตฺตรสหิ ธาตูหิ อสงฺคหิตํ.}

\proseitem{55}{จกฺขุนฺทฺริยญฺจ โสตินฺทฺริยญฺจ เอเกน ขนฺเธน ทฺวีหายตเนหิ ทฺวีหิ ธาตูหิ สงฺคหิตา. กติหิ อสงฺคหิตา, จตูหิ ขนฺเธหิ ทสหายตเนหิ}

\proseitem{56}{จกฺขุนฺทฺริยญฺจ ฆานินฺทฺริยญฺจ. จกฺขุนฺทฺริยญฺจ ชิวฺหินฺทฺริยญฺจ. จกฺขุนฺทฺริยญฺจ กายินฺทฺริยญฺจ. จกฺขุนฺทฺริยญฺจ อิตฺถินฺทฺริยญฺจ. จกฺขุนฺทฺริยญฺจ ปุริสินฺทฺริยญฺจ เอเกน ขนฺเธน ทฺวีหายเนหิ ทฺวีหิ ธาตูหิ สงฺคหิตา. กติหิ อสงฺคหิตา, จตูหิ ขนฺเธหิ ทสหายตเนหิ}

\proseitem{57}{จกฺขุนฺทฺริยญฺจ มนินฺทฺริยญฺจ ทฺวีหิ ขนฺเธหิ ทฺวีหายตเนหิ อฏฺฐหิ ธาตูหิ สงฺคหิตา. กติหิ อสงฺคหิตา, ตีหิ ขนฺเธหิ ทสหายตเนหิ ทสหิ ธาตูหิ อสงฺคหิตา.}

\proseitem{58}{จกฺขุนฺทฺริยญฺจ ชีวิตินฺทฺริยญฺจ ทฺวีหิ ขนฺเธหิ ทฺวีหายตเนหิ ทฺวีหิ ธาตูหิ สงฺคหิตา. กติหิ อสงฺคหิตา, ตีหิ ขนฺเธหิ ทสหายตเนหิ}

\proseitem{59}{จกฺขุนฺทฺริยญฺจ สุขินฺทฺริยญฺจ. จกฺขุนฺทฺริยญฺจ ทุกฺขินฺทฺริยญฺจ. จกฺขุนฺทฺริยญฺจ โสมนสฺสินฺทฺริยญฺจ. จกฺขุนฺทฺริยญฺจ โทมนสฺสินฺทฺริยญฺจ. จกฺขุนฺทฺริยญฺจ อุเปกฺขินฺทฺริยญฺจ. จกฺขุนฺทฺริยญฺจ สทฺธินฺทฺริยญฺจ. จกฺขุนฺทฺริยญฺจ วีริยินฺทฺริยญฺจ. จกฺขุนฺทฺริยญฺจ สตินฺทฺริยญฺจ. จกฺขุนฺทฺริยญฺจ สมาธินฺทฺริยญฺจ. จกฺขุนฺทฺริยญฺจ ปญฺญินฺทฺริยญฺจ. จกฺขุนฺทฺริยญฺจ อนญฺญาตญฺญสฺสามีตินฺทฺริยญฺจ. จกฺขุนฺทฺริยญฺจ อญฺญินฺทฺริยญฺจ. จกฺขุนฺทฺริยญฺจ อญฺญาตาวินฺทฺริยญฺจ ทฺวีหิ ขนฺเธหิ ทฺวีหายตเนหิ ทฺวีหิ ธาตูหิ สงฺคหิตา. กติหิ อสงฺคหิตา, ตีหิ ขนฺเธหิ ทสหายตเนหิ โสฬสหิ ธาตูหิ อสงฺคหิตา.}

\proseitem{60}{\csromanfolio{12}พาวีสติ\-นฺทฺริยานิ กติหิ ขนฺเธหิ กติหายตเนหิ กติหิ ธาตูหิ สงฺคหิตานิ. พาวีสติ\-นฺทฺริยานิ จตูหิ ขนฺเธหิ สตฺตหา\-ยตเนหิ เตรสหิ ธาตูหิ สงฺคหิตานิ. กติหิ อสงฺคหิตานิ, เอเกน ขนฺเธน ปญฺจหา\-ยตเนหิ ปญฺจหิ ธาตูหิ อสงฺคหิตานิ.}

\csromanmatikaanchor{mtk.14}
\csromantocmarkref{subsection}{6. ปฏิจฺจสมุปฺปาทาทิ}{mtk.14}
\bookmark[dest={mtk.14},level=2]{6. ปฏิจฺจสมุปฺปาทาทิ}
\csromanheader{2}{6. \textbf{ปฏิจฺจสมุปฺปาทาทิ}}
\proseitem{61}{อวิชฺชา เอเกน ขนฺเธน เอเกนายตเนน เอกาย ธาตุยา สงฺคหิตา. กติหิ อสงฺคหิตา, จตูหิ ขนฺเธหิ เอกาทสหา\-ยตเนหิ สตฺตรสหิ ธาตูหิ}

\proseitem{62}{อวิชฺชาปจฺจยา สงฺขารา เอเกน ขนฺเธน เอเกนายตเนน เอกาย ธาตุยา สงฺคหิตา. กติหิ อสงฺคหิตา, จตูหิ ขนฺเธหิ เอกาทสหา\-ยตเนหิ สตฺตรสหิ ธาตูหิ อสงฺคหิตา.}

\proseitem{63}{สงฺขาร\-ปจฺจยา วิญฺญาณํ เอเกน ขนฺเธน เอเกนายตเนน สตฺตหิ ธาตูหิ สงฺคหิตํ. กติหิ อสงฺคหิตํ, จตูหิ ขนฺเธหิ เอกาทสหา\-ยตเนหิ เอกาทสหิ ธาตูหิ อสงฺคหิตํ.}

\proseitem{64}{วิญฺญาณปจฺจยา นามรูปํ จตูหิ ขนฺเธหิ เอกาทสหา\-ยตเนหิ เอกาทสหิ ธาตูหิ สงฺคหิตํ. กติหิ อสงฺคหิตํ, เอเกน ขนฺเธน เอเกนายตเนน สตฺตหิ ธาตูหิ อสงฺคหิตํ.}

\proseitem{65}{นามรูป\-ปจฺจยา สฬายตนํ ทฺวีหิ ขนฺเธหิ ฉหายตเนหิ ทฺวาทสหิ ธาตูหิ สงฺคหิตํ. กติหิ อสงฺคหิตํ, ตีหิ ขนฺเธหิ ฉหายตเนหิ ฉหิ ธาตูหิ อสงฺคหิตํ.}

\proseitem{66}{สฬายตน\-ปจฺจยา ผสฺโส. ผสฺสปจฺจยา เวทนา. เวทนาปจฺจยา ตณฺหา. ตณฺหาปจฺจยา อุปาทานํ. กมฺมภโว\csromansharedfootnote{6544f97e1ad0d78d}{อุปาทานปจฺจยา กมฺมภโว (สี, สฺยา)} เอเกน ขนฺเธน เอเกนายตเนน เอกาย ธาตุยา สงฺคหิโต. กติหิ อสงฺคหิโต, จตูหิ ขนฺเธหิ เอกาทสหา\-ยตเนหิ สตฺตรสหิ ธาตูหิ อสงฺคหิโต.}

\proseitem{67}{\csromanfolio{13}อุปปตฺติภโว. กามภโว. สญฺญาภโว. ปญฺจโวการภโว ปญฺจหิ ขนฺเธหิ เอกาทสหา\-ยตเนหิ สตฺตรสหิ ธาตูหิ สงฺคหิโต. กติหิ อสงฺคโต, น เกหิจิ ขนฺเธหิ เอเกนายตเนน เอกาย ธาตุยา อสงฺคหิโต.}

\proseitem{68}{รูปภโว ปญฺจหิ ขนฺเธหิ ปญฺจหายเนหิ อฏฺฐหิ ธาตูหิ สงฺคหิโต. กติหิ อสงฺคหิโต, น เกหิจิ ขนฺเธหิ สตฺตหา\-ยตเนหิ ทสหิ ธาตูหิ อสงฺคโต.}

\proseitem{69}{อรูปภโว. เนว\-สญฺญา\-นา\-สญฺญาภโว. จตุโวการภโว จตูหิ ขนฺเธหิ ทฺวีหายตเนหิ ทฺวีหิ ธาตูหิ สงฺคหิโต. กติหิ อสงฺคหิโต, เอเกน ขนฺเธน ทสหายตเนหิ โสฬสหิ ธาตูหิ อสงฺคหิโต.}

\proseitem{70}{อสญฺญาภโว. เอกโวการภโว เอเกน ขนฺเธน ทฺวีหายตเนหิ ทฺวีหิ ธาตูหิ สงฺคหิโต. กติหิ อสงฺคหิโต, จตูหิ ขนฺเธหิ ทสหายตเนหิ โสฬสหิ ธาตูหิ อสงฺคหิโต.}

\proseitem{71}{ชาติ ทฺวีหิ ขนฺเธหิ. ชรา ทฺวีหิ ขนฺเธหิ. มรณํ ทฺวีหิ ขนฺเธหิ เอเกนายตเนน เอกาย ธาตุยา สงฺคหิตํ. กติหิ อสงฺคหิตํ, ตีหิ ขนฺเธหิ เอกาทสหา\-ยตเนหิ สตฺตรสหิ ธาตูหิ อสงฺคหิตํ.}

\proseitem{72}{โสโก. ปริเทโว. ทุกฺขํ. โทมนสฺสํ. อุปายาโส. สติปฏฺฐานํ. สมฺมปฺปธานํ เอเกน ขนฺเธน เอเกนายตเนน เอกาย ธาตุยา สงฺคหิตํ. กติหิ อสงฺคหิตํ, จตูหิ ขนฺเธหิ เอกาทสหา\-ยตเนหิ}

\proseitem{73}{อิทฺธิปาโท ทฺวีหิ ขนฺเธหิ ทฺวีหายตเนหิ ทฺวีหิ ธาตูหิ สงฺคหิโต. กติหิ อสงฺคหิโต, ตีหิ ขนฺเธหิ ทสหายตเนหิ โสฬสหิ ธาตูหิ อสงฺคหิโต.}

\proseitem{74}{ฌานํ ทฺวีหิ ขนฺเธหิ เอเกนายตเนน เอกาย ธาตุยา สงฺคหิตํ. กติหิ อสงฺคหิตํ, ตีหิ ขนฺเธหิ เอกาทสหา\-ยตเนหิ สตฺตรสหิ ธาตูหิ อสงฺคหิตํ.}

\proseitem{75}{\csromanfolio{14}อปฺปมญฺญา. ปญฺจินฺทฺริยานิ. ปญฺจ พลานิ. สตฺต โพชฺฌงฺคา. อริโย อฏฺฐงฺคิโก มคฺโค. ผสฺโส. เวทนา. สญฺญา. เจตนา. อธิโมกฺโข. มนสิกาโร เอเกน ขนฺเธน เอเกนายตเนน เอกาย ธาตุยา สงฺคหิโต. กติหิ อสงฺคหิโต, จตูหิ ขนฺเธหิ เอกาทสหา\-ยตเนหิ สตฺตรสหิ ธาตูหิ อสงฺคหิโต.}

\proseitem{76}{จิตฺตํ เอเกน ขนฺเธน เอเกนายตเนน สตฺตหิ ธาตูหิ สงฺคหิตํ. กติหิ อสงฺคหิตํ, จตูหิ ขนฺเธหิ เอกาทสหา\-ยตเนหิ เอกาทสหิ ธาตูหิ อสงฺคหิตํ.}

\csromanmatikaanchor{mtk.15}
\csromantocmarkref{subsection}{7. ติก}{mtk.15}
\bookmark[dest={mtk.15},level=2]{7. ติก}
\csromanheader{2}{7. \textbf{ติก}}
\proseitem{77}{กุสลา ธมฺมา. อกุสลา ธมฺมา กติหิ ขนฺเธหิ กติหายตเนหิ กติหิ ธาตูหิ สงฺคหิตา. กุสลา ธมฺมา. อกุสลา ธมฺมา จตูหิ ขนฺเธหิ ทฺวีหายตเนหิ ทฺวีหิ ธาตูหิ สงฺคหิตา. กติหิ อสงฺคหิตา, เอเกน ขนฺเธน ทสหายตเนหิ โสฬสหิ ธาตูหิ อสงฺคหิตา.}

\proseitem{78}{อพฺยากตา ธมฺมา อสงฺขตํ ขนฺธโต ฐเปตฺวา ปญฺจหิ ขนฺเธหิ ทฺวาทสหา\-ยตเนหิ อฏฺฐารสหิ ธาตูหิ สงฺคหิตา. กติหิ อสงฺคหิตา, น เกหิจิ ขนฺเธหิ น เกหิจิ อายตเนหิ น กาหิจิ ธาตูหิ}

\proseitem{79}{สุขาย เวทนาย สมฺปยุตฺตา ธมฺมา. ทุกฺขาย เวทนาย สมฺปยุตฺตา ธมฺมา ตีหิ ขนฺเธหิ ทฺวีหายตเนหิ ตีหิ ธาตูหิ สงฺคหิตา. กติหิ อสงฺคหิตา, ทฺวีหิ ขนฺเธหิ ทสหายตเนหิ ปนฺนรสหิ ธาตูหิ}

\proseitem{80}{อทุกฺขม\-สุขาย เวทนาย สมฺปยุตฺตา ธมฺมา ตีหิ ขนฺเธหิ ทฺวีหายตเนหิ สตฺตหิ ธาตูหิ สงฺคหิตา. กติหิ อสงฺคหิตา, ทฺวีหิ ขนฺเธหิ ทสหายตเนหิ เอกาทสติ ธาตูหิ อสงฺคหิตา.}

\proseitem{81}{\csromanfolio{15}วิปากา ธมฺมา จตูหิ ขนฺเธหิ ทฺวีหายตเนหิ อฏฺฐหิ ธาตูหิ สงฺคหิตา. กติหิ อสงฺคหิตา, เอเกน ขนฺเธน ทสหายตเนหิ ทสหิ ธาตูหิ อสงฺคหิตา.}

\proseitem{82}{วิปาก\-ธมฺม\-ธมฺมา. สํกิลิฏฺฐ\-สํกิเลสิกา ธมฺมา จตูหิ ขนฺเธหิ ทฺวีหายตเนหิ ทฺวีหิ ธาตูหิ สงฺคหิตา. กติหิ อสงฺคหิตา, เอเกน ขนฺเธน ทสหายตเนหิ โสฬสหิ ธาตูหิ อสงฺคหิตา.}

\proseitem{83}{เนว\-วิปาก\-น\-วิปากธมฺม\-ธมฺมา อสงฺขตํ ขนฺธโต ฐเปตฺวา ปญฺจหิ ขนฺเธหิ ทฺวาทสหา\-ยตเนหิ เตรสหิ ธาตูหิ สงฺคหิตา. กติหิ อสงฺคหิตา, น เกหิจิ ขนฺเธหิ น เกหิจิ อายตเนหิ ปญฺจหิ ธาตูหิ}

\proseitem{84}{อุปาทินฺนุ\-ปาทานิยา ธมฺมา ปญฺจหิ ขนฺเธหิ เอกาทสหา\-ยตเนหิ สตฺตรสหิ ธาตูหิ สงฺคหิตา. กติหิ อสงฺคหิตา, น เกหิจิ ขนฺเธหิ เอเกนายตเนน เอกาย ธาตุยา อสงฺคหิตา.}

\proseitem{85}{อนุปาทินฺนุ\-ปาทานิยา ธมฺมา ปญฺจหิ ขนฺเธหิ สตฺตหา\-ยตเนหิ อฏฺฐหิ ธาตูหิ สงฺคหิตา. กติหิ อสงฺคหิตา, น เกหิจิ ขนฺเธหิ ปญฺจหา\-ยตเนหิ ทสหิ ธาตูหิ อสงฺคหิตา.}

\proseitem{86}{อนุปาทินฺน\-อนุปาทานิยา ธมฺมา. อสํกิลิฏฺฐ\-อสํกิเลสิกา ธมฺมา อสงฺขตํ ขนฺธโต ฐเปตฺวา จตูหิ ขนฺเธหิ ทฺวีหายตเนหิ ทฺวีหิ ธาตูหิ สงฺคหิตา. กติหิ อสงฺคหิตา, เอเกน ขนฺเธน ทสหายตเนหิ}

\proseitem{87}{อสํกิลิฏฺฐ\-สํกิเลสิกา ธมฺมา ปญฺจหิ ขนฺเธหิ ทฺวาทสหา\-ยตเนหิ อฏฺฐารสหิ ธาตูหิ สงฺคหิตา. กติหิ อสงฺคหิตา, น เกหิจิ ขนฺเธหิ น เกหิจิ อายตเนหิ น กาหิจิ ธาตูหิ อสงฺคหิตา.}

\proseitem{88}{สวิตกฺก\-สวิจารา ธมฺมา จตูหิ ขนฺเธหิ ทฺวีหายตเนหิ ตีหิ ธาตูหิ สงฺคหิตา. กติหิ อสงฺคหิตา, เอเกน ขนฺเธน ทสหายตเนหิ ปนฺนรสหิ ธาตูหิ อสงฺคหิตา.}

\proseitem{89}{\csromanfolio{16}อวิตกฺก\-วิจารมตฺตา ธมฺมา. ปีติสหคตา ธมฺมา จตูหิ ขนฺเธหิ ทฺวีหายตเนหิ ทฺวีหิ ธาตูหิ สงฺคหิตา. กติหิ อสงฺคหิตา, เอเกน ขนฺเธน ทสหายตเนหิ โสฬสหิ ธาตูหิ อสงฺคหิตา.}

\proseitem{90}{อวิตกฺก\-อวิจารา ธมฺมา อสงฺขตํ ขนฺธโต ฐเปตฺวา ปญฺจหิ ขนฺเธหิ ทฺวาทสหา\-ยตเนหิ สตฺตรสหิ ธาตูหิ สงฺคหิตา. กติหิ อสงฺคหิตา, น เกหิจิ ขนฺเธหิ น เกหิจิ อายตเนหิ เอกาย ธาตุยา}

\proseitem{91}{สุขสหคตา ธมฺมา ตีหิ ขนฺเธหิ ทฺวีหายตเนหิ ตีหิ ธาตูหิ สงฺคหิตา. กติหิ อสงฺคหิตา, ทฺวีหิ ขนฺเธหิ ทสหายตเนหิ ปนฺนรสหิ ธาตูหิ อสงฺคหิตา.}

\proseitem{92}{อุเปกฺขา\-สหคตา ธมฺมา ตีหิ ขนฺเธหิ ทฺวีหายตเนหิ สตฺตหิ ธาตูหิ สงฺคหิตา. กติหิ อสงฺคหิตา, ทฺวีหิ ขนฺเธหิ ทสหายตเนหิ เอกาทสหิ ธาตูหิ อสงฺคหิตา.}

\proseitem{93}{ทสฺสเนน ปหาตพฺพา ธมฺมา. ภาวนาย ปหาตพฺพา ธมฺมา. ทสฺสเนน ปหาตพฺพ\-เหตุกา ธมฺมา. ภาวนาย ปหาตพฺพ\-เหตุกา ธมฺมา. อาจยคามิโน ธมฺมา. อปจยคามิโน ธมฺมา. เสกฺขา ธมฺมา. อเสกฺขา ธมฺมา. มหคฺคตา ธมฺมา จตูหิ ขนฺเธหิ ทฺวีหายตเนหิ ทฺวีหิ ธาตูหิ สงฺคหิตา. กติหิ อสงฺคหิตา, เอเกน ขนฺเธน ทสหายตเนหิ}

\proseitem{94}{เนว ทสฺสเนน น ภาวนาย ปหาตพฺพา ธมฺมา. เนว ทสฺสเนน น ภาวนาย ปหาตพฺพ\-เหตุกา ธมฺมา. เนวาจยคามิ นาปจยคามิโน ธมฺมา. เนว\-เสกฺข\-นาเสกฺขา ธมฺมา อสงฺขตํ ขนฺธโต ฐเปตฺวา ปญฺจหิ ขนฺเธหิ ทฺวาทสหา\-ยตเนหิ อฏฺฐารสหิ ธาตูหิ สงฺคหิตา. กติหิ อสงฺคหิตา, น เกหิจิ ขนฺเธหิ น เกหิจิ อายตเนหิ น กาหิจิ ธาตูหิ}

\proseitem{95}{ปริตฺตา ธมฺมา ปญฺจหิ ขนฺเธหิ ทฺวาทสหา\-ยตเนหิ อฏฺฐารสหิ ธาตูหิ สงฺคหิตา. กติหิ อสงฺคหิตา, น เกหิจิ ขนฺเธหิ น เกหิจิ อายตเนหิ น กาหิจิ ธาตูหิ อสงฺคหิตา.}

\proseitem{96}{\csromanfolio{17}อปฺปมาณา ธมฺมา. ปณีตา ธมฺมา อสงฺขตํ ขนฺธโต ฐเปตฺวา จตูหิ ขนฺเธหิ ทฺวีหายตเนหิ ทฺวีหิ ธาตูหิ สงฺคหิตา. กติหิ อสงฺคหิตา, เอเกน ขนฺเธน ทสหายตเนหิ โสฬสหิ ธาตูหิ อสงฺคหิตา.}

\proseitem{97}{ปริตฺตารมฺมณา\csromansharedfootnote{f7305b9b54f5cd67}{ปริตฺตารมณา (?)} ธมฺมา จตูหิ ขนฺเธหิ ทฺวีหายตเนหิ อฏฺฐหิ ธาตูหิ สงฺคหิตา. กติหิ อสงฺคหิตา, เอเกน ขนฺเธน ทสหายตเนหิ ทสหิ ธาตูหิ อสงฺคหิตา.}

\proseitem{98}{มหคฺคตา\-รมฺมณา ธมฺมา. อปฺปมาณา\-รมฺมณา ธมฺมา. หีนา ธมฺมา. มิจฺฉตฺต\-นิยตา ธมฺมา. สมฺมตฺตนิยตา ธมฺมา, มคฺคารมฺมณา ธมฺมา. มคฺคเหตุกา ธมฺมา. มคฺคาธิปติโน ธมฺมา จตูหิ ขนฺเธหิ ทฺวีหายตเนหิ ทฺวีหิ ธาตูหิ สงฺคหิตา. กติหิ อสงฺคหิตา, เอเกน ขนฺเธน ทสหายตเนหิ โสฬสหิ ธาตูหิ อสงฺคหิตา.}

\proseitem{99}{มชฺฌิมา ธมฺมา ปญฺจหิ ขนฺเธหิ ทฺวาทสหา\-ยตเนหิ อฏฺฐารสหิ ธาตูหิ สงฺคหิตา. กติหิ อสงฺคหิตา, น เกหิจิ ขนฺเธหิ น เกหิจิ อายตเนหิ น กาหิจิ ธาตูหิ อสงฺคหิตา.}

\proseitem{100}{อนิยตา ธมฺมา อสงฺขตํ ขนฺธโต ฐเปตฺวา ปญฺจหิ ขนฺเธหิ ทฺวาทสหา\-ยตเนหิ อฏฺฐารสหิ ธาตูหิ สงฺคหิตา. กติหิ อสงฺคหิตา, น เกหิจิ ขนฺเธหิ น เกหิจิ อายตเนหิ น กาหิจิ ธาตูหิ}

\proseitem{101}{อุปฺปนฺนา ธมฺมา ปญฺจหิ ขนฺเธหิ ทฺวาทสหา\-ยตเนหิ อฏฺฐารสหิ ธาตูหิ สงฺคหิตา. กติหิ อสงฺคหิตา, น เกหิจิ ขนฺเธหิ น เกหิจิ อายตเนหิ น กาหิจิ ธาตูหิ อสงฺคหิตา.}

\proseitem{102}{อนุปฺปนฺนา ธมฺมา ปญฺจหิ ขนฺเธหิ สตฺตหา\-ยตเนหิ อฏฺฐหิ ธาตูหิ สงฺคหิตา. กติหิ อสงฺคหิตา, น เกหิจิ ขนฺเธหิ ปญฺจหา\-ยตเนหิ ทสหิ ธาตูหิ อสงฺคหิตา.}

\proseitem{103}{\csromanfolio{18}อุปฺปาทิโน ธมฺมา ปญฺจหิ ขนฺเธหิ เอกาทสหา\-ยตเนหิ สตฺตรสหิ ธาตูหิ สงฺคหิตา. กติหิ อสงฺคหิตา, น เกหิจิ ขนฺเธหิ เอเกนายตเนน เอกาย ธาตุยา อสงฺคหิตา.}

\proseitem{104}{อตีตา ธมฺมา. อนาคตา ธมฺมา. ปจฺจุปฺปนฺนา ธมฺมา. อชฺฌตฺตา ธมฺมา อชฺฌตฺต\-พหิทฺธา ธมฺมา ปญฺจหิ ขนฺเธหิ ทฺวาทสหา\-ยตเนหิ อฏฺฐารสหิ ธาตูหิ สงฺคหิตา. กติหิ อสงฺคหิตา, น เกหิจิ ขนฺเธหิ น เกหิจิ อายตเนหิ น กาหิจิ ธาตูหิ อสงฺคหิตา.}

\proseitem{105}{พหิทฺธา ธมฺมา อสงฺขตํ ขนฺธโต ฐเปตฺวา ปญฺจหิ ขนฺเธหิ ทฺวาทสหา\-ยตเนหิ อฏฺฐารสหิ ธาตูหิ สงฺคหิตา. กติหิ อสงฺคหิตา, น เกหิจิ ขนฺเธหิ น เกหิจิ อายตเนหิ น กาหิจิ ธาตูหิ}

\proseitem{106}{อตีตารมฺมณา ธมฺมา. อนาคตารมฺมณา ธมฺมา จตูหิ ขนฺเธหิ ทฺวีหายตเนหิ ทฺวีหิ ธาตูหิ สงฺคหิตา. กติหิ อสงฺคหิตา, เอเกน ขนฺเธน ทสหายตเนหิ โสฬสหิ ธาตูหิ อสงฺคหิตา.}

\proseitem{107}{ปจฺจุปฺปนฺนา\-รมฺมณา ธมฺมา. อชฺฌตฺตา\-รมฺมณา ธมฺมา. พหิทฺธา\-รมฺมณา ธมฺมา. อชฺฌตฺตพหิทฺธา\-รมฺมณา ธมฺมา จตูหิ ขนฺเธหิ ทฺวีหายตเนหิ อฏฺฐหิ ธาตูหิ สงฺคหิตา. กติหิ อสงฺคหิตา, เอเกน ขนฺเธน ทสหายตเนหิ ทสหิ ธาตูหิ อสงฺคหิตา.}

\proseitem{108}{สนิทสฺสน\-สปฺปฏิฆา ธมฺมา เอเกน ขนฺเธน เอเกนายตเนน เอกาย ธาตุยา สงฺคหิตา. กติหิ อสงฺคหิตา, จตูหิ ขนฺเธหิ เอกาทสหา\-ยตเนหิ สตฺตรสหิ ธาตูหิ อสงฺคหิตา.}

\proseitem{109}{อนิทสฺสน\-สปฺปฏิฆา ธมฺมา เอเกน ขนฺเธน นวหายตเนหิ นวหิ ธาตูหิ สงฺคหิตา. กติหิ อสงฺคหิตา, จตูหิ ขนฺเธหิ ตีหายตเนหิ นวหิ ธาตูหิ อสงฺคหิตา.}

\proseitem{110}{อนิทสฺสน\-อปฺปฏิฆา ธมฺมา อสงฺขตํ ขนฺธโต ฐเปตฺวา ปญฺจหิ ขนฺเธหิ ทฺวีหายตเนหิ อฏฺฐหิ ธาตูหิ สงฺคหิตา. กติหิ อสงฺคหิตา, น เกหิจิ ขนฺเธหิ ทสหายตเนหิ ทสหิ ธาตูหิ อสงฺคหิตา.}

\csromanmatikaanchor{mtk.16}
\csromantocmarkref{subsection}{8. ทุก}{mtk.16}
\bookmark[dest={mtk.16},level=2]{8. ทุก}
\csromanheader{2}{8. \textbf{ทุก}}
\proseitem{111}{\csromanfolio{19}เหตู ธมฺมา. เหตู เจว สเหตุกา จ ธมฺมา. เหตู เจวเหตุสมฺปยุตฺตา จ ธมฺมา เอเกน ขนฺเธน เอเกนายตเนน เอกาย ธาตุยา สงฺคหิตา. กติหิ อสงฺคหิตา, จตูหิ ขนฺเธหิ เอกาทสหา\-ยตเนหิ สตฺตรสหิ ธาตูหิ อสงฺคหิตา.}

\proseitem{112}{น เหตู ธมฺมา. อเหตุกา ธมฺมา. เหตุวิปฺปยุตฺตา ธมฺมา. น เหตุ อเหตุกา\csromansharedfootnote{31d60ec7a94e1b9a}{น เหตู อเหตุกา (สฺยา, ก) อภิ 2 ทุกปญฺหาปุจฺฉเกปิ.} ธมฺมา อสงฺขตํ ขนฺธโต ฐเปตฺวา ปญฺจหิ ขนฺเธหิ ทฺวาทสหา\-ยตเนหิ อฏฺฐารสหิ ธาตูหิ สงฺคหิตา. กติหิ อสงฺคหิตา, น เกหิจิ ขนฺเธหิ น เกหิจิ อายตเนหิ น กาหิจิ ธาตูหิ อสงฺคหิตา.}

\proseitem{113}{สเหตุกา ธมฺมา. เหตุสมฺปยุตฺตา ธมฺมา. สเหตุกา เจว น จ เหตู ธมฺมา. เหตุสมฺปยุตฺตา เจว น จ เหตู ธมฺมา. น เหตุ สเหตุกา ธมฺมา จตูหิ ขนฺเธหิ ทฺวีหายตเนหิ ทฺวีหิ ธาตูหิ สงฺคหิตา. กติหิ อสงฺคหิตา, เอเกน ขนฺเธน ทสหายตเนหิ โสฬสหิ ธาตูหิ อสงฺคหิตา.}

\proseitem{114}{สปฺปจฺจยา ธมฺมา. สงฺขตา ธมฺมา ปญฺจหิ ขนฺเธหิ ทฺวาทสหา\-ยตเนหิ อฏฺฐารสหิ ธาตูหิ สงฺคหิตา. กติหิ อสงฺคหิตา, น เกหิจิ ขนฺเธหิ น เกหิจิ อายตเนหิ น กาหิจิ ธาตูหิ อสงฺคหิตา.}

\proseitem{115}{อปฺปจฺจยา ธมฺมา. อสงฺขตา ธมฺมา น เกหิจิ ขนฺเธหิ เอเกนายตเนน เอกาย ธาตุยา สงฺคหิตา. กติหิ อสงฺคหิตา, ปญฺจหิ ขนฺเธหิ เอกาทสหา\-ยตเนหิ สตฺตรสหิ ธาตูหิ อสงฺคหิตา.}

\proseitem{116}{สนิทสฺสนา ธมฺมา เอเกน ขนฺเธน เอเกนายตเนน เอกาย ธาตุยา สงฺคหิตา. กติหิ อสงฺคหิตา, จตูหิ ขนฺเธหิ เอกาทสหา\-ยตเนหิ สตฺตรสหิ ธาตูหิ อสงฺคหิตา.}

\proseitem{117}{\csromanfolio{20}อนิทสฺสนา ธมฺมา อสงฺขตํ ขนฺธโต ฐเปตฺวา ปญฺจหิ ขนฺเธหิ เอกาทสหา\-ยตเนหิ สตฺตรสหิ ธาตูหิ สงฺคหิตา. กติหิ อสงฺคหิตา, น เกหิจิ ขนฺเธหิ เอเกนายตเนน เอกาย ธาตุยา อสงฺคหิตา.}

\proseitem{118}{สปฺปฏิฆา ธมฺมา เอเกน ขนฺเธน ทสหายตเนหิ ทสหิ ธาตูหิ สงฺคหิตา. กติหิ อสงฺคหิตา, จตูหิ ขนฺเธหิ ทฺวีหายตเนหิ อฏฺฐหิ ธาตูหิ อสงฺคหิตา.}

\proseitem{119}{อปฺปฏิฆา ธมฺมา อสงฺขตํ ขนฺธโต ฐเปตฺวา ปญฺจหิ ขนฺเธหิ ทฺวีหายตเนหิ อฏฺฐหิ ธาตูหิ สงฺคหิตา. กติหิ อสงฺคหิตา, น เกหิจิ ขนฺเธหิ ทสหายตเนหิ ทสหิ ธาตูหิ อสงฺคหิตา.}

\proseitem{120}{รูปิโน ธมฺมา เอเกน ขนฺเธน เอกาทสหา\-ยตเนหิ เอกาทสหิ ธาตูหิ สงฺคหิตา. กติหิ อสงฺคหิตา, จตูหิ ขนฺเธหิ เอเกนายตเนน สตฺตหิ ธาตูหิ อสงฺคหิตา.}

\proseitem{121}{อรูปิโน ธมฺมา อสงฺขตํ ขนฺธโต ฐเปตฺวา จตูหิ ขนฺเธหิ ทฺวีหายตเนหิ อฏฺฐหิ ธาตูหิ สงฺคหิตา. กติหิ อสงฺคหิตา เอเกน ขนฺเธน ทสหายตเนหิ ทสหิ ธาตูหิ อสงฺคหิตา.}

\proseitem{122}{โลกิยา ธมฺมา ปญฺจหิ ขนฺเธหิ ทฺวาทสหา\-ยตเนหิ อฏฺฐารสหิ ธาตูหิ สงฺคหิตา. กติหิ อสงฺคหิตา, น เกหิจิ ขนฺเธหิ น เกหิจิ อายตเนหิ น กาหิจิ ธาตูหิ อสงฺคหิตา.}

\proseitem{123}{โลกุตฺตรา ธมฺมา อสงฺขตํ ขนฺธโต ฐเปตฺวา จตูหิ ขนฺเธหิ ทฺวีหายตเนหิ ทฺวีหิ ธาตูหิ สงฺคหิตา. กติหิ อสงฺคหิตา เอเกน ขนฺเธน ทสหายตเนหิ โสฬสหิ ธาตูหิ อสงฺคหิตา.}

\proseitem{124}{เกนจิ วิญฺเญยฺยา ธมฺมา. เกนจิ น วิญฺเญยฺยา ธมฺมา อสงฺขตํ ขนฺธโต ฐเปตฺวา ปญฺจหิ ขนฺเธหิ ทฺวาทสหา\-ยตเนหิ อฏฺฐารสหิ ธาตูหิ สงฺคหิตา. กติหิ อสงฺคหิตา, น เกหิจิ ขนฺเธหิ น เกหิจิ อายตเนหิ น กาหิจิ ธาตูหิ อสงฺคหิตา.}

\proseitem{125}{\csromanfolio{21}อาสวา ธมฺมา. อาสวา เจว สาสวา จ ธมฺมา. อาสวา เจว อาสว\-สมฺปยุตฺตา จ ธมฺมา เอเกน ขนฺเธน เอเกนายตเนน เอกาย ธาตุยา สงฺคหิตา. กติหิ อสงฺคหิตา, จตูหิ ขนฺเธหิ เอกาทสหา\-ยตเนหิ สตฺตรสหิ ธาตูหิ อสงฺคหิตา.}

\proseitem{126}{โน อาสวา ธมฺมา. อาสว\-วิปฺปยุตฺตา ธมฺมา อสงฺขตํ ขนฺธโต ฐเปตฺวา ปญฺจหิ ขนฺเธหิ ทฺวาทสหา\-ยตเนหิ อฏฺฐารสหิ ธาตูหิ สงฺคหิตา. กติหิ อสงฺคหิตา, น เกหิจิ ขนฺเธหิ น เกหิจิ อายตเนหิ น กาหิจิ ธาตูหิ อสงฺคหิตา.}

\proseitem{127}{สาสวา ธมฺมา. สาสวา เจว โน จ อาสวา ธมฺมา. อาสว\-วิปฺปยุตฺตา สาสวา ธมฺมา ปญฺจหิ ขนฺเธหิ ทฺวาทสหา\-ยตเนหิ อฏฺฐารสหิ ธาตูหิ สงฺคหิตา. กติหิ อสงฺคหิตา, น เกหิจิ ขนฺเธหิ น เกหิจิ อายตเนหิ น กาหิจิ ธาตูหิ อสงฺคหิตา.}

\proseitem{128}{อนาสวา ธมฺมา. อาสว\-วิปฺปยุตฺตา อนาสวา ธมฺมา อสงฺขตํ ขนฺธโต ฐเปตฺวา จตูหิ ขนฺเธหิ ทฺวีหายตเนหิ ทฺวีหิ ธาตูหิ สงฺคหิตา. กติหิ อสงฺคหิตา, เอเกน ขนฺเธน ทสหายตเนหิ โสฬสหิ ธาตูหิ}

\proseitem{129}{อาสว\-สมฺปยุตฺตา ธมฺมา. อาสว\-สมฺปยุตฺตา เจว โน จ อาสวา ธมฺมา จตูหิ ขนฺเธหิ ทฺวีหายตเนหิ ทฺวีหิ ธาตูหิ สงฺคหิตา. กติหิ อสงฺคหิตา, เอเกน ขนฺเธน ทสหายตเนหิ โสฬสหิ ธาตูหิ อสงฺคหิตา.}

\proseitem{130}{สํโยชนา ธมฺมา. คนฺถา ธมฺมา. โอฆา ธมฺมา. โยคา ธมฺมา. นีวรณา ธมฺมา. ปรามาสา ธมฺมา. ปรามาสา เจว ปรามฏฺฐา จ ธมฺมา เอเกน ขนฺเธน เอเกนายตเนน เอกาย ธาตุยา สงฺคหิตา. กติหิ อสงฺคหิตา, จตูหิ ขนฺเธหิ เอกาทสหา\-ยตเนหิ สตฺตรสหิ ธาตูหิ}

\proseitem{131}{โน ปรามาสา ธมฺมา. ปรามาส\-วิปฺปยุตฺตา ธมฺมา อสงฺขตํ ขนฺธโต ฐเปตฺวา ปญฺจหิ ขนฺเธหิ ทฺวาทสหา\-ยตเนหิ อฏฺฐารสหิ ธาตูหิ \csromanfolio{22}สงฺคหิตา. กติหิ อสงฺคหิตา, น เกหิจิ ขนฺเธหิ น เกหิจิ อายตเนหิ น กาหิจิ ธาตูหิ อสงฺคหิตา.}

\proseitem{132}{ปรามฏฺฐา ธมฺมา. ปรามฏฺฐา เจว โน จ ปรามาสา ธมฺมา. ปรามาส\-วิปฺปยุตฺตา ปรามฏฺฐา ธมฺมา ปญฺจหิ ขนฺเธหิ ทฺวาทสหา\-ยตเนหิ อฏฺฐารสหิ ธาตูหิ สงฺคหิตา. กติหิ อสงฺคหิตา, น เกหิจิ ขนฺเธหิ น เกหิจิ อายตเนหิ น กาหิจิ ธาตูหิ อสงฺคหิตา.}

\proseitem{133}{อปรามฏฺฐา ธมฺมา. ปรามาส\-วิปฺปยุตฺตา อปรามฏฺฐา ธมฺมา อสงฺขตํ ขนฺธโต ฐเปตฺวา จตูหิ ขนฺเธหิ ทฺวีหายตเนหิ ทฺวีหิ ธาตูหิ สงฺคหิตา. กติหิ อสงฺคหิตา, เอเกน ขนฺเธน ทสหายตเนหิ}

\proseitem{134}{ปรามาส\-สมฺปยุตฺตา ธมฺมา จตูหิ ขนฺเธหิ ทฺวีหายตเนหิ ทฺวีหิ ธาตูหิ สงฺคหิตา. กติหิ อสงฺคหิตา, เอเกน ขนฺเธน ทสหายตเนหิ}

\proseitem{135}{สารมฺมณา ธมฺมา จตูหิ ขนฺเธหิ ทฺวีหายตเนหิ อฏฺฐหิ ธาตูหิ สงฺคหิตา. กติหิ อสงฺคหิตา, เอเกน ขนฺเธน ทสหายตเนหิ ทสหิ ธาตูหิ อสงฺคหิตา.}

\proseitem{136}{อนารมฺมณา ธมฺมา อสงฺขตํ ขนฺธโต ฐเปตฺวา เอเกน ขนฺเธน เอกาทสหา\-ยตเนหิ เอกาทสหิ ธาตูหิ สงฺคหิตา. กติหิ อสงฺคหิตา, จตูหิ ขนฺเธหิ เอเกนายตเนน สตฺตหิ ธาตูหิ อสงฺคหิตา.}

\proseitem{137}{จิตฺตา ธมฺมา เอเกน ขนฺเธน เอเกนายตเนน สตฺตหิ ธาตูหิ สงฺคหิตา. กติหิ อสงฺคหิตา, จตูหิ ขนฺเธหิ เอกาทสหา\-ยตเนหิ เอกาทสหิ ธาตูหิ อสงฺคหิตา.}

\proseitem{138}{โน จิตฺตา ธมฺมา อสงฺขตํ ขนฺธโต ฐเปตฺวา จตูหิ ขนฺเธหิ เอกาทสหา\-ยตเนหิ เอกาทสหิ ธาตูหิ สงฺคหิตา. กติหิ อสงฺคหิตา, เอเกน ขนฺเธน เอเกนายตเนน สตฺตหิ ธาตูหิ อสงฺคหิตา.}

\proseitem{139}{\csromanfolio{23}เจตสิกา ธมฺมา. จิตฺต\-สมฺปยุตฺตา ธมฺมา. จิตฺตสํสฏฺฐา ธมฺมา ตีหิ ขนฺเธหิ เอเกนายตเนน เอกาย ธาตุยา สงฺคหิตา. กติหิ อสงฺคหิตา, ทฺวีหิ ขนฺเธหิ เอกาทสหา\-ยตเนหิ สตฺตรสหิ ธาตูหิ}

\proseitem{140}{อเจตสิกา ธมฺมา อสงฺขตํ ขนฺธโต ฐเปตฺวา ทฺวีหิ ขนฺเธหิ ทฺวาทสหา\-ยตเนหิ อฏฺฐารสหิ ธาตูหิ สงฺคหิตา. กติหิ อสงฺคหิตา, ตีหิ ขนฺเธหิ น เกหิจิ อายตเนหิ น กาหิจิ ธาตูหิ อสงฺคหิตา.}

\proseitem{141}{จิตฺต\-วิปฺปยุตฺตา ธมฺมา. จิตฺต\-วิสํสฏฺฐา ธมฺมา อสงฺขตํ ขนฺธโต ฐเปตฺวา เอเกน ขนฺเธน เอกาทสหา\-ยตเนหิ เอกาทสหิ ธาตูหิ สงฺคหิตา. กติหิ อสงฺคหิตา, จตูหิ ขนฺเธหิ เอเกนายตเนน สตฺตหิ ธาตูหิ อสงฺคหิตา.}

\proseitem{142}{จิตฺต\-สมุฏฺฐานา ธมฺมา จตูหิ ขนฺเธหิ ฉหายตเนหิ ฉหิ ธาตูหิ สงฺคหิตา. กติหิ อสงฺคหิตา, เอเกน ขนฺเธน ฉหายตเนหิ ทฺวาทสหิ ธาตูหิ อสงฺคหิตา.}

\proseitem{143}{โน จิตฺต\-สมุฏฺฐานา ธมฺมา. โน จิตฺตสหภุโน ธมฺมา. โน จิตฺตา\-นุปริวตฺติโน ธมฺมา อสงฺขตํ ขนฺธโต ฐเปตฺวา ทฺวีหิ ขนฺเธหิ ทฺวาทสหา\-ยตเนหิ อฏฺฐารสหิ ธาตูหิ สงฺคหิตา. กติหิ อสงฺคหิตา, ตีหิ ขนฺเธหิ น เกหิจิ อายตเนหิ น กาหิจิ ธาตูหิ อสงฺคหิตา.}

\proseitem{144}{จิตฺตสหภุโน ธมฺมา. จิตฺตา\-นุปริวตฺติโน ธมฺมา จตูหิ ขนฺเธหิ เอเกนายตเนน เอกาย ธาตุยา สงฺคหิตา. กติหิ อสงฺคหิตา, เอเกน ขนฺเธน เอกาทสหา\-ยตเนหิ สตฺตรสหิ ธาตูหิ อสงฺคหิตา.}

\proseitem{145}{จิตฺต\-สํสฏฺฐ\-สมุฏฺฐานา ธมฺมา. จิตฺต\-สํสฏฺฐ\-สมุฏฺฐาน\-สหภุโน ธมฺมา. จิตฺต\-สํสฏฺฐ\-สมุฏฺฐานา\-นุปริวตฺติโน ธมฺมา ตีหิ ขนฺเธหิ เอเกนายตเนน เอกาย ธาตุยา สงฺคหิตา. กติหิ อสงฺคหิตา, ทฺวีหิ ขนฺเธหิ เอกาทสหา\-ยตเนหิ สตฺตรสหิ ธาตูหิ อสงฺคหิตา.}

\proseitem{146}{\csromanfolio{24}โน จิตฺต\-สํสฏฺฐ\-สมุฏฺฐานา ธมฺมา. โน จิตฺต\-สํสฏฺฐ\-สมุฏฺฐาน\-สหภุโน ธมฺมา. โน จิตฺต\-สํสฏฺฐ\-สมุฏฺฐานา\-นุปริวตฺติโน ธมฺมา อสงฺขตํ ขนฺธโต ฐเปตฺวา ทฺวีหิ ขนฺเธหิ ทฺวาทสหา\-ยตเนหิ อฏฺฐารสหิ ธาตูหิ สงฺคหิตา. กติหิ อสงฺคหิตา, ตีหิ ขนฺเธหิ น เกหิจิ อายตเนหิ น กาหิจิ ธาตูหิ}

\proseitem{147}{อชฺฌตฺติกา ธมฺมา ทฺวีหิ ขนฺเธหิ ฉหายตเนหิ ทฺวาทสหิ ธาตูหิ สงฺคหิตา. กติหิ อสงฺคหิตา, ตีหิ ขนฺเธหิ ฉหายตเนหิ ฉหิ ธาตูหิ อสงฺคหิตา.}

\proseitem{148}{พาหิรา ธมฺมา อสงฺขตํ ขนฺธโต ฐเปตฺวา จตูหิ ขนฺเธหิ ฉหายตเนหิ ฉหิ ธาตูหิ สงฺคหิตา. กติหิ อสงฺคหิตา, เอเกน ขนฺเธน ฉหายตเนหิ ทฺวาทสหิ ธาตูหิ อสงฺคหิตา.}

\proseitem{149}{อุปาทา ธมฺมา เอเกน ขนฺเธน ทสหายตเนหิ ทสหิ ธาตูหิ สงฺคหิตา. กติหิ อสงฺคหิตา, จตูหิ ขนฺเธหิ ทฺวีหายตเนหิ อฏฺฐหิ ธาตูหิ}

\proseitem{150}{โน อุปาทา ธมฺมา อสงฺขตํ ขนฺธโต ฐเปตฺวา ปญฺจหิ ขนฺเธหิ ตีหายตเนหิ นวหิ ธาตูหิ สงฺคหิตา. กติหิ อสงฺคหิตา, น เกหิจิ ขนฺเธหิ นวหายตเนหิ นวหิ ธาตูหิ อสงฺคหิตา.}

\proseitem{151}{อุปาทินฺนา ธมฺมา ปญฺจหิ ขนฺเธหิ เอกาทสหา\-ยตเนหิ สตฺตรสหิ ธาตูหิ สงฺคหิตา. กติหิ อสงฺคหิตา, น เกหิจิ ขนฺเธหิ เอเกนายตเนน เอกาย ธาตุยา อสงฺคหิตา.}

\proseitem{152}{อนุปาทินฺนา ธมฺมา อสงฺขตํ ขนฺธโต ฐเปตฺวา ปญฺจหิ ขนฺเธหิ สตฺตหา\-ยตเนหิ อฏฺฐหิ ธาตูหิ สงฺคหิตา. กติหิ อสงฺคหิตา, น เกหิจิ ขนฺเธหิ ปญฺจหา\-ยตเนหิ ทสหิ ธาตูหิ อสงฺคหิตา.}

\proseitem{153}{อุปาทานา ธมฺมา. กิเลสา ธมฺมา. กิเลสา เจว สํกิเลสิกา จ ธมฺมา. กิเลสา เจว สํกิลิฏฺฐา จ ธมฺมา. กิเลสา เจว กิเลส\-สมฺปยุตฺตา จ ธมฺมา เอเกน ขนฺเธน เอเกนายตเนน เอกาย \csromanfolio{25}ธาตุยา สงฺคหิตา. กติหิ อสงฺคหิตา, จตูหิ ขนฺเธหิ เอกาทสหา\-ยตเนหิ สตฺตรสหิ ธาตูหิ อสงฺคหิตา.}

\proseitem{154}{โน กิเลสา ธมฺมา. อสํกิลิฏฺฐา ธมฺมา. กิเลส\-วิปฺปยุตฺตา ธมฺมา อสงฺขตํ ขนฺธโต ฐเปตฺวา ปญฺจหิ ขนฺเธหิ ทฺวาทสหา\-ยตเนหิ อฏฺฐารสหิ ธาตูหิ สงฺคหิตา. กติหิ อสงฺคหิตา, น เกหิจิ ขนฺเธหิ น เกหิจิ อายตเนหิ น กาหิจิ ธาตูหิ อสงฺคหิตา.}

\proseitem{155}{สํกิเลสิกา ธมฺมา. สํกิเลสิกา เจว โน จ กิเลสา ธมฺมา. กิเลส\-วิปฺปยุตฺตา สํกิเลสิกา ธมฺมา ปญฺจหิ ขนฺเธหิ ทฺวาทสหา\-ยตเนหิ อฏฺฐารสหิ ธาตูหิ สงฺคหิตา. กติหิ อสงฺคหิตา, น เกหิจิ ขนฺเธหิ น เกหิจิ อายตเนหิ น กาหิจิ ธาตูหิ อสงฺคหิตา.}

\proseitem{156}{อสํกิเลสิกา ธมฺมา. กิเลส\-วิปฺปยุตฺตา อสํกิเลสิกา ธมฺมา อสงฺขตํ ขนฺธโต ฐเปตฺวา จตูหิ ขนฺเธหิ ทฺวีหายตเนหิ ทฺวีหิ ธาตูหิ สงฺคหิตา. กติหิ อสงฺคหิตา, เอเกน ขนฺเธน ทสหายตเนหิ}

\proseitem{157}{สํกิลิฏฺฐา ธมฺมา. กิเลส\-สมฺปยุตฺตา ธมฺมา. สํกิลิฏฺฐา เจว โน จ กิเลสา ธมฺมา. กิเลส\-สมฺปยุตฺตา เจว โน จ กิเลสา ธมฺมา จตูหิ ขนฺเธหิ ทฺวีหายตเนหิ ทฺวีหิ ธาตูหิ สงฺคหิตา. กติหิ อสงฺคหิตา, เอเกน ขนฺเธน ทสหายตเนหิ โสฬสหิ ธาตูหิ อสงฺคหิตา.}

\proseitem{158}{ทสฺสเนน ปหาตพฺพา ธมฺมา. ภาวนาย ปหาตพฺพา ธมฺมา. ทสฺสเนน ปหาตพฺพ\-เหตุกา ธมฺมา. ภาวนาย ปหาตพฺพ\-เหตุกา ธมฺมา จตูหิ ขนฺเธหิ ทฺวีหายตเนหิ ทฺวีหิ ธาตูหิ สงฺคหิตา. กติหิ อสงฺคหิตา, เอเกน ขนฺเธน ทสหายตเนหิ โสฬสหิ ธาตูหิ อสงฺคหิตา.}

\proseitem{159}{น ทสฺสเนน ปหาตพฺพา ธมฺมา. น ภาวนาย ปหาตพฺพา ธมฺมา. น ทสฺสเนน ปหาตพฺพ\-เหตุกา ธมฺมา. น ภาวนาย ปหาตพฺพ\-เหตุกา ธมฺมา อสงฺขตํ ขนฺธโต ฐเปตฺวา ปญฺจหิ ขนฺเธหิ ทฺวาทสหา\-ยตเนหิ \csromanfolio{26}อฏฺฐารสหิ ธาตูหิ สงฺคหิตา. กติหิ อสงฺคหิตา, น เกหิจิ ขนฺเธหิ น เกหิจิ อายตเนหิ น กาหิจิ ธาตูหิ อสงฺคหิตา.}

\proseitem{160}{สวิตกฺกา ธมฺมา. สวิจารา ธมฺมา จตูหิ ขนฺเธหิ ทฺวีหายตเนหิ ตีหิ ธาตูหิ สงฺคหิตา. กติหิ อสงฺคหิตา, เอเกน ขนฺเธน ทสหายตเนหิ ปนฺนรสหิ ธาตูหิ อสงฺคหิตา.}

\proseitem{161}{อวิตกฺกา ธมฺมา. อวิจารา ธมฺมา อสงฺขตํ ขนฺธโต ฐเปตฺวา ปญฺจหิ ขนฺเธหิ ทฺวาทสหา\-ยตเนหิ สตฺตรสหิ ธาตูหิ สงฺคหิตา. กติหิ อสงฺคหิตา, น เกหิจิ ขนฺเธหิ น เกหิจิ อายตเนหิ เอกาย ธาตุยา อสงฺคหิตา.}

\proseitem{162}{สปฺปีติกา ธมฺมา. ปีติสหคตา ธมฺมา จตูหิ ขนฺเธหิ ทฺวีหายตเนหิ ทฺวีหิ ธาตูหิ สงฺคหิตา. กติหิ อสงฺคหิตา, เอเกน ขนฺเธน ทสหายตเนหิ โสฬสหิ ธาตูหิ อสงฺคหิตา.}

\proseitem{163}{อปฺปีติกา ธมฺมา. น ปีติสหคตา ธมฺมา. น สุขสหคตา ธมฺมา อสงฺขตํ ขนฺธโต ฐเปตฺวา ปญฺจหิ ขนฺเธหิ ทฺวาทสหา\-ยตเนหิ อฏฺฐารสหิ ธาตูหิ สงฺคหิตา. กติหิ อสงฺคหิตา, น เกหิจิ ขนฺเธหิ น เกหิจิ อายตเนหิ น กาหิจิ ธาตูหิ อสงฺคหิตา.}

\proseitem{164}{สุขสหคตา ธมฺมา ตีหิ ขนฺเธหิ ทฺวีหายตเนหิ ตีหิ ธาตูหิ สงฺคหิตา. กติหิ อสงฺคหิตา, ทฺวีหิ ขนฺเธหิ ทสหายตเนหิ ปนฺนรสหิ ธาตูหิ อสงฺคหิตา.}

\proseitem{165}{อุเปกฺขา\-สหคตา ธมฺมา ตีหิ ขนฺเธหิ ทฺวีหายตเนหิ สตฺตหิ ธาตูหิ สงฺคหิตา. กติหิ อสงฺคหิตา, ทฺวีหิ ขนฺเธหิ ทสหายตเนหิ เอกาทสหิ ธาตูหิ อสงฺคหิตา.}

\proseitem{166}{น อุเปกฺขา\-สหคตา ธมฺมา อสงฺขตํ ขนฺธโต ฐเปตฺวา ปญฺจหิ ขนฺเธหิ ทฺวาทสหา\-ยตเนหิ เตรสหิ ธาตูหิ สงฺคหิตา. กติหิ อสงฺคหิตา, น เกหิจิ ขนฺเธหิ น เกหิจิ อายตเนหิ ปญฺจหิ ธาตูหิ}

\proseitem{167}{\csromanfolio{27}กามาวจรา ธมฺมา. ปริยาปนฺนา ธมฺมา. สอุตฺตรา ธมฺมา ปญฺจหิ ขนฺเธหิ ทฺวาทสหา\-ยตเนหิ อฏฺฐารสหิ ธาตูหิ สงฺคหิตา. กติหิ อสงฺคหิตา, น เกหิจิ ขนฺเธหิ น เกหิจิ อายตเนหิ น กาหิจิ ธาตูหิ}

\proseitem{168}{น กามาวจรา ธมฺมา. อปริยาปนฺนา ธมฺมา. อนุตฺตรา ธมฺมา อสงฺขตํ ขนฺธโต ฐเปตฺวา จตูหิ ขนฺเธหิ ทฺวีหายตเนหิ ทฺวีหิ ธาตูหิ สงฺคหิตา. กติหิ อสงฺคหิตา, เอเกน ขนฺเธน ทสหายตเนหิ}

\proseitem{169}{รูปาวจรา ธมฺมา. อรูปาวจรา ธมฺมา. นิยฺยานิกา ธมฺมา. นิยตา ธมฺมา. สรณา ธมฺมา จตูหิ ขนฺเธหิ ทฺวีหายตเนหิ ทฺวีหิ ธาตูหิ สงฺคหิตา. กติหิ อสงฺคหิตา, เอเกน ขนฺเธน ทสหายตเนหิ}

\proseitem{170}{น รูปาวจรา ธมฺมา. น อรูปาวจรา ธมฺมา. อนิยฺยานิกา ธมฺมา. อนิยตา ธมฺมา. อรณา ธมฺมา กติหิ ขนฺเธหิ กติหายตเนหิ กติหิ ธาตูหิ สงฺคหิตา. อรณา ธมฺมา อสงฺขตํ ขนฺธโต ฐเปตฺวา ปญฺจหิ ขนฺเธหิ ทฺวาทสหา\-ยตเนหิ อฏฺฐารสหิ ธาตูหิ สงฺคหิตา. กติหิ อสงฺคหิตา, น เกหิจิ ขนฺเธหิ น เกหิจิ อายตเนหิ น กาหิจิ ธาตูหิ}

\vspace{\tipitakaparskip}
\csromancenter{สงฺคหา\-สงฺคห\-ปท\-นิทฺเทโส ปฐโม.}
\csromanmatikaanchor{mtk.17}
\csromantocmarknopage{section}{2. ทุติยนย}{mtk.17}
\bookmark[dest={mtk.17},level=1]{2. ทุติยนย}
\csromanheader{1}{2. \textbf{ทุติยนย}}
\csromanmatikaanchor{mtk.18}
\csromantocmarkref{subsection}{2. สงฺคหิเตนอสงฺคหิตปทนิทฺเทส}{mtk.18}
\bookmark[dest={mtk.18},level=2]{2. สงฺคหิเตนอสงฺคหิตปทนิทฺเทส}
\csromanheader{2}{2. สงฺคหิเตน\-อสงฺคหิต\-ปท\-นิทฺเทส}
\proseitem{171}{\csromanfolio{28}จกฺขายตเนน เย ธมฺมา ฯเปฯ. โผฏฺฐพฺพา\-ยตเนน เย ธมฺมา. จกฺขุธาตุยา เย ธมฺมา ฯเปฯ. โผฏฺฐพฺพ\-ธาตุยา เย ธมฺมา ขนฺธ\-สงฺคเหน สงฺคหิตา อายตน\-สงฺคเหน อสงฺคหิตา ธาตุสงฺคเหน อสงฺคหิตา. เต ธมฺมา กติหิ ขนฺเธหิ กติหายตเนหิ กติหิ ธาตูหิ อสงฺคหิตา, เต ธมฺมา จตูหิ ขนฺเธหิ ทฺวีหายตเนหิ อฏฺฐหิ ธาตูหิ}

\proseitem{172}{จกฺขุวิญฺญาณ\-ธาตุยา เย ธมฺมา. โสตวิญฺญาณ\-ธาตุยา เย ธมฺมา. ฆานวิญฺญาณ\-ธาตุยา เย ธมฺมา. ชิวฺหาวิญฺญาณ\-ธาตุยา เย ธมฺมา. กายวิญฺญาณ\-ธาตุยา เย ธมฺมา. มโนธาตุยา เย ธมฺมา. มโนวิญฺญาณ\-ธาตุยา เย ธมฺมา ขนฺธ\-สงฺคเหน สงฺคหิตา อายตน\-สงฺคเหน อสงฺคหิตา ธาตุสงฺคเหน อสงฺคหิตา ฯเปฯ เต ธมฺมา จตูหิ ขนฺเธหิ เอกาทสหา\-ยตเนหิ ทฺวาทสหิ ธาตูหิ อสงฺคหิตา.}

\proseitem{173}{จกฺขุนฺทฺริเยน เย ธมฺมา. โสตินฺทฺริเยน เย ธมฺมา. ฆานินฺทฺริเยน เย ธมฺมา. ชิวฺหินฺทฺริเยน เย ธมฺมา. กายินฺทฺริเยน เย ธมฺมา. อิตฺถินฺทฺริเยน เย ธมฺมา. ปุริสินฺทฺริเยน เย ธมฺมา ขนฺธ\-สงฺคเหน สงฺคหิตา อายตน\-สงฺคเหน อสงฺคหิตา ธาตุสงฺคเหน อสงฺคหิตา ฯเปฯ เต ธมฺมา จตูหิ ขนฺเธหิ ทฺวีหายตเนหิ อฏฺฐหิ ธาตูหิ}

\proseitem{174}{อสญฺญาภเวน เย ธมฺมา. เอกโวการ\-ภเวน เย ธมฺมา ขนฺธ\-สงฺคเหน สงฺคหิตา อายตน\-สงฺคเหน อสงฺคหิตา ธาตุสงฺคเหน อสงฺคหิตา ฯเปฯ เต ธมฺมา จตูหิ ขนฺเธหิ ตีหายตเนหิ นวหิ ธาตูหิ}

\proseitem{175}{ปริเทเวน เย ธมฺมา. สนิทสฺสน\-สปฺปฏิเฆหิ ธมฺเมหิ เย ธมฺมา ขนฺธ\-สงฺคเหน สงฺคหิตา อายตน\-สงฺคเหน อสงฺคหิตา ธาตุสงฺคเหน อสงฺคหิตา ฯเปฯ เต ธมฺมา จตูหิ ขนฺเธหิ ทฺวีหายตเนหิ อฏฺฐหิ ธาตูหิ อสงฺคหิตา.}

\csromanheader{2}{2. ทุติยนย}
\proseitem{176}{\csromanfolio{29}อนิทสฺสน\-สปฺปฏิเฆหิ ธมฺเมหิ เย ธมฺมา ขนฺธ\-สงฺคเหน สงฺคหิตา อายตน\-สงฺคเหน อสงฺคหิตา ธาตุสงฺคเหน อสงฺคหิตา ฯเปฯ เต ธมฺมา จตูหิ ขนฺเธหิ ทสหายตเนหิ โสฬสหิ ธาตูหิ อสงฺคหิตา.}

\proseitem{177}{สนิทสฺสเนหิ ธมฺเมหิ เย ธมฺมา ขนฺธ\-สงฺคเหน สงฺคหิตา อายตน\-สงฺคเหน อสงฺคหิตา ธาตุสงฺคเหน อสงฺคหิตา ฯเปฯ เต ธมฺมา จตูหิ ขนฺเธหิ ทฺวีหายตเนหิ อฏฺฐหิ ธาตูหิ อสงฺคหิตา.}

\proseitem{178}{สปฺปฏิเฆหิ ธมฺเมหิ เย ธมฺมา. อุปาทาธมฺเมหิ เย ธมฺมา ขนฺธ\-สงฺคเหน สงฺคหิตา อายตน\-สงฺคเหน อสงฺคหิตา ธาตุสงฺคเหน อสงฺคหิตา. เต ธมฺมา กติหิ ขนฺเธหิ กติหายตเนหิ กติหิ ธาตูหิ อสงฺคหิตา. เต ธมฺมา จตูหิ ขนฺเธหิ เอกาทสหา\-ยตเนหิ สตฺตรสหิ ธาตูหิ อสงฺคหิตา.}

\setlength{\csromangathapagewidth}{0pt}
\setlength{\csromangathawakwidth}{0pt}
\csromangathameasuregroup{}{ทสายตนา สตฺตรส ธาตุโย, \\ สตฺตินฺทฺริยา อสญฺญาภโว เอกโวการภโว. \\ ปริเทโว สนิทสฺสนสปฺปฏิฆํ, \\ อนิทสฺสนํ ปุนเทว\textsuperscript{1} สปฺปฏิฆํ อุปาทาติ. \\ \textbf{สงฺคหิเตน อสงฺคหิตปทนิทฺเทโส ทุติโย.}}
\csromangathameasurewak{ทสายตนา สตฺตรส ธาตุโย, \\ สตฺตินฺทฺริยา อสญฺญาภโว เอกโวการภโว. \\ ปริเทโว สนิทสฺสนสปฺปฏิฆํ, \\ อนิทสฺสนํ ปุนเทว\textsuperscript{1} สปฺปฏิฆํ อุปาทาติ.}
\setlength{\csromangathaleftcol}{0pt}
\csromangathagroup{\csromangathastanza{\csromangathawak{ทสายตนา สตฺตรส ธาตุโย, \\ สตฺตินฺทฺริยา อสญฺญาภโว เอกโวการภโว. \\ ปริเทโว สนิทสฺสน\-สปฺปฏิฆํ, \\ อนิทสฺสนํ ปุนเทว\csromansharedfootnote{be8c753a9828e054}{ปุนเรว (อิ)} สปฺปฏิฆํ อุปาทาติ.}}\csromangathastanza{\textbf{สงฺคหิเตน อสงฺคหิตปทนิทฺเทโส ทุติโย.}}}

\csromanmatikaanchor{mtk.19}
\csromantocmarknopage{section}{3. ตติยนย}{mtk.19}
\bookmark[dest={mtk.19},level=1]{3. ตติยนย}
\csromanheader{1}{3. \textbf{ตติยนย}}
\csromanmatikaanchor{mtk.20}
\csromantocmarkref{subsection}{3. อสงฺคหิเตนสงฺคหิตปทนิทฺเทส}{mtk.20}
\bookmark[dest={mtk.20},level=2]{3. อสงฺคหิเตนสงฺคหิตปทนิทฺเทส}
\csromanheader{2}{3. อสงฺคหิเตน\-สงฺคหิต\-ปท\-นิทฺเทส}
\proseitem{179}{\csromanfolio{30}เวทนา\-กฺขนฺเธน เย ธมฺมา. สญฺญา\-กฺขนฺเธน เย ธมฺมา. สงฺขาร\-กฺขนฺเธน เย ธมฺมา. สมุทย\-สจฺเจน เย ธมฺมา. มคฺคสจฺเจน เย ธมฺมา ขนฺธ\-สงฺคเหน อสงฺคหิตา อายตน\-สงฺคเหน สงฺคหิตา ธาตุสงฺคเหน สงฺคหิตา. เต ธมฺมา กติหิ ขนฺเธหิ กติหายตเนหิ กติหิ ธาตูหิ สงฺคหิตา, เต ธมฺมา อสงฺขตํ ขนฺธโต ฐเปตฺวา ตีหิ ขนฺเธหิ เอเกนายตเนน เอกาย ธาตุยา สงฺคหิตา.}

\proseitem{180}{นิโรธสจฺเจน เย ธมฺมา ขนฺธ\-สงฺคเหน อสงฺคหิตา อายตน\-สงฺคเหน สงฺคหิตา ธาตุสงฺคเหน สงฺคหิตา ฯเปฯ เต ธมฺมา จตูหิ ขนฺเธหิ เอเกนายตเนน เอกาย ธาตุยา สงฺคหิตา.}

\proseitem{181}{ชีวิตินฺทฺริเยน เย ธมฺมา ขนฺธ\-สงฺคเหน อสงฺคหิตา อายตน\-สงฺคเหน สงฺคหิตา ธาตุสงฺคเหน สงฺคหิตา ฯเปฯ เต ธมฺมา อสงฺขตํ ขนฺธโต ฐเปตฺวา ทฺวีหิ ขนฺเธหิ เอเกนายตเนน เอกาย ธาตุยา สงฺคหิตา.}

\proseitem{182}{อิตฺถินฺทฺริเยน เย ธมฺมา. ปุริสินฺทฺริเยน เย ธมฺมา. สุขินฺทฺริเยน เย ธมฺมา. ทุกฺขินฺทฺริเยน เย ธมฺมา. โสมนสฺสิ\-นฺทฺริเยน เย ธมฺมา. โทมนสฺสิ\-นฺทฺริเยน เย ธมฺมา. อุเปกฺขิ\-นฺทฺริเยน เย ธมฺมา. สทฺธินฺทฺริเยน เย ธมฺมา. วีริยินฺทฺริเยน เย ธมฺมา. สตินฺทฺริเยน เย ธมฺมา. สมาธิ\-นฺทฺริเยน เย ธมฺมา. ปญฺญินฺทฺริเยน เย ธมฺมา. อนญฺญาต\-ญฺญสฺสามี\-ตินฺทฺริเยน เย ธมฺมา. อญฺญินฺทฺริเยน เย ธมฺมา. อญฺญาตาวิ\-นฺทฺริเยน เย ธมฺมา. อวิชฺชาย เย ธมฺมา. อวิชฺชาปจฺจยา สงฺขาเรน เย ธมฺมา. สฬายตน\-ปจฺจยา ผสฺเสน เย ธมฺมา. ผสฺสปจฺจยา เวทนาย เย ธมฺมา. เวทนาปจฺจยา ตณฺหาย เย ธมฺมา. ตณฺหาปจฺจยา อุปาทาเนน เย ธมฺมา. กมฺมภเวน\csromansharedfootnote{f976a084b08c7509}{อุปาทานปจฺจยา กมฺมภเวน (สฺยา)} เย ธมฺมา ขนฺธ\-สงฺคเหน อสงฺคหิตา อายตน\-สงฺคเหน สงฺคหิตา ธาตุสงฺคเหน สงฺคหิตา ฯเปฯ เต ธมฺมา อสงฺขตํ ขนฺธโต ฐเปตฺวา ตีหิ ขนฺเธหิ เอเกนายตเนน เอกาย ธาตุยา สงฺคหิตา.}

\csromanheader{2}{3. ตติยนย}
\proseitem{183}{\csromanfolio{31}ชาติยา เย ธมฺมา. ชราย เย ธมฺมา. มรเณน เย ธมฺมา. ฌาเนน เย ธมฺมา ขนฺธ\-สงฺคเหน อสงฺคหิตา อายตน\-สงฺคเหน สงฺคหิตา ธาตุสงฺคเหน สงฺคหิตา ฯเปฯ เต ธมฺมา อสงฺขตํ ขนฺธโต ฐเปตฺวา ทฺวีหิ ขนฺเทหิ เอเกนายตเนน เอกาย ธาตุยา สงฺคหิตา.}

\proseitem{184}{โสเกน เย ธมฺมา. ทุกฺเขน เย ธมฺมา. โทมนสฺเสน เย ธมฺมา. อุปายาเสน เย ธมฺมา. สติปฏฺฐาเนน เย ธมฺมา. สมฺม\-ปฺปธาเนน เย ธมฺมา. อปฺปมญฺญาย เย ธมฺมา. ปญฺจหิ อินฺทฺริเยหิ เย ธมฺมา. ปญฺจหิ พเลหิ เย ธมฺมา. สตฺตหิ โพชฺฌงฺเคหิ เย ธมฺมา. อริเยน อฏฺฐงฺคิเกน มคฺเคน เย ธมฺมา. ผสฺเสน เย ธมฺมา. เวทนาย เย ธมฺมา. สญฺญาย เย ธมฺมา. เจตนาย เย ธมฺมา. อธิโมกฺเขน เย ธมฺมา. มนสิกาเรน เย ธมฺมา. เหตูหิ ธมฺเมหิ เย ธมฺมา. เหตูหิ เจว สเหตุเกหิ จ ธมฺเมหิ เย ธมฺมา. เหตูหิ เจว เหตุ\-สมฺปยุตฺเตหิ จ ธมฺเมหิ เย ธมฺมา ขนฺธ\-สงฺคเหน อสงฺคหิตา อายตน\-สงฺคเหน สงฺคหิตา ธาตุสงฺคเหน สงฺคหิตา ฯเปฯ เต ธมฺมา อสงฺขตํ ขนฺธโต ฐเปตฺวา ตีหิ ขนฺเธหิ เอเกนายตเนน เอกาย ธาตุยา สงฺคหิตา.}

\proseitem{185}{อปฺปจฺจเยหิ ธมฺเมหิ เย ธมฺมา. อสงฺขเตหิ ธมฺเมหิ เย ธมฺมา ขนฺธ\-สงฺคเหน อสงฺคหิตา อายตน\-สงฺคเหน สงฺคหิตา ธาตุสงฺคเหน สงฺคหิตา ฯเปฯ เต ธมฺมา จตูหิ ขนฺเธหิ เอเกนายตเนน เอกาย ธาตุยา สงฺคหิตา.}

\proseitem{186}{อาสเวหิ ธมฺเมหิ เย ธมฺมา. อาสเวหิ เจว สาสเวหิ จ ธมฺเมหิ เย ธมฺมา. อาสเวหิ เจว อาสว\-สมฺปยุตฺเตหิ จ ธมฺเมหิ เย ธมฺมา ขนฺธ\-สงฺคเหน อสงฺคหิตา อายตน\-สงฺคเหน สงฺคหิตา ธาตุสงฺคเหน สงฺคหิตา. เต ธมฺมา อสงฺขตํ ขนฺธโต ฐเปตฺวา ตีหิ ขนฺเธหิ เอเกนายตเนน เอกาย ธาตุยา สงฺคหิตา.}

\proseitem{187}{สํโยชเนหิ. คนฺเถหิ. โอเฆหิ. โยเคหิ. นีวรเณหิ. ปรามาเสหิ ธมฺเมหิ เย ธมฺมา. ปรามาเสหิ เจว ปรามฏฺเฐหิ จ ธมฺเมหิ เย ธมฺมา ขนฺธ\-สงฺคเหน อสงฺคหิตา อายตน\-สงฺคเหน สงฺคหิตา ธาตุสงฺคเหน สงฺคหิตา. เต ธมฺมา อสงฺขตํ ขนฺธโต ฐเปตฺวา ตีหิ ขนฺเธหิ เอเกนายตเนน เอกาย ธาตุยา สงฺคหิตา.}

\proseitem{188}{\csromanfolio{32}เจตสิเกหิ ธมฺเมหิ เย ธมฺมา. จิตฺต\-สมฺปยุตฺเตหิ ธมฺเมหิ เย ธมฺมา. จิตฺต\-สํสฏฺเฐหิ ธมฺเมหิ เย ธมฺมา. จิตฺต\-สํสฏฺฐ\-สมุฏฺฐาเนหิ ธมฺเมหิ เย ธมฺมา. จิตฺต\-สํสฏฺฐ\-สมุฏฺฐาน\-สหภูหิ ธมฺเมหิ เย ธมฺมา. จิตฺต\-สํสฏฺฐ\-สมุฏฺฐานา\-นุปริวตฺตีหิ ธมฺเมหิ เย ธมฺมา ขนฺธ\-สงฺคเหน อสงฺคหิตา อายตน\-สงฺคเหน สงฺคหิตา ธาตุสงฺคเหน สงฺคหิตา ฯเปฯ เต ธมฺมา อสงฺขตํ ขนฺธโต ฐเปตฺวา เอเกน ขนฺเธน เอเกนายตเนน เอกาย ธาตุยา สงฺคหิตา.}

\proseitem{189}{จิตฺตสหภูหิ ธมฺเมหิ เย ธมฺมา. จิตฺตา\-นุปริวตฺตีหิ ธมฺเมหิ เย ธมฺมา ขนฺธ\-สงฺคเหน อสงฺคหิตา อายตน\-สงฺคเหน สงฺคหิตา ธาตุสงฺคเหน สงฺคหิตา ฯเปฯ เต ธมฺมา น เกหิจิ ขนฺเธหิ เอเกนายตเนน เอกาย ธาตุยา สงฺคหิตา.}

\proseitem{190}{อุปาทาเนหิ ธมฺเมหิ เย ธมฺมา. กิเลเสหิ ธมฺเมหิ เย ธมฺมา. กิเลเสหิ เจว สํกิเลสิเกหิ จ ธมฺเมหิ เย ธมฺมา. กิเลเสหิ เจว สํกิลิฏฺเฐหิ จ ธมฺเมหิ เย ธมฺมา. กิเลเสหิ เจว กิเลส\-สมฺปยุตฺเตหิ จ ธมฺเมหิ เย ธมฺมา ขนฺธ\-สงฺคเหน อสงฺคหิตา อายตน\-สงฺคเหน สงฺคหิตา ธาตุสงฺคเหน สงฺคหิตา. เต ธมฺมา กติหิ ขนฺเธหิ กติหายตเนหิ กติหิ ธาตูหิ สงฺคหิตา, เต ธมฺมา อสงฺขตํ ขนฺธโต ฐเปตฺวา ตีหิ ขนฺเธหิ เอเกนายตเนน เอกาย ธาตุยา สงฺคหิตา.}

\setlength{\csromangathapagewidth}{0pt}
\setlength{\csromangathawakwidth}{0pt}
\csromangathameasuregroup{ตโย ขนฺธา ตถา สจฺจา, \\ ปทานิ ปจฺจยากาเร, \\ สมติํส ปทา โหนฺติ, \\ ทุเว จูฬนฺตรทุกา\textsuperscript{1},}{\csromangathabat{ตโย ขนฺธา ตถา สจฺจา,}{อินฺทฺริยานิ จ โสฬส.} \\ \csromangathabat{ปทานิ ปจฺจยากาเร,}{จุทฺทสูปริ จุทฺทส.} \\ \csromangathabat{สมติํส ปทา โหนฺติ,}{โคจฺฉเกสุ ทสสฺวถ.} \\ \csromangathabat{ทุเว จูฬนฺตรทุกา\textsuperscript{1},}{อฏฺฐ โหนฺติ มหนฺตราติ.}}
\csromangathasetleft{ตโย ขนฺธา ตถา สจฺจา, \\ ปทานิ ปจฺจยากาเร, \\ สมติํส ปทา โหนฺติ, \\ ทุเว จูฬนฺตรทุกา\textsuperscript{1},}
\csromangathagroup{\csromangathastanza{\csromangathabat{ตโย ขนฺธา ตถา สจฺจา,}{อินฺทฺริยานิ จ โสฬส.} \\ \csromangathabat{ปทานิ ปจฺจยากาเร,}{จุทฺทสูปริ จุทฺทส.}}\csromangathastanza{\csromangathabat{สมติํส ปทา โหนฺติ,}{โคจฺฉเกสุ ทสสฺวถ.} \\ \csromangathabat{ทุเว จูฬนฺตรทุกา\csromansharedfootnote{5c658f8ea90dae1d}{จุลฺลนฺตรทุกา (สี)},}{อฏฺฐ โหนฺติ มหนฺตราติ.}}}

\vspace{\tipitakaparskip}
\csromancenter{อสงฺคหิเตน สงฺคหิต\-ปท\-นิทฺเทโส ตติโย.}
\csromanmatikaanchor{mtk.21}
\csromantocmarknopage{section}{4. จตุตฺถนย}{mtk.21}
\bookmark[dest={mtk.21},level=1]{4. จตุตฺถนย}
\csromanheader{1}{4. \textbf{จตุตฺถนย}}
\csromanmatikaanchor{mtk.22}
\csromantocmarkref{subsection}{4. สงฺคหิเตนสงฺคหิตปทนิทฺเทส}{mtk.22}
\bookmark[dest={mtk.22},level=2]{4. สงฺคหิเตนสงฺคหิตปทนิทฺเทส}
\csromanheader{2}{4. สงฺคหิเตน\-สงฺคหิต\-ปท\-นิทฺเทส}
\proseitem{191}{\csromanfolio{33}สมุทย\-สจฺเจน เย ธมฺมา. มคฺคสจฺเจน เย ธมฺมา ขนฺธ\-สงฺคเหน สงฺคหิตา อายตน\-สงฺคเหน สงฺคหิตา ธาตุสงฺคเหน สงฺคหิตา, เตหิ ธมฺเมหิ เย ธมฺมา ขนฺธ\-สงฺคเหน สงฺคหิตา อายตน\-สงฺคเหน สงฺคหิตา ธาตุสงฺคเหน สงฺคหิตา. เต ธมฺมา กติหิ ขนฺเธหิ กติหายตเนหิ กติหิ ธาตูหิ สงฺคหิตา, เต ธมฺมา เอเกน ขนฺเธน เอเกนายตเนน เอกาย ธาตุยา สงฺคหิตา.}

\proseitem{192}{อิตฺถินฺทฺริเยน เย ธมฺมา. ปุริสินฺทฺริเยน เย ธมฺมา. สุขินฺทฺริเยน เย ธมฺมา. ทุกฺขินฺทฺริเยน เย ธมฺมา. โสมนสฺสิ\-นฺทฺริเยน เย ธมฺมา. โทมนสฺสิ\-นฺทฺริเยน เย ธมฺมา. อุเปกฺขิ\-นฺทฺริเยน เย ธมฺมา. สทฺธินฺทฺริเยน เย ธมฺมา. วีริยินฺทฺริเยน เย ธมฺมา. สตินฺทฺริเยน เย ธมฺมา. สมาธิ\-นฺทฺริเยน เย ธมฺมา. ปญฺญินฺทฺริเยน เย ธมฺมา. อนญฺญาต\-ญฺญสฺสามี\-ตินฺทฺริเยน เย ธมฺมา. อญฺญินฺทฺริเยน เย ธมฺมา. อญฺญาตาวิ\-นฺทฺริเยน เย ธมฺมา.}

\prose{อวิชฺชาย เย ธมฺมา. อวิชฺชาปจฺจยา สงฺขาเรน เย ธมฺมา. สฬายตน\-ปจฺจยา ผสฺเสน เย ธมฺมา. เวทนาปจฺจยา ตณฺหาย เย ธมฺมา. ตณฺหาปจฺจยา อุปาทาเนน เย ธมฺมา. กมฺมภเวน เย ธมฺมา. โสเกน เย ธมฺมา. ปริเทเวน เย ธมฺมา. ทุกฺเขน เย ธมฺมา. โทมนสฺเสน เย ธมฺมา. อุปายาเสน เย ธมฺมา.}

\prose{สติปฏฺฐาเนน เย ธมฺมา. สมฺม\-ปฺปธาเนน เย ธมฺมา. อปฺปมญฺญาย เย ธมฺมา. ปญฺจหิ อินฺทฺริเยหิ เย ธมฺมา. ปญฺจหิ พเลหิ เย ธมฺมา. สตฺตหิ โพชฺฌงฺเคหิ เย ธมฺมา. อริเยน อฏฺฐงฺคิเกน มคฺเคน เย ธมฺมา. ผสฺเสน เย ธมฺมา. เจตนาย เย ธมฺมา. อธิโมกฺเขน เย ธมฺมา. มนสิกาเรน เย ธมฺมา.}

\prose{เหตูหิ ธมฺเมหิ เย ธมฺมา. เหตูหิ เจว สเหตุเกหิ จ ธมฺเมหิ เย ธมฺมา. เหตูหิ เจว เหตุ\-สมฺปยุตฺเตหิ จ ธมฺเมหิ เย ธมฺมา. อาสเวหิ. สํโยชเนหิ. คนฺเถหิ. โอเฆหิ. โยเคหิ. นีวรเณหิ. ปรามาเสหิ. อุปาทาเนหิ. กิเลเสหิ ธมฺเมหิ เย ธมฺมา. กิเลเสหิ \csromanfolio{34}เจว สํกิเลสิเกหิ จ ธมฺเมหิ เย ธมฺมา. กิเลเสหิ เจว สํกิลิฏฺเฐหิ จ ธมฺเมหิ เย ธมฺมา. กิเลเสหิ เจว กิเลส\-สมฺปยุตฺเตหิ จ ธมฺเมหิ เย ธมฺมา ขนฺธ\-สงฺคเหน สงฺคหิตา อายตน\-สงฺคเหน สงฺคหิตา ธาตุสงฺคเหน สงฺคหิตา. เตหิ ธมฺเมหิ เย ธมฺมา ขนฺธ\-สงฺคเหน สงฺคหิตา อายตน\-สงฺคเหน สงฺคหิตา ธาตุสงฺคเหน สงฺคหิตา. เต ธมฺมา กติหิ ขนฺเธหิ กติหายตเนหิ กติหิ ธาตูหิ สงฺคหิตา, เต ธมฺมา เอเกน ขนฺเธน เอเกนายตเนน เอกาย ธาตุยา สงฺคหิตา.}

\setlength{\csromangathapagewidth}{0pt}
\setlength{\csromangathawakwidth}{0pt}
\csromangathameasuregroup{เทฺว สจฺจา ปนฺนรสินฺทฺริยา, \\ อุทฺธํ ปุน เอกาทส,}{\csromangathabat{เทฺว สจฺจา ปนฺนรสินฺทฺริยา,}{เอกาทส ปฏิจฺจปทา.} \\ \csromangathabat{อุทฺธํ ปุน เอกาทส,}{โคจฺฉกปทเมตฺถ ติํสวิธาติ\textsuperscript{1}.}}
\csromangathasetleft{เทฺว สจฺจา ปนฺนรสินฺทฺริยา, \\ อุทฺธํ ปุน เอกาทส,}
\csromangathagroup{\csromangathastanza{\csromangathabat{เทฺว สจฺจา ปนฺนรสิ\-นฺทฺริยา,}{เอกาทส ปฏิจฺจปทา.} \\ \csromangathabat{อุทฺธํ ปุน เอกาทส,}{โคจฺฉกปทเมตฺถ ติํสวิธาติ\csromansharedfootnote{018d26250ffa500c}{ติํสวิธนฺติ (อิ)}.}}}

\vspace{\tipitakaparskip}
\csromancenter{\textbf{สงฺคหิเตน สงฺคหิตปสนิทฺเทโส จตุตฺโถ.}}
\csromanmatikaanchor{mtk.23}
\csromantocmarknopage{section}{5. ปญฺจมนย}{mtk.23}
\bookmark[dest={mtk.23},level=1]{5. ปญฺจมนย}
\csromanheader{1}{5. \textbf{ปญฺจมนย}}
\csromanmatikaanchor{mtk.24}
\csromantocmarkref{subsection}{5. อสงฺคหิเตนอสงฺคหิตปทนิทฺเทส}{mtk.24}
\bookmark[dest={mtk.24},level=2]{5. อสงฺคหิเตนอสงฺคหิตปทนิทฺเทส}
\csromanheader{2}{5. \textbf{อสงฺคหิเตน อสงฺคหิตปทนิทฺเทส}}
\proseitem{193}{\csromanfolio{35}รูปกฺขนฺเธน เย ธมฺมา ขนฺธ\-สงฺคเหน อสงฺคหิตา อายตน\-สงฺคเหน อสงฺคหิตา ธาตุสงฺคเหน อสงฺคหิตา, เตหิ ธมฺเมหิ เย ธมฺมา ขนฺธ\-สงฺคเหน อสงฺคหิตา อายตน\-สงฺคเหน อสงฺคหิตา ธาตุสงฺคเหน อสงฺคหิตา. เต ธมฺมา กติหิ ขนฺเธหิ กติหายตเนหิ กติหิ ธาตูหิ อสงฺคหิตา, เต ธมฺมา เอเกน ขนฺเธน เอเกนายตเนน สตฺตหิ ธาตูหิ อสงฺคหิตา.}

\proseitem{194}{เวทนา\-กฺขนฺเธน เย ธมฺมา. สญฺญา\-กฺขนฺเธน เย ธมฺมา. สงฺขาร\-กฺขนฺเธน เย ธมฺมา ขนฺธ\-สงฺคเหน อสงฺคหิตา อายตน\-สงฺคเหน อสงฺคหิตา ธาตุสงฺคเหน อสงฺคหิตา, เตหิ ธมฺเมหิ เย ธมฺมา ฯเปฯ เต ธมฺมา ทฺวีหิ ขนฺเธหิ เอกาทสหา\-ยตเนหิ สตฺตรสหิ ธาตูหิ อสงฺคหิตา.}

\proseitem{195}{วิญฺญาณ\-กฺขนฺเธน เย ธมฺมา. มนายตเนน เย ธมฺมา. จกฺขุวิญฺญาณ\-ธาตุยา เย ธมฺมา ฯเปฯ. มโนธาตุยา เย ธมฺมา. มโนวิญฺญาณ\-ธาตุยา เย ธมฺมา. มนินฺทฺริเยน เย ธมฺมา ขนฺธ\-สงฺคเหน อสงฺคหิตา อายตน\-สงฺคเหน อสงฺคหิตา ธาตุสงฺคเหน อสงฺคหิตา, เตหิ ธมฺเมหิ เย ธมฺมา ฯเปฯ เต ธมฺมา จตูหิ ขนฺเธหิ เอกาทสหา\-ยตเนหิ เอกาทสหิ ธาตูหิ อสงฺคหิตา.}

\proseitem{196}{จกฺขายตเนน เย ธมฺมา ฯเปฯ. โผฏฺฐพฺพา\-ยตเนน เย ธมฺมา. จกฺขุธาตุยา เย ธมฺมา ฯเปฯ. โผฏฺฐพฺพ\-ธาตุยา เย ธมฺมา ขนฺธ\-สงฺคเหน อสงฺคหิตา อายตน\-สงฺคเหน อสงฺคหิตา ธาตุสงฺคเหน อสงฺคหิตา, เตหิ ธมฺเมหิ เย ธมฺมา ฯเปฯ เต ธมฺมา จตูหิ ขนฺเธหิ ทฺวีหายตเนหิ อฏฺฐหิ ธาตูหิ อสงฺคหิตา.}

\proseitem{197}{ธมฺมายตเนน เย ธมฺมา. ธมฺมธาตุยา เย ธมฺมา. อิตฺถินฺทฺริเยน เย ธมฺมา. ปุริสินฺทฺริเยน เย ธมฺมา. ชีวิตินฺทฺริเยน เย ธมฺมา ขนฺธ\-สงฺคเหน อสงฺคหิตา อายตน\-สงฺคเหน อสงฺคหิตา ธาตุสงฺคเหน อสงฺคหิตา, เตหิ ธมฺเมหิ เย ธมฺมา ฯเปฯ เต ธมฺมา เอเกน ขนฺเธน เอเกนายตเนน สตฺตหิ ธาตูหิ อสงฺคหิตา.}

\proseitem{198}{\csromanfolio{36}สมุทย\-สจฺเจน เย ธมฺมา. มคฺคสจฺเจน เย ธมฺมา. นิโรธสจฺเจน เย ธมฺมา ขนฺธ\-สงฺคเหน อสงฺคหิตา อายตน\-สงฺคเหน อสงฺคหิตา. ธาตุสงฺคเหน อสงฺคหิตา, เตหิ ธมฺเมหิ เย ธมฺมา ฯเปฯ เต ธมฺมา ทฺวีหิ ขนฺเธหิ เอกาทสหา\-ยตเนหิ สตฺตรสหิ ธาตูหิ อสงฺคหิตา.}

\proseitem{199}{จกฺขุนฺทฺริเยน เย ธมฺมา ฯเปฯ. กายินฺทฺริเยน เย ธมฺมา ขนฺธ\-สงฺคเหน อสงฺคหิตา อายตน\-สงฺคเหน อสงฺคหิตา ธาตุสงฺคเหน อสงฺคหิตา, เตหิ ธมฺเมหิ เย ธมฺมา ฯเปฯ เต ธมฺมา จตูหิ ขนฺเธหิ ทฺวีหายตเนหิ อฏฺฐหิ ธาตูหิ อสงฺคหิตา.}

\proseitem{200}{สุขินฺทฺริเยน เย ธมฺมา. ทุกฺขินฺทฺริเยน เย ธมฺมา. โสมนสฺสิ\-นฺทฺริเยน เย ธมฺมา. โทมนสฺสิ\-นฺทฺริเยน เย ธมฺมา. อุเปกฺขิ\-นฺทฺริเยน เย ธมฺมา. สทฺธินฺทฺริเยน เย ธมฺมา. วีริยินฺทฺริเยน เย ธมฺมา. สตินฺทฺริเยน เย ธมฺมา. สมาธิ\-นฺทฺริเยน เย ธมฺมา. ปญฺญินฺทฺริเยน เย ธมฺมา. อนญฺญาต\-ญฺญสฺสามี\-ตินฺทฺริเยน เย ธมฺมา. อญฺญินฺทฺริเยน เย ธมฺมา. อญฺญาตาวิ\-นฺทฺริเยน เย ธมฺมา. อวิชฺชาย เย ธมฺมา. อวิชฺชาปจฺจยา สงฺขาเรน เย ธมฺมา ขนฺธ\-สงฺคเหน อสงฺคหิตา อายตน\-สงฺคเหน อสงฺคหิตา ธาตุสงฺคเหน อสงฺคหิตา, เตหิ ธมฺเมหิ เย ธมฺมา ฯเปฯ เต ธมฺมา ทฺวีหิ ขนฺเธหิ เอกาทสหา\-ยตเนหิ สตฺตรสหิ ธาตูหิ อสงฺคหิตา.}

\proseitem{201}{สงฺขาร\-ปจฺจยา วิญฺญาเณน เย ธมฺมา ขนฺธ\-สงฺคเหน อสงฺคหิตา อายตน\-สงฺคเหน อสงฺคหิตา ธาตุสงฺคเหน อสงฺคหิตา, เตหิ ธมฺเมหิ เย ธมฺมา ฯเปฯ เต ธมฺมา จตูหิ ขนฺเธหิ เอกาทสหา\-ยตเนหิ เอกาทสหิ ธาตูหิ อสงฺคหิตา.}

\proseitem{202}{วิญฺญาณปจฺจยา นามรูเปน เย ธมฺมา ขนฺธ\-สงฺคเหน อสงฺคหิตา อายตน\-สงฺคเหน อสงฺคหิตา ธาตุสงฺคเหน อสงฺคหิตา, เตหิ ธมฺเมหิ เย ธมฺมา ฯเปฯ เต ธมฺมา เอเกน ขนฺเธน เอเกนายตเนน สตฺตหิ ธาตูหิ อสงฺคหิตา.}

\proseitem{203}{นามรูป\-ปจฺจยา สฬายตเนน เย ธมฺมา ขนฺธ\-สงฺคเหน อสงฺคหิตา อายตน\-สงฺคเหน อสงฺคหิตา ธาตุสงฺคเหน อสงฺคหิตา, เตหิ ธมฺเมหิ เย ธมฺมา ฯเปฯ เต ธมฺมา ตีหิ ขนฺเธหิ เอเกนายตเนน เอกาย ธาตุยา อสงฺคหิตา.}

\proseitem{204}{\csromanfolio{37}สฬายตน\-ปจฺจยา ผสฺเสน เย ธมฺมา. ผสฺสปจฺจยา เวทนาย เย ธมฺมา. เวทนาปจฺจยา ตณฺหาย เย ธมฺมา. ตณฺหาปจฺจยา อุปาทาเนน เย ธมฺมา. กมฺมภเวน เย ธมฺมา ขนฺธ\-สงฺคเหน อสงฺคหิตา อายตน\-สงฺคเหน อสงฺคหิตา ธาตุสงฺคเหน อสงฺคหิตา, เตหิ ธมฺเมหิ เย ธมฺมา ฯเปฯ เต ธมฺมา ทฺวีหิ ขนฺเธหิ เอกาทสหา\-ยตเนหิ สตฺตรสหิ ธาตูหิ อสงฺคหิตา.}

\proseitem{205}{อรูปภเวน เย ธมฺมา. เนว\-สญฺญา\-นา\-สญฺญาภเวน เย ธมฺมา. จตุโวการภเวน เย ธมฺมา. อิทฺธิปาเทน เย ธมฺมา ขนฺธ\-สงฺคเหน อสงฺคหิตา อายตน\-สงฺคเหน อสงฺคหิตา ธาตุสงฺคเหน อสงฺคหิตา, เตหิ ธมฺเมหิ เย ธมฺมา ฯเปฯ เต ธมฺมา เอเกน ขนฺเธน ทสหายตเนหิ ทสหิ ธาตูหิ อสงฺคหิตา.}

\proseitem{206}{อสญฺญาภเวน เย ธมฺมา. เอกโวการ\-ภเวน เย ธมฺมา. ชาติยา เย ธมฺมา. ชราย เย ธมฺมา. มรเณน เย ธมฺมา ขนฺธ\-สงฺคเหน อสงฺคหิตา อายตน\-สงฺคเหน อสงฺคหิตา ธาตุสงฺคเหน อสงฺคหิตา, เตหิ ธมฺเมหิ เย ธมฺมา ฯเปฯ เต ธมฺมา เอเกน ขนฺเธน เอเกนายตเนน สตฺตหิ ธาตูหิ อสงฺคหิตา.}

\proseitem{207}{ปริเทเวน เย ธมฺมา ขนฺธ\-สงฺคเหน อสงฺคหิตา อายตน\-สงฺคเหน อสงฺคหิตา ธาตุสงฺคเหน อสงฺคหิตา, เตหิ ธมฺเมหิ เย ธมฺมา ฯเปฯ เต ธมฺมา จตูหิ ขนฺเธหิ ทฺวีหายตเนหิ อฏฺฐหิ ธาตูหิ}

\proseitem{208}{โสเกน เย ธมฺมา. ทุกฺเขน เย ธมฺมา. โทมนสฺเสน เย ธมฺมา. อุปายาเสน เย ธมฺมา. สติปฏฺฐาเนน เย ธมฺมา. สมฺม\-ปฺปธาเนน เย ธมฺมา. ฌาเนน เย ธมฺมา. อปฺปมญฺญาย เย ธมฺมา. ปญฺจหิ อินฺทฺริเยหิ เย ธมฺมา. ปญฺจหิ พเลหิ เย ธมฺมา. สตฺตหิ โพชฺฌงฺเคหิ เย ธมฺมา. อริเยน อฏฺฐงฺคิเกน มคฺเคน เย ธมฺมา. ผสฺเสน เย ธมฺมา. เวทนาย เย ธมฺมา. สญฺญาย เย ธมฺมา. เจตนาย เย ธมฺมา. อธิโมกฺเขน เย ธมฺมา. มนสิกาเรน เย ธมฺมา ขนฺธ\-สงฺคเหน อสงฺคหิตา อายตน\-สงฺคเหน อสงฺคหิตา ธาตุสงฺคเหน อสงฺคหิตา, เตหิ ธมฺเมหิ เย ธมฺมา ฯเปฯ เต ธมฺมา ทฺวีหิ ขนฺเธหิ เอกาทสหา\-ยตเนหิ สตฺตรสหิ ธาตูหิ อสงฺคหิตา.}

\proseitem{209}{\csromanfolio{38}จิตฺเตน เย ธมฺมา ขนฺธ\-สงฺคเหน อสงฺคหิตา อายตน\-สงฺคเหน อสงฺคหิตา ธาตุสงฺคเหน อสงฺคหิตา, เตหิ ธมฺเมหิ เย ธมฺมา ฯเปฯ เต ธมฺมา จตูหิ ขนฺเธหิ เอกาทสหา\-ยตเนหิ เอกาทสหิ ธาตูหิ อสงฺคหิตา.}

\csromanmatikaanchor{mtk.25}
\csromantocmarkref{subsection}{1. ติก}{mtk.25}
\bookmark[dest={mtk.25},level=2]{1. ติก}
\csromanheader{2}{1. \textbf{ติก}}
\proseitem{210}{กุสเลหิ ธมฺเมหิ เย ธมฺมา. อกุสเลหิ ธมฺเมหิ เย ธมฺมา. สุขาย เวทนาย สมฺปยุตฺเตหิ ธมฺเมหิ เย ธมฺมา. ทุกฺขาย เวทนาย สมฺปยุตฺเตหิ ธมฺเมหิ เย ธมฺมา. อทุกฺขม\-สุขาย เวทนาย สมฺปยุตฺเตหิ ธมฺเมหิ เย ธมฺมา. วิปาเกหิ ธมฺเมหิ เย ธมฺมา. วิปาก\-ธมฺม\-ธมฺเมหิ เย ธมฺมา. อนุปาทินฺน\-อนุปาทานิเยหิ ธมฺเมหิ เย ธมฺมา. สํกิลิฏฺฐ\-สํกิเลสิเกหิ ธมฺเมหิ เย ธมฺมา. อสํกิลิฏฺฐ- อสํกิเลสิเกหิ ธมฺเมหิ เย ธมฺมา. สวิตกฺก\-สวิจาเรหิ ธมฺเมหิ เย ธมฺมา. อวิตกฺก\-วิจารมตฺเตหิ ธมฺเมหิ เย ธมฺมา. ปีติสหคเตหิ ธมฺเมหิ เย ธมฺมา. สุขสหคเตหิ ธมฺเมหิ เย ธมฺมา. อุเปกฺขา\-สหคเตหิ ธมฺเมหิ เย ธมฺมา. ทสฺสเนน ปหาตพฺเพหิ ธมฺเมหิ เย ธมฺมา. ภาวนาย ปหาตพฺเพหิ ธมฺเมหิ เย ธมฺมา. ทสฺสเนน ปหาตพฺพ\-เหตุเกหิ ธมฺเมหิ เย ธมฺมา. ภาวนาย ปหาตพฺพ\-เหตุเกหิ ธมฺเมหิ เย ธมฺมา. อาจยคามีหิ ธมฺเมหิ เย ธมฺมา. อปจยคามีหิ ธมฺเมหิ เย ธมฺมา. เสกฺเขหิ ธมฺเมหิ เย ธมฺมา. อเสกฺเขหิ ธมฺเมหิ เย ธมฺมา. มหคฺคเตหิ ธมฺเมหิ เย ธมฺมา. อปฺปมาเณหิ ธมฺเมหิ เย ธมฺมา. ปริตฺตา\-รมฺมเณหิ ธมฺเมหิ เย ธมฺมา. มหคฺคตา\-รมฺมเณหิ ธมฺเมหิ เย ธมฺมา. อปฺปมาณา\-รมฺมเณหิ ธมฺเมหิ เย ธมฺมา. หีเนหิ ธมฺเมหิ เย ธมฺมา. ปณีเตหิ ธมฺเมหิ เย ธมฺมา. มิจฺฉตฺต\-นิยเตหิ ธมฺเมหิ เย ธมฺมา. สมฺมตฺต\-นิยเตหิ ธมฺเมหิ เย ธมฺมา. มคฺคารมฺมเณหิ ธมฺเมหิ เย ธมฺมา. มคฺคเหตุเกหิ ธมฺเมหิ เย ธมฺมา. มคฺคาธิปตีหิ ธมฺเมหิ เย ธมฺมา. อตีตารมฺมเณหิ ธมฺเมหิ เย ธมฺมา. อนาคตา\-รมฺมเณหิ ธมฺเมหิ เย ธมฺมา. ปจฺจุปฺปนฺนา\-รมฺมเณหิ ธมฺเมหิ เย ธมฺมา. อชฺฌตฺตา\-รมฺมเณหิ ธมฺเมหิ เย ธมฺมา. พหิทฺธา\-รมฺมเณหิ ธมฺเมหิ เย ธมฺมา. อชฺฌตฺตพหิทฺธา\-รมฺมเณหิ ธมฺเมหิ เย ธมฺมา ขนฺธ\-สงฺคเหน อสงฺคหิตา อายตน\-สงฺคเหน อสงฺคหิตา ธาตุสงฺคเหน อสงฺคหิตา, เตหิ ธมฺเมหิ เย ธมฺมา ฯเปฯ เต ธมฺมา เอเกน ขนฺเธน ทสหายตเนหิ ทสหิ ธาตูหิ อสงฺคหิตา.}

\proseitem{211}{\csromanfolio{39}สนิทสฺสน\-สปฺปฏิเฆหิ ธมฺเมหิ เย ธมฺมา. อนิทสฺสน\-สปฺปฏิเฆหิ ธมฺเมหิ เย ธมฺมา ขนฺธ\-สงฺคเหน อสงฺคหิตา อายตน\-สงฺคเหน อสงฺคหิตา ธาตุสงฺคเหน อสงฺคหิตา, เตหิ ธมฺเมหิ เย ธมฺมา ฯเปฯ เต ธมฺมา จตูหิ ขนฺเธหิ ทฺวีหายตเนหิ อฏฺฐหิ ธาตูหิ อสงฺคหิตา.}

\csromanmatikaanchor{mtk.26}
\csromantocmarkref{subsection}{2. ทุก}{mtk.26}
\bookmark[dest={mtk.26},level=2]{2. ทุก}
\csromanheader{2}{2. \textbf{ทุก}}
\proseitem{212}{เหตูหิ ธมฺเมหิ เย ธมฺมา. เหตูหิ เจว สเหตุเกหิ จ ธมฺเมหิ เย ธมฺมา. เหตูหิ เจว เหตุ\-สมฺปยุตฺเตหิ จ ธมฺเมหิ เย ธมฺมา. ขนฺธ\-สงฺคเหน อสงฺคหิตา อายตน\-สงฺคเหน อสงฺคหิตา ธาตุสงฺคเหน อสงฺคหิตา, เตหิ ธมฺเมหิ เย ธมฺมา ฯเปฯ เต ธมฺมา ทฺวีหิ ขนฺเธหิ เอกาทสหา\-ยตเนหิ สตฺตรสหิ ธาตูหิ อสงฺคหิตา.}

\proseitem{213}{สเหตุเกหิ ธมฺเมหิ เย ธมฺมา. เหตุ\-สมฺปยุตฺเตหิ ธมฺเมหิ เย ธมฺมา. สเหตุเกหิ เจว น จ เหตูหิ ธมฺเมหิ เย ธมฺมา. เหตุ\-สมฺปยุตฺเตหิ เจว น จ เหตูหิ ธมฺเมหิ เย ธมฺมา. น เหตุสเหตุเกหิ\csromansharedfootnote{58098e21b9126315}{น เหตู สเหตุเกหิ (สี), น เหตูหิ สเหตุเกหิ (สฺยา, ก)} ธมฺเมหิ เย ธมฺมา ขนฺธ\-สงฺคเหน อสงฺคหิตา อายตน\-สงฺคเหน อสงฺคหิตา ธาตุสงฺคเหน อสงฺคหิตา, เตหิ ธมฺเมหิ เย ธมฺมา ฯเปฯ เต ธมฺมา เอเกน ขนฺเธน ทสหายตเนหิ ทสหิ ธาตูหิ}

\proseitem{214}{อปฺปจฺจเยหิ ธมฺเมหิ เย ธมฺมา. อสงฺขเตหิ ธมฺเมหิ เย ธมฺมา ขนฺธ\-สงฺคเหน อสงฺคหิตา อายตน\-สงฺคเหน อสงฺคหิตา ธาตุสงฺคเหน อสงฺคหิตา, เตหิ ธมฺเมหิ เย ธมฺมา ฯเปฯ เต ธมฺมา ทฺวีหิ ขนฺเธหิ เอกาทสหา\-ยตเนหิ สตฺตรสหิ ธาตูหิ อสงฺคหิตา.}

\proseitem{215}{สนิทสฺสเนหิ ธมฺเมหิ เย ธมฺมา. สปฺปฏิเฆหิ ธมฺเมหิ เย ธมฺมา ขนฺธ\-สงฺคเหน อสงฺคหิตา อายตน\-สงฺคเหน อสงฺคหิตา ธาตุสงฺคเหน อสงฺคหิตา, เตหิ ธมฺเมหิ เย ธมฺมา ฯเปฯ เต ธมฺมา จตูหิ ขนฺเธหิ ทฺวีหายตเนหิ อฏฺฐหิ ธาตูหิ อสงฺคหิตา.}

\proseitem{216}{รูปีหิ ธมฺเมหิ เย ธมฺมา ขนฺธ\-สงฺคเหน อสงฺคหิตา อายตน\-สงฺคเหน อสงฺคหิตา ธาตุสงฺคเหน อสงฺคหิตา, เตหิ ธมฺเมหิ \csromanfolio{40}เย ธมฺมา ฯเปฯ เต ธมฺมา เอเกน ขนฺเธน เอเกนายตเนน สตฺตหิ ธาตูหิ}

\proseitem{217}{อรูปีหิ ธมฺเมหิ เย ธมฺมา. โลกุตฺตเรหิ ธมฺเมหิ เย ธมฺมา ขนฺธ\-สงฺคเหน อสงฺคหิตา อายตน\-สงฺคเหน อสงฺคหิตา ธาตุสงฺคเหน อสงฺคหิตา, เตหิ ธมฺเมหิ เย ธมฺมา ฯเปฯ เต ธมฺมา เอเกน ขนฺเธน ทสหายตเนหิ ทสหิ ธาตูหิ อสงฺคหิตา.}

\proseitem{218}{อาสเวหิ ธมฺเมหิ เย ธมฺมา. อาสเวหิ เจว สาสเวหิ จ ธมฺเมหิ เย ธมฺมา. อาสเวหิ เจว อาสว\-สมฺปยุตฺเตหิ จ ธมฺเมหิ เย ธมฺมา ขนฺธ\-สงฺคเหน อสงฺคหิตา อายตน\-สงฺคเหน อสงฺคหิตา ธาตุสงฺคเหน อสงฺคหิตา, เตหิ ธมฺเมหิ เย ธมฺมา ฯเปฯ เต ธมฺมา ทฺวีหิ ขนฺเธหิ เอกาทสหา\-ยตเนหิ สตฺตรสหิ ธาตูหิ อสงฺคหิตา.}

\proseitem{219}{อนาสเวหิ ธมฺเมหิ เย ธมฺมา. อาสว\-สมฺปยุตฺเตหิ ธมฺเมหิ เย ธมฺมา. อาสว\-สมฺปยุตฺเตหิ เจว โน จ อาสเวหิ ธมฺเมหิ เย ธมฺมา. อาสว\-วิปฺปยุตฺเตหิ อนาสเวหิ ธมฺเมหิ เย ธมฺมา ขนฺธ\-สงฺคเหน อสงฺคหิตา อายตน\-สงฺคเหน อสงฺคหิตา ธาตุสงฺคเหน อสงฺคหิตา, เตหิ ธมฺเมหิ เย ธมฺมา ฯเปฯ เต ธมฺมา เอเกน ขนฺเธน ทสหายตเนหิ ทสหิ ธาตูหิ อสงฺคหิตา.}

\proseitem{220}{สํโยชเนหิ ธมฺเมหิ เย ธมฺมา. คนฺเถหิ ธมฺเมหิ เย ธมฺมา. โอเฆหิ ธมฺเมหิ เย ธมฺมา. โยเคหิ ธมฺเมหิ เย ธมฺมา. นีวรเณหิ ธมฺเมหิ เย ธมฺมา. ปรามาเสหิ ธมฺเมหิ เย ธมฺมา. ปรามาเสหิ เจว ปรามฏฺเฐหิ จ ธมฺเมหิ เย ธมฺมา ขนฺธ\-สงฺคเหน อสงฺคหิตา อายตน\-สงฺคเหน อสงฺคหิตา ธาตุสงฺคเหน อสงฺคหิตา, เตหิ ธมฺเมหิ เย ธมฺมา ฯเปฯ เต ธมฺมา ทฺวีหิ ขนฺเธหิ เอกาทสหา\-ยตเนหิ สตฺตรสหิ ธาตูหิ อสงฺคหิตา.}

\proseitem{221}{อปรามฏฺเฐหิ ธมฺเมหิ เย ธมฺมา. ปรามาส\-สมฺปยุตฺเตหิ ธมฺเมหิ เย ธมฺมา. ปรามาส\-วิปฺปยุตฺเตหิ อปรามฏฺเฐหิ ธมฺเมหิ เย ธมฺมา. สารมฺมเณหิ ธมฺเมหิ เย ธมฺมา ขนฺธ\-สงฺคเหน อสงฺคหิตา อายตน\-สงฺคเหน อสงฺคหิตา ธาตุสงฺคเหน อสงฺคหิตา, เตหิ ธมฺเมหิ เย ธมฺมา ฯเปฯ เต ธมฺมา เอเกน ขนฺเธน ทสหายตเนหิ ทสหิ ธาตูหิ อสงฺคหิตา.}

\proseitem{222}{\csromanfolio{41}อนารมฺมเณหิ ธมฺเมหิ เย ธมฺมา. โน จิตฺเตหิ ธมฺเมหิ เย ธมฺมา. จิตฺต\-วิปฺปยุตฺเตหิ ธมฺเมหิ เย ธมฺมา. จิตฺต\-วิสํสฏฺเฐหิ ธมฺเมหิ เย ธมฺมา. จิตฺต\-สมุฏฺฐาเนหิ ธมฺเมหิ เย ธมฺมา. จิตฺตสหภูหิ ธมฺเมหิ เย ธมฺมา. จิตฺตา\-นุปริวตฺตีหิ ธมฺเมหิ เย ธมฺมา. พาหิเรหิ ธมฺเมหิ เย ธมฺมา. อุปาทาธมฺเมหิ เย ธมฺมา ขนฺธ\-สงฺคเหน อสงฺคหิตา อายตน\-สงฺคเหน อสงฺคหิตา ธาตุสงฺคเหน อสงฺคหิตา, เตหิ ธมฺเมหิ เย ธมฺมา ฯเปฯ เต ธมฺมา เอเกน ขนฺเธน เอเกนายตเนน สตฺตหิ ธาตูหิ อสงฺคหิตา.}

\proseitem{223}{จิตฺเตหิ ธมฺเมหิ เย ธมฺมา ขนฺธ\-สงฺคเหน อสงฺคหิตา อายตน\-สงฺคเหน อสงฺคหิตา ธาตุสงฺคเหน อสงฺคหิตา, เตหิ ธมฺเมหิ เย ธมฺมา ฯเปฯ เต ธมฺมา จตูหิ ขนฺเธหิ เอกาทสหา\-ยตเนหิ เอกาทสหิ ธาตูหิ อสงฺคหิตา.}

\proseitem{224}{เจตสิเกหิ ธมฺเมหิ เย ธมฺมา. จิตฺต\-สมฺปยุตฺเตหิ ธมฺเมหิ เย ธมฺมา. จิตฺต\-สํสฏฺเฐหิ ธมฺเมหิ เย ธมฺมา. จิตฺต\-สํสฏฺฐ\-สมุฏฺฐาเนหิ ธมฺเมหิ เย ธมฺมา. จิตฺต\-สํสฏฺฐ\-สมุฏฺฐาน\-สหภูหิ ธมฺเมหิ เย ธมฺมา. จิตฺต\-สํสฏฺฐ\-สมุฏฺฐานา\-นุปริวตฺตีหิ ธมฺเมหิ เย ธมฺมา ขนฺธ\-สงฺคเหน อสงฺคหิตา อายตน\-สงฺคเหน อสงฺคหิตา ธาตุสงฺคเหน อสงฺคหิตา, เตหิ ธมฺเมหิ เย ธมฺมา ฯเปฯ เต ธมฺมา ทฺวีหิ ขนฺเธหิ เอกาทสหา\-ยตเนหิ สตฺตรสหิ ธาตูหิ อสงฺคหิตา.}

\proseitem{225}{อชฺฌตฺติเกหิ ธมฺเมหิ เย ธมฺมา ขนฺธ\-สงฺคเหน อสงฺคหิตา อายตน\-สงฺคเหน อสงฺคหิตา ธาตุสงฺคเหน อสงฺคหิตา, เตหิ ธมฺเมหิ เย ธมฺมา ฯเปฯ เต ธมฺมา ตีหิ ขนฺเธหิ เอเกนายตเนน เอกาย ธาตุยา}

\proseitem{226}{อุปาทาเนหิ ธมฺเมหิ เย ธมฺมา. กิเลเสหิ ธมฺเมหิ เย ธมฺมา. กิเลเสหิ เจว สํกิเลสิเกหิ จ ธมฺเมหิ เย ธมฺมา. กิเลเสหิ เจว สํกิลิฏฺเฐหิ จ ธมฺเมหิ เย ธมฺมา. กิเลเสหิ เจว กิเลส\-สมฺปยุตฺเตหิ จ ธมฺเมหิ เย ธมฺมา ขนฺธ\-สงฺคเหน อสงฺคหิตา อายตน\-สงฺคเหน อสงฺคหิตา ธาตุสงฺคเหน อสงฺคหิตา, เตหิ ธมฺเมหิ เย ธมฺมา ฯเปฯ เต ธมฺมา ทฺวีหิ ขนฺเธหิ เอกาทสหา\-ยตเนหิ สตฺตรสหิ ธาตูหิ อสงฺคหิตา.}

\proseitem{227}{\csromanfolio{42}อสํกิเลสิเกหิ ธมฺเมหิ เย ธมฺมา. สํกิลิฏฺเฐหิ ธมฺเมหิ เย ธมฺมา. กิเลส\-สมฺปยุตฺเตหิ ธมฺเมหิ เย ธมฺมา. สํกิลิฏฺเฐหิ เจว โน จ กิเลเสหิ ธมฺเมหิ เย ธมฺมา. กิเลส\-สมฺปยุตฺเตหิ เจว โน จ กิเลเสหิ ธมฺเมหิ เย ธมฺมา. กิเลส\-วิปฺปยุตฺเตหิ อสํกิเลสิเกหิ ธมฺเมหิ เย ธมฺมา. ทสฺสเนน ปหาตพฺเพหิ ธมฺเมหิ เย ธมฺมา. ภาวนาย ปหาตพฺเพหิ ธมฺเมหิ เย ธมฺมา. ทสฺสเนน ปหาตพฺพ\-เหตุเกหิ ธมฺเมหิ เย ธมฺมา. ภาวนาย ปหาตพฺพ\-เหตุเกหิ ธมฺเมหิ เย ธมฺมา. สวิตกฺเกหิ ธมฺเมหิ เย ธมฺมา. สวิจาเรหิ ธมฺเมหิ เย ธมฺมา. สปฺปีติเกหิ ธมฺเมหิ เย ธมฺมา. ปีติสหคเตหิ ธมฺเมหิ เย ธมฺมา. สุขสหคเตหิ ธมฺเมหิ เย ธมฺมา. อุเปกฺขา\-สหคเตหิ ธมฺเมหิ เย ธมฺมา. น กามาวจเรหิ ธมฺเมหิ เย ธมฺมา. รูปาวจเรหิ ธมฺเมหิ เย ธมฺมา. อรูปาวจเรหิ ธมฺเมหิ เย ธมฺมา. อปริยาปนฺเนหิ ธมฺเมหิ เย ธมฺมา. นิยฺยานิเกหิ ธมฺเมหิ เย ธมฺมา. นิยเตหิ ธมฺเมหิ เย ธมฺมา. อนุตฺตเรหิ ธมฺเมหิ เย ธมฺมา. สรเณหิ ธมฺเมหิ เย ธมฺมา ขนฺธ\-สงฺคเหน อสงฺคหิตา อายตน\-สงฺคเหน อสงฺคหิตา ธาตุสงฺคเหน อสงฺคหิตา, เตหิ ธมฺเมหิ เย ธมฺมา ขนฺธ\-สงฺคเหน อสงฺคหิตา อายตน\-สงฺคเหน อสงฺคหิตา ธาตุสงฺคเหน อสงฺคหิตา. เต ธมฺมา กติหิ ขนฺเธหิ กติหายตเนหิ กติหิ ธาตูหิ อสงฺคหิตา, เต ธมฺมา เอเกน ขนฺเธน ทสหายตเนหิ ทสหิ ธาตูหิ อสงฺคหิตา.}

\setlength{\csromangathapagewidth}{0pt}
\setlength{\csromangathawakwidth}{0pt}
\csromangathameasuregroup{}{รูปญฺจ ธมฺมายตนํ ธมฺมธาตุ, \\ อิตฺถิปุมํ ชีวิตํ นามรูปํ. \\ เทฺว ภวา ชาติชรา มจฺจุรูปํ, \\ อนารมฺมณํ โน จิตฺตํ จิตฺเตน วิปฺปยุตฺตํ. \\ วิสํสฏฺฐํ สมุฏฺฐานสหภุ อนุปริวตฺติ, \\ พาหิรํ อุปาทา เทฺว วิสโย\textsuperscript{1} เอสนโย สุพุทฺโธ. \\ \textbf{อสงฺคหิเตน อสงฺคหิตปทนิทฺเทโส ปญฺจโม.}}
\csromangathameasurewak{รูปญฺจ ธมฺมายตนํ ธมฺมธาตุ, \\ อิตฺถิปุมํ ชีวิตํ นามรูปํ. \\ เทฺว ภวา ชาติชรา มจฺจุรูปํ, \\ อนารมฺมณํ โน จิตฺตํ จิตฺเตน วิปฺปยุตฺตํ. \\ วิสํสฏฺฐํ สมุฏฺฐานสหภุ อนุปริวตฺติ, \\ พาหิรํ อุปาทา เทฺว วิสโย\textsuperscript{1} เอสนโย สุพุทฺโธ. \\ \textbf{อสงฺคหิเตน อสงฺคหิตปทนิทฺเทโส ปญฺจโม.}}
\setlength{\csromangathaleftcol}{0pt}
\csromangathagroup{\csromangathastanza{\csromangathawak{รูปญฺจ ธมฺมายตนํ ธมฺมธาตุ, \\ อิตฺถิปุมํ ชีวิตํ นามรูปํ. \\ เทฺว ภวา ชาติชรา มจฺจุรูปํ, \\ อนารมฺมณํ โน จิตฺตํ จิตฺเตน วิปฺปยุตฺตํ.}}\csromangathastanza{\csromangathawak{วิสํสฏฺฐํ สมุฏฺฐานสหภุ อนุปริวตฺติ, \\ พาหิรํ อุปาทา เทฺว วิสโย\csromansharedfootnote{4d233fe981ea0fd3}{เทฺววีสติ (สฺยา)} เอสนโย สุพุทฺโธ. \\ \textbf{อสงฺคหิเตน อสงฺคหิตปทนิทฺเทโส ปญฺจโม.}}}}

\csromanmatikaanchor{mtk.27}
\csromantocmarknopage{section}{6. ฉฏฺฐนย}{mtk.27}
\bookmark[dest={mtk.27},level=1]{6. ฉฏฺฐนย}
\csromanheader{1}{6. \textbf{ฉฏฺฐนย}}
\csromanmatikaanchor{mtk.28}
\csromantocmarkref{subsection}{6. สมฺปโยควิปฺปโยคปทนิทฺเทส}{mtk.28}
\bookmark[dest={mtk.28},level=2]{6. สมฺปโยควิปฺปโยคปทนิทฺเทส}
\csromanheader{2}{6. สมฺปโยค\-วิปฺปโยค\-ปท\-นิทฺเทส}
\csromanmatikaanchor{mtk.29}
\csromantocmarkref{subsection}{1. ขนฺธ}{mtk.29}
\bookmark[dest={mtk.29},level=2]{1. ขนฺธ}
\csromanheader{2}{1. \textbf{ขนฺธ}}
\proseitem{228}{\csromanfolio{43}รูปกฺขนฺโธ กติหิ ขนฺเธหิ กติหายตเนหิ กติหิ ธาตูหิ สมฺปยุตฺโตติ นตฺถิ. กติหิ วิปฺปยุตฺโต, จตูหิ ขนฺเธหิ เอเกนายตเนน สตฺตหิ ธาตูหิ วิปฺปยุตฺโต, เอเกนายตเนน เอกาย ธาตุยา เกหิจิ วิปฺปยุตฺโต.}

\proseitem{229}{เวทนากฺขนฺโธ. สญฺญากฺขนฺโธ. สงฺขาร\-กฺขนฺโธ ตีหิ ขนฺเธหิ เอเกนายตเนน สตฺตหิ ธาตูหิ สมฺปยุตฺโต, เอเกนายตเนน เอกาย ธาตุยา เกหิจิ สมฺปยุตฺโต. กติหิ วิปฺปยุตฺโต, เอเกน ขนฺเธน ทสหายตเนหิ ทสหิ ธาตูหิ วิปฺปยุตฺโต, เอเกนายตเนน เอกาย ธาตุยา เกหิจิ วิปฺปยุตฺโต.}

\proseitem{230}{วิญฺญาณ\-กฺขนฺโธ ตีหิ ขนฺเธหิ สมฺปยุตฺโต, เอเกนายตเนน เอกาย ธาตุยา เกหิจิ สมฺปยุตฺโต. กติหิ วิปฺปยุตฺโต, เอเกน ขนฺเธน ทสหายตเนหิ ทสหิ ธาตูหิ วิปฺปยุตฺโต, เอเกนายตเนน เอกาย ธาตุยา เกหิจิ วิปฺปยุตฺโต.}

\csromanmatikaanchor{mtk.30}
\csromantocmarkref{subsection}{2. อายตน}{mtk.30}
\bookmark[dest={mtk.30},level=2]{2. อายตน}
\csromanheader{2}{2. \textbf{อายตน}}
\proseitem{231}{จกฺขายตนํ ฯเปฯ. โผฏฺฐพฺพา\-ยตนํ ฯเปฯ สมฺปยุตฺตนฺติ นตฺถิ. กติหิ วิปฺปยุตฺตํ, จตูหิ ขนฺเธหิ เอเกนายตเนน สตฺตหิ ธาตูหิ วิปฺปยุตฺตํ, เอเกนายตเนน เอกาย ธาตุยา เกหิจิ วิปฺปยุตฺตํ.}

\proseitem{232}{มนายตนํ ตีหิ ขนฺเธหิ สมฺปยุตฺตํ, เอเกนายตเนน เอกาย ธาตุยา เกหิจิ สมฺปยุตฺตํ. กติหิ วิปฺปยุตฺตํ, เอเกน ขนฺเธน ทสหายตเนหิ ทสหิ ธาตูหิ วิปฺปยุตฺตํ, เอเกนายตเนน เอกาย ธาตุยา เกหิจิ วิปฺปยุตฺตํ.}

\csromanmatikaanchor{mtk.31}
\csromantocmarkref{subsection}{3. ธาตุ}{mtk.31}
\bookmark[dest={mtk.31},level=2]{3. ธาตุ}
\csromanheader{2}{3. ธาตุ}
\proseitem{233}{\csromanfolio{44}จกฺขุธาตุ ฯเปฯ. โผฏฺฐพฺพ\-ธาตุ ฯเปฯ สมฺปยุตฺตาติ นตฺถิ. กติหิ วิปฺปยุตฺตา, จตูหิ ขนฺเธหิ เอเกนายตเนน สตฺตหิ ธาตูหิ วิปฺปยุตฺตา, เอเกนายตเนน เอกาย ธาตุยา เกหิจิ วิปฺปยุตฺตา.}

\proseitem{234}{จกฺขุวิญฺญาณ\-ธาตุ ฯเปฯ. มโนธาตุ. มโนวิญฺญาณ\-ธาตุ ตีหิ ขนฺเธหิ สมฺปยุตฺตา, เอเกนายตเนน เอกาย ธาตุยา เกหิจิ สมฺปยุตฺตา. กติหิ วิปฺปยุตฺตา, เอเกน ขนฺเธน ทสหายตเนหิ โสฬสหิ ธาตูหิ วิปฺปยุตฺตา, เอเกนายตเนน เอกาย ธาตุยา เกหิจิ วิปฺปยุตฺตา.}

\csromanmatikaanchor{mtk.32}
\csromantocmarkref{subsection}{4. สจฺจาทิ}{mtk.32}
\bookmark[dest={mtk.32},level=2]{4. สจฺจาทิ}
\csromanheader{2}{4. \textbf{สจฺจาทิ}}
\proseitem{235}{สมุทยสจฺจํ. มคฺคสจฺจํ ตีหิ ขนฺเธหิ เอเกนายตเนน เอกาย ธาตุยา สมฺปยุตฺตํ, เอเกน ขนฺเธน เอเกนายตเนน เอกาย ธาตุยา เกหิจิ สมฺปยุตฺตํ. กติหิ วิปฺปยุตฺตํ, เอเกน ขนฺเธน ทสหายตเนหิ โสฬสหิ ธาตูหิ วิปฺปยุตฺตํ, เอเกนายตเนน เอกาย ธาตุยา เกหิจิ วิปฺปยุตฺตํ.}

\proseitem{236}{นิโรธสจฺจํ. จกฺขุนฺทฺริยํ ฯเปฯ. กายินฺทฺริยํ. อิตฺถินฺทฺริยํ. ปุริสินฺทฺริยํ ฯเปฯ สมฺปยุตฺตนฺติ นตฺถิ. กติหิ วิปฺปยุตฺตํ, จตูหิ ขนฺเธหิ เอเกนายตเนน สตฺตหิ ธาตูหิ วิปฺปยุตฺตํ, เอเกนายตเนน เอกาย ธาตุยา เกหิจิ วิปฺปยุตฺตํ.}

\proseitem{237}{มนินฺทฺริยํ ตีหิ ขนฺเธหิ สมฺปยุตฺตํ, เอเกนายตเนน เอกาย ธาตุยา เกหิจิ สมฺปยุตฺตํ. กติหิ วิปฺปยุตฺตํ, เอเกน ขนฺเธน ทสหายตเนหิ ทสหิ ธาตูหิ วิปฺปยุตฺตํ, เอเกนายตเนน เอกาย ธาตุยา เกหิจิ วิปฺปยุตฺตํ.}

\proseitem{238}{สุขินฺทฺริยํ. ทุกฺขินฺทฺริยํ. โสมนสฺสิ\-นฺทฺริยํ. โทมนสฺสิ\-นฺทฺริยํ ตีหิ ขนฺเธหิ เอเกนายตเนน เอกาย ธาตุยา สมฺปยุตฺตํ, เอเกนายตเนน เอกาย ธาตุยา เกหิจิ สมฺปยุตฺตํ. กติหิ วิปฺปยุตฺตํ, เอเกน ขนฺเธน ทสหายตเนหิ โสฬสหิ ธาตูหิ วิปฺปยุตฺตํ, เอเกนายตเนน เอกาย ธาตุยา เกหิจิ วิปฺปยุตฺตํ.}

\proseitem{239}{\csromanfolio{45}อุเปกฺขินฺทฺริยํ ตีหิ ขนฺเธหิ เอเกนายตเนน ฉหิ ธาตูหิ สมฺปยุตฺตํ, เอเกนายตเนน เอกาย ธาตุยา เกหิจิ สมฺปยุตฺตํ. กติหิ วิปฺปยุตฺตํ, เอเกน ขนฺเธน ทสหายตเนหิ เอกาทสหิ ธาตูหิ วิปฺปยุตฺตํ, เอเกนายตเนน เอกาย ธาตุยา เกหิจิ วิปฺปยุตฺตํ.}

\proseitem{240}{สทฺธินฺทฺริยํ. วีริยินฺทฺริยํ. สตินฺทฺริยํ. สมาธินฺทฺริยํ. ปญฺญินฺทฺริยํ. อนญฺญาต\-ญฺญสฺสามี\-ตินฺทฺริยํ. อญฺญินฺทฺริยํ. อญฺญาตาวิ\-นฺทฺริยํ. อวิชฺชา. อวิชฺชาปจฺจยา สงฺขารา ตีหิ ขนฺเธหิ เอเกนายตเนน เอกาย ธาตุยา สมฺปยุตฺตา, เอเกน ขนฺเธน เอเกนายตเนน เอกาย ธาตุยา เกหิจิ สมฺปยุตฺตา. กติหิ วิปฺปยุตฺตา, เอเกน ขนฺเธน ทสหายตเนหิ โสฬสหิ ธาตูหิ วิปฺปยุตฺตา, เอเกนายตเนน เอกาย ธาตุยา เกหิจิ วิปฺปยุตฺตา.}

\proseitem{241}{สงฺขาร\-ปจฺจยา วิญฺญาณํ ตีหิ ขนฺเธหิ สมฺปยุตฺตํ, เอเกนายตเนน เอกาย ธาตุยา เกหิจิ สมฺปยุตฺตํ. กติหิ วิปฺปยุตฺตํ, เอเกน ขนฺเธน ทสหายตเนหิ ทสหิ ธาตูหิ วิปฺปยุตฺตํ, เอเกนายตเนน เอกาย ธาตุยา เกหิจิ วิปฺปยุตฺตํ.}

\proseitem{242}{สฬายตน\-ปจฺจยา ผสฺโส ตีหิ ขนฺเธหิ เอเกนายตเนน สตฺตหิ ธาตูหิ สมฺปยุตฺโต, เอเกน ขนฺเธน เอเกนายตเนน เอกาย ธาตุยา เกหิจิ สมฺปยุตฺโต. กติหิ วิปฺปยุตฺโต, เอเกน ขนฺเธน ทสหายตเนหิ ทสหิ ธาตูหิ วิปฺปยุตฺโต, เอเกนายตเนน เอกาย ธาตุยา เกหิจิ วิปฺปยุตฺโต.}

\proseitem{243}{ผสฺสปจฺจยา เวทนา ตีหิ ขนฺเธหิ เอเกนายตเนน สตฺตหิ ธาตูหิ สมฺปยุตฺตา, เอเกนายตเนน เอกาย ธาตุยา เกหิจิ สมฺปยุตฺตา. กติหิ วิปฺปยุตฺตา, เอเกน ขนฺเธน ทสหายตเนหิ ทสหิ ธาตูหิ วิปฺปยุตฺตา, เอเกนายตเนน เอกาย ธาตุยา เกหิจิ วิปฺปยุตฺตา.}

\proseitem{244}{เวทนาปจฺจยา ตณฺหา. ตณฺหาปจฺจยา อุปาทานํ. กมฺมภโว ตีหิ ขนฺเธหิ เอเกนายตเนน เอกาย ธาตุยา สมฺปยุตฺโต, เอเกน ขนฺเธน เอเกนายตเนน เอกาย ธาตุยา เกหิจิ สมฺปยุตฺโต. กติหิ \csromanfolio{46}วิปฺปยุตฺโต, เอเกน ขนฺเธน ทสหายตเนหิ โสฬสหิ ธาตูหิ วิปฺปยุตฺโต, เอเกนายตเนน เอกาย ธาตุยา เกหิจิ วิปฺปยุตฺโต.}

\proseitem{245}{รูปภโว ฯเปฯ สมฺปยุตฺโตติ นตฺถิ. กติหิ วิปฺปยุตฺโต, น เกหิจิ ขนฺเธหิ น เกหิจิ อายตเนหิ ตีหิ ธาตูหิ วิปฺปยุตฺโต.}

\proseitem{246}{อรูปภโว. เนว\-สญฺญา\-นา\-สญฺญาภโว. จตุโวการภโว ฯเปฯ สมฺปยุตฺโตติ นตฺถิ. กติหิ วิปฺปยุตฺโต, เอเกน ขนฺเธน ทสหายตเนหิ โสฬสหิ ธาตูหิ วิปฺปยุตฺโต, เอเกนายตเนน เอกาย ธาตุยา เกหิจิ วิปฺปยุตฺโต.}

\proseitem{247}{อสญฺญาภโว. เอกโวการภโว. ปริเทโว ฯเปฯ สมฺปยุตฺโตติ นตฺถิ. กติหิ วิปฺปยุตฺโต, จตูหิ ขนฺเธหิ เอเกนายตเนน สตฺตหิ ธาตูหิ วิปฺปยุตฺโต, เอเกนายตเนน เอกาย ธาตุยา เกหิจิ วิปฺปยุตฺโต.}

\proseitem{248}{โสโก. ทุกฺขํ. โทมนสฺสํ ตีหิ ขนฺเธหิ เอเกนายตเนน เอกาย ธาตุยา สมฺปยุตฺตํ, เอเกนายตเนน เอกาย ธาตุยา เกหิจิ สมฺปยุตฺตํ. กติหิ วิปฺปยุตฺตํ, เอเกน ขนฺเธน ทสหายตเนหิ โสฬสหิ ธาตูหิ วิปฺปยุตฺตํ, เอเกนายตเนน เอกาย ธาตุยา เกหิจิ วิปฺปยุตฺตํ.}

\proseitem{249}{อุปายาโส. สติปฏฺฐานํ. สมฺมปฺปธานํ ตีหิ ขนฺเธหิ เอเกนายตเนน เอกาย ธาตุยา สมฺปยุตฺตํ, เอเกน ขนฺเธน เอเกนายตเนน เอกาย ธาตุยา เกหิจิ สมฺปยุตฺตํ. กติหิ วิปฺปยุตฺตํ, เอเกน ขนฺเธน ทสหายตเนหิ โสฬสหิ ธาตูหิ วิปฺปยุตฺตํ, เอเกนายตเนน เอกาย ธาตุยา เกหิจิ วิปฺปยุตฺตํ.}

\proseitem{250}{อิทฺธิปาโท ทฺวีหิ ขนฺเธหิ สมฺปยุตฺโต, เอเกน ขนฺเธน เอเกนายตเนน เอกาย ธาตุยา เกหิจิ สมฺปยุตฺโต. กติหิ วิปฺปยุตฺโต, เอเกน ขนฺเธน ทสหายตเนหิ โสฬสหิ ธาตูหิ วิปฺปยุตฺโต, เอเกนายตเนน เอกาย ธาตุยา เกหิจิ วิปฺปยุตฺโต.}

\proseitem{251}{\csromanfolio{47}ฌานํ ทฺวีหิ ขนฺเธหิ เอเกนายตเนน เอกาย ธาตุยา สมฺปยุตฺตํ, เอเกน ขนฺเธน เอเกนายตเนน เอกาย ธาตุยา เกหิจิ สมฺปยุตฺตํ. กติหิ วิปฺปยุตฺตํ, เอเกน ขนฺเธน ทสหายตเนหิ โสฬสหิ ธาตูหิ วิปฺปยุตฺตํ, เอเกนายตเนน เอกาย ธาตุยา เกหิจิ วิปฺปยุตฺตํ.}

\proseitem{252}{อปฺปมญฺญา. ปญฺจินฺทฺริยานิ. ปญฺจ พลานิ. สตฺต โพชฺฌงฺคา. อริโย อฏฺฐงฺคิโก มคฺโค ตีหิ ขนฺเธหิ เอเกนายตเนน เอกาย ธาตุยา สมฺปยุตฺโต, เอเกน ขนฺเธน เอเกนายตเนน เอกาย ธาตุยา เกหิจิ สมฺปยุตฺโต. กติหิ วิปฺปยุตฺโต, เอเกน ขนฺเธน ทสหายตเนหิ โสฬสหิ ธาตูหิ วิปฺปยุตฺโต, เอเกนายตเนน เอกาย ธาตุยา เกหิจิ วิปฺปยุตฺโต.}

\proseitem{253}{ผสฺโส. เจตนา. มนสิกาโร ตีหิ ขนฺเธหิ เอเกนายตเนน สตฺตหิ ธาตูหิ สมฺปยุตฺโต, เอเกน ขนฺเธน เอเกนายตเนน เอกาย ธาตุยา เกหิจิ สมฺปยุตฺโต. กติหิ วิปฺปยุตฺโต, เอเกน ขนฺเธน ทสหายตเนหิ ทสหิ ธาตูหิ วิปฺปยุตฺโต, เอเกนายตเนน เอกาย ธาตุยา เกหิจิ วิปฺปยุตฺโต.}

\proseitem{254}{เวทนา. สญฺญา ตีหิ ขนฺเธหิ เอเกนายตเนน สตฺตหิ ธาตูหิ สมฺปยุตฺตา, เอเกนายตเนน เอกาย ธาตุยา เกหิจิ สมฺปยุตฺตา. กติหิ วิปฺปยุตฺตา, เอเกน ขนฺเธน ทสหายตเนหิ ทสหิ ธาตูหิ วิปฺปยุตฺตา, เอเกนายตเนน เอกาย ธาตุยา เกหิจิ วิปฺปยุตฺตา.}

\proseitem{255}{จิตฺตํ ตีหิ ขนฺเธหิ สมฺปยุตฺตํ, เอเกนายตเนน เอกาย ธาตุยา เกหิจิ สมฺปยุตฺตํ. กติหิ วิปฺปยุตฺตํ, เอเกน ขนฺเธน ทสหายตเนหิ ทสหิ ธาตูหิ วิปฺปยุตฺตํ, เอเกนายตเนน เอกาย ธาตุยา เกหิจิ วิปฺปยุตฺตํ.}

\proseitem{256}{อธิโมกฺโข ตีหิ ขนฺเธหิ เอเกนายตเนน ทฺวีหิ ธาตูหิ สมฺปยุตฺโต, เอเกน ขนฺเธน เอเกนายตเนน เอกาย ธาตุยา เกหิจิ สมฺปยุตฺโต. กติหิ วิปฺปยุตฺโต, เอเกน ขนฺเธน ทสหายตเนหิ ปนฺนรสหิ ธาตุหิ วิปฺปยุตฺโต, เอเกนายตเนน เอกาย ธาตุยา เกหิจิ วิปฺปยุตฺโต.}

\csromanmatikaanchor{mtk.33}
\csromantocmarkref{subsection}{5. ติก}{mtk.33}
\bookmark[dest={mtk.33},level=2]{5. ติก}
\csromanheader{2}{5. \textbf{ติก}}
\proseitem{257}{\csromanfolio{48}กุสลา ธมฺมา. อกุสลา ธมฺมา กติหิ ขนฺเธหิ กติหายตเนหิ กติหิ ธาตูหิ สมฺปยุตฺตาติ นตฺถิ. กติหิ วิปฺปยุตฺตา, เอเกน ขนฺเธน ทสหายตเนหิ โสฬสหิ ธาตูหิ วิปฺปยุตฺตา, เอเกนายตเนน เอกาย ธาตุยา เกหิจิ วิปฺปยุตฺตา.}

\proseitem{258}{สุขาย เวทนาย สมฺปยุตฺตา ธมฺมา. ทุกฺขาย เวทนาย สมฺปยุตฺตา ธมฺมา เอเกน ขนฺเธน สมฺปยุตฺตา, เอเกนายตเนน เอกาย ธาตุยา เกหิจิ สมฺปยุตฺตา. กติหิ วิปฺปยุตฺตา, เอเกน ขนฺเธน ทสหายตเนหิ ปนฺนรสหิ ธาตูหิ วิปฺปยุตฺตา, เอเกนายตเนน เอกาย ธาตุยา เกหิจิ วิปฺปยุตฺตา.}

\proseitem{259}{อทุกฺขม\-สุขาย เวทนาย สมฺปยุตฺตา ธมฺมา เอเกน ขนฺเธน สมฺปยุตฺตา, เอเกนายตเนน เอกาย ธาตุยา เกหิจิ สมฺปยุตฺตา. กติหิ วิปฺปยุตฺตา, เอเกน ขนฺเธน ทสหายตเนหิ เอกาทสหิ ธาตูหิ วิปฺปยุตฺตา, เอเกนายตเนน เอกาย ธาตุยา เกหิจิ วิปฺปยุตฺตา.}

\proseitem{260}{วิปากา ธมฺมา ฯเปฯ สมฺปยุตฺตาติ นตฺถิ. กติหิ วิปฺปยุตฺตา, เอเกน ขนฺเธน ทสหายตเนหิ ทสหิ ธาตูหิ วิปฺปยุตฺตา, เอเกนายตเนน เอกาย ธาตุยา เกหิจิ วิปฺปยุตฺตา.}

\proseitem{261}{วิปาก\-ธมฺม\-ธมฺมา. สํกิลิฏฺฐ\-สํกิเลสิกา ธมฺมา ฯเปฯ สมฺปยุตฺตาติ นตฺถิ. กติหิ วิปฺปยุตฺตา, เอเกน ขนฺเธน ทสหายตเนหิ โสฬสหิ ธาตูหิ วิปฺปยุตฺตา, เอเกนายตเนน เอกาย ธาตุยา เกหิจิ วิปฺปยุตฺตา.}

\proseitem{262}{เนว\-วิปาก\-น\-วิปากธมฺม\-ธมฺมา. อนุปาทินฺนุ\-ปาทานิยา ธมฺมา ฯเปฯ สมฺปยุตฺตาติ นตฺถิ. กติหิ วิปฺปยุตฺตา, น เกหิจิ ขนฺเธหิ น เกหิจิ อายตเนหิ ปญฺจหิ ธาตูหิ วิปฺปยุตฺตา.}

\proseitem{263}{อนุปาทินฺน\-อนุปาทานิยา ธมฺมา. อสํกิลิฏฺฐ\-อสํกิเลสิกา ธมฺมา ฯเปฯ สมฺปยุตฺตาติ นตฺถิ. กติหิ วิปฺปยุตฺตา, น เกหิจิ ขนฺเธหิ น เกหิจิ อายตเนหิ ฉหิ ธาตูหิ วิปฺปยุตฺตา.}

\proseitem{264}{\csromanfolio{49}สวิตกฺก\-สวิจารา ธมฺมา เอเกน ขนฺเธน เอเกนายตเนน เอกาย ธาตุยา เกหิจิ สมฺปยุตฺตา. กติหิ วิปฺปยุตฺตา, เอเกน ขนฺเธน ทสหายตเนหิ ปนฺนรสหิ ธาตูหิ วิปฺปยุตฺตา, เอเกนายตเนน เอกาย ธาตุยา เกหิจิ วิปฺปยุตฺตา.}

\proseitem{265}{อวิตกฺก\-วิจารมตฺตา ธมฺมา. ปีติสหคตา ธมฺมา เอเกน ขนฺเธน เอเกนายตเนน เอกาย ธาตุยา เกหิจิ สมฺปยุตฺตา. กติหิ วิปฺปยุตฺตา, เอเกน ขนฺเธน ทสหายตเนหิ โสฬสหิ ธาตูหิ วิปฺปยุตฺตา, เอเกนายตเนน เอกาย ธาตุยา เกหิจิ วิปฺปยุตฺตา.}

\proseitem{266}{อวิตกฺก\-อวิจารา ธมฺมา ฯเปฯ สมฺปยุตฺตาติ นตฺถิ. กติหิ วิปฺปยุตฺตา, น เกหิจิ ขนฺเธหิ น เกหิจิ อายตเนหิ เอกาย ธาตุยา วิปฺปยุตฺตา.}

\proseitem{267}{สุขสหคตา ธมฺมา เอเกน ขนฺเธน สมฺปยุตฺตา, เอเกนายตเนน เอกาย ธาตุยา เกหิจิ สมฺปยุตฺตา. กติหิ วิปฺปยุตฺตา, เอเกน ขนฺเธน ทสหายตเนหิ ปนฺนรสหิ ธาตูหิ วิปฺปยุตฺตา, เอเกนายตเนน เอกาย ธาตุยา เกหิจิ วิปฺปยุตฺตา.}

\proseitem{268}{อุเปกฺขา\-สหคตา ธมฺมา เอเกน ขนฺเธน สมฺปยุตฺตา, เอเกนายตเนน เอกาย ธาตุยา เกหิจิ สมฺปยุตฺตา. กติหิ วิปฺปยุตฺตา, เอเกน ขนฺเธน ทสหายตเนหิ เอกาทสหิ ธาตูหิ วิปฺปยุตฺตา, เอเกนายตเนน เอกาย ธาตุยา เกหิจิ วิปฺปยุตฺตา.}

\proseitem{269}{ทสฺสเนน ปหาตพฺพา ธมฺมา. ภาวนาย ปหาตพฺพา ธมฺมา. ทสฺสเนน ปหาตพฺพ\-เหตุกา ธมฺมา. ภาวนาย ปหาตพฺพ\-เหตุกา ธมฺมา. อาจยคามิโน ธมฺมา. อปจยคามิโน ธมฺมา. เสกฺขา ธมฺมา. อเสกฺขา ธมฺมา. มหคฺคตา ธมฺมา ฯเปฯ สมฺปยุตฺตาติ นตฺถิ. กติหิ วิปฺปยุตฺตา, เอเกน ขนฺเธน ทสหายตเนหิ โสฬสหิ ธาตูหิ วิปฺปยุตฺตา, เอเกนายตเนน เอกาย ธาตุยา เกหิจิ วิปฺปยุตฺตา.}

\proseitem{270}{อปฺปมาณา ธมฺมา. ปณีตา ธมฺมา ฯเปฯ สมฺปยุตฺตาติ นตฺถิ. กติหิ วิปฺปยุตฺตา, น เกหิจิ ขนฺเธหิ น เกหิจิ อายตเนหิ ฉหิ ธาตูหิ วิปฺปยุตฺตา.}

\proseitem{271}{\csromanfolio{50}ปริตฺตารมฺมณา ธมฺมา ฯเปฯ สมฺปยุตฺตาติ นตฺถิ. กติหิ วิปฺปยุตฺตา, เอเกน ขนฺเธน ทสหายตเนหิ ทสหิ ธาตูหิ วิปฺปยุตฺตา, เอเกนายตเนน เอกาย ธาตุยา เกหิจิ วิปฺปยุตฺตา.}

\proseitem{272}{มหคฺคตา\-รมฺมณา ธมฺมา. อปฺปมาณา\-รมฺมณา ธมฺมา. หีนา ธมฺมา. มิจฺฉตฺต\-นิยตา ธมฺมา. สมฺมตฺตนิยตา ธมฺมา. มคฺคารมฺมณา ธมฺมา. มคฺคเหตุกา ธมฺมา. มคฺคาธิปติโน ธมฺมา ฯเปฯ สมฺปยุตฺตาติ นตฺถิ. กติหิ วิปฺปยุตฺตา, เอเกน ขนฺเธน ทสหายตเนหิ โสฬสหิ ธาตูหิ วิปฺปยุตฺตา, เอเกนายตเนน เอกาย ธาตุยา เกหิจิ วิปฺปยุตฺตา.}

\proseitem{273}{อนุปฺปนฺนา ธมฺมา ฯเปฯ สมฺปยุตฺตาติ นตฺถิ. กติหิ วิปฺปยุตฺตา, น เกหิจิ ขนฺเธหิ น เกหิจิ อายตเนหิ ปญฺจหิ ธาตูหิ วิปฺปยุตฺตา.}

\proseitem{274}{อตีตารมฺมณา ธมฺมา. อนาคตารมฺมณา ธมฺมา ฯเปฯ สมฺปยุตฺตาติ นตฺถิ. กติหิ วิปฺปยุตฺตา, เอเกน ขนฺเธน ทสหายตเนหิ โสฬสหิ ธาตูหิ วิปฺปยุตฺตา, เอเกนายตเนน เอกาย ธาตุยา เกหิจิ วิปฺปยุตฺตา.}

\proseitem{275}{ปจฺจุปฺปนฺนา\-รมฺมณา ธมฺมา. อชฺฌตฺตา\-รมฺมณา ธมฺมา. พหิทฺธา\-รมฺมณา ธมฺมา. อชฺฌตฺตพหิทฺธา\-รมฺมณา ธมฺมา ฯเปฯ สมฺปยุตฺตาติ นตฺถิ. กติหิ วิปฺปยุตฺตา, เอเกน ขนฺเธน ทสหายตเนหิ ทสหิ ธาตูหิ วิปฺปยุตฺตา, เอเกนายตเนน เอกาย ธาตุยา เกหิจิ วิปฺปยุตฺตา.}

\proseitem{276}{สนิทสฺสน\-สปฺปฏิฆา ธมฺมา. อนิทสฺสน\-สปฺปฏิฆา ธมฺมา ฯเปฯ สมฺปยุตฺตาติ นตฺถิ. กติหิ วิปฺปยุตฺตา, จตูหิ ขนฺเธหิ เอเกนายตเนน สตฺตหิ ธาตูหิ วิปฺปยุตฺตา, เอเกนายตเนน เอกาย ธาตุยา เกหิจิ วิปฺปยุตฺตา.}

\csromanmatikaanchor{mtk.34}
\csromantocmarkref{subsection}{6. ทุก}{mtk.34}
\bookmark[dest={mtk.34},level=2]{6. ทุก}
\csromanheader{2}{6. \textbf{ทุก}}
\proseitem{277}{เหตู ธมฺมา. เหตู เจว สเหตุกา จ ธมฺมา. เหตู เจว เหตุสมฺปยุตฺตา จ ธมฺมา ตีหิ ขนฺเธหิ เอเกนายตเนน เอกาย ธาตุยา สมฺปยุตฺตา, เอเกน ขนฺเธน เอเกนายตเนน เอกาย ธาตุยา เกหิจิ สมฺปยุตฺตา. กติหิ วิปฺปยุตฺตา, เอเกน ขนฺเธน ทสหายตเนหิ \csromanfolio{51}โสฬสหิ ธาตูหิ วิปฺปยุตฺตา, เอเกนายตเนน เอกาย ธาตุยา เกหิจิ วิปฺปยุตฺตา.}

\proseitem{278}{สเหตุกา ธมฺมา. เหตุสมฺปยุตฺตา ธมฺมา ฯเปฯ สมฺปยุตฺตาติ นตฺถิ. กติหิ วิปฺปยุตฺตา, เอเกน ขนฺเธน ทสหายตเนหิ โสฬสหิ ธาตูหิ วิปฺปยุตฺตา, เอเกนายตเนน เอกาย ธาตุยา เกหิจิ วิปฺปยุตฺตา.}

\proseitem{279}{สเหตุกา เจว น จ เหตู ธมฺมา. เหตุสมฺปยุตฺตา เจว น จ เหตู ธมฺมา. น เหตุสเหตุกา ธมฺมา เอเกน ขนฺเธน เอเกนายตเนน เอกาย ธาตุยา เกหิจิ สมฺปยุตฺตา. กติหิ วิปฺปยุตฺตา, เอเกน ขนฺเธน ทสหายตเนหิ โสฬสหิ ธาตูหิ วิปฺปยุตฺตา, เอเกนายตเนน เอกาย ธาตุยา เกหิจิ วิปฺปยุตฺตา.}

\proseitem{280}{อปฺปจฺจยา ธมฺมา. อสงฺขตา ธมฺมา. สนิทสฺสนา ธมฺมา. สปฺปฏิฆา ธมฺมา. รูปิโน ธมฺมา ฯเปฯ สมฺปยุตฺตาติ นตฺถิ. กติหิ วิปฺปยุตฺตา, จตูหิ ขนฺเธหิ เอเกนายตเนน สตฺตหิ ธาตูหิ วิปฺปยุตฺตา, เอเกนายตเนน เอกาย ธาตุยา เกหิจิ วิปฺปยุตฺตา.}

\proseitem{281}{โลกุตฺตรา ธมฺมา ฯเปฯ สมฺปยุตฺตาติ นตฺถิ. กติหิ วิปฺปยุตฺตา, น เกหิจิ ขนฺเธหิ น เกหิจิ อายตเนหิ ฉหิ ธาตูหิ วิปฺปยุตฺตา.}

\proseitem{282}{อาสวา ธมฺมา. อาสวา เจว สาสวา จ ธมฺมา. อาสวา เจว อาสว\-สมฺปยุตฺตา จ ธมฺมา ตีหิ ขนฺเธหิ เอเกนายตเนน เอกาย ธาตุยา สมฺปยุตฺตา, เอเกน ขนฺเธน เอเกนายตเนน เอกาย ธาตุยา เกหิจิ สมฺปยุตฺตา. กติหิ วิปฺปยุตฺตา, เอเกน ขนฺเธน ทสหายตเนหิ โสฬสหิ ธาตูหิ วิปฺปยุตฺตา, เอเกนายตเนน เอกาย ธาตุยา เกหิจิ วิปฺปยุตฺตา.}

\proseitem{283}{อนาสวา ธมฺมา. อาสว\-วิปฺปยุตฺตา อนาสวา ธมฺมา ฯเปฯ สมฺปยุตฺตาติ นตฺถิ. กติหิ วิปฺปยุตฺตา, น เกหิจิ ขนฺเธหิ น เกหิจิ อายตเนหิ ฉหิ ธาตูหิ วิปฺปยุตฺตา.}

\proseitem{284}{\csromanfolio{52}อาสว\-สมฺปยุตฺตา ธมฺมา ฯเปฯ สมฺปยุตฺตาติ นตฺถิ. กติหิ วิปฺปยุตฺตา, เอเกน ขนฺเธน ทสหายตเนหิ โสฬสหิ ธาตูหิ วิปฺปยุตฺตา, เอเกนายตเนน เอกาย ธาตุยา เกหิจิ วิปฺปยุตฺตา.}

\proseitem{285}{อาสว\-สมฺปยุตฺตา เจว โน จ อาสวา ธมฺมา เอเกน ขนฺเธน เอเกนายตเนน เอกาย ธาตุยา เกหิจิ สมฺปยุตฺตา. กติหิ วิปฺปยุตฺตา, เอเกน ขนฺเธน ทสหายตเนหิ โสฬสหิ ธาตูหิ วิปฺปยุตฺตา, เอเกนายตเนน เอกาย ธาตุยา เกหิจิ วิปฺปยุตฺตา.}

\proseitem{286}{สํโยชนา ธมฺมา. คนฺถา ธมฺมา. โอฆา ธมฺมา. โยคา ธมฺมา. นีวรณา ธมฺมา. ปรามาสา ธมฺมา. ปรามาสา เจว ปรามฏฺฐา จ ธมฺมา ตีหิ ขนฺเธหิ เอเกนายตเนน เอกาย ธาตุยา สมฺปยุตฺตา, เอเกน ขนฺเธน เอเกนายตเนน เอกาย ธาตุยา เกหิจิ สมฺปยุตฺตา. กติหิ วิปฺปยุตฺตา, เอเกน ขนฺเธน ทสหายตเนหิ โสฬสหิ ธาตูหิ วิปฺปยุตฺตา, เอเกนายตเนน เอกาย ธาตุยา เกหิจิ วิปฺปยุตฺตา.}

\proseitem{287}{อปรามฏฺฐา ธมฺมา. ปรามาส\-วิปฺปยุตฺตา อปรามฏฺฐา ธมฺมา สมฺปยุตฺตาติ นตฺถิ. กติหิ วิปฺปยุตฺตา, น เกหิจิ ขนฺเธหิ น เกหิจิ อายตเนหิ ฉหิ ธาตูหิ วิปฺปยุตฺตา.}

\proseitem{288}{ปรามาส\-สมฺปยุตฺตา ธมฺมา เอเกน ขนฺเธน เอเกนายตเนน เอกาย ธาตุยา เกหิจิ สมฺปยุตฺตา. กติหิ วิปฺปยุตฺตา, เอเกน ขนฺเธน ทสหายตเนหิ โสฬสหิ ธาตูหิ วิปฺปยุตฺตา, เอเกนายตเนน เอกาย ธาตุยา เกหิจิ วิปฺปยุตฺตา.}

\proseitem{289}{สารมฺมณา ธมฺมา ฯเปฯ สมฺปยุตฺตาติ นตฺถิ. กติหิ วิปฺปยุตฺตา, เอเกน ขนฺเธน ทสหายตเนหิ ทสหิ ธาตูหิ วิปฺปยุตฺตา, เอเกนายตเนน เอกาย ธาตุยา เกหิจิ วิปฺปยุตฺตา.}

\proseitem{290}{อนารมฺมณา ธมฺมา. จิตฺต\-วิปฺปยุตฺตา ธมฺมา. จิตฺต\-วิสํสฏฺฐา ธมฺมา. อุปาทา ธมฺมา ฯเปฯ สมฺปยุตฺตาติ นตฺถิ. กติหิ วิปฺปยุตฺตา, จตูหิ ขนฺเธหิ เอเกนายตเนน สตฺตหิ ธาตูหิ วิปฺปยุตฺตา, เอเกนายตเนน เอกาย ธาตุยา เกหิจิ วิปฺปยุตฺตา.}

\proseitem{291}{\csromanfolio{53}จิตฺตา ธมฺมา ตีหิ ขนฺเธหิ สมฺปยุตฺตา, เอเกนายตเนน เอกาย ธาตุยา เกหิจิ สมฺปยุตฺตา. กติหิ วิปฺปยุตฺตา, เอเกน ขนฺเธน ทสหายตเนหิ ทสหิ ธาตูหิ วิปฺปยุตฺตา, เอเกนายตเนน เอกาย ธาตุยา เกหิจิ วิปฺปยุตฺตา.}

\proseitem{292}{เจตสิกา ธมฺมา. จิตฺต\-สมฺปยุตฺตา ธมฺมา. จิตฺตสํสฏฺฐา ธมฺมา. จิตฺต\-สํสฏฺฐ\-สมุฏฺฐานา ธมฺมา. จิตฺต\-สํสฏฺฐ\-สมุฏฺฐาน\-สหภุโน ธมฺมา. จิตฺต\-สํสฏฺฐ\-สมุฏฺฐานา\-นุปริวตฺติโน ธมฺมา เอเกน ขนฺเธน เอเกนายตเนน สตฺตหิ ธาตูหิ สมฺปยุตฺตา. กติหิ วิปฺปยุตฺตา, เอเกน ขนฺเธน ทสหายตเนหิ ทสหิ ธาตูหิ วิปฺปยุตฺตา, เอเกนายตเนน เอกาย ธาตุยา เกหิจิ วิปฺปยุตฺตา.}

\proseitem{293}{อนุปาทินฺนา ธมฺมา ฯเปฯ สมฺปยุตฺตาติ นตฺถิ. กติหิ วิปฺปยุตฺตา, น เกหิจิ ขนฺเธหิ น เกหิจิ อายตเนหิ ปญฺจหิ ธาตูหิ วิปฺปยุตฺตา.}

\proseitem{294}{อุปาทานา ธมฺมา. กิเลสา ธมฺมา. กิเลสา เจว สํกิเลสิกา จ ธมฺมา. กิเลสา เจว สํกิลิฏฺฐา จ ธมฺมา. กิเลสา เจว กิเลส\-สมฺปยุตฺตา จ ธมฺมา ตีหิ ขนฺเธหิ เอเกนายตเนน เอกาย ธาตุยา สมฺปยุตฺตา, เอเกน ขนฺเธน เอเกนายตเนน เอกาย ธาตุยา เกหิจิ สมฺปยุตฺตา. กติหิ วิปฺปยุตฺตา, เอเกน ขนฺเธน ทสหายตเนหิ โสฬสหิ ธาตูหิ วิปฺปยุตฺตา, เอเกนายตเนน เอกาย ธาตุยา เกหิจิ วิปฺปยุตฺตา.}

\proseitem{295}{อสํกิเลสิกา ธมฺมา. กิเลส\-วิปฺปยุตฺตา อสํกิเลสิกา ธมฺมา ฯเปฯ สมฺปยุตฺตาติ นตฺถิ. กติหิ วิปฺปยุตฺตา, น เกหิจิ ขนฺเธหิ น เกหิจิ อายตเนหิ ฉหิ ธาตูหิ วิปฺปยุตฺตา.}

\proseitem{296}{สํกิลิฏฺฐา ธมฺมา. กิเลส\-สมฺปยุตฺตา ธมฺมา ฯเปฯ สมฺปยุตฺตาติ นตฺถิ. กติหิ วิปฺปยุตฺตา, เอเกน ขนฺเธน ทสหายตเนหิ โสฬสหิ ธาตูหิ วิปฺปยุตฺตา, เอเกนายตเนน เอกาย ธาตุยา เกหิจิ วิปฺปยุตฺตา.}

\proseitem{297}{สํกิลิฏฺฐา เจว โน จ กิเลสา ธมฺมา. กิเลส\-สมฺปยุตฺตา เจว โน จ กิเลสา ธมฺมา เอเกน ขนฺเธน เอเกนายตเนน เอกาย \csromanfolio{54}ธาตุยา เกหิจิ สมฺปยุตฺตา. กติหิ วิปฺปยุตฺตา, เอเกน ขนฺเธน ทสหายตเนหิ โสฬสหิ ธาตูหิ วิปฺปยุตฺตา, เอเกนายตเนน เอกาย ธาตุยา เกหิจิ วิปฺปยุตฺตา.}

\proseitem{298}{ทสฺสเนน ปหาตพฺพา ธมฺมา. ภาวนาย ปหาตพฺพา ธมฺมา ทสฺสเนน ปหาตพฺพ\-เหตุกา ธมฺมา. ภาวนาย ปหาตพฺพ\-เหตุกา ธมฺมา สมฺปยุตฺตาติ นตฺถิ. กติหิ วิปฺปยุตฺตา, เอเกน ขนฺเธน ทสหายตเนหิ โสฬสหิ ธาตูหิ วิปฺปยุตฺตา เอเกนายตเนน เอกาย ธาตุยา เกหิจิ วิปฺปยุตฺตา.}

\proseitem{299}{สวิตกฺกา ธมฺมา. สวิจารา ธมฺมา เอเกน ขนฺเธน เอเกนายตเนน เอกาย ธาตุยา เกหิจิ สมฺปยุตฺตา. กติหิ วิปฺปยุตฺตา, เอเกน ขนฺเธน ทสหายตเนหิ ปนฺนรสหิ ธาตูหิ วิปฺปยุตฺตา, เอเกนายตเนน เอกาย ธาตุยา เกหิจิ วิปฺปยุตฺตา.}

\proseitem{300}{อวิตกฺกา ธมฺมา. อวิจารา ธมฺมา ฯเปฯ สมฺปยุตฺตาติ นตฺถิ. กติหิ วิปฺปยุตฺตา, น เกหิจิ ขนฺเธหิ น เกหิจิ อายตเนหิ เอกาย ธาตุยา วิปฺปยุตฺตา.}

\proseitem{301}{สปฺปีติกา ธมฺมา. ปีติสหคตา ธมฺมา เอเกน ขนฺเธน เอเกนายตเนน เอกาย ธาตุยา เกหิจิ สมฺปยุตฺตา. กติหิ วิปฺปยุตฺตา, เอเกน ขนฺเธน ทสหายตเนหิ โสฬสหิ ธาตูหิ วิปฺปยุตฺตา, เอเกนายตเนน เอกาย ธาตุยา เกหิจิ วิปฺปยุตฺตา.}

\proseitem{302}{สุขสหคตา ธมฺมา เอเกน ขนฺเธน สมฺปยุตฺตา, เอเกนายตเนน เอกาย ธาตุยา เกหิจิ สมฺปยุตฺตา. กติหิ วิปฺปยุตฺตา, เอเกน ขนฺเธน ทสหายตเนหิ ปนฺนรสหิ ธาตูหิ วิปฺปยุตฺตา, เอเกนายตเนน เอกาย ธาตุยา เกหิจิ วิปฺปยุตฺตา.}

\proseitem{303}{อุเปกฺขา\-สหคตา ธมฺมา เอเกน ขนฺเธน สมฺปยุตฺตา, เอเกนายตเนน เอกาย ธาตุยา เกหิจิ สมฺปยุตฺตา. กติหิ วิปฺปยุตฺตา, เอเกน ขนฺเธน ทสหายตเนหิ เอกาทสหิ ธาตูหิ วิปฺปยุตฺตา, เอเกนายตเนน เอกาย ธาตุยา เกหิจิ วิปฺปยุตฺตา.}

\proseitem{304}{\csromanfolio{55}น กามาวจรา ธมฺมา. อปริยาปนฺนา ธมฺมา. อนุตฺตรา ธมฺมา ฯเปฯ สมฺปยุตฺตาติ นตฺถิ. กติหิ วิปฺปยุตฺตา, น เกหิจิ ขนฺเธหิ น เกหิจิ อายตเนหิ ฉหิ ธาตูหิ วิปฺปยุตฺตา.}

\proseitem{305}{รูปาวจรา ธมฺมา. อรูปาวจรา ธมฺมา. นิยฺยานิกา ธมฺมา. นิยตา ธมฺมา. สรณา ธมฺมา กติหิ ขนฺเธหิ กติหายตเนหิ กติหิ ธาตูหิ สมฺปยุตฺตาติ นตฺถิ. กติหิ วิปฺปยุตฺตา, เอเกน ขนฺเธน ทสหายตเนหิ โสฬสหิ ธาตูหิ วิปฺปยุตฺตา, เอเกนายตเนน เอกาย ธาตุยา เกหิจิ วิปฺปยุตฺตา.}

\setlength{\csromangathapagewidth}{0pt}
\setlength{\csromangathawakwidth}{0pt}
\csromangathameasuregroup{ธมฺมายตนํ ธมฺมธาตุ, \\ สฬายตนํ นามรูปํ, \\ ชาติ ชรา จ มรณํ, \\ โคจฺฉเกสุ จ ปญฺญาส, \\ มหนฺตเร ปนฺนรส, \\ เตวีส ปทสตํ เอตํ,}{\csromangathabat{ธมฺมายตนํ ธมฺมธาตุ,}{ทุกฺขสจฺจญฺจ ชีวิตํ.} \\ \csromangathabat{สฬายตนํ นามรูปํ,}{จตฺตาโร จ มหาภวา.} \\ \csromangathabat{ชาติ ชรา จ มรณํ,}{ติเกเสฺวกูนวีสติ.} \\ \csromangathabat{โคจฺฉเกสุ จ ปญฺญาส,}{อฏฺฐ จูฬนฺตเร ปทา.} \\ \csromangathabat{มหนฺตเร ปนฺนรส,}{อฏฺฐารส ตโต ปเร.} \\ \csromangathabat{เตวีส ปทสตํ เอตํ,}{สมฺปโยเค น ลพฺภตีติ.}}
\csromangathasetleft{ธมฺมายตนํ ธมฺมธาตุ, \\ สฬายตนํ นามรูปํ, \\ ชาติ ชรา จ มรณํ, \\ โคจฺฉเกสุ จ ปญฺญาส, \\ มหนฺตเร ปนฺนรส, \\ เตวีส ปทสตํ เอตํ,}
\csromangathagroup{\csromangathastanza{\csromangathabat{ธมฺมายตนํ ธมฺมธาตุ,}{ทุกฺขสจฺจญฺจ ชีวิตํ.} \\ \csromangathabat{สฬายตนํ นามรูปํ,}{จตฺตาโร จ มหาภวา.}}\csromangathastanza{\csromangathabat{ชาติ ชรา จ มรณํ,}{ติเกเสฺวกูนวีสติ.} \\ \csromangathabat{โคจฺฉเกสุ จ ปญฺญาส,}{อฏฺฐ จูฬนฺตเร ปทา.}}\csromangathastanza{\csromangathabat{มหนฺตเร ปนฺนรส,}{อฏฺฐารส ตโต ปเร.} \\ \csromangathabat{เตวีส ปทสตํ เอตํ,}{สมฺปโยเค น ลพฺภตีติ.}}}

\vspace{\tipitakaparskip}
\csromancenter{สมฺปโยค\-วิปฺปโยค\-ปท\-นิทฺเทโส ฉฏฺโฐ.}
\csromanmatikaanchor{mtk.35}
\csromantocmarknopage{section}{7. สตฺตมนย}{mtk.35}
\bookmark[dest={mtk.35},level=1]{7. สตฺตมนย}
\csromanheader{1}{7. \textbf{สตฺตมนย}}
\csromanmatikaanchor{mtk.36}
\csromantocmarkref{subsection}{7. สมฺปยุตฺเตนวิปฺปยุตฺตปทนิทฺเทส}{mtk.36}
\bookmark[dest={mtk.36},level=2]{7. สมฺปยุตฺเตนวิปฺปยุตฺตปทนิทฺเทส}
\csromanheader{2}{7. สมฺปยุตฺเตน\-วิปฺปยุตฺต\-ปท\-นิทฺเทส}
\proseitem{306}{\csromanfolio{56}เวทนา\-กฺขนฺเธน เย ธมฺมา. สญฺญา\-กฺขนฺเธน เย ธมฺมา. สงฺขาร\-กฺขนฺเธน เย ธมฺมา. วิญฺญาณ\-กฺขนฺเธน เย ธมฺมา. มนายตเนน เย ธมฺมา สมฺปยุตฺตา, เตหิ ธมฺเมหิ เย ธมฺมา วิปฺปยุตฺตา. เต ธมฺมา กติหิ ขนฺเธหิ กติหายตเนหิ กติหิ ธาตูหิ วิปฺปยุตฺตา, เต ธมฺมา จตูหิ ขนฺเธหิ เอเกนายตเนน สตฺตหิ ธาตูหิ วิปฺปยุตฺตา, เอเกนายตเนน เอกาย ธาตุยา เกหิจิ วิปฺปยุตฺตา.}

\proseitem{307}{จกฺขุวิญฺญาณ\-ธาตุยา เย ธมฺมา ฯเปฯ. มโนธาตุยา เย ธมฺมา. มโนวิญฺญาณ\-ธาตุยา เย ธมฺมา สมฺปยุตฺตา, เตหิ ธมฺเมหิ เย ธมฺมา วิปฺปยุตฺตา. เต ธมฺมา น เกหิจิ ขนฺเธหิ น เกหิจิ อายตเนหิ เอกาย ธาตุยา วิปฺปยุตฺตา.}

\proseitem{308}{มนินฺทฺริเยน เย ธมฺมา สมฺปยุตฺตา, เตหิ ธมฺเมหิ เย ธมฺมา วิปฺปยุตฺตา. เต ธมฺมา จตูหิ ขนฺเธหิ เอเกนายตเนน สตฺตหิ ธาตูหิ วิปฺปยุตฺตา, เอเกนายตเนน เอกาย ธาตุยา เกหิจิ วิปฺปยุตฺตา.}

\proseitem{309}{อุเปกฺขิ\-นฺทฺริเยน เย ธมฺมา สมฺปยุตฺตา, เตหิ ธมฺเมหิ เย ธมฺมา วิปฺปยุตฺตา. เต ธมฺมา น เกหิจิ ขนฺเธหิ น เกหิจิ อายตเนหิ ปญฺจหิ ธาตูหิ วิปฺปยุตฺตา.}

\proseitem{310}{สงฺขาร\-ปจฺจยา วิญฺญาเณน เย ธมฺมา. สฬายตน\-ปจฺจยา ผสฺเสน เย ธมฺมา. ผสฺสปจฺจยา เวทนาย เย ธมฺมา. ผสฺเสน เย ธมฺมา. เวทนาย เย ธมฺมา. สญฺญาย เย ธมฺมา. เจตนาย เย ธมฺมา. จิตฺเตน เย ธมฺมา. มนสิกาเรน เย ธมฺมา สมฺปยุตฺตา, เตหิ ธมฺเมหิ เย ธมฺมา วิปฺปยุตฺตา. เต ธมฺมา จตูหิ ขนฺเธหิ เอเกนายตเนน สตฺตหิ ธาตูหิ วิปฺปยุตฺตา, เอเกนายตเนน เอกาย ธาตุยา เกหิจิ วิปฺปยุตฺตา.}

\proseitem{311}{อธิโมกฺเขน เย ธมฺมา สมฺปยุตฺตา, เตหิ ธมฺเมหิ เย ธมฺมา วิปฺปยุตฺตา. เต ธมฺมา น เกหิจิ ขนฺเธหิ น เกหิจิ อายตเนหิ เอกาย ธาตุยา วิปฺปยุตฺตา.}

\csromanheader{2}{7. สตฺตมนย}
\proseitem{312}{\csromanfolio{57}อทุกฺขม\-สุขาย เวทนาย สมฺปยุตฺเตหิ ธมฺเมหิ เย ธมฺมา. อุเปกฺขา\-สหคเตหิ ธมฺเมหิ เย ธมฺมา สมฺปยุตฺตา, เตหิ ธมฺเมหิ เย ธมฺมา วิปฺปยุตฺตา. เต ธมฺมา น เกหิจิ ขนฺเธหิ น เกหิจิ อายตเนหิ ปญฺจหิ ธาตูหิ วิปฺปยุตฺตา.}

\proseitem{313}{สวิตกฺก\-สวิจาเรหิ ธมฺเมหิ เย ธมฺมา สมฺปยุตฺตา, เตหิ ธมฺเมหิ เย ธมฺมา วิปฺปยุตฺตา. เต ธมฺมา น เกหิจิ ขนฺเธหิ น เกหิจิ อายตเนหิ เอกาย ธาตุยา วิปฺปยุตฺตา.}

\proseitem{314}{จิตฺเตหิ ธมฺเมหิ เย ธมฺมา. เจตสิเกหิ ธมฺเมหิ เย ธมฺมา. จิตฺต\-สมฺปยุตฺเตหิ ธมฺเมหิ เย ธมฺมา. จิตฺต\-สํสฏฺเฐหิ ธมฺเมหิ เย ธมฺมา. จิตฺต\-สํสฏฺฐ\-สมุฏฺฐาเนหิ ธมฺเมหิ เย ธมฺมา. จิตฺต\-สํสฏฺฐ\-สมุฏฺฐาน\-สหภูหิ ธมฺเมหิ เย ธมฺมา. จิตฺต\-สํสฏฺฐ\-สมุฏฺฐานา\-นุปริวตฺตีหิ ธมฺเมหิ เย ธมฺมา สมฺปยุตฺตา, เตหิ ธมฺเมหิ เย ธมฺมา วิปฺปยุตฺตา. เต ธมฺมา จตูหิ ขนฺเธหิ เอเกนายตเนน สตฺตหิ ธาตูหิ วิปฺปยุตฺตา, เอเกนายตเนน เอกาย ธาตุยา เกหิจิ วิปฺปยุตฺตา.}

\proseitem{315}{สวิตกฺเกหิ ธมฺเมหิ เย ธมฺมา. สวิจาเรหิ ธมฺเมหิ เย ธมฺมา สมฺปยุตฺตา, เตหิ ธมฺเมหิ เย ธมฺมา วิปฺปยุตฺตา. เต ธมฺมา น เกหิจิ ขนฺเธหิ น เกหิจิ อายตเนหิ เอกาย ธาตุยา วิปฺปยุตฺตา.}

\proseitem{316}{อุเปกฺขา\-สหคเตหิ ธมฺเมหิ เย ธมฺมา สมฺปยุตฺตา, เตหิ ธมฺเมหิ เย ธมฺมา วิปฺปยุตฺตา. เต ธมฺมา กติหิ ขนฺเธหิ กติหายตเนหิ กติหิ ธาตูหิ วิปฺปยุตฺตา, เต ธมฺมา น เกหิจิ ขนฺเธหิ น เกหิจิ อายตเนหิ ปญฺจหิ ธาตูหิ วิปฺปยุตฺตา.}

\setlength{\csromangathapagewidth}{0pt}
\setlength{\csromangathawakwidth}{0pt}
\csromangathameasuregroup{}{ขนฺธา จตุโร อายตนญฺจ เมกํ, \\ ธาตูสุ สตฺต เทฺวปิ จ อินฺทฺริยโต. \\ ตโย ปฏิจฺจ ตถริว ผสฺสปญฺจมา, \\ อธิมุจฺจนา มนสิ ติเกสุ ตีณิ. \\ สตฺตนฺตรา เทฺว จ มเนน ยุตฺตา, \\ วิตกฺกวิจารณา อุเปกฺขกาย จาติ. \\ สมฺปยุตฺเตน วิปฺปยุตฺตปทนิทฺเทโส สตฺตโม.}
\csromangathameasurewak{ขนฺธา จตุโร อายตนญฺจ เมกํ, \\ ธาตูสุ สตฺต เทฺวปิ จ อินฺทฺริยโต. \\ ตโย ปฏิจฺจ ตถริว ผสฺสปญฺจมา, \\ อธิมุจฺจนา มนสิ ติเกสุ ตีณิ. \\ สตฺตนฺตรา เทฺว จ มเนน ยุตฺตา, \\ วิตกฺกวิจารณา อุเปกฺขกาย จาติ. \\ สมฺปยุตฺเตน วิปฺปยุตฺตปทนิทฺเทโส สตฺตโม.}
\setlength{\csromangathaleftcol}{0pt}
\csromangathagroup{\csromangathastanza{\csromangathawak{ขนฺธา จตุโร อายตนญฺจ เมกํ, \\ ธาตูสุ สตฺต เทฺวปิ จ อินฺทฺริยโต. \\ ตโย ปฏิจฺจ ตถริว ผสฺสปญฺจมา, \\ อธิมุจฺจนา มนสิ ติเกสุ ตีณิ.}}\csromangathastanza{\csromangathawak{สตฺตนฺตรา เทฺว จ มเนน ยุตฺตา, \\ วิตกฺกวิจารณา อุเปกฺขกาย จาติ. \\ สมฺปยุตฺเตน วิปฺปยุตฺต\-ปท\-นิทฺเทโส สตฺตโม.}}}

\csromanmatikaanchor{mtk.37}
\csromantocmarknopage{section}{8. อฏฺฐมนย}{mtk.37}
\bookmark[dest={mtk.37},level=1]{8. อฏฺฐมนย}
\csromanheader{1}{8. \textbf{อฏฺฐมนย}}
\csromanmatikaanchor{mtk.38}
\csromantocmarkref{subsection}{8. วิปฺปยุตฺเตนสมฺปยุตฺตปทนิทฺเทส}{mtk.38}
\bookmark[dest={mtk.38},level=2]{8. วิปฺปยุตฺเตนสมฺปยุตฺตปทนิทฺเทส}
\csromanheader{2}{8. วิปฺปยุตฺเตน\-สมฺปยุตฺต\-ปท\-นิทฺเทส}
\proseitem{317}{\csromanfolio{58}รูปกฺขนฺเธน เย ธมฺมา วิปฺปยุตฺตา, เต ธมฺมา กติหิ ขนฺเธหิ กติหายตเนหิ กติหิ ธาตูหิ สมฺปยุตฺตาติ นตฺถิ.}

\proseitem{318}{เวทนา\-กฺขนฺเธน เย ธมฺมา. สญฺญา\-กฺขนฺเธน เย ธมฺมา. สงฺขาร\-กฺขนฺเธน เย ธมฺมา. วิญฺญาณ\-กฺขนฺเธน เย ธมฺมา ฯเปฯ. สรเณหิ ธมฺเมหิ เย ธมฺมา. อรเณหิ ธมฺเมหิ เย ธมฺมา วิปฺปยุตฺตา, เต ธมฺมา กติหิ ขนฺเธหิ กติหายตเนหิ กติหิ ธาตูหิ สมฺปยุตฺตาติ นตฺถิ.}

\setlength{\csromangathapagewidth}{0pt}
\setlength{\csromangathawakwidth}{0pt}
\csromangathameasuregroup{ธมฺมายตนํ ธมฺมธาตุ, \\ สฬายตนํ ชาติชรามตํ, \\ ปฐมนฺตเร สตฺต จ, \\ จุทฺทส ฉ จ มตฺถเก, \\ สมุจฺเฉเท น ลพฺภนฺติ,}{\csromangathabat{ธมฺมายตนํ ธมฺมธาตุ,}{อถ ชีวิตํ นามรูปํ.} \\ \csromangathabat{สฬายตนํ ชาติชรามตํ,}{เทฺว จ ติเก น ลพฺภเร.} \\ \csromangathabat{ปฐมนฺตเร สตฺต จ,}{โคจฺฉเก ทส อปรนฺเต.} \\ \csromangathabat{จุทฺทส ฉ จ มตฺถเก,}{อิจฺเจเต สตฺตจตฺตาลีส ธมฺมา.} \\ \csromangathabat{สมุจฺเฉเท น ลพฺภนฺติ,}{โมฆปุจฺฉเกน จาติ.}}
\csromangathasetleft{ธมฺมายตนํ ธมฺมธาตุ, \\ สฬายตนํ ชาติชรามตํ, \\ ปฐมนฺตเร สตฺต จ, \\ จุทฺทส ฉ จ มตฺถเก, \\ สมุจฺเฉเท น ลพฺภนฺติ,}
\csromangathagroup{\csromangathastanza{\csromangathabat{ธมฺมายตนํ ธมฺมธาตุ,}{อถ ชีวิตํ นามรูปํ.} \\ \csromangathabat{สฬายตนํ ชาติชรามตํ,}{เทฺว จ ติเก น ลพฺภเร.}}\csromangathastanza{\csromangathabat{ปฐมนฺตเร สตฺต จ,}{โคจฺฉเก ทส อปรนฺเต.} \\ \csromangathabat{จุทฺทส ฉ จ มตฺถเก,}{อิจฺเจเต สตฺตจตฺตาลีส ธมฺมา.} \\ \csromangathabat{สมุจฺเฉเท น ลพฺภนฺติ,}{โมฆปุจฺฉเกน จาติ.}}}

\vspace{\tipitakaparskip}
\csromancenter{วิปฺปยุตฺเตน สมฺปยุตฺต\-ปท\-นิทฺเทโส อฏฺฐโม.}
\csromanmatikaanchor{mtk.39}
\csromantocmarknopage{section}{9. นวมนย}{mtk.39}
\bookmark[dest={mtk.39},level=1]{9. นวมนย}
\csromanheader{1}{9. \textbf{นวมนย}}
\csromanmatikaanchor{mtk.40}
\csromantocmarkref{subsection}{9. สมฺปยุตฺเตนสมฺปยุตฺตปทนิทฺเทส}{mtk.40}
\bookmark[dest={mtk.40},level=2]{9. สมฺปยุตฺเตนสมฺปยุตฺตปทนิทฺเทส}
\csromanheader{2}{9. สมฺปยุตฺเตน\-สมฺปยุตฺต\-ปท\-นิทฺเทส}
\proseitem{319}{\csromanfolio{59}เวทนา\-กฺขนฺเธน เย ธมฺมา. สญฺญา\-กฺขนฺเธน เย ธมฺมา. สงฺขาร\-กฺขนฺเธน เย ธมฺมา สมฺปยุตฺตา, เตหิ ธมฺเมหิ เย ธมฺมา สมฺปยุตฺตา. เต ธมฺมา กติหิ ขนฺเธหิ กติหายตเนหิ กติหิ ธาตูหิ สมฺปยุตฺตา, เต ธมฺมา ตีหิ ขนฺเธหิ เอเกนายตเนน สตฺตหิ ธาตูหิ สมฺปยุตฺตา, เอเกนายตเนน เอกาย ธาตุยา เกหิจิ สมฺปยุตฺตา.}

\proseitem{320}{วิญฺญาณ\-กฺขนฺเธน เย ธมฺมา. มนายตเนน เย ธมฺมา. จกฺขุวิญฺญาณ\-ธาตุยา เย ธมฺมา ฯเปฯ. มโนธาตุยา เย ธมฺมา. มโนวิญฺญาณ\-ธาตุยา เย ธมฺมา สมฺปยุตฺตา, เตหิ ธมฺเมหิ เย ธมฺมา สมฺปยุตฺตา ฯเปฯ เต ธมฺมา ตีหิ ขนฺเธหิ สมฺปยุตฺตา, เอเกนายตเนน เอกาย ธาตุยา เกหิจิ สมฺปยุตฺตา.}

\proseitem{321}{สมุทย\-สจฺเจน เย ธมฺมา. มคฺคสจฺเจน เย ธมฺมา สมฺปยุตฺตา, เตหิ ธมฺเมหิ เย ธมฺมา สมฺปยุตฺตา. เต ธมฺมา ตีหิ ขนฺเธหิ เอเกนายตเนน เอกาย ธาตุยา สมฺปยุตฺตา, เอเกน ขนฺเธน เอเกนายตเนน เอกาย ธาตุยา เกหิจิ สมฺปยุตฺตา.}

\proseitem{322}{มนินฺทฺริเยน เย ธมฺมา สมฺปยุตฺตา, เตหิ ธมฺเมหิ เย ธมฺมา สมฺปยุตฺตา. เต ธมฺมา ตีหิ ขนฺเธหิ สมฺปยุตฺตา, เอเกนายตเนน เอกาย ธาตุยา เกหิจิ สมฺปยุตฺตา.}

\proseitem{323}{สุขินฺทฺริเยน เย ธมฺมา. ทุกฺขินฺทฺริเยน เย ธมฺมา. โสมนสฺสิ\-นฺทฺริเยน เย ธมฺมา. โทมนสฺสิ\-นฺทฺริเยน เย ธมฺมา สมฺปยุตฺตา, เตหิ ธมฺเมหิ เย ธมฺมา สมฺปยุตฺตา. เต ธมฺมา ตีหิ ขนฺเธหิ เอเกนายตเนน เอกาย ธาตุยา สมฺปยุตฺตา, เอเกนายตเนน เอกาย ธาตุยา เกหิจิ สมฺปยุตฺตา.}

\proseitem{324}{อุเปกฺขิ\-นฺทฺริเยน เย ธมฺมา สมฺปยุตฺตา, เตหิ ธมฺเมหิ เย ธมฺมา สมฺปยุตฺตา. เต ธมฺมา ตีหิ ขนฺเธหิ เอเกนายตเนน ฉหิ ธาตูหิ สมฺปยุตฺตา, เอเกนายตเนน เอกาย ธาตุยา เกหิจิ สมฺปยุตฺตา.}

\proseitem{325}{\csromanfolio{60}สทฺธินฺทฺริเยน เย ธมฺมา. วีริยินฺทฺริเยน เย ธมฺมา. สตินฺทฺริเยน เย ธมฺมา. สมาธิ\-นฺทฺริเยน เย ธมฺมา. ปญฺญินฺทฺริเยน เย ธมฺมา. อนญฺญาต\-ญฺญสฺสามี\-ตินฺทฺริเยน เย ธมฺมา. อญฺญินฺทฺริเยน เย ธมฺมา. อญฺญาตาวิ\-นฺทฺริเยน เย ธมฺมา. อวิชฺชาย เย ธมฺมา. อวิชฺชาปจฺจยา สงฺขาเรหิ เย ธมฺมา สมฺปยุตฺตา, เตหิ ธมฺเมหิ เย ธมฺมา สมฺปยุตฺตา. เต ธมฺมา ตีหิ ขนฺเธหิ เอเกนายตเนน เอกาย ธาตุยา สมฺปยุตฺตา, เอเกน ขนฺเธน เอเกนายตเนน เอกาย ธาตุยา เกหิจิ สมฺปยุตฺตา.}

\proseitem{326}{สงฺขาร\-ปจฺจยา วิญฺญาเณน เย ธมฺมา สมฺปยุตฺตา, เตหิ ธมฺเมหิ เย ธมฺมา สมฺปยุตฺตา. เต ธมฺมา ตีหิ ขนฺเธหิ สมฺปยุตฺตา, เอเกนายตเนน เอกาย ธาตุยา เกหิจิ สมฺปยุตฺตา.}

\proseitem{327}{สฬายตน\-ปจฺจยา ผสฺเสน เย ธมฺมา สมฺปยุตฺตา, เตหิ ธมฺเมหิ เย ธมฺมา สมฺปยุตฺตา. เต ธมฺมา ตีหิ ขนฺเธหิ เอเกนายตเนน สตฺตหิ ธาตูหิ สมฺปยุตฺตา, เอเกน ขนฺเธน เอเกนายตเนน เอกาย ธาตุยา เกหิจิ สมฺปยุตฺตา.}

\proseitem{328}{ผสฺสปจฺจยา เวทนาย เย ธมฺมา สมฺปยุตฺตา, เตหิ ธมฺเมหิ เย ธมฺมา สมฺปยุตฺตา. เต ธมฺมา ตีหิ ขนฺเธหิ เอเกนายตเนน สตฺตหิ ธาตูหิ สมฺปยุตฺตา, เอเกนายตเนน เอกาย ธาตุยา เกหิจิ สมฺปยุตฺตา.}

\proseitem{329}{เวทนาปจฺจยา ตณฺหาย เย ธมฺมา. ตณฺหาปจฺจยา อุปาทาเนน เย ธมฺมา. กมฺมภเวน เย ธมฺมา สมฺปยุตฺตา, เตหิ ธมฺเมหิ เย ธมฺมา สมฺปยุตฺตา. เต ธมฺมา ตีหิ ขนฺเธหิ เอเกนายตเนน เอกาย ธาตุยา สมฺปยุตฺตา, เอเกน ขนฺเธน เอเกนายตเนน เอกาย ธาตุยา เกหิจิ สมฺปยุตฺตา.}

\proseitem{330}{โสเกน เย ธมฺมา. ทุกฺเขน เย ธมฺมา. โทมนสฺเสน เย ธมฺมา สมฺปยุตฺตา, เตหิ ธมฺเมหิ เย ธมฺมา สมฺปยุตฺตา. เต ธมฺมา ตีหิ ขนฺเธหิ เอเกนายตเนน เอกาย ธาตุยา สมฺปยุตฺตา, เอเกนายตเนน เอกาย ธาตุยา เกหิจิ สมฺปยุตฺตา.}

\proseitem{331}{\csromanfolio{61}อุปายาเสน เย ธมฺมา. สติปฏฺฐาเนน เย ธมฺมา. สมฺม\-ปฺปธาเนน เย ธมฺมา สมฺปยุตฺตา, เตหิ ธมฺเมหิ เย ธมฺมา สมฺปยุตฺตา. เต ธมฺมา ตีหิ ขนฺเธหิ เอเกนายตเนน เอกาย ธาตุยา สมฺปยุตฺตา, เอเกน ขนฺเธน เอเกนายตเนน เอกาย ธาตุยา เกหิจิ สมฺปยุตฺตา.}

\proseitem{332}{อิทฺธิปาเทน เย ธมฺมา สมฺปยุตฺตา, เตหิ ธมฺเมหิ เย ธมฺมา สมฺปยุตฺตา. เต ธมฺมา ทฺวีหิ ขนฺเธหิ สมฺปยุตฺตา, เอเกนายตเนน เอกาย ธาตุยา เกหิจิ สมฺปยุตฺตา.}

\proseitem{333}{ฌาเนน เย ธมฺมา สมฺปยุตฺตา, เตหิ ธมฺเมหิ เย ธมฺมา สมฺปยุตฺตา. เต ธมฺมา ทฺวีหิ ขนฺเธหิ เอเกนายตเนน เอกาย ธาตุยา สมฺปยุตฺตา, เอเกน ขนฺเธน เอเกนายตเนน เอกาย ธาตุยา เกหิจิ สมฺปยุตฺตา.}

\proseitem{334}{อปฺปมญฺญาย เย ธมฺมา. ปญฺจหิ อินฺทฺริเยหิ เย ธมฺมา. ปญฺจหิ พเลหิ เย ธมฺมา. สตฺตหิ โพชฺฌงฺเคหิ เย ธมฺมา. อริเยน อฏฺฐงฺคิเกน มคฺเคน เย ธมฺมา สมฺปยุตฺตา, เตหิ ธมฺเมหิ เย ธมฺมา สมฺปยุตฺตา. เต ธมฺมา ตีหิ ขนฺเธหิ เอเกนายตเนน เอกาย ธาตุยา สมฺปยุตฺตา, เอเกน ขนฺเธน เอเกนายตเนน เอกาย ธาตุยา เกหิจิ สมฺปยุตฺตา.}

\proseitem{335}{ผสฺเสน เย ธมฺมา. เจตนาย เย ธมฺมา. มนสิกาเรน เย ธมฺมา สมฺปยุตฺตา, เตหิ ธมฺเมหิ เย ธมฺมา สมฺปยุตฺตา. เต ธมฺมา ตีหิ ขนฺเธหิ เอเกนายตเนน สตฺตหิ ธาตูหิ สมฺปยุตฺตา, เอเกน ขนฺเธน เอเกนายตเนน เอกาย ธาตุยา เกหิจิ สมฺปยุตฺตา.}

\proseitem{336}{เวทนาย เย ธมฺมา. สญฺญาย เย ธมฺมา สมฺปยุตฺตา, ติหิ ธมฺเมหิ เย ธมฺมา สมฺปยุตฺตา, เต ธมฺมา ตีหิ ขนฺเธหิ เอเกนายตเนน สตฺตหิ ธาตูหิ สมฺปยุตฺตา, เอเกนายตเนน เอกาย ธาตุยา เกหิจิ สมฺปยุตฺตา.}

\proseitem{337}{จิตฺเตน เย ธมฺมา สมฺปยุตฺตา, เตหิ ธมฺเมหิ เย ธมฺมา สมฺปยุตฺตา. เต ธมฺมา ตีหิ ขนฺเธหิ สมฺปยุตฺตา, เอเกนายตเนน เอกาย ธาตุยา เกหิจิ สมฺปยุตฺตา.}

\proseitem{338}{\csromanfolio{62}อธิโมกฺเขน เย ธมฺมา สมฺปยุตฺตา, เตหิ ธมฺเมหิ เย ธมฺมา สมฺปยุตฺตา. เต ธมฺมา ตีหิ ขนฺเธหิ เอเกนายตเนน ทฺวีหิ ธาตูหิ สมฺปยุตฺตา, เอเกน ขนฺเธน เอเกนายตเนน เอกาย ธาตุยา เกหิจิ สมฺปยุตฺตา.}

\proseitem{339}{สุขาย เวทนาย สมฺปยุตฺเตหิ ธมฺเมหิ เย ธมฺมา. ทุกฺขาย เวทนาย สมฺปยุตฺเตหิ ธมฺเมหิ เย ธมฺมา. อทุกฺขม\-สุขาย เวทนาย สมฺปยุตฺเตหิ ธมฺเมหิ เย ธมฺมา สมฺปยุตฺตา, เตหิ ธมฺเมหิ เย ธมฺมา สมฺปยุตฺตา. เต ธมฺมา เอเกน ขนฺเธน สมฺปยุตฺตา, เอเกนายตเนน เอกาย ธาตุยา เกหิจิ สมฺปยุตฺตา.}

\proseitem{340}{สวิตกฺก\-สวิจาเรหิ ธมฺเมหิ เย ธมฺมา. อวิตกฺก\-วิจารมตฺเตหิ ธมฺเมหิ เย ธมฺมา. ปีติสหคเตหิ ธมฺเมหิ เย ธมฺมา สมฺปยุตฺตา, เตหิ ธมฺเมหิ เย ธมฺมา สมฺปยุตฺตา. เต ธมฺมา เอเกน ขนฺเธน เอเกนายตเนน เอกาย ธาตุยา เกหิจิ สมฺปยุตฺตา.}

\proseitem{341}{สุขสหคเตหิ ธมฺเมหิ เย ธมฺมา, อุเปกฺขา\-สหคเตหิ ธมฺเมหิ เย ธมฺมา สมฺปยุตฺตา, เตหิ ธมฺเมหิ เย ธมฺมา สมฺปยุตฺตา. เต ธมฺมา เอเกน ขนฺเธน สมฺปยุตฺตา, เอเกนายตเนน เอกาย ธาตุยา เกหิจิ สมฺปยุตฺตา.}

\proseitem{342}{เหตูหิ ธมฺเมหิ เย ธมฺมา. เหตูหิ เจว สเหตุเกหิ จ ธมฺเมหิ เย ธมฺมา. เหตูหิ เจว เหตุ\-สมฺปยุตฺเตหิ จ ธมฺเมหิ เย ธมฺมา สมฺปยุตฺตา, เตหิ ธมฺเมหิ เย ธมฺมา สมฺปยุตฺตา. เต ธมฺมา ตีหิ ขนฺเธหิ เอเกนายตเนน เอกาย ธาตุยา สมฺปยุตฺตา, เอเกน ขนฺเธน เอเกนายตเนน เอกาย ธาตุยา เกหิจิ สมฺปยุตฺตา.}

\proseitem{343}{สเหตุเกหิ เจว น จ เหตูหิ ธมฺเมหิ เย ธมฺมา. เหตุ\-สมฺปยุตฺเตหิ เจว น จ เหตูหิ ธมฺเมหิ เย ธมฺมา. น เหตุสเหตุเกหิ\csromansharedfootnote{466303ca518780e9}{น เหตูหิ สเหตุเกหิ (พหูสุ)} ธมฺเมหิ เย ธมฺมา สมฺปยุตฺตา, เตหิ ธมฺเมหิ เย ธมฺมา สมฺปยุตฺตา, เต ธมฺมา เอเกน ขนฺเธน เอเกนายตเนน เอกาย ธาตุยา เกหิจิ สมฺปยุตฺตา.}

\proseitem{344}{\csromanfolio{63}อาสเวหิ ธมฺเมหิ เย ธมฺมา. อาสเวหิ เจว สาสเวหิ จ ธมฺเมหิ เย ธมฺมา. อาสเวหิ เจว อาสว\-สมฺปยุตฺเตหิ จ ธมฺเมหิ เย ธมฺมา สมฺปยุตฺตา, เตหิ ธมฺเมหิ เย ธมฺมา สมฺปยุตฺตา. เต ธมฺมา ตีหิ ขนฺเธหิ เอเกนายตเนน เอกาย ธาตุยา สมฺปยุตฺตา, เอเกน ขนฺเธน เอเกนายตเนน เอกาย ธาตุยา เกหิจิ สมฺปยุตฺตา.}

\proseitem{345}{อาสว\-สมฺปยุตฺเตหิ เจว โน จ อาสเวหิ ธมฺเมหิ เย ธมฺมา สมฺปยุตฺตา, เตหิ ธมฺเมหิ เย ธมฺมา สมฺปยุตฺตา. เต ธมฺมา เอเกน ขนฺเธน เอเกนายตเนน เอกาย ธาตุยา เกหิจิ สมฺปยุตฺตา.}

\proseitem{346}{สํโยชเนหิ. คนฺเถหิ. โอเฆหิ. โยเคหิ. นีวรเณหิ. ปรามาเสหิ ธมฺเมหิ เย ธมฺมา. ปรามาเสหิ เจว ปรามฏฺเฐหิ จ ธมฺเมหิ เย ธมฺมา สมฺปยุตฺตา, เตหิ ธมฺเมหิ เย ธมฺมา สมฺปยุตฺตา. เต ธมฺมา ตีหิ ขนฺเธหิ เอเกนายตเนน เอกาย ธาตุยา สมฺปยุตฺตา, เอเกน ขนฺเธน เอเกนายตเนน เอกาย ธาตุยา เกหิจิ สมฺปยุตฺตา.}

\proseitem{347}{ปรามาส\-สมฺปยุตฺเตหิ ธมฺเมหิ เย ธมฺมา สมฺปยุตฺตา, เตหิ ธมฺเมหิ เย ธมฺมา สมฺปยุตฺตา. เต ธมฺมา เอเกน ขนฺเธน เอเกนายตเนน เอกาย ธาตุยา เกหิจิ สมฺปยุตฺตา.}

\proseitem{348}{จิตฺเตหิ ธมฺเมหิ เย ธมฺมา สมฺปยุตฺตา, เตหิ ธมฺเมหิ เย ธมฺมา สมฺปยุตฺตา. เต ธมฺมา ตีหิ ขนฺเธหิ สมฺปยุตฺตา, เอเกนายตเนน เอกาย ธาตุยา เกหิจิ สมฺปยุตฺตา.}

\proseitem{349}{เจตสิเกหิ ธมฺเมหิ เย ธมฺมา. จิตฺต\-สมฺปยุตฺเตหิ ธมฺเมหิ เย ธมฺมา. จิตฺต\-สํสฏฺเฐหิ ธมฺเมหิ เย ธมฺมา. จิตฺต\-สํสฏฺฐ\-สมุฏฺฐาเนหิ ธมฺเมหิ เย ธมฺมา. จิตฺต\-สํสฏฺฐ\-สมุฏฺฐาน\-สหภูหิ ธมฺเมหิ เย ธมฺมา. จิตฺต\-สํสฏฺฐ\-สมุฏฺฐานา\-นุปริวตฺตีหิ ธมฺเมหิ เย ธมฺมา สมฺปยุตฺตา, เตหิ ธมฺเมหิ เย ธมฺมา สมฺปยุตฺตา. เต ธมฺมา เอเกน ขนฺเธน เอเกนายตเนน สตฺตหิ ธาตูหิ สมฺปยุตฺตา.}

\proseitem{350}{อุปาทาเนหิ ธมฺเมหิ เย ธมฺมา. กิเลเสหิ ธมฺเมหิ เย ธมฺมา. กิเลเสหิ เจว สํกิเลสิเกหิ จ ธมฺเมหิ เย ธมฺมา. กิเลเสหิ เจว สํกิลิฏฺเฐหิ จ ธมฺเมหิ เย ธมฺมา. กิเลเสหิ เจว \csromanfolio{64}กิเลส\-สมฺปยุตฺเตหิ จ ธมฺเมหิ เย ธมฺมา สมฺปยุตฺตา, เตหิ ธมฺเมหิ เย ธมฺมา สมฺปยุตฺตา. เต ธมฺมา ตีหิ ขนฺเธหิ เอเกนายตเนน เอกาย ธาตุยา สมฺปยุตฺตา, เอเกน ขนฺเธน เอเกนายตเนน เอกาย ธาตุยา เกหิจิ สมฺปยุตฺตา.}

\proseitem{351}{สํกิลิฏฺเฐหิ เจว โน จ กิเลเสหิ ธมฺเมหิ เย ธมฺมา. กิเลส\-สมฺปยุตฺเตหิ เจว โน จ กิเลเสหิ ธมฺเมหิ เย ธมฺมา. สวิตกฺเกหิ ธมฺเมหิ เย ธมฺมา. สวิจาเรหิ ธมฺเมหิ เย ธมฺมา. สปฺปีติเกหิ ธมฺเมหิ เย ธมฺมา. ปีติสหคเตหิ ธมฺเมหิ เย ธมฺมา สมฺปยุตฺตา, เตหิ ธมฺเมหิ เย ธมฺมา สมฺปยุตฺตา. เต ธมฺมา เอเกน ขนฺเธน เอเกนายตเนน เอกาย ธาตุยา เกหิจิ สมฺปยุตฺตา.}

\proseitem{352}{สุขสหคเตหิ ธมฺเมหิ เย ธมฺมา. อุเปกฺขา\-สหคเตหิ ธมฺเมหิ เย ธมฺมา สมฺปยุตฺตา, เตหิ ธมฺเมหิ เย ธมฺมา สมฺปยุตฺตา. เต ธมฺมา กติหิ ขนฺเธหิ กติหายตเนหิ กติหิ ธาตูหิ สมฺปยุตฺตา, เต ธมฺมา เอเกน ขนฺเธน สมฺปยุตฺตา, เอเกนายตเนน เอกาย ธาตุยา เกหิจิ สมฺปยุตฺตา.}

\setlength{\csromangathapagewidth}{0pt}
\setlength{\csromangathawakwidth}{0pt}
\csromangathameasuregroup{อรูปกฺขนฺธา จตฺตาโร, \\ วิญฺญาณธาตุโย สตฺต, \\ ปจฺจเย ทฺวาทส ปทา, \\ ติเกสุ อฏฺฐ โคจฺฉเก, \\ มหนฺตรทุเก สตฺต, \\ นวมสฺส ปทสฺเสเต,}{\csromangathabat{อรูปกฺขนฺธา จตฺตาโร,}{มนายตนเมว จ.} \\ \csromangathabat{วิญฺญาณธาตุโย สตฺต,}{เทฺว สจฺจา จุทฺทสินฺทฺริยา.} \\ \csromangathabat{ปจฺจเย ทฺวาทส ปทา,}{ตโต อุปริ โสฬส.} \\ \csromangathabat{ติเกสุ อฏฺฐ โคจฺฉเก,}{เตจตฺตาลีสเมว จ.} \\ \csromangathabat{มหนฺตรทุเก สตฺต,}{ปทา ปิฏฺฐิ ทุเกสุ ฉ.} \\ \csromangathabat{นวมสฺส ปทสฺเสเต,}{นิทฺเทเส สงฺคหํ คตาติ.}}
\csromangathasetleft{อรูปกฺขนฺธา จตฺตาโร, \\ วิญฺญาณธาตุโย สตฺต, \\ ปจฺจเย ทฺวาทส ปทา, \\ ติเกสุ อฏฺฐ โคจฺฉเก, \\ มหนฺตรทุเก สตฺต, \\ นวมสฺส ปทสฺเสเต,}
\csromangathagroup{\csromangathastanza{\csromangathabat{อรูปกฺขนฺธา จตฺตาโร,}{มนายตนเมว จ.} \\ \csromangathabat{วิญฺญาณธาตุโย สตฺต,}{เทฺว สจฺจา จุทฺทสินฺทฺริยา.}}\csromangathastanza{\csromangathabat{ปจฺจเย ทฺวาทส ปทา,}{ตโต อุปริ โสฬส.} \\ \csromangathabat{ติเกสุ อฏฺฐ โคจฺฉเก,}{เตจตฺตาลีสเมว จ.}}\csromangathastanza{\csromangathabat{มหนฺตรทุเก สตฺต,}{ปทา ปิฏฺฐิ ทุเกสุ ฉ.} \\ \csromangathabat{นวมสฺส ปทสฺเสเต,}{นิทฺเทเส สงฺคหํ คตาติ.}}}

\vspace{\tipitakaparskip}
\csromancenter{สมฺปยุตฺเตน สมฺปยุตฺต\-ปท\-นิทฺเทโส นวโม.}
\csromanmatikaanchor{mtk.41}
\csromantocmarknopage{section}{10. ทสมนย}{mtk.41}
\bookmark[dest={mtk.41},level=1]{10. ทสมนย}
\csromanheader{1}{10. \textbf{ทสมนย}}
\csromanmatikaanchor{mtk.42}
\csromantocmarkref{subsection}{10. วิปฺปยุตฺเตนวิปฺปยุตฺตปทนิทฺเทส}{mtk.42}
\bookmark[dest={mtk.42},level=2]{10. วิปฺปยุตฺเตนวิปฺปยุตฺตปทนิทฺเทส}
\csromanheader{2}{10. วิปฺปยุตฺเตน\-วิปฺปยุตฺต\-ปท\-นิทฺเทส}
\proseitem{353}{\csromanfolio{65}รูปกฺขนฺเธน เย ธมฺมา วิปฺปยุตฺตา, เตหิ ธมฺเมหิ เย ธมฺมา วิปฺปยุตฺตา. เต ธมฺมา กติหิ ขนฺเธหิ กติหายตเนหิ กติหิ ธาตูหิ วิปฺปยุตฺตา, เต ธมฺมา จตูหิ ขนฺเธหิ เอเกนายตเนน สตฺตหิ ธาตูหิ วิปฺปยุตฺตา, เอเกนายตเนน เอกาย ธาตุยา เกหิจิ วิปฺปยุตฺตา.}

\proseitem{354}{เวทนา\-กฺขนฺเธน เย ธมฺมา. สญฺญา\-กฺขนฺเธน เย ธมฺมา. สงฺขาร\-กฺขนฺเธน เย ธมฺมา. วิญฺญาณ\-กฺขนฺเธน เย ธมฺมา. มนายตเนน เย ธมฺมา วิปฺปยุตฺตา, เตหิ ธมฺเมหิ เย ธมฺมา วิปฺปยุตฺตา ฯเปฯ เต ธมฺมา เอเกน ขนฺเธน ทสหายตเนหิ ทสหิ ธาตูหิ วิปฺปยุตฺตา, เอเกนายตเนน เอกาย ธาตุยา เกหิจิ วิปฺปยุตฺตา.}

\proseitem{355}{จกฺขายตเนน เย ธมฺมา ฯเปฯ. โผฏฺฐพฺพา\-ยตเนน เย ธมฺมา. จกฺขุธาตุยา เย ธมฺมา ฯเปฯ. โผฏฺฐพฺพ\-ธาตุยา เย ธมฺมา วิปฺปยุตฺตา, เตหิ ธมฺเมหิ เย ธมฺมา วิปฺปยุตฺตา. เต ธมฺมา จตูหิ ขนฺเธหิ เอเกนายตเนน สตฺตหิ ธาตูหิ วิปฺปยุตฺตา, เอเกนายตเนน เอกาย ธาตูยา เกหิจิ วิปฺปยุตฺตา.}

\proseitem{356}{จกฺขุวิญฺญาณ\-ธาตุยา เย ธมฺมา ฯเปฯ. มโนวิญฺญาณ\-ธาตุยา เย ธมฺมา. สมุทย\-สจฺเจน เย ธมฺมา. มคฺคสจฺเจน เย ธมฺมา วิปฺปยุตฺตา, เตหิ ธมฺเมหิ เย ธมฺมา วิปฺปยุตฺตา. เต ธมฺมา เอเกน ขนฺเธน ทสหายตเนหิ โสฬสหิ ธาตูหิ วิปฺปยุตฺตา, เอเกนายตเนน เอกาย ธาตุยา เกหิจิ วิปฺปยุตฺตา.}

\proseitem{357}{นิโรธสจฺเจน เย ธมฺมา. จกฺขุนฺทฺริเยน เย ธมฺมา ฯเปฯ. กายินฺทฺริเยน เย ธมฺมา. อิตฺถินฺทฺริเยน เย ธมฺมา. ปุริสินฺทฺริเยน เย ธมฺมา วิปฺปยุตฺตา, เตหิ ธมฺเมหิ เย ธมฺมา วิปฺปยุตฺตา. เต ธมฺมา จตูหิ ขนฺเธหิ เอเกนายตเนน สตฺตหิ ธาตูหิ วิปฺปยุตฺตา, เอเกนายตเนน เอกาย ธาตุยา เกหิจิ วิปฺปยุตฺตา.}

\proseitem{358}{\csromanfolio{66}มนินฺทฺริเยน เย ธมฺมา วิปฺปยุตฺตา, เตหิ ธมฺเมหิ เย ธมฺมา วิปฺปยุตฺตา. เต ธมฺมา เอเกน ขนฺเธน ทสหายตเนหิ ทสหิ ธาตูหิ วิปฺปยุตฺตา, เอเกนายตเนน เอกาย ธาตุยา เกหิจิ วิปฺปยุตฺตา.}

\proseitem{359}{สุขินฺทฺริเยน เย ธมฺมา. ทุกฺขินฺทฺริเยน เย ธมฺมา. โสมนสฺสิ\-นฺทฺริเยน เย ธมฺมา. โทมนสฺสิ\-นฺทฺริเยน เย ธมฺมา วิปฺปยุตฺตา, เตหิ ธมฺเมหิ เย ธมฺมา วิปฺปยุตฺตา. เต ธมฺมา เอเกน ขนฺเธน ทสหายตเนหิ โสฬสหิ ธาตูหิ วิปฺปยุตฺตา, เอเกนายตเนน เอกาย ธาตุยา เกหิจิ วิปฺปยุตฺตา.}

\proseitem{360}{อุเปกฺขิ\-นฺทฺริเยน เย ธมฺมา วิปฺปยุตฺตา, เตหิ ธมฺเมหิ เย ธมฺมา วิปฺปยุตฺตา. เต ธมฺมา เอเกน ขนฺเธน ทสหายตเนหิ เอกาทสหิ ธาตูหิ วิปฺปยุตฺตา, เอเกนายตเนน เอกาย ธาตุยา เกหิจิ วิปฺปยุตฺตา.}

\proseitem{361}{สทฺธินฺทฺริเยน เย ธมฺมา. วีริยินฺทฺริเยน เย ธมฺมา. สตินฺทฺริเยน เย ธมฺมา. สมาธิ\-นฺทฺริเยน เย ธมฺมา. ปญฺญินฺทฺริเยน เย ธมฺมา. อนญฺญาต\-ญฺญสฺสามี\-ตินฺทฺริเยน เย ธมฺมา. อญฺญินฺทฺริเยน เย ธมฺมา. อญฺญาตาวิ\-นฺทฺริเยน เย ธมฺมา. อวิชฺชาย เย ธมฺมา. อวิชฺชาปจฺจยา สงฺขาเรหิ เย ธมฺมา วิปฺปยุตฺตา, เตหิ ธมฺเมหิ เย ธมฺมา วิปฺปยุตฺตา. เต ธมฺมา เอเกน ขนฺเธน ทสหายตเนหิ โสฬสหิ ธาตูหิ วิปฺปยุตฺตา, เอเกนายตเนน เอกาย ธาตุยา เกหิจิ วิปฺปยุตฺตา.}

\proseitem{362}{สงฺขาร\-ปจฺจยา วิญฺญาเณน เย ธมฺมา. สฬายตน\-ปจฺจยา ผสฺเสน เย ธมฺมา. ผสฺสปจฺจยา เวทนาย เย ธมฺมา วิปฺปยุตฺตา, เตหิ ธมฺเมหิ เย ธมฺมา วิปฺปยุตฺตา. เต ธมฺมา เอเกน ขนฺเธน ทสหายตเนหิ ทสหิ ธาตูหิ วิปฺปยุตฺตา, เอเกนายตเนน เอกาย ธาตุยา เกหิจิ วิปฺปยุตฺตา.}

\proseitem{363}{เวทนาปจฺจยา ตณฺหาย เย ธมฺมา. ตณฺหาปจฺจยา อุปาทาเนน เย ธมฺมา. กมฺมภเวน เย ธมฺมา วิปฺปยุตฺตา, เตหิ ธมฺเมหิ เย ธมฺมา วิปฺปยุตฺตา. เต ธมฺมา เอเกน ขนฺเธน ทสหายตเนหิ โสฬสหิ ธาตูหิ วิปฺปยุตฺตา, เอเกนายตเนน เอกาย ธาตุยา เกหิจิ วิปฺปยุตฺตา.}

\proseitem{364}{\csromanfolio{67}รูปภเวน เย ธมฺมา วิปฺปยุตฺตา, เตหิ ธมฺเมหิ เย ธมฺมา วิปฺปยุตฺตา. เต ธมฺมา น เกหิจิ ขนฺเธหิ น เกหิจิ อายตเนหิ ตีหิ ธาตูหิ วิปฺปยุตฺตา.}

\proseitem{365}{อสญฺญาภเวน เย ธมฺมา. เอกโวการ\-ภเวน เย ธมฺมา. ปริเทเวน เย ธมฺมา วิปฺปยุตฺตา, เตหิ ธมฺเมหิ เย ธมฺมา วิปฺปยุตฺตา. เต ธมฺมา จตูหิ ขนฺเธหิ เอเกนายตเนน สตฺตหิ ธาตูหิ วิปฺปยุตฺตา, เอเกนายตเนน เอกาย ธาตุยา เกหิจิ วิปฺปยุตฺตา.}

\proseitem{366}{อรูปภเวน เย ธมฺมา. เนว\-สญฺญา\-นา\-สญฺญาภเวน เย ธมฺมา. จตุโวการภเวน เย ธมฺมา. โสเกน เย ธมฺมา. ทุกฺเขน เย ธมฺมา. โทมนสฺเสน เย ธมฺมา. อุปายาเสน เย ธมฺมา. สติปฏฺฐาเนน เย ธมฺมา. สมฺม\-ปฺปธาเนน เย ธมฺมา. อิทฺธิปาเทน เย ธมฺมา. ฌาเนน เย ธมฺมา. อปฺปมญฺญาย เย ธมฺมา. ปญฺจหิ อินฺทฺริเยหิ เย ธมฺมา. ปญฺจหิ พเลหิ เย ธมฺมา. สตฺตหิ โพชฺฌงฺเคหิ เย ธมฺมา. อริเยน อฏฺฐงฺคิเกน มคฺเคน เย ธมฺมา วิปฺปยุตฺตา, เตหิ ธมฺเมหิ เย ธมฺมา วิปฺปยุตฺตา. เต ธมฺมา เอเกน ขนฺเธน ทสหายตเนหิ โสฬสหิ ธาตูหิ วิปฺปยุตฺตา, เอเกนายตเนน เอกาย ธาตุยา เกหิจิ วิปฺปยุตฺตา.}

\proseitem{367}{ผสฺเสน เย ธมฺมา. เวทนาย เย ธมฺมา. สญฺญาย เย ธมฺมา. เจตนาย เย ธมฺมา. จิตฺเตน เย ธมฺมา. มนสิกาเรน เย ธมฺมา วิปฺปยุตฺตา, เตหิ ธมฺเมหิ เย ธมฺมา วิปฺปยุตฺตา. เต ธมฺมา เอเกน ขนฺเธน ทสหายตเนหิ ทสหิ ธาตูหิ วิปฺปยุตฺตา, เอเกนายตเนน เอกาย ธาตุยา เกหิจิ วิปฺปยุตฺตา.}

\proseitem{368}{อธิโมกฺเขน เย ธมฺมา วิปฺปยุตฺตา, เตหิ ธมฺเมหิ เย ธมฺมา วิปฺปยุตฺตา. เต ธมฺมา เอเกน ขนฺเธน ทสหายตเนหิ ปนฺนรสหิ ธาตูหิ วิปฺปยุตฺตา, เอเกนายตเนน เอกาย ธาตุยา เกหิจิ วิปฺปยุตฺตา.}

\csromanmatikaanchor{mtk.43}
\csromantocmarkref{subsection}{1. ติก}{mtk.43}
\bookmark[dest={mtk.43},level=2]{1. ติก}
\csromanheader{2}{1. \textbf{ติก}}
\proseitem{369}{กุสเลหิ ธมฺเมหิ เย ธมฺมา. อกุสเลหิ ธมฺเมหิ เย ธมฺมา วิปฺปยุตฺตา, เตหิ ธมฺเมหิ เย ธมฺมา วิปฺปยุตฺตา. เต ธมฺมา เอเกน ขนฺเธน ทสหายตเนหิ โสฬสหิ ธาตูหิ วิปฺปยุตฺตา, เอเกนายตเนน เอกาย ธาตุยา เกหิจิ วิปฺปยุตฺตา.}

\proseitem{370}{\csromanfolio{68}สุขาย เวทนาย สมฺปยุตฺเตหิ ธมฺเมหิ เย ธมฺมา. ทุกฺขาย เวทนาย สมฺปยุตฺเตหิ ธมฺเมหิ เย ธมฺมา วิปฺปยุตฺตา, เตหิ ธมฺเมหิ เย ธมฺมา วิปฺปยุตฺตา. เต ธมฺมา เอเกน ขนฺเธน ทสหายตเนหิ ปนฺนรสหิ ธาตูหิ วิปฺปยุตฺตา, เอเกนายตเนน เอกาย ธาตุยา เกหิจิ วิปฺปยุตฺตา.}

\proseitem{371}{อทุกฺขม\-สุขาย เวทนาย สมฺปยุตฺเตหิ ธมฺเมหิ เย ธมฺมา วิปฺปยุตฺตา, เตหิ ธมฺเมหิ เย ธมฺมา วิปฺปยุตฺตา. เต ธมฺมา เอเกน ขนฺเธน ทสหายตเนหิ เอกาทสหิ ธาตูหิ วิปฺปยุตฺตา, เอเกนายตเนน เอกาย ธาตุยา เกหิจิ วิปฺปยุตฺตา.}

\proseitem{372}{วิปาเกหิ ธมฺเมหิ เย ธมฺมา วิปฺปยุตฺตา, เตหิ ธมฺเมหิ เย ธมฺมา วิปฺปยุตฺตา. เต ธมฺมา เอเกน ขนฺเธน ทสหายตเนหิ ทสหิ ธาตูหิ วิปฺปยุตฺตา, เอเกนายตเนน เอกาย ธาตุยา เกหิจิ วิปฺปยุตฺตา.}

\proseitem{373}{วิปาก\-ธมฺม\-ธมฺเมหิ เย ธมฺมา. สํกิลิฏฺฐ\-สํกิเลสิเกหิ ธมฺเมหิ เย ธมฺมา วิปฺปยุตฺตา, เตหิ ธมฺเมหิ เย ธมฺมา วิปฺปยุตฺตา. เต ธมฺมา เอเกน ขนฺเธน ทสหายตเนหิ โสฬสหิ ธาตูหิ วิปฺปยุตฺตา, เอเกนายตเนน เอกาย ธาตุยา เกหิจิ วิปฺปยุตฺตา.}

\proseitem{374}{เนว\-วิปาก\-น\-วิปากธมฺม\-ธมฺเมหิ เย ธมฺมา. อนุปาทินฺนุ\-ปาทานิเยหิ ธมฺเมหิ เย ธมฺมา วิปฺปยุตฺตา, เตหิ ธมฺเมหิ เย ธมฺมา วิปฺปยุตฺตา. เต ธมฺมา น เกหิจิ ขนฺเธหิ น เกหิจิ อายตเนหิ ปญฺจหิ ธาตูหิ วิปฺปยุตฺตา.}

\proseitem{375}{อนุปาทินฺนอนุปาทาเยหิ ธมฺเมหิ เย ธมฺมา. อสํกิลิฏฺฐ- อสํกิเลสิเกหิ ธมฺเมหิ เย ธมฺมา วิปฺปยุตฺตา, เตหิ ธมฺเมหิ เย ธมฺมา วิปฺปยุตฺตา. เต ธมฺมา น เกหิจิ ขนฺเธหิ น เกหิจิ อายตเนหิ ฉหิ ธาตูหิ วิปฺปยุตฺตา.}

\proseitem{376}{สวิตกฺก\-สวิจาเรหิ ธมฺเมหิ เย ธมฺมา วิปฺปยุตฺตา, เตหิ ธมฺเมหิ เย ธมฺมา วิปฺปยุตฺตา. เต ธมฺมา เอเกน ขนฺเธน ทสหายตเนหิ ปนฺนรสหิ ธาตูหิ วิปฺปยุตฺตา, เอเกนายตเนน เอกาย ธาตุยา เกหิจิ วิปฺปยุตฺตา.}

\proseitem{377}{\csromanfolio{69}อวิตกฺก\-วิจารมตฺเตหิ ธมฺเมหิ เย ธมฺมา. ปีติสหคเตหิ ธมฺเมหิ เย ธมฺมา วิปฺปยุตฺตา, เตหิ ธมฺเมหิ เย ธมฺมา วิปฺปยุตฺตา. เต ธมฺมา เอเกน ขนฺเธน ทสหายตเนหิ โสฬสหิ ธาตูหิ วิปฺปยุตฺตา, เอเกนายตเนน เอกาย ธาตุยา เกหิจิ วิปฺปยุตฺตา.}

\proseitem{378}{อวิตกฺก\-อวิจาเรหิ ธมฺเมหิ เย ธมฺมา วิปฺปยุตฺตา, เตหิ ธมฺเมหิ เย ธมฺมา วิปฺปยุตฺตา. เต ธมฺมา น เกหิจิ ขนฺเธหิ น เกหิจิ อายตเนหิ เอกาย ธาตุยา วิปฺปยุตฺตา.}

\proseitem{379}{สุขสหคเตหิ ธมฺเมหิ เย ธมฺมา วิปฺปยุตฺตา, เตหิ ธมฺเมหิ เย ธมฺมา วิปฺปยุตฺตา. เต ธมฺมา เอเกน ขนฺเธน ทสหายตเนหิ ปนฺนรสหิ ธาตูหิ วิปฺปยุตฺตา, เอเกนายตเนน เอกาย ธาตุยา เกหิจิ วิปฺปยุตฺตา.}

\proseitem{380}{อุเปกฺขา\-สหคเตหิ ธมฺเมหิ เย ธมฺมา วิปฺปยุตฺตา, เตหิ ธมฺเมหิ เย ธมฺมา วิปฺปยุตฺตา. เต ธมฺมา เอเกน ขนฺเธน ทสหายตเนหิ เอกาทสหิ ธาตูหิ วิปฺปยุตฺตา, เอเกนายตเนน เอกาย ธาตุยา เกหิจิ วิปฺปยุตฺตา.}

\proseitem{381}{ทสฺสเนน ปหาตพฺเพหิ ธมฺเมหิ เย ธมฺมา. ภาวนาย ปหาตพฺเพหิ ธมฺเมหิ เย ธมฺมา. ทสฺสเนน ปหาตพฺพ\-เหตุเกหิ ธมฺเมหิ เย ธมฺมา. ภาวนาย ปหาตพฺพ\-เหตุเกหิ ธมฺเมหิ เย ธมฺมา. อาจยคามีหิ ธมฺเมหิ เย ธมฺมา. อปจยคามีหิ ธมฺเมหิ เย ธมฺมา. เสกฺเขหิ ธมฺเมหิ เย ธมฺมา. อเสกฺเขหิ ธมฺเมหิ เย ธมฺมา. มหคฺคเตหิ ธมฺเมหิ เย ธมฺมา วิปฺปยุตฺตา, เตหิ ธมฺเมหิ เย ธมฺมา วิปฺปยุตฺตา. เต ธมฺมา เอเกน ขนฺเธน ทสหายตเนหิ โสฬสหิ ธาตูหิ วิปฺปยุตฺตา, เอเกนายตเนน เอกาย ธาตุยา เกหิจิ วิปฺปยุตฺตา.}

\proseitem{382}{อปฺปมาเณหิ ธมฺเมหิ เย ธมฺมา. ปณีเตหิ ธมฺเมหิ เย ธมฺมา วิปฺปยุตฺตา, เตหิ ธมฺเมหิ เย ธมฺมา วิปฺปยุตฺตา. เต ธมฺมา น เกหิจิ ขนฺเธหิ น เกหิจิ อายตเนหิ ฉหิ ธาตูหิ วิปฺปยุตฺตา.}

\proseitem{383}{ปริตฺตา\-รมฺมเณหิ ธมฺเมหิ เย ธมฺมา วิปฺปยุตฺตา, เตหิ ธมฺเมหิ เย ธมฺมา วิปฺปยุตฺตา. เต ธมฺมา เอเกน ขนฺเธน ทสหายตเนหิ \csromanfolio{70}ทสหิ ธาตูหิ วิปฺปยุตฺตา, เอเกนายตเนน เอกาย ธาตุยา เกหิจิ วิปฺปยุตฺตา.}

\proseitem{384}{มหคฺคตา\-รมฺมเณหิ ธมฺเมหิ เย ธมฺมา. อปฺปมาณา\-รมฺมเณหิ ธมฺเมหิ เย ธมฺมา. หีเนหิ ธมฺเมหิ เย ธมฺมา. มิจฺฉตฺต\-นิยเตหิ ธมฺเมหิ เย ธมฺมา. สมฺมตฺต\-นิยเตหิ ธมฺเมหิ เย ธมฺมา. มคฺคารมฺมเณหิ ธมฺเมหิ เย ธมฺมา. มคฺคเหตุเกหิ ธมฺเมหิ เย ธมฺมา. มคฺคาธิปตีหิ ธมฺเมหิ เย ธมฺมา วิปฺปยุตฺตา, เตหิ ธมฺเมหิ เย ธมฺมา วิปฺปยุตฺตา. เต ธมฺมา เอเกน ขนฺเธน ทสหายตเนหิ โสฬสหิ ธาตูหิ วิปฺปยุตฺตา, เอเกนายตเนน เอกาย ธาตุยา เกหิจิ วิปฺปยุตฺตา.}

\proseitem{385}{อนุปฺปนฺเนหิ ธมฺเมหิ เย ธมฺมา วิปฺปยุตฺตา, เตหิ ธมฺเมหิ เย ธมฺมา วิปฺปยุตฺตา. เต ธมฺมา น เกหิจิ ขนฺเธหิ น เกหิจิ อายตเนหิ ปญฺจหิ ธาตูหิ วิปฺปยุตฺตา.}

\proseitem{386}{อตีตารมฺมเณหิ ธมฺเมหิ เย ธมฺมา. อนาคตา\-รมฺมเณหิ ธมฺเมหิ เย ธมฺมา วิปฺปยุตฺตา, เตหิ ธมฺเมหิ เย ธมฺมา วิปฺปยุตฺตา. เต ธมฺมา เอเกน ขนฺเธน ทสหายตเนหิ โสฬสหิ ธาตูหิ วิปฺปยุตฺตา, เอเกนายตเนน เอกาย ธาตุยา เกหิจิ วิปฺปยุตฺตา.}

\proseitem{387}{ปจฺจุปฺปนฺนา\-รมฺมเณหิ ธมฺเมหิ เย ธมฺมา. อชฺฌตฺตา\-รมฺมเณหิ ธมฺเมหิ เย ธมฺมา. พหิทฺธา\-รมฺมเณหิ ธมฺเมหิ เย ธมฺมา. อชฺฌตฺตพหิทฺธา\-รมฺมเณหิ ธมฺเมหิ เย ธมฺมา วิปฺปยุตฺตา, เตหิ ธมฺเมหิ เย ธมฺมา วิปฺปยุตฺตา. เต ธมฺมา เอเกน ขนฺเธน ทสหายตเนหิ ทสหิ ธาตูหิ วิปฺปยุตฺตา, เอเกนายตเนน เอกาย ธาตุยา เกหิจิ วิปฺปยุตฺตา.}

\proseitem{388}{สนิทสฺสน\-สปฺปฏิเฆหิ ธมฺเมหิ เย ธมฺมา. อนิทสฺสน\-สปฺปฏิเฆหิ ธมฺเมหิ เย ธมฺมา วิปฺปยุตฺตา, เตหิ ธมฺเมหิ เย ธมฺมา วิปฺปยุตฺตา. เต ธมฺมา จตูหิ ขนฺเธหิ เอเกนายตเนน สตฺตหิ ธาตูหิ วิปฺปยุตฺตา, เอเกนายตเนน เอกาย ธาตุยา เกหิจิ วิปฺปยุตฺตา.}

\csromanmatikaanchor{mtk.44}
\csromantocmarkref{subsection}{2. ทุก}{mtk.44}
\bookmark[dest={mtk.44},level=2]{2. ทุก}
\csromanheader{2}{2. \textbf{ทุก}}
\proseitem{389}{เหตูหิ ธมฺเมหิ เย ธมฺมา. สเหตุเกหิ ธมฺเมหิ เย ธมฺมา. เหตุ\-สมฺปยุตฺเตหิ ธมฺเมหิ เย ธมฺมา. เหตูหิ เจว สเหตุเกหิ จ \csromanfolio{71}ธมฺเมหิ เย ธมฺมา. สเหตุเกหิ เจว น จ เหตูหิ ธมฺเมหิ เย ธมฺมา. เหตูหิ เจว เหตุ\-สมฺปยุตฺเตหิ จ ธมฺเมหิ เย ธมฺมา. เหตุ\-สมฺปยุตฺเตหิ เจว น จ เหตูหิ ธมฺเมหิ เย ธมฺมา. น เหตุสเหตุเกหิ ธมฺเมหิ เย ธมฺมา วิปฺปยุตฺตา, เตหิ ธมฺเมหิ เย ธมฺมา วิปฺปยุตฺตา. เต ธมฺมา เอเกน ขนฺเธน ทสหายตเนหิ โสฬสหิ ธาตูหิ วิปฺปยุตฺตา, เอเกนายตเนน เอกาย ธาตุยา เกหิจิ วิปฺปยุตฺตา.}

\proseitem{390}{อปฺปจฺจเยหิ ธมฺเมหิ เย ธมฺมา. อสงฺขเตหิ ธมฺเมหิ เย ธมฺมา. สนิทสฺสเนหิ ธมฺเมหิ เย ธมฺมา. สปฺปฏิเฆหิ ธมฺเมหิ เย ธมฺมา. รูปีหิ ธมฺเมหิ เย ธมฺมา วิปฺปยุตฺตา, เตหิ ธมฺเมหิ เย ธมฺมา วิปฺปยุตฺตา. เต ธมฺมา จตูหิ ขนฺเธหิ เอเกนายตเนน สตฺตหิ ธาตูหิ วิปฺปยุตฺตา, เอเกนายตเนน เอกาย ธาตุยา เกหิจิ วิปฺปยุตฺตา.}

\proseitem{391}{โลกุตฺตเรหิ ธมฺเมหิ เย ธมฺมา วิปฺปยุตฺตา, เตหิ ธมฺเมหิ เย ธมฺมา วิปฺปยุตฺตา. เต ธมฺมา น เกหิจิ ขนฺเธหิ น เกหิจิ อายตเนหิ ฉหิ ธาตูหิ วิปฺปยุตฺตา.}

\proseitem{392}{อาสเวหิ ธมฺเมหิ เย ธมฺมา. อาสว\-สมฺปยุตฺเตหิ ธมฺเมหิ เย ธมฺมา. อาสเวหิ เจว สาสเวหิ จ ธมฺเมหิ เย ธมฺมา. อาสเวหิ เจว อาสว\-สมฺปยุตฺเตหิ จ ธมฺเมหิ เย ธมฺมา. อาสว\-สมฺปยุตฺเตหิ เจว โน จ อาสเวหิ ธมฺเมหิ เย ธมฺมา วิปฺปยุตฺตา, เตหิ ธมฺเมหิ เย ธมฺมา วิปฺปยุตฺตา. เต ธมฺมา เอเกน ขนฺเธน ทสหายตเนหิ โสฬสหิ ธาตูหิ วิปฺปยุตฺตา, เอเกนายตเนน เอกาย ธาตุยา เกหิจิ วิปฺปยุตฺตา.}

\proseitem{393}{อนาสเวหิ ธมฺเมหิ เย ธมฺมา. อาสว\-วิปฺปยุตฺเตหิ อนาสเวหิ ธมฺเมหิ เย ธมฺมา วิปฺปยุตฺตา, เตหิ ธมฺเมหิ เย ธมฺมา วิปฺปยุตฺตา. เต ธมฺมา น เกหิจิ ขนฺเธหิ น เกหิจิ อายตเนหิ ฉหิ ธาตูหิ วิปฺปยุตฺตา.}

\proseitem{394}{สญฺโญชเนหิ ธมฺเมหิ เย ธมฺมา. คนฺเถหิ ธมฺเมหิ เย ธมฺมา. โอเฆหิ ธมฺเมหิ เย ธมฺมา. โยเคหิ ธมฺเมหิ เย ธมฺมา. นีวรเณหิ ธมฺเมหิ เย ธมฺมา. ปรามาเสหิ ธมฺเมหิ เย ธมฺมา. ปรามาส\-สมฺปยุตฺเตหิ \csromanfolio{72}ธมฺเมหิ เย ธมฺมา. ปรามาเสหิ เจว ปรามฏฺเฐหิ จ ธมฺเมหิ เย ธมฺมา วิปฺปยุตฺตา, เตหิ ธมฺเมหิ เย ธมฺมา วิปฺปยุตฺตา. เต ธมฺมา เอเกน ขนฺเธน ทสหายตเนหิ โสฬสหิ ธาตูหิ วิปฺปยุตฺตา, เอเกนายตเนน เอกาย ธาตุยา เกหิจิ วิปฺปยุตฺตา.}

\proseitem{395}{อปรามฏฺเฐหิ ธมฺเมหิ เย ธมฺมา. ปรามาส\-วิปฺปยุตฺเตหิ อปรามฏฺเฐหิ ธมฺเมหิ เย ธมฺมา วิปฺปยุตฺตา, เตหิ ธมฺเมหิ เย ธมฺมา วิปฺปยุตฺตา. เต ธมฺมา น เกหิจิ ขนฺเธหิ น เกหิจิ อายตเนหิ ฉหิ ธาตูหิ วิปฺปยุตฺตา.}

\proseitem{396}{สารมฺมเณหิ ธมฺเมหิ เย ธมฺมา. จิตฺเตหิ ธมฺเมหิ เย ธมฺมา. เจตสิเกหิ ธมฺเมหิ เย ธมฺมา. จิตฺต\-สมฺปยุตฺเตหิ ธมฺเมหิ เย ธมฺมา. จิตฺต\-สํสฏฺเฐหิ ธมฺเมหิ เย ธมฺมา. จิตฺต\-สํสฏฺฐ\-สมุฏฺฐาเนหิ ธมฺเมหิ เย ธมฺมา. จิตฺต\-สํสฏฺฐ\-สมุฏฺฐาน\-สหภูหิ ธมฺเมหิ เย ธมฺมา. จิตฺต\-สํสฏฺฐ\-สมุฏฺฐานา\-นุปริวตฺตีหิ ธมฺเมหิ เย ธมฺมา วิปฺปยุตฺตา, เตหิ ธมฺเมหิ เย ธมฺมา วิปฺปยุตฺตา. เต ธมฺมา เอเกน ขนฺเธน ทสหายตเนหิ ทสหิ ธาตูหิ วิปฺปยุตฺตา, เอเกนายตเนน เอกาย ธาตุยา เกหิจิ วิปฺปยุตฺตา.}

\proseitem{397}{อนารมฺมเณหิ ธมฺเมหิ เย ธมฺมา. จิตฺต\-วิปฺปยุตฺเตหิ ธมฺเมหิ เย ธมฺมา. จิตฺต\-วิสํสฏฺเฐหิ ธมฺเมหิ เย ธมฺมา. อุปาทาธมฺเมหิ เย ธมฺมา วิปฺปยุตฺตา, เตหิ ธมฺเมหิ เย ธมฺมา วิปฺปยุตฺตา. เต ธมฺมา จตูหิ ขนฺเธหิ เอเกนายตเนน สตฺตหิ ธาตูหิ วิปฺปยุตฺตา, เอเกนายตเนน เอกาย ธาตุยา เกหิจิ วิปฺปยุตฺตา.}

\proseitem{398}{อนุปาทินฺเนหิ\csromansharedfootnote{7824cd9ba3be2934}{อนุปาทิณฺเณหิ (สี, ก)} ธมฺเมหิ เย ธมฺมา วิปฺปยุตฺตา, เตหิ ธมฺเมหิ เย ธมฺมา วิปฺปยุตฺตา, เต ธมฺมา น เกหิจิ ขนฺเธหิ น เกหิจิ อายตเนหิ ปญฺจหิ ธาตูหิ วิปฺปยุตฺตา.}

\proseitem{399}{อุปาทาเนหิ ธมฺเมหิ เย ธมฺมา. กิเลเสหิ ธมฺเมหิ เย ธมฺมา. สํกิลิฏฺเฐหิ ธมฺเมหิ เย ธมฺมา. กิเลส\-สมฺปยุตฺเตหิ ธมฺเมหิ เย ธมฺมา. กิเลเสหิ เจว สํกิเลสิเกหิ จ ธมฺเมหิ เย ธมฺมา. กิเลเสหิ เจว สํกิลิฏฺเฐหิ จ ธมฺเมหิ เย ธมฺมา. สํกิลิฏฺเฐหิ เจว \csromanfolio{73}โน จ กิเลเสหิ ธมฺเมหิ เย ธมฺมา. กิเลเสหิ เจว กิเลส\-สมฺปยุตฺเตหิ จ ธมฺเมหิ เย ธมฺมา. กิเลส\-สมฺปยุตฺเตหิ เจว โน จ กิเลเสหิ ธมฺเมหิ เย ธมฺมา วิปฺปยุตฺตา, เตหิ ธมฺเมหิ เย ธมฺมา วิปฺปยุตฺตา. เต ธมฺมา เอเกน ขนฺเธน ทสหายตเนหิ โสฬสหิ ธาตูหิ วิปฺปยุตฺตา, เอเกนายตเนน เอกาย ธาตุยา เกหิจิ วิปฺปยุตฺตา.}

\proseitem{400}{อสํกิเลสิเกหิ ธมฺเมหิ เย ธมฺมา,. กิเลส\-วิปฺปยุตฺเตหิ อสํกิเลสิเกหิ ธมฺเมหิ เย ธมฺมา วิปฺปยุตฺตา, เตหิ ธมฺเมหิ เย ธมฺมา วิปฺปยุตฺตา. เต ธมฺมา น เกหิจิ ขนฺเธหิ น เกหิจิ อายตเนหิ ฉหิ ธาตูหิ วิปฺปยุตฺตา.}

\proseitem{401}{ทสฺสเนน ปหาตพฺเพหิ ธมฺเมหิ เย ธมฺมา. ภาวนาย ปหาตพฺเพหิ ธมฺเมหิ เย ธมฺมา. ทสฺสเนน ปหาตพฺพ\-เหตุเกหิ ธมฺเมหิ เย ธมฺมา. ภาวนาย ปหาตพฺพ\-เหตุเกหิ ธมฺเมหิ เย ธมฺมา วิปฺปยุตฺตา, เตหิ ธมฺเมหิ เย ธมฺมา วิปฺปยุตฺตา. เต ธมฺมา เอเกน ขนฺเธน ทสหายตเนหิ โสฬสหิ ธาตูหิ วิปฺปยุตฺตา, เอเกนายตเนน เอกาย ธาตุยา เกหิจิ วิปฺปยุตฺตา.}

\proseitem{402}{สวิตกฺเกหิ ธมฺเมหิ เย ธมฺมา. สวิจาเรหิ ธมฺเมหิ เย ธมฺมา วิปฺปยุตฺตา, เตหิ ธมฺเมหิ เย ธมฺมา วิปฺปยุตฺตา. เต ธมฺมา เอเกน ขนฺเธน ทสหายตเนหิ ปนฺนรสหิ ธาตูหิ วิปฺปยุตฺตา, เอเกนายตเนน เอกาย ธาตุยา เกหิจิ วิปฺปยุตฺตา.}

\proseitem{403}{อวิตกฺเกหิ ธมฺเมหิ เย ธมฺมา. อวิจาเรหิ ธมฺเมหิ เย ธมฺมา วิปฺปยุตฺตา, เตหิ ธมฺเมหิ เย ธมฺมา วิปฺปยุตฺตา. เต ธมฺมา น เกหิจิ ขนฺเธหิ น เกหิจิ อายตเนหิ เอกาย ธาตุยา วิปฺปยุตฺตา.}

\proseitem{404}{สปฺปีติเกหิ ธมฺเมหิ เย ธมฺมา. ปีติสหคเตหิ ธมฺเมหิ เย ธมฺมา วิปฺปยุตฺตา, เตหิ ธมฺเมหิ เย ธมฺมา วิปฺปยุตฺตา. เต ธมฺมา เอเกน ขนฺเธน ทสหายตเนหิ โสฬสหิ ธาตูหิ วิปฺปยุตฺตา, เอเกนายตเนน เอกาย ธาตุยา เกหิจิ วิปฺปยุตฺตา.}

\proseitem{405}{สุขสหคเตหิ ธมฺเมหิ เย ธมฺมา วิปฺปยุตฺตา, เตหิ ธมฺเมหิ เย ธมฺมา วิปฺปยุตฺตา. เต ธมฺมา เอเกน ขนฺเธน ทสหายตเนหิ ปนฺนรสหิ \csromanfolio{74}ธาตูหิ วิปฺปยุตฺตา, เอเกนายตเนน เอกาย ธาตุยา เกหิจิ วิปฺปยุตฺตา.}

\proseitem{406}{อุเปกฺขา\-สหคเตหิ ธมฺเมหิ เย ธมฺมา วิปฺปยุตฺตา, เตหิ ธมฺเมหิ เย ธมฺมา วิปฺปยุตฺตา. เต ธมฺมา เอเกน ขนฺเธน ทสหายตเนหิ เอกาทสหิ ธาตูหิ วิปฺปยุตฺตา, เอเกนายตเนน เอกาย ธาตุยา เกหิจิ วิปฺปยุตฺตา.}

\proseitem{407}{น กามาวจเรหิ ธมฺเมหิ เย ธมฺมา. อปริยาปนฺเนหิ ธมฺเมหิ เย ธมฺมา. อนุตฺตเรหิ ธมฺเมหิ เย ธมฺมา วิปฺปยุตฺตา, เตหิ ธมฺเมหิ เย ธมฺมา วิปฺปยุตฺตา. เต ธมฺมา น เกหิจิ ขนฺเธหิ น เกหิจิ อายตเนหิ ฉหิ ธาตูหิ วิปฺปยุตฺตา.}

\proseitem{408}{รูปาวจเรหิ ธมฺเมหิ เย ธมฺมา. อรูปาวจเรหิ ธมฺเมหิ เย ธมฺมา. นิยฺยานิเกหิ ธมฺเมหิ เย ธมฺมา. นิยเตหิ ธมฺเมหิ เย ธมฺมา. สรเณหิ ธมฺเมหิ เย ธมฺมา วิปฺปยุตฺตา, เตหิ ธมฺเมหิ เย ธมฺมา วิปฺปยุตฺตา. เต ธมฺมา กติหิ ขนฺเธหิ กติหายตเนหิ กติหิ ธาตูหิ วิปฺปยุตฺตา, เต ธมฺมา เอเกน ขนฺเธน ทสหายตเนหิ โสฬสหิ ธาตูหิ วิปฺปยุตฺตา, เอเกนายตเนน เอกาย ธาตุยา เกหิจิ วิปฺปยุตฺตา.}

\setlength{\csromangathapagewidth}{0pt}
\setlength{\csromangathawakwidth}{0pt}
\csromangathameasuregroup{ธมฺมายตนํ ธมฺมธาตุ, \\ สฬายตนํ นามรูปํ, \\ ชาติ ชรา จ มรณํ, \\ โคจฺฉเกสุ จ ปญฺญาส, \\ มหนฺตเร ปนฺนรส, \\ เตวีส ปทสตํ เอตํ,}{\csromangathabat{ธมฺมายตนํ ธมฺมธาตุ,}{ทุกฺขสจฺจญฺจ ชีวิตํ.} \\ \csromangathabat{สฬายตนํ นามรูปํ,}{จตฺตาโร จ มหาภวา.} \\ \csromangathabat{ชาติ ชรา จ มรณํ,}{ติเกเสฺวกูนวีสติ.} \\ \csromangathabat{โคจฺฉเกสุ จ ปญฺญาส,}{อฏฺฐ จูฬนฺตเร ปทา.} \\ \csromangathabat{มหนฺตเร ปนฺนรส,}{อฏฺฐารส ตโต ปเร.} \\ \csromangathabat{เตวีส ปทสตํ เอตํ,}{สมฺปโยเค น ลพฺภตีติ.}}
\csromangathasetleft{ธมฺมายตนํ ธมฺมธาตุ, \\ สฬายตนํ นามรูปํ, \\ ชาติ ชรา จ มรณํ, \\ โคจฺฉเกสุ จ ปญฺญาส, \\ มหนฺตเร ปนฺนรส, \\ เตวีส ปทสตํ เอตํ,}
\csromangathagroup{\csromangathastanza{\csromangathabat{ธมฺมายตนํ ธมฺมธาตุ,}{ทุกฺขสจฺจญฺจ ชีวิตํ.} \\ \csromangathabat{สฬายตนํ นามรูปํ,}{จตฺตาโร จ มหาภวา.}}\csromangathastanza{\csromangathabat{ชาติ ชรา จ มรณํ,}{ติเกเสฺวกูนวีสติ.} \\ \csromangathabat{โคจฺฉเกสุ จ ปญฺญาส,}{อฏฺฐ จูฬนฺตเร ปทา.}}\csromangathastanza{\csromangathabat{มหนฺตเร ปนฺนรส,}{อฏฺฐารส ตโต ปเร.} \\ \csromangathabat{เตวีส ปทสตํ เอตํ,}{สมฺปโยเค น ลพฺภตีติ.}}}

\vspace{\tipitakaparskip}
\csromancenter{วิปฺปยุตฺเตน วิปฺปยุตฺต\-ปท\-นิทฺเทโส ทสโม.}
\csromanmatikaanchor{mtk.45}
\csromantocmarknopage{section}{11. เอกาทสมนย}{mtk.45}
\bookmark[dest={mtk.45},level=1]{11. เอกาทสมนย}
\csromanheader{1}{11. \textbf{เอกาทสมนย}}
\csromanmatikaanchor{mtk.46}
\csromantocmarkref{subsection}{11. สงฺคหิเตนสมฺปยุตฺตวิปฺปยุตฺตปทนิทฺเทส}{mtk.46}
\bookmark[dest={mtk.46},level=2]{11. สงฺคหิเตนสมฺปยุตฺตวิปฺปยุตฺตปทนิทฺเทส}
\csromanheader{2}{11. สงฺคหิเตน\-สมฺปยุตฺต\-วิปฺปยุตฺต\-ปท\-นิทฺเทส}
\proseitem{409}{\csromanfolio{75}สมุทย\-สจฺเจน เย ธมฺมา. มคฺคสจฺเจน เย ธมฺมา ขนฺธ\-สงฺคเหน สงฺคหิตา อายตนสงฺคหิตา สงฺคหิตา ธาตุสงฺคเหน สงฺคหิตา. เต ธมฺมา กติหิ ขนฺเธหิ กติหายตเนหิ กติหิ ธาตูหิ สมฺปยุตฺตา, เต ธมฺมา ตีหิ ขนฺเธหิ เอเกนายตเนน สตฺตหิ ธาตูหิ สมฺปยุตฺตา, เอเกน ขนฺเธน เอเกนายตเนน เอกาย ธาตุยา เกหิจิ สมฺปยุตฺตา. กติหิ วิปฺปยุตฺตา, เอเกน ขนฺเธน ทสหายตเนหิ ทสหิ ธาตูหิ วิปฺปยุตฺตา, เอเกนายตเนน เอกาย ธาตุยา เกหิจิ วิปฺปยุตฺตา.}

\proseitem{410}{อิตฺถินฺทฺริเยน เย ธมฺมา. ปุริสินฺทฺริเยน เย ธมฺมา ขนฺธ\-สงฺคเหน สงฺคหิตา อายตน\-สงฺคเหน สงฺคหิตา ธาตุสงฺคเหน สงฺคหิตา. เต ธมฺมา กติหิ ขนฺเธหิ กติหายตเนหิ กติหิ ธาตูหิ สมฺปยุตฺตาติ นตฺถิ. กติหิ วิปฺปยุตฺตา, จตูหิ ขนฺเธหิ เอเกนายตเนน สตฺตหิ ธาตูหิ วิปฺปยุตฺตา, เอเกนายตเนน เอกาย ธาตุยา เกหิจิ วิปฺปยุตฺตา.}

\proseitem{411}{สุขินฺทฺริเยน เย ธมฺมา. ทุกฺขินฺทฺริเยน เย ธมฺมา. โสมนสฺสิ\-นฺทฺริเยน เย ธมฺมา. โทมนสฺสิ\-นฺทฺริเยน เย ธมฺมา ขนฺธ\-สงฺคเหน สงฺคหิตา อายตน\-สงฺคเหน สงฺคหิตา ธาตุสงฺคเหน สงฺคหิตา ฯเปฯ เต ธมฺมา ตีหิ ขนฺเธหิ เอเกนายตเนน สตฺตหิ ธาตูหิ สมฺปยุตฺตา, เอเกนายตเนน เอกาย ธาตุยา เกหิจิ สมฺปยุตฺตา. กติหิ วิปฺปยุตฺตา, เอเกน ขนฺเธน ทสหายตเนหิ ทสหิ ธาตูหิ วิปฺปยุตฺตา, เอเกนายตเนน เอกาย ธาตุยา เกหิจิ วิปฺปยุตฺตา.}

\proseitem{412}{อุเปกฺขิ\-นฺทฺริเยน เย ธมฺมา ขนฺธ\-สงฺคเหน สงฺคหิตา อายตน\-สงฺคเหน สงฺคหิตา ธาตุสงฺคเหน สงฺคหิตา. เต ธมฺมา ตีหิ ขนฺเธหิ เอเกนายตเนน ทฺวีหิ ธาตูหิ สมฺปยุตฺตา, เอเกนายตเนน เอกาย ธาตุยา เกหิจิ สมฺปยุตฺตา. กติหิ วิปฺปยุตฺตา, เอเกน ขนฺเธน ทสหายตเนหิ ปนฺนรสหิ ธาตูหิ วิปฺปยุตฺตา, เอเกนายตเนน เอกาย ธาตุยา เกหิจิ วิปฺปยุตฺตา.}

\proseitem{413}{\csromanfolio{76}สทฺธินฺทฺริเยน เย ธมฺมา. วีริยินฺทฺริเยน เย ธมฺมา. สตินฺทฺริเยน เย ธมฺมา. สมาธิ\-นฺทฺริเยน เย ธมฺมา. ปญฺญินฺทฺริเยน เย ธมฺมา. อนญฺญาต\-ญฺญสฺสามี\-ตินฺทฺริเยน เย ธมฺมา. อญฺญินฺทฺริเยน เย ธมฺมา. อญฺญาตาวิ\-นฺทฺริเยน เย ธมฺมา. อวิชฺชาย เย ธมฺมา. อวิชฺชาปจฺจยา สงฺขาเรหิ เย ธมฺมา. สฬายตน\-ปจฺจยา ผสฺเสน เย ธมฺมา. เวทนาปจฺจยา ตณฺหาย เย ธมฺมา. ตณฺหาปจฺจยา อุปาทาเนน เย ธมฺมา. กมฺมภเวน เย ธมฺมา ขนฺธ\-สงฺคเหน สงฺคหิตา อายตน\-สงฺคเหน สงฺคหิตา ธาตุสงฺคเหน สงฺคหิตา. เต ธมฺมา ตีหิ ขนฺเธหิ เอเกนายตเนน สตฺตหิ ธาตูหิ สมฺปยุตฺตา, เอเกน ขนฺเธน เอเกนายตเนน เอกาย ธาตุยา เกหิจิ สมฺปยุตฺตา. กติหิ วิปฺปยุตฺตา, เอเกน ขนฺเธน ทสหายตเนหิ ทสหิ ธาตูหิ วิปฺปยุตฺตา, เอเกนายตเนน เอกาย ธาตุยา เกหิจิ วิปฺปยุตฺตา.}

\proseitem{414}{ปริเทเวน เย ธมฺมา ขนฺธ\-สงฺคเหน สงฺคหิตา อายตน\-สงฺคเหน สงฺคหิตา ธาตุสงฺคเหน สงฺคหิตา. เต ธมฺมา กติหิ ขนฺเธหิ กติหายตเนหิ กติหิ ธาตูหิ สมฺปยุตฺตาติ นตฺถิ. กติหิ วิปฺปยุตฺตา, จตูหิ ขนฺเธหิ เอเกนายตเนน สตฺตหิ ธาตูหิ วิปฺปยุตฺตา, เอเกนายตเนน เอกาย ธาตุยา เกหิจิ วิปฺปยุตฺตา.}

\proseitem{415}{โสเกน เย ธมฺมา. ทุกฺเขน เย ธมฺมา. โทมนสฺเสน เย ธมฺมา ขนฺธ\-สงฺคเหน สงฺคหิตา อายตน\-สงฺคเหน สงฺคหิตา ธาตุสงฺคเหน สงฺคหิตา. เต ธมฺมา ตีหิ ขนฺเธหิ เอเกนายตเนน สตฺตหิ ธาตูหิ สมฺปยุตฺตา, เอเกนายตเนน เอกาย ธาตุยา เกหิจิ สมฺปยุตฺตา. กติหิ วิปฺปยุตฺตา, เอเกน ขนฺเธน ทสหายตเนหิ ทสหิ ธาตูหิ วิปฺปยุตฺตา, เอเกนายตเนน เอกาย ธาตุยา เกหิจิ วิปฺปยุตฺตา.}

\proseitem{416}{อุปายาเสน เย ธมฺมา. สติปฏฺฐาเนน เย ธมฺมา. สมฺม\-ปฺปธาเนน เย ธมฺมา. อปฺปมญฺญาย เย ธมฺมา. ปญฺจหิ อินฺทฺริเยหิ เย ธมฺมา. ปญฺจหิ พเลหิ เย ธมฺมา, สตฺตหิ โพชฺฌงฺเคหิ เย ธมฺมา. อริเยน อฏฺฐงฺคิเกน มคฺเคน เย ธมฺมา. ผสฺเสน เย ธมฺมา. เจตนาย เย ธมฺมา. อธิโมกฺเขน เย ธมฺมา. มนสิกาเรน เย ธมฺมา. เหตูหิ ธมฺเมหิ เย ธมฺมา. เหตูหิ เจว สเหตุเกหิ จ ธมฺเมหิ เย ธมฺมา. เหตูหิ เจว เหตุ\-สมฺปยุตฺเตหิ จ ธมฺเมหิ เย ธมฺมา. อาสเวหิ ธมฺเมหิ เย ธมฺมา. อาสเวหิ เจว สาสเวหิ จ ธมฺเมหิ เย ธมฺมา. อาสเวหิ เจว}

\vspace{\tipitakaparskip}
\csromancenter{เอกาทสมนย อาสว\-สมฺปยุตฺเตหิ จ ธมฺเมหิ เย ธมฺมา. สญฺโญชเนหิ ธมฺเมหิ เย ธมฺมา. คนฺเถหิ ธมฺเมหิ เย ธมฺมา. โอเฆหิ ธมฺเมหิ เย ธมฺมา. โยเคหิ ธมฺเมหิ เย ธมฺมา. นีวรเณหิ ธมฺเมหิ เย ธมฺมา. ปรามาเสหิ ธมฺเมหิ เย ธมฺมา. อุปาทาเนหิ ธมฺเมหิ เย ธมฺมา. กิเลเสหิ ธมฺเมหิ เย ธมฺมา. กิเลเสหิ เจว สํกิเลสิเกหิ จ ธมฺเมหิ เย ธมฺมา. กิเลเสหิ เจว สํกิลิฏฺเฐหิ จ ธมฺเมหิ เย ธมฺมา. กิเลเสหิ เจว กิเลส\-สมฺปยุตฺเตหิ จ ธมฺเมหิ เย ธมฺมา ขนฺธ\-สงฺคเหน สงฺคหิตา อายตน\-สงฺคเหน สงฺคหิตา ธาตุสงฺคเหน สงฺคหิตา. เต ธมฺมา กติหิ ขนฺเธหิ กติหายตเนหิ กติหิ ธาตูหิ สมฺปยุตฺตา, เต ธมฺมา ตีหิ ขนฺเธหิ เอเกนายตเนน สตฺตหิ ธาตูหิ สมฺปยุตฺตา, เอเกน ขนฺเธน เอเกนายตเนน เอกาย ธาตุยา เกหิจิ สมฺปยุตฺตา. กติหิ วิปฺปยุตฺตา, เอเกน ขนฺเธน ทสหายตเนหิ ทสหิ ธาตูหิ วิปฺปยุตฺตา, เอเกนายตเนน เอกาย ธาตุยา เกหิจิ วิปฺปยุตฺตา.}
\setlength{\csromangathapagewidth}{0pt}
\setlength{\csromangathawakwidth}{0pt}
\csromangathameasuregroup{เทฺว สจฺจา ปนฺนรสินฺทฺริยา, \\ อุทฺธํ ปุน เอกาทส,}{\csromangathabat{เทฺว สจฺจา ปนฺนรสินฺทฺริยา,}{เอกาทส ปฏิจฺจปทา.} \\ \csromangathabat{อุทฺธํ ปุน เอกาทส,}{โคจฺฉกปทเมตฺถ ติํสวิธาติ.}}
\csromangathasetleft{เทฺว สจฺจา ปนฺนรสินฺทฺริยา, \\ อุทฺธํ ปุน เอกาทส,}
\csromangathagroup{\csromangathastanza{\csromanfolio{77}\csromangathabat{เทฺว สจฺจา ปนฺนรสิ\-นฺทฺริยา,}{เอกาทส ปฏิจฺจปทา.} \\ \csromangathabat{อุทฺธํ ปุน เอกาทส,}{โคจฺฉกปทเมตฺถ ติํสวิธาติ.}}}

\prose{สงฺคหิเตน สมฺปยุตฺต\-วิปฺปยุตฺต\-ปท\-นิทฺเทโส เอกาทสโม.}

\csromanmatikaanchor{mtk.47}
\csromantocmarknopage{section}{12. ทฺวาทสมนย}{mtk.47}
\bookmark[dest={mtk.47},level=1]{12. ทฺวาทสมนย}
\csromanheader{1}{12. \textbf{ทฺวาทสมนย}}
\csromanheader{2}{12. \textbf{สมฺปยุตฺเตน สงฺคหิตาสงฺคหิตปทนิทฺเทส}}
\csromanmatikaanchor{mtk.48}
\csromantocmarkref{subsection}{12. สมฺปยุตฺเตนสงฺคหิตาสงฺคหิตปทนิทฺเทส}{mtk.48}
\bookmark[dest={mtk.48},level=2]{12. สมฺปยุตฺเตนสงฺคหิตาสงฺคหิตปทนิทฺเทส}
\proseitem{417}{\csromanfolio{78}เวทนา\-กฺขนฺเธน เย ธมฺมา. สญฺญา\-กฺขนฺเธน เย ธมฺมา. สงฺขาร\-กฺขนฺเธน เย ธมฺมา สมฺปยุตฺตา, เต ธมฺมา กติหิ ขนฺเธหิ กติหายตเนหิ กติหิ ธาตูหิ สงฺคหิตา, เต ธมฺมา ตีหิ ขนฺเธหิ ทฺวีหายตเนหิ อฏฺฐหิ ธาตูหิ สงฺคหิตา. กติหิ อสงฺคหิตา, ทฺวีหิ ขนฺเธหิ ทสหายตเนหิ ทสหิ ธาตูหิ อสงฺคหิตา.}

\proseitem{418}{วิญฺญาณ\-กฺขนฺเธน เย ธมฺมา. มนายตเนน เย ธมฺมา. จกฺขุวิญฺญาณ\-ธาตุยา เย ธมฺมา ฯเปฯ. มโนธาตุยา เย ธมฺมา. มโนวิญฺญาณ\-ธาตุยา เย ธมฺมา สมฺปยุตฺตา, เต ธมฺมา ตีหิ ขนฺเธหิ เอเกนายตเนน เอกาย ธาตุยา สงฺคหิตา. กติหิ อสงฺคหิตา, ทฺวีหิ ขนฺเธหิ เอกาทสหา\-ยตเนหิ สตฺตรสหิ ธาตูหิ อสงฺคหิตา.}

\proseitem{419}{สมุทย\-สจฺเจน เย ธมฺมา. มคฺคสจฺเจน เย ธมฺมา สมฺปยุตฺตา, เต ธมฺมา จตูหิ ขนฺเธหิ ทฺวีหายตเนหิ ทฺวีหิ ธาตูหิ สงฺคหิตา. กติหิ อสงฺคหิตา, เอเกน ขนฺเธน ทสหายตเนหิ โสฬสหิ ธาตูหิ อสงฺคหิตา.}

\proseitem{420}{มนินฺทฺริเยน เย ธมฺมา สมฺปยุตฺตา, เต ธมฺมา ตีหิ ขนฺเธหิ เอเกนายตเนน เอกาย ธาตุยา สงฺคหิตา. กติหิ อสงฺคหิตา, ทฺวีหิ ขนฺเธหิ เอกาทสหา\-ยตเนหิ สตฺตรสหิ ธาตูหิ อสงฺคหิตา.}

\proseitem{421}{สุขินฺทฺริเยน เย ธมฺมา. ทุกฺขินฺทฺริเยน เย ธมฺมา. โสมนสฺสิ\-นฺทฺริเยน เย ธมฺมา. โทมนสฺสิ\-นฺทฺริเยน เย ธมฺมา สมฺปยุตฺตา, เต ธมฺมา ตีหิ ขนฺเธหิ ทฺวีหายตเนหิ ทฺวีหิ ธาตูหิ สงฺคหิตา. กติหิ อสงฺคหิตา, ทฺวีหิ ขนฺเธหิ ทสหายตเนหิ โสฬสหิ ธาตูหิ อสงฺคหิตา.}

\proseitem{422}{อุเปกฺขิ\-นฺทฺริเยน เย ธมฺมา สมฺปยุตฺตา, เต ธมฺมา ตีหิ ขนฺเธหิ ทฺวีหายตเนหิ สตฺตหิ ธาตูหิ สงฺคหิตา. กติหิ อสงฺคหิตา, ทฺวีหิ ขนฺเธหิ ทสหายตเนหิ เอกาทสหิ ธาตูหิ อสงฺคหิตา.}

\proseitem{423}{\csromanfolio{79}สทฺธินฺทฺริเยน เย ธมฺมา. วีริยินฺทฺริเยน เย ธมฺมา. สตินฺทฺริเยน เย ธมฺมา. สมาธิ\-นฺทฺริเยน เย ธมฺมา. ปญฺญินฺทฺริเยน เย ธมฺมา. อนญฺญาต\-ญฺญสฺสามี\-ตินฺทฺริเยน เย ธมฺมา. อญฺญินฺทฺริเยน เย ธมฺมา. อญฺญาตาวิ\-นฺทฺริเยน เย ธมฺมา. อวิชฺชาย เย ธมฺมา. อวิชฺชาปจฺจยา สงฺขาเรหิ เย ธมฺมา สมฺปยุตฺตา, เต ธมฺมา จตูหิ ขนฺเธหิ ทฺวีหายตเนหิ ทฺวีหิ ธาตูหิ สงฺคหิตา. กติหิ อสงฺคหิตา, เอเกน ขนฺเธน ทสหายตเนหิ โสฬสหิ ธาตูหิ อสงฺคหิตา.}

\proseitem{424}{สงฺขาร\-ปจฺจยา วิญฺญาเณน เย ธมฺมา สมฺปยุตฺตา, เต ธมฺมา ตีหิ ขนฺเธหิ เอเกนายตเนน เอกาย ธาตุยา สงฺคหิตา. กติหิ อสงฺคหิตา, ทฺวีหิ ขนฺเธหิ เอกาทสหา\-ยตเนหิ สตฺตรสหิ ธาตูหิ อสงฺคหิตา.}

\proseitem{425}{สฬายตน\-ปจฺจยา ผสฺเสน เย ธมฺมา สมฺปยุตฺตา, เต ธมฺมา จตูหิ ขนฺเธหิ ทฺวีหายตเนหิ อฏฺฐหิ ธาตูหิ สงฺคหิตา. กติหิ อสงฺคหิตา, เอเกน ขนฺเธน ทสหายตเนหิ ทสหิ ธาตูหิ อสงฺคหิตา.}

\proseitem{426}{ผสฺสปจฺจยา เวทนาย เย ธมฺมา สมฺปยุตฺตา, เต ธมฺมา ตีหิ ขนฺเธหิ ทฺวีหายตเนหิ อฏฺฐหิ ธาตูหิ สงฺคหิตา. กติหิ อสงฺคหิตา, ทฺวีหิ ขนฺเธหิ ทสหายตเนหิ ทสหิ ธาตูหิ อสงฺคหิตา.}

\proseitem{427}{เวทนาปจฺจยา ตณฺหาย เย ธมฺมา. ตณฺหาปจฺจยา อุปาทาเนน เย ธมฺมา. กมฺมภเวน เย ธมฺมา สมฺปยุตฺตา, เต ธมฺมา จตูหิ ขนฺเธหิ ทฺวีหายตเนหิ ทฺวีหิ ธาตูหิ สงฺคหิตา. กติหิ อสงฺคหิตา, เอเกน ขนฺเธน ทสหายตเนหิ โสฬสหิ ธาตูหิ อสงฺคหิตา.}

\proseitem{428}{โสเกน เย ธมฺมา. ทุกฺเขน เย ธมฺมา. โทมนสฺเสน เย ธมฺมา สมฺปยุตฺตา, เต ธมฺมา ตีหิ ขนฺเธหิ ทฺวีหายตเนหิ ทฺวีหิ ธาตูหิ สงฺคหิตา. กติหิ อสงฺคหิตา, ทฺวีหิ ขนฺเธหิ ทสหายตเนหิ โสฬสหิ ธาตูหิ อสงฺคหิตา.}

\proseitem{429}{อุปายาเสน เย ธมฺมา. สติปฏฺฐาเนน เย ธมฺมา. สมฺม\-ปฺปธาเนน เย ธมฺมา สมฺปยุตฺตา, เต ธมฺมา จตูหิ ขนฺเธหิ ทฺวีหายตเนหิ ทฺวีหิ \csromanfolio{80}ธาตูหิ สงฺคหิตา. กติหิ อสงฺคหิตา, เอเกน ขนฺเธน ทสหายตเนหิ}

\proseitem{430}{อิทฺธิปาเทน เย ธมฺมา สมฺปยุตฺตา, เต ธมฺมา ตีหิ ขนฺเธหิ เอเกนายตเนน เอกาย ธาตุยา สงฺคหิตา. กติหิ อสงฺคหิตา, ทฺวีหิ ขนฺเธหิ เอกาทสหา\-ยตเนหิ สตฺตรสหิ ธาตูหิ อสงฺคหิตา.}

\proseitem{431}{ฌาเนน เย ธมฺมา สมฺปยุตฺตา, เต ธมฺมา ตีหิ ขนฺเธหิ ทฺวีหายตเนหิ ทฺวีหิ ธาตูหิ สงฺคหิตา. กติหิ อสงฺคหิตา, ทฺวีหิ ขนฺเธหิ ทสหายตเนหิ โสฬสหิ ธาตูหิ อสงฺคหิตา.}

\proseitem{432}{อปฺปมญฺญาย เย ธมฺมา. ปญฺจหิ อินฺทฺริเยหิ เย ธมฺมา. ปญฺจหิ พเลหิ เย ธมฺมา. สตฺตหิ โพชฺฌงฺเคหิ เย ธมฺมา. อริเยน อฏฺฐงฺคิเกน มคฺเคน เย ธมฺมา สมฺปยุตฺตา, เต ธมฺมา จตูหิ ขนฺเธหิ ทฺวีหายตเนหิ ทฺวีหิ ธาตูหิ สงฺคหิตา. กติหิ อสงฺคหิตา, เอเกน ขนฺเธน ทสหายตเนหิ โสฬสหิ ธาตูหิ อสงฺคหิตา.}

\proseitem{433}{ผสฺเสน เย ธมฺมา. เจตนาย เย ธมฺมา. มนสิกาเรน เย ธมฺมา สมฺปยุตฺตา, เต ธมฺมา จตูหิ ขนฺเธหิ ทฺวีหายตเนหิ อฏฺฐหิ ธาตูหิ สงฺคหิตา. กติหิ อสงฺคหิตา, เอเกน ขนฺเธน ทสหายตเนหิ ทสหิ ธาตูหิ อสงฺคหิตา.}

\proseitem{434}{เวทนาย เย ธมฺมา. สญฺญาย เย ธมฺมา สมฺปยุตฺตา, เต ธมฺมา ตีหิ ขนฺเธหิ ทฺวีหายตเนหิ อฏฺฐหิ ธาตูหิ สงฺคหิตา. กติหิ อสงฺคหิตา, ทฺวีหิ ขนฺเธหิ ทสหายตเนหิ ทสหิ ธาตูหิ อสงฺคหิตา.}

\proseitem{435}{จิตฺเตน เย ธมฺมา สมฺปยุตฺตา, เต ธมฺมา ตีหิ ขนฺเธหิ เอเกนายตเนน เอกาย ธาตุยา สงฺคหิตา. กติหิ อสงฺคหิตา, ทฺวีหิ ขนฺเธหิ เอกาทสหา\-ยตเนหิ สตฺตรสหิ ธาตูหิ อสงฺคหิตา.}

\proseitem{436}{อธิโมกฺเขน เย ธมฺมา สมฺปยุตฺตา, เต ธมฺมา จตูหิ ขนฺเธหิ ทฺวีหายตเนหิ ตีหิ ธาตูหิ สงฺคหิตา. กติหิ อสงฺคหิตา, เอเกน ขนฺเธน ทสหายตเนหิ ปนฺนรสหิ ธาตูหิ อสงฺคหิตา.}

\proseitem{437}{\csromanfolio{81}สุขาย เวทนาย สมฺปยุตฺเตหิ ธมฺเมหิ เย ธมฺมา. ทุกฺขาย เวทนาย สมฺปยุตฺเตหิ ธมฺเมหิ เย ธมฺมา. อทุกฺขม\-สุขาย เวทนาย สมฺปยุตฺเตหิ ธมฺเมหิ เย ธมฺมา. สวิตกฺก\-สวิจาเรหิ ธมฺเมหิ เย ธมฺมา. อวิตกฺก\-วิจารมตฺเตหิ ธมฺเมหิ เย ธมฺมา. ปีติสหคเตหิ ธมฺเมหิ เย ธมฺมา. สุขสหคเตหิ ธมฺเมหิ เย ธมฺมา. อุเปกฺขา\-สหคเตหิ ธมฺเมหิ เย ธมฺมา สมฺปยุตฺตา, เต ธมฺมา เอเกน ขนฺเธน เอเกนายตเนน เอกาย ธาตุยา สงฺคหิตา. กติหิ อสงฺคหิตา, จตูหิ ขนฺเธหิ เอกาทสหา\-ยตเนหิ สตฺตรสหิ ธาตูหิ}

\proseitem{438}{เหตูหิ ธมฺเมหิ เย ธมฺมา. เหตูหิ เจว สเหตุเกหิ จ ธมฺเมหิ เย ธมฺมา. เหตูหิ เจว เหตุ\-สมฺปยุตฺเตหิ จ ธมฺเมหิ เย ธมฺมา สมฺปยุตฺตา, เต ธมฺมา จตูหิ ขนฺเธหิ ทฺวีหายตเนหิ ทฺวีหิ ธาตูหิ สงฺคหิตา. กติหิ อสงฺคหิตา, เอเกน ขนฺเธน ทสหายตเนหิ}

\proseitem{439}{สเหตุเกหิ เจว น จ เหตูหิ ธมฺเมหิ เย ธมฺมา. เหตุ\-สมฺปยุตฺเตหิ เจว น จ เหตูหิ ธมฺเมหิ เย ธมฺมา. น เหตุสเหตุเกหิ ธมฺเมหิ เย ธมฺมา สมฺปยุตฺตา, เต ธมฺมา เอเกน ขนฺเธน เอเกนายตเนน เอกาย ธาตุยา สงฺคหิตา. กติหิ อสงฺคหิตา, จตูหิ ขนฺเธหิ เอกาทสหา\-ยตเนหิ สตฺตรสหิ ธาตูหิ อสงฺคหิตา.}

\proseitem{440}{อาสเวหิ ธมฺเมหิ เย ธมฺมา. อาสเวหิ เจว สาสเวหิ จ ธมฺเมหิ เย ธมฺมา. อาสเวหิ เจว อาสว\-สมฺปยุตฺเตหิ จ ธมฺเมหิ เย ธมฺมา สมฺปยุตฺตา, เต ธมฺมา จตูหิ ขนฺเธหิ ทฺวีหายตเนหิ ทฺวีหิ ธาตูหิ สงฺคหิตา. กติหิ อสงฺคหิตา, เอเกน ขนฺเธน ทสหายตเนหิ}

\proseitem{441}{อาสว\-สมฺปยุตฺเตหิ เจว โน จ อาสเวหิ ธมฺเมหิ เย ธมฺมา สมฺปยุตฺตา, เต ธมฺมา เอเกน ขนฺเธน เอเกนายตเนน เอกาย ธาตุยา สงฺคหิตา. กติหิ อสงฺคหิตา, จตูหิ ขนฺเธหิ เอกาทสหา\-ยตเนหิ สตฺตรสหิ ธาตูหิ อสงฺคหิตา.}

\proseitem{442}{สญฺโญชเนหิ ธมฺเมหิ เย ธมฺมา. คนฺเถหิ ธมฺเมหิ เย ธมฺมา. โอเฆหิ ธมฺเมหิ เย ธมฺมา. โยเคหิ ธมฺเมหิ เย ธมฺมา. นีวรเณหิ \csromanfolio{82}ธมฺเมหิ เย ธมฺมา. ปรามาเสหิ ธมฺเมหิ เย ธมฺมา. ปรามาเสหิ เจว ปรามฏฺเฐหิ จ ธมฺเมหิ เย ธมฺมา สมฺปยุตฺตา, เต ธมฺมา จตูหิ ขนฺเธหิ ทฺวีหายตเนหิ ทฺวีหิ ธาตูหิ สงฺคหิตา. กติหิ อสงฺคหิตา, เอเกน ขนฺเธน ทสหายตเนหิ โสฬสหิ ธาตูหิ อสงฺคหิตา.}

\proseitem{443}{ปรามาส\-สมฺปยุตฺเตหิ ธมฺเมหิ เย ธมฺมา สมฺปยุตฺตา, เต ธมฺมา เอเกน ขนฺเธน เอเกนายตเนน เอกาย ธาตุยา สงฺคหิตา. กติหิ อสงฺคหิตา, จตูหิ ขนฺเธหิ เอกาทสหา\-ยตเนหิ สตฺตรสหิ ธาตูหิ}

\proseitem{444}{จิตฺเตหิ ธมฺเมหิ เย ธมฺมา สมฺปยุตฺตา, เต ธมฺมา ตีหิ ขนฺเธหิ เอเกนายตเนน เอกาย ธาตุยา สงฺคหิตา. กติหิ อสงฺคหิตา, ทฺวีหิ ขนฺเธหิ เอกาทสหา\-ยตเนหิ สตฺตรสหิ ธาตูหิ อสงฺคหิตา.}

\proseitem{445}{เจตสิเกหิ ธมฺเมหิ เย ธมฺมา. จิตฺต\-สมฺปยุตฺเตหิ ธมฺเมหิ เย ธมฺมา. จิตฺต\-สํสฏฺเฐหิ ธมฺเมหิ เย ธมฺมา. จิตฺต\-สํสฏฺฐ\-สมุฏฺฐาเนหิ ธมฺเมหิ เย ธมฺมา. จิตฺต\-สํสฏฺฐ\-สมุฏฺฐาน\-สหภูหิ ธมฺเมหิ เย ธมฺมา. จิตฺต\-สํสฏฺฐ\-สมุฏฺฐานา\-นุปริวตฺตีหิ ธมฺเมหิ เย ธมฺมา สมฺปยุตฺตา, เต ธมฺมา เอเกน ขนฺเธน เอเกนายตเนน สตฺตหิ ธาตูหิ สงฺคหิตา. กติหิ อสงฺคหิตา, จตูหิ ขนฺเธหิ เอกาทสหา\-ยตเนหิ เอกาทสหิ ธาตูหิ}

\proseitem{446}{อุปาทาเนหิ ธมฺเมหิ เย ธมฺมา. กิเลเสหิ ธมฺเมหิ เย ธมฺมา. กิเลเสหิ เจว สํกิเลสิเกหิ จ ธมฺเมหิ เย ธมฺมา. กิเลเสหิ เจว สํกิลิฏฺเฐหิ จ ธมฺเมหิ เย ธมฺมา. กิเลเสหิ เจว กิเลส\-สมฺปยุตฺเตหิ จ ธมฺเมหิ เย ธมฺมา สมฺปยุตฺตา, เต ธมฺมา จตูหิ ขนฺเธหิ ทฺวีหายตเนหิ ทฺวีหิ ธาตูหิ สงฺคหิตา. กติหิ อสงฺคหิตา, เอเกน ขนฺเธน ทสหายตเนหิ โสฬสหิ ธาตูหิ อสงฺคหิตา.}

\proseitem{447}{สํกิลิฏฺเฐหิ เจว โน จ กิเลเสหิ ธมฺเมหิ เย ธมฺมา. กิเลส\-สมฺปยุตฺเตหิ เจว โน จ กิเลเสหิ ธมฺเมหิ เย ธมฺมา. สวิตกฺเกหิ ธมฺเมหิ เย ธมฺมา. สวิจาเรหิ ธมฺเมหิ เย ธมฺมา. สปฺปีติเกหิ ธมฺเมหิ เย ธมฺมา. ปีติสหคเตหิ ธมฺเมหิ เย ธมฺมา. สุขสหคเตหิ \csromanfolio{83}ธมฺเมหิ เย ธมฺมา. อุเปกฺขา\-สหคเตหิ ธมฺเมหิ เย ธมฺมา สมฺปยุตฺตา, เต ธมฺมา กติหิ ขนฺเธหิ กติหายตเนหิ กติหิ ธาตูหิ สงฺคหิตา, เต ธมฺมา เอเกน ขนฺเธน เอเกนายตเนน เอกาย ธาตุยา สงฺคหิตา. กติหิ อสงฺคหิตา, จตูหิ ขนฺเธหิ เอกาทสหา\-ยตเนหิ สตฺตรสหิ ธาตูหิ อสงฺคหิตา.}

\setlength{\csromangathapagewidth}{0pt}
\setlength{\csromangathawakwidth}{0pt}
\csromangathameasuregroup{อรูปกฺขนฺธา จตฺตาโร, \\ วิญฺญาณธาตุโย สตฺต, \\ ปจฺจเย ทฺวาทส ปทา, \\ ติเกสุ อฏฺฐ โคจฺฉเก, \\ มหนฺตรทุเก สตฺต, \\ นวมสฺส ปทสฺเสเต,}{\csromangathabat{อรูปกฺขนฺธา จตฺตาโร,}{มนายตนเมว จ.} \\ \csromangathabat{วิญฺญาณธาตุโย สตฺต,}{เทฺว สจฺจา จุทฺทสินฺทฺริยา.} \\ \csromangathabat{ปจฺจเย ทฺวาทส ปทา,}{ตโต อุปริ โสฬส.} \\ \csromangathabat{ติเกสุ อฏฺฐ โคจฺฉเก,}{เตจตฺตาลีสเมว จ.} \\ \csromangathabat{มหนฺตรทุเก สตฺต,}{ปทา ปิฏฺฐิทุเกสุ ฉ.} \\ \csromangathabat{นวมสฺส ปทสฺเสเต,}{นิทฺเทเส สงฺคหํ คตาติ.}}
\csromangathasetleft{อรูปกฺขนฺธา จตฺตาโร, \\ วิญฺญาณธาตุโย สตฺต, \\ ปจฺจเย ทฺวาทส ปทา, \\ ติเกสุ อฏฺฐ โคจฺฉเก, \\ มหนฺตรทุเก สตฺต, \\ นวมสฺส ปทสฺเสเต,}
\csromangathagroup{\csromangathastanza{\csromangathabat{อรูปกฺขนฺธา จตฺตาโร,}{มนายตนเมว จ.} \\ \csromangathabat{วิญฺญาณธาตุโย สตฺต,}{เทฺว สจฺจา จุทฺทสินฺทฺริยา.}}\csromangathastanza{\csromangathabat{ปจฺจเย ทฺวาทส ปทา,}{ตโต อุปริ โสฬส.} \\ \csromangathabat{ติเกสุ อฏฺฐ โคจฺฉเก,}{เตจตฺตาลีสเมว จ.}}\csromangathastanza{\csromangathabat{มหนฺตรทุเก สตฺต,}{ปทา ปิฏฺฐิทุเกสุ ฉ.} \\ \csromangathabat{นวมสฺส ปทสฺเสเต,}{นิทฺเทเส สงฺคหํ คตาติ.}}}

\prose{สมฺปยุตฺเตน สงฺคหิตา\-สงฺคหิต\-ปท\-นิทฺเทโส ทฺวาทสโม.}

\csromanmatikaanchor{mtk.49}
\csromantocmarknopage{section}{13. เตรสมนย}{mtk.49}
\bookmark[dest={mtk.49},level=1]{13. เตรสมนย}
\csromanheader{1}{13. \textbf{เตรสมนย}}
\csromanmatikaanchor{mtk.50}
\csromantocmarkref{subsection}{13. อสงฺคหิเตนสมฺปยุตฺตวิปฺปยุตฺตปทนิทฺเทส}{mtk.50}
\bookmark[dest={mtk.50},level=2]{13. อสงฺคหิเตนสมฺปยุตฺตวิปฺปยุตฺตปทนิทฺเทส}
\csromanheader{2}{13. อสงฺคหิเตน\-สมฺปยุตฺต\-วิปฺปยุตฺต\-ปท\-นิทฺเทส}
\proseitem{448}{\csromanfolio{84}รูปกฺขนฺเธน เย ธมฺมา ขนฺธ\-สงฺคเหน อสงฺคหิตา อายตน\-สงฺคเหน อสงฺคหิตา ธาตุสงฺคเหน อสงฺคหิตา. เต ธมฺมา กติหิ ขนฺเธหิ กติหายตเนหิ กติหิ ธาตูหิ สมฺปยุตฺตา, เต ธมฺมา ตีหิ ขนฺเธหิ สมฺปยุตฺตา, เอเกนายตเนน เอกาย ธาตุยา เกหิจิ สมฺปยุตฺตา. กติหิ วิปฺปยุตฺตา, เอเกน ขนฺเธน ทสหายตเนหิ ทสหิ ธาตูหิ วิปฺปยุตฺตา, เอเกนายตเนน เอกาย ธาตุยา เกหิจิ วิปฺปยุตฺตา.}

\proseitem{449}{ธมฺมายตเนน เย ธมฺมา. ธมฺมธาตุยา เย ธมฺมา. อิตฺถินฺทฺริเยน เย ธมฺมา. ปุริสินฺทฺริเยน เย ธมฺมา. ชีวิตินฺทฺริเยน เย ธมฺมา. วิญฺญาณปจฺจยา นามรูเปน เย ธมฺมา. อสญฺญาภเวน เย ธมฺมา. เอกโวการ\-ภเวน เย ธมฺมา. ชาติยา เย ธมฺมา. ชราย เย ธมฺมา. มรเณน เย ธมฺมา ขนฺธ\-สงฺคเหน อสงฺคหิตา อายตน\-สงฺคเหน อสงฺคหิตา ธาตุสงฺคเหน อสงฺคหิตา ฯเปฯ เต ธมฺมา ตีหิ ขนฺเธหิ สมฺปยุตฺตา, เอเกนายตเนน เอกาย ธาตุยา เกหิจิ สมฺปยุตฺตา. กติหิ วิปฺปยุตฺตา, เอเกน ขนฺเธน ทสหายตเนหิ ทสหิ ธาตูหิ วิปฺปยุตฺตา, เอเกนายตเนน เอกาย ธาตุยา เกหิจิ วิปฺปยุตฺตา.}

\proseitem{450}{อรูปภเวน เย ธมฺมา. เนว\-สญฺญา\-นา\-สญฺญาภเวน เย ธมฺมา. จตุโวการภเวน เย ธมฺมา. อิทฺธิปาเทน เย ธมฺมา ขนฺธ\-สงฺคเหน อสงฺคหิตา อายตน\-สงฺคเหน อสงฺคหิตา ธาตุสงฺคเหน อสงฺคหิตา. เต ธมฺมา กติหิ ขนฺเธหิ กติหายตเนหิ กติหิ ธาตูหิ สมฺปยุตฺตาติ นตฺถิ. กติหิ วิปฺปยุตฺตา, จตูหิ ขนฺเธหิ เอเกนายตเนน สตฺตหิ ธาตูหิ วิปฺปยุตฺตา, เอเกนายตเนน เอกาย ธาตุยา เกหิจิ วิปฺปยุตฺตา.}

\proseitem{451}{กุสเลหิ ธมฺเมหิ เย ธมฺมา. อกุสเลหิ ธมฺเมหิ เย ธมฺมา. สุขาย เวทนาย สมฺปยุตฺเตหิ ธมฺเมหิ เย ธมฺมา. ทุกฺขาย เวทนาย สมฺปยุตฺเตหิ ธมฺเมหิ เย ธมฺมา. อทุกฺขม\-สุขาย เวทนาย สมฺปยุตฺเตหิ ธมฺเมหิ เย ธมฺมา. วิปาเกหิ ธมฺเมหิ เย ธมฺมา. วิปาก\-ธมฺม\-ธมฺเมหิ เย ธมฺมา.}

\prose{\csromanfolio{85}อนุปาทินฺน\-อนุปาทานิเยหิ ธมฺเมหิ เย ธมฺมา. สํกิลิฏฺฐ\-สํกิเลสิเกหิ ธมฺเมหิ เย ธมฺมา. อสํกิลิฏฺฐ- อสํกิเลสิเกหิ ธมฺเมหิ เย ธมฺมา. สวิตกฺก\-สวิจาเรหิ ธมฺเมหิ เย ธมฺมา. อวิตกฺก\-วิจารมตฺเตหิ ธมฺเมหิ เย ธมฺมา. ปีติสหคเตหิ ธมฺเมหิ เย ธมฺมา. สุขสหคเตหิ ธมฺเมหิ เย ธมฺมา. อุเปกฺขา\-สหคเตหิ ธมฺเมหิ เย ธมฺมา. ทสฺสเนน ปหาตพฺเพหิ ธมฺเมหิ เย ธมฺมา. ภาวนาย ปหาตพฺเพหิ ธมฺเมหิ เย ธมฺมา. ทสฺสเนน ปหาตพฺพ\-เหตุเกหิ ธมฺเมหิ เย ธมฺมา. ภาวนาย ปหาตพฺพ\-เหตุเกหิ ธมฺเมหิ เย ธมฺมา. อาจยคามีหิ ธมฺเมหิ เย ธมฺมา. อปจยคามีหิ ธมฺเมหิ เย ธมฺมา. เสกฺเขหิ ธมฺเมหิ เย ธมฺมา. อเสกฺเขหิ ธมฺเมหิ เย ธมฺมา. มหคฺคเตหิ ธมฺเมหิ เย ธมฺมา. อปฺปมาเณหิ ธมฺเมหิ เย ธมฺมา. ปริตฺตา\-รมฺมเณหิ ธมฺเมหิ เย ธมฺมา. มหคฺคตา\-รมฺมเณหิ ธมฺเมหิ เย ธมฺมา. อปฺปมาณา\-รมฺมเณหิ ธมฺเมหิ เย ธมฺมา. หีเนหิ ธมฺเมหิ เย ธมฺมา. ปณีเตหิ ธมฺเมหิ เย ธมฺมา. มิจฺฉตฺต\-นิยเตหิ ธมฺเมหิ เย ธมฺมา. สมฺมตฺต\-นิยเตหิ ธมฺเมหิ เย ธมฺมา. มคฺคารมฺมเณหิ ธมฺเมหิ เย ธมฺมา. มคฺคเหตุเกหิ ธมฺเมหิ เย ธมฺมา. มคฺคาธิปตีหิ ธมฺเมหิ เย ธมฺมา. อตีตารมฺมเณหิ ธมฺเมหิ เย ธมฺมา. อนาคตา\-รมฺมเณหิ ธมฺเมหิ เย ธมฺมา. ปจฺจุปฺปนฺนา\-รมฺมเณหิ ธมฺเมหิ เย ธมฺมา. อชฺฌตฺตา\-รมฺมเณหิ ธมฺเมหิ เย ธมฺมา. พหิทฺธา\-รมฺมเณหิ ธมฺเมหิ เย ธมฺมา. อชฺฌตฺตพหิทฺธา\-รมฺมเณหิ ธมฺเมหิ เย ธมฺมา.}

\prose{สเหตุเกหิ ธมฺเมหิ เย ธมฺมา. เหตุ\-สมฺปยุตฺเตหิ ธมฺเมหิ เย ธมฺมา. สเหตุเกหิ เจว น จ เหตูหิ ธมฺเมหิ เย ธมฺมา. เหตุ\-สมฺปยุตฺเตหิ เจว น จ เหตูหิ ธมฺเมหิ เย ธมฺมา. น เหตุสเหตุเกหิ ธมฺเมหิ เย ธมฺมา ขนฺธ\-สงฺคเหน อสงฺคหิตา อายตน\-สงฺคเหน อสงฺคหิตา ธาตุสงฺคเหน อสงฺคหิตา. เต ธมฺมา กติหิ ขนฺเธหิ กติหายตเนหิ กติหิ ธาตูหิ สมฺปยุตฺตาติ นตฺถิ. กติหิ วิปฺปยุตฺตา, จตูหิ ขนฺเธหิ เอเกนายตเนน สตฺตหิ ธาตูหิ วิปฺปยุตฺตา, เอเกนายตเนน เอกาย ธาตุยา เกหิจิ วิปฺปยุตฺตา.}

\proseitem{452}{รูปีหิ ธมฺเมหิ เย ธมฺมา ขนฺธ\-สงฺคเหน อสงฺคหิตา อายตน\-สงฺคเหน อสงฺคหิตา ธาตุสงฺคเหน อสงฺคหิตา. เต ธมฺมา ตีหิ ขนฺเธหิ สมฺปยุตฺตา, เอเกนายตเนน เอกาย ธาตุยา เกหิจิ สมฺปยุตฺตา. กติหิ วิปฺปยุตฺตา, เอเกน ขนฺเธน ทสหายตเนหิ ทสหิ \csromanfolio{86}ธาตูหิ วิปฺปยุตฺตา, เอเกนายตเนน เอกาย ธาตุยา เกหิจิ วิปฺปยุตฺตา.}

\proseitem{453}{อรูปีหิ ธมฺเมหิ เย ธมฺมา. โลกุตฺตเรหิ ธมฺเมหิ เย ธมฺมา. อนาสเวหิ ธมฺเมหิ เย ธมฺมา. อาสว\-สมฺปยุตฺเตหิ ธมฺเมหิ เย ธมฺมา. อาสว\-สมฺปยุตฺเตหิ เจว โน จ อาสเวหิ ธมฺเมหิ เย ธมฺมา. อาสว\-วิปฺปยุตฺเตหิ อนาสเวหิ ธมฺเมหิ เย ธมฺมา. อสํโยชนิเยหิ ธมฺเมหิ เย ธมฺมา. อคนฺถนิเยหิ ธมฺเมหิ เย ธมฺมา. อโนฆนิเยหิ ธมฺเมหิ เย ธมฺมา. อโยคนิเยหิ ธมฺเมหิ เย ธมฺมา. อนีวรณิเยหิ ธมฺเมหิ เย ธมฺมา. อปรามฏฺเฐหิ ธมฺเมหิ เย ธมฺมา. ปรามาส\-สมฺปยุตฺเตหิ ธมฺเมหิ เย ธมฺมา. ปรามาส\-วิปฺปยุตฺเตหิ อปรามฏฺเฐหิ ธมฺเมหิ เย ธมฺมา. สารมฺมเณหิ ธมฺเมหิ เย ธมฺมา ขนฺธ\-สงฺคเหน อสงฺคหิตา อายตน\-สงฺคเหน อสงฺคหิตา ธาตุสงฺคเหน อสงฺคหิตา. เต ธมฺมา กติหิ ขนฺเธหิ กติหายตเนหิ กติหิ ธาตูหิ สมฺปยุตฺตาติ นตฺถิ. กติหิ วิปฺปยุตฺตา, จตูหิ ขนฺเธหิ เอเกนายตเนน สตฺตหิ ธาตูหิ วิปฺปยุตฺตา, เอเกนายตเนน เอกาย ธาตุยา เกหิจิ วิปฺปยุตฺตา.}

\proseitem{454}{อนารมฺมเณหิ ธมฺเมหิ เย ธมฺมา. โน จิตฺเตหิ ธมฺเมหิ เย ธมฺมา. จิตฺต\-วิปฺปยุตฺเตหิ ธมฺเมหิ เย ธมฺมา. จิตฺต\-วิสํสฏฺเฐหิ ธมฺเมหิ เย ธมฺมา. จิตฺต\-สมุฏฺฐาเนหิ ธมฺเมหิ เย ธมฺมา. จิตฺตสหภูหิ ธมฺเมหิ เย ธมฺมา. จิตฺตา\-นุปริวตฺตีหิ ธมฺเมหิ เย ธมฺมา. พาหิเรหิ ธมฺเมหิ เย ธมฺมา. อุปาทาธมฺเมหิ เย ธมฺมา ขนฺธ\-สงฺคเหน อสงฺคหิตา อายตน\-สงฺคเหน อสงฺคหิตา ธาตุสงฺคเหน อสงฺคหิตา. เต ธมฺมา กติหิ ขนฺเธหิ กติหายตเนหิ กติหิ ธาตูหิ สมฺปยุตฺตา, เต ธมฺมา ตีหิ ขนฺเธหิ สมฺปยุตฺตา, เอเกนายตเนน เอกาย ธาตุยา เกหิจิ สมฺปยุตฺตา. กติหิ วิปฺปยุตฺตา, เอเกน ขนฺเธน ทสหายตเนหิ ทสหิ ธาตูหิ วิปฺปยุตฺตา, เอเกนายตเนน เอกาย ธาตุยา เกหิจิ วิปฺปยุตฺตา.}

\proseitem{455}{อนุปาทานิเยหิ ธมฺเมหิ เย ธมฺมา. อุปาทาน\-สมฺปยุตฺเตหิ ธมฺเมหิ เย ธมฺมา. อุปาทาน\-สมฺปยุตฺเตหิ เจว โน จ อุปาทาเนหิ ธมฺเมหิ เย ธมฺมา. อุปาทาน\-วิปฺปยุตฺเตหิ อนุปาทานิเยหิ ธมฺเมหิ เย ธมฺมา. อสํกิเลสิเกหิ ธมฺเมหิ เย ธมฺมา. อสํกิลิฏฺเฐหิ ธมฺเมหิ เย ธมฺมา. กิเลส\-สมฺปยุตฺเตหิ ธมฺเมหิ เย ธมฺมา. สํกิลิฏฺเฐหิ เจว โน จ \csromanfolio{87}กิเลเสหิ ธมฺเมหิ เย ธมฺมา. กิเลส\-สมฺปยุตฺเตหิ เจว โน จ กิเลเสหิ ธมฺเมหิ เย ธมฺมา. กิเลส\-วิปฺปยุตฺเตหิ อสํกิเลสิเกหิ ธมฺเมหิ เย ธมฺมา. ทสฺสเนน ปหาตพฺเพหิ ธมฺเมหิ เย ธมฺมา. ภาวนาย ปหาตพฺเพหิ ธมฺเมหิ เย ธมฺมา. ทสฺสเนน ปหาตพฺพ\-เหตุเกหิ ธมฺเมหิ เย ธมฺมา. ภาวนาย ปหาตพฺพ\-เหตุเกหิ ธมฺเมหิ เย ธมฺมา. สวิตกฺเกหิ ธมฺเมหิ เย ธมฺมา. สวิจาเรหิ ธมฺเมหิ เย ธมฺมา. สปฺปีติเกหิ ธมฺเมหิ เย ธมฺมา. ปีติสหคเตหิ ธมฺเมหิ เย ธมฺมา. สุขสหคเตหิ ธมฺเมหิ เย ธมฺมา. อุเปกฺขา\-สหคเตหิ ธมฺเมหิ เย ธมฺมา. น กามาวจเรหิ ธมฺเมหิ เย ธมฺมา. รูปาวจเรหิ ธมฺเมหิ เย ธมฺมา. อรูปาวจเรหิ ธมฺเมหิ เย ธมฺมา. อปริยาปนฺเนหิ ธมฺเมหิ เย ธมฺมา. นิยฺยานิเกหิ ธมฺเมหิ เย ธมฺมา. นิยเตหิ ธมฺเมหิ เย ธมฺมา. อนุตฺตเรหิ ธมฺเมหิ เย ธมฺมา. สรเณหิ ธมฺเมหิ เย ธมฺมา ขนฺธ\-สงฺคเหน อสงฺคหิตา, อายตน\-สงฺคเหน อสงฺคหิตา, ธาตุสงฺคเหน อสงฺคหิตา, เต ธมฺมา กติหิ ขนฺเธหิ กติหายตเนหิ กติหิ ธาตูหิ สมฺปยุตฺตาติ นตฺถิ. กติหิ วิปฺปยุตฺตา, จตูหิ ขนฺเธหิ เอเกนายตเนน สตฺตหิ ธาตูหิ วิปฺปยุตฺตา, เอเกนายตเนน เอกาย ธาตุยา เกหิจิ วิปฺปยุตฺตา.}

\setlength{\csromangathapagewidth}{0pt}
\setlength{\csromangathawakwidth}{0pt}
\csromangathameasuregroup{วิสํสฏฺฐํ สมุฏฺฐานสหภุ,}{รูปญฺจ ธมฺมายตนํ ธมฺมธาตุ, \\ อิตฺถิปุมํ ชีวิตํ นามรูปํ. \\ เทฺว ภวา ชาติชรา มจฺจุรูปํ, \\ อนารมฺมณํ โน จิตฺตํ จิตฺเตน วิปฺปยุตฺตํ. \\ \csromangathabat{วิสํสฏฺฐํ สมุฏฺฐานสหภุ,}{อนุปริวตฺติ พาหิรํ อุปาทา.} \\ เทฺว วิสโย เอสนโย สุพุทฺโธติ. \\ อสงฺคหิเตน สมฺปยุตฺตวิปฺปยุตฺตปทนิทฺเทโส เตรสโม.}
\csromangathameasurewak{รูปญฺจ ธมฺมายตนํ ธมฺมธาตุ, \\ อิตฺถิปุมํ ชีวิตํ นามรูปํ. \\ เทฺว ภวา ชาติชรา มจฺจุรูปํ, \\ อนารมฺมณํ โน จิตฺตํ จิตฺเตน วิปฺปยุตฺตํ.}
\csromangathasetleft{วิสํสฏฺฐํ สมุฏฺฐานสหภุ,}
\csromangathagroup{\csromangathastanza{\csromangathawak{รูปญฺจ ธมฺมายตนํ ธมฺมธาตุ, \\ อิตฺถิปุมํ ชีวิตํ นามรูปํ. \\ เทฺว ภวา ชาติชรา มจฺจุรูปํ, \\ อนารมฺมณํ โน จิตฺตํ จิตฺเตน วิปฺปยุตฺตํ.}}\csromangathastanza{\csromangathabat{วิสํสฏฺฐํ สมุฏฺฐานสหภุ,}{อนุปริวตฺติ พาหิรํ อุปาทา.} \\ เทฺว วิสโย เอสนโย สุพุทฺโธติ. \\ อสงฺคหิเตน สมฺปยุตฺต\-วิปฺปยุตฺต\-ปท\-นิทฺเทโส เตรสโม.}}

\csromanmatikaanchor{mtk.51}
\csromantocmarknopage{section}{14. จุทฺทสมนย}{mtk.51}
\bookmark[dest={mtk.51},level=1]{14. จุทฺทสมนย}
\csromanheader{1}{14. \textbf{จุทฺทสมนย}}
\csromanmatikaanchor{mtk.52}
\csromantocmarkref{subsection}{14. วิปฺปยุตฺเตนสงฺคหิตาสงฺคหิตปทนิทฺเทส}{mtk.52}
\bookmark[dest={mtk.52},level=2]{14. วิปฺปยุตฺเตนสงฺคหิตาสงฺคหิตปทนิทฺเทส}
\csromanheader{2}{14. วิปฺปยุตฺเตน\-สงฺคหิตา\-สงฺคหิต\-ปท\-นิทฺเทส}
\csromanmatikaanchor{mtk.53}
\csromantocmarkref{subsection}{1. ขนฺธาทิ}{mtk.53}
\bookmark[dest={mtk.53},level=2]{1. ขนฺธาทิ}
\csromanheader{2}{1. \textbf{ขนฺธาทิ}}
\proseitem{456}{\csromanfolio{88}รูปกฺขนฺเธน เย ธมฺมา วิปฺปยุตฺตา. เต ธมฺมา กติหิ ขนฺเธหิ กติหายตเนหิ กติหิ ธาตูหิ สงฺคหิตา, เต ธมฺมา จตูหิ ขนฺเธหิ ทฺวีหายตเนหิ อฏฺฐหิ ธาตูหิ สงฺคหิตา. กติหิ อสงฺคหิตา, เอเกน ขนฺเธน ทสหายตเนหิ ทสหิ ธาตูหิ อสงฺคหิตา.}

\proseitem{457}{เวทนา\-กฺขนฺเธน เย ธมฺมา. สญฺญา\-กฺขนฺเธน เย ธมฺมา. สงฺขาร\-กฺขนฺเธน เย ธมฺมา. วิญฺญาณ\-กฺขนฺเธน เย ธมฺมา. มนายตเนน เย ธมฺมา. มนินฺทฺริเยน เย ธมฺมา วิปฺปยุตฺตา. เต ธมฺมา กติหิ ขนฺเธหิ กติหายตเนหิ กติหิ ธาตูหิ สงฺคหิตา, เต ธมฺมา อสงฺขตํ ขนฺธโต ฐเปตฺวา เอเกน ขนฺเธน เอกาทสหา\-ยตเนหิ เอกาทสหิ ธาตูหิ สงฺคหิตา. กติหิ อสงฺคหิตา, จตูหิ ขนฺเธหิ เอเกนายตเนน สตฺตหิ ธาตูหิ อสงฺคหิตา.}

\proseitem{458}{จกฺขายตเนน เย ธมฺมา ฯเปฯ. โผฏฺฐพฺพา\-ยตเนน เย ธมฺมา. จกฺขุธาตุยา เย ธมฺมา ฯเปฯ. โผฏฺฐพฺพ\-ธาตุยา เย ธมฺมา วิปฺปยุตฺตา ฯเปฯ เต ธมฺมา จตูหิ ขนฺเธหิ ทฺวีหายตเนหิ อฏฺฐหิ ธาตูหิ สงฺคหิตา. กติหิ อสงฺคหิตา, เอเกน ขนฺเธน ทสหายตเนหิ ทสหิ ธาตูหิ อสงฺคหิตา.}

\proseitem{459}{จกฺขุวิญฺญาณ\-ธาตุยา เย ธมฺมา. โสตวิญฺญาณ\-ธาตุยา เย ธมฺมา. ฆานวิญฺญาณ\-ธาตุยา เย ธมฺมา. ชิวฺหาวิญฺญาณ\-ธาตุยา เย ธมฺมา. กายวิญฺญาณ\-ธาตุยา เย ธมฺมา. มโนธาตุยา เย ธมฺมา. มโนวิญฺญาณ\-ธาตุยา เย ธมฺมา วิปฺปยุตฺตา. เต ธมฺมา อสงฺขตํ ขนฺธโต ฐเปตฺวา ปญฺจหิ ขนฺเธหิ ทฺวาทสหา\-ยตเนหิ สตฺตรสหิ ธาตูหิ สงฺคหิตา. กติหิ อสงฺคหิตา, น เกหิจิ ขนฺเธหิ น เกหิจิ อายตเนหิ เอกาย ธาตุยา อสงฺคหิตา.}

\csromanmatikaanchor{mtk.54}
\csromantocmarkref{subsection}{2. สจฺจาทิ}{mtk.54}
\bookmark[dest={mtk.54},level=2]{2. สจฺจาทิ}
\csromanmarkh{14. จุทฺทสมนย}\csromanmarkhii{2. สจฺจาทิ}\csromanheadernomark{2}{14. จุทฺทสมนย \textbf{2. สจฺจาทิ}}
\proseitem{460}{\csromanfolio{89}ทุกฺขสจฺเจน เย ธมฺมา วิปฺปยุตฺตา. เต ธมฺมา จตูหิ ขนฺเธหิ ทฺวีหายตเนหิ ทฺวีหิ ธาตูหิ สงฺคหิตา. กติหิ อสงฺคหิตา, เอเกน ขนฺเธน ทสหายตเนหิ โสฬสหิ ธาตูหิ อสงฺคหิตา.}

\proseitem{461}{สมุทย\-สจฺเจน เย ธมฺมา. มคฺคสจฺเจน เย ธมฺมา วิปฺปยุตฺตา. เต ธมฺมา อสงฺขตํ ขนฺธโต ฐเปตฺวา ปญฺจหิ ขนฺเธหิ ทฺวาทสหา\-ยตเนหิ อฏฺฐารสหิ ธาตูหิ สงฺคหิตา. กติหิ อสงฺคหิตา, น เกหิจิ ขนฺเธหิ น เกหิจิ อายตเนหิ น กาหิจิ ธาตูหิ อสงฺคหิตา.}

\proseitem{462}{นิโรธสจฺเจน เย ธมฺมา. จกฺขุนฺทฺริเยน เย ธมฺมา ฯเปฯ. กายินฺทฺริเยน เย ธมฺมา. อิตฺถินฺทฺริเยน เย ธมฺมา. ปุริสินฺทฺริเยน เย ธมฺมา วิปฺปยุตฺตา. เต ธมฺมา จตูหิ ขนฺเธหิ ทฺวีหายตเนหิ อฏฺฐหิ ธาตูหิ สงฺคหิตา. กติหิ อสงฺคหิตา, เอเกน ขนฺเธน ทสหายตเนหิ ทสหิ ธาตูหิ อสงฺคหิตา.}

\proseitem{463}{สุขินฺทฺริเยน เย ธมฺมา. ทุกฺขินฺทฺริเยน เย ธมฺมา. โสมนสฺสิ\-นฺทฺริเยน เย ธมฺมา. โทมนสฺสิ\-นฺทฺริเยน เย ธมฺมา วิปฺปยุตฺตา. เต ธมฺมา อสงฺขตํ ขนฺธโต ฐเปตฺวา ปญฺจหิ ขนฺเธหิ ทฺวาทสหา\-ยตเนหิ อฏฺฐารสหิ ธาตูหิ สงฺคหิตา. กติหิ อสงฺคหิตา, น เกหิจิ ขนฺเธหิ น เกหิจิ อายตเนหิ น กาหิจิ ธาตูหิ อสงฺคหิตา.}

\proseitem{464}{อุเปกฺขิ\-นฺทฺริเยน เย ธมฺมา วิปฺปยุตฺตา. เต ธมฺมา อสงฺขตํ ขนฺธโต ฐเปตฺวา ปญฺจหิ ขนฺเธหิ ทฺวาทสหา\-ยตเนหิ เตรสหิ ธาตูหิ สงฺคหิตา. กติหิ อสงฺคหิตา, น เกหิจิ ขนฺเธหิ น เกหิจิ อายตเนหิ ปญฺจหิ ธาตูหิ อสงฺคหิตา.}

\proseitem{465}{สทฺธินฺทฺริเยน เย ธมฺมา. วีริยินฺทฺริเยน เย ธมฺมา. สตินฺทฺริเยน เย ธมฺมา. สมาธิ\-นฺทฺริเยน เย ธมฺมา. ปญฺญินฺทฺริเยน เย ธมฺมา. อนญฺญาต\-ญฺญสฺสามี\-ตินฺทฺริเยน เย ธมฺมา. อญฺญินฺทฺริเยน เย ธมฺมา. อญฺญาตาวิ\-นฺทฺริเยน เย ธมฺมา. อวิชฺชาย เย ธมฺมา. อวิชฺชาปจฺจยา สงฺขาเรหิ เย ธมฺมา วิปฺปยุตฺตา. เต ธมฺมา อสงฺขตํ ขนฺธโต ฐเปตฺวา ปญฺจหิ ขนฺเธหิ ทฺวาทสหา\-ยตเนหิ \csromanfolio{90}อฏฺฐารสหิ ธาตูหิ สงฺคหิตา. กติหิ อสงฺคหิตา, น เกหิจิ ขนฺเธหิ น เกหิจิ อายตเนหิ น กาหิจิ ธาตูหิ อสงฺคหิตา.}

\proseitem{466}{สงฺขาร\-ปจฺจยา วิญฺญาเณน เย ธมฺมา. สฬายตน\-ปจฺจยา ผสฺเสน เย ธมฺมา. ผสฺสปจฺจยา เวทนาย เย ธมฺมา วิปฺปยุตฺตา. เต ธมฺมา อสงฺขตํ ขนฺธโต ฐเปตฺวา เอเกน ขนฺเธน เอกาทสหา\-ยตเนหิ เอกาทสหิ ธาตูหิ สงฺคหิตา. กติหิ อสงฺคหิตา, จตูหิ ขนฺเธหิ เอเกนายตเนน สตฺตหิ ธาตูหิ อสงฺคหิตา.}

\proseitem{467}{เวทนาปจฺจยา ตณฺหาย เย ธมฺมา. ตณฺหาปจฺจยา อุปาทาเนน เย ธมฺมา. กมฺมภเวน เย ธมฺมา วิปฺปยุตฺตา. เต ธมฺมา อสงฺขตํ ขนฺธโต ฐเปตฺวา ปญฺจหิ ขนฺเธหิ ทฺวาทสหา\-ยตเนหิ อฏฺฐารสหิ ธาตูหิ สงฺคหิตา. กติหิ อสงฺคหิตา, น เกหิจิ ขนฺเธหิ น เกหิจิ อายตเนหิ น กาหิจิ ธาตูหิ อสงฺคหิตา.}

\proseitem{468}{อุปปตฺติภเวน เย ธมฺมา. สญฺญาภเวน เย ธมฺมา. ปญฺจโวการภเวน เย ธมฺมา วิปฺปยุตฺตา. เต ธมฺมา จตูหิ ขนฺเธหิ ทฺวีหายตเนหิ ตีหิ ธาตูหิ สงฺคหิตา. กติหิ อสงฺคหิตา, เอเกน ขนฺเธน ทสหายตเนหิ ปนฺนรสหิ ธาตูหิ อสงฺคหิตา.}

\proseitem{469}{กามภเวน เย ธมฺมา วิปฺปยุตฺตา. เต ธมฺมา จตูหิ ขนฺเธหิ ทฺวีหายตเนหิ ปญฺจหิ ธาตูหิ สงฺคหิตา. กติหิ อสงฺคหิตา. เอเกน ขนฺเธน ทสหายตเนหิ เตรสหิ ธาตูหิ อสงฺคหิตา.}

\proseitem{470}{รูปภเวน เย ธมฺมา. อสญฺญาภเวน เย ธมฺมา. เอกโวการ\-ภเวน เย ธมฺมา. ปริเทเวน เย ธมฺมา วิปฺปยุตฺตา. เต ธมฺมา จตูหิ ขนฺเธหิ ทฺวีหายตเนหิ อฏฺฐหิ ธาตูหิ สงฺคหิตา. กติหิ อสงฺคหิตา, เอเกน ขนฺเธน ทสหายตเนหิ ทสหิ ธาตูหิ อสงฺคหิตา.}

\proseitem{471}{อรูปภเวน เย ธมฺมา. เนว\-สญฺญา\-นา\-สญฺญาภเวน เย ธมฺมา. จตุโวการภเวน เย ธมฺมา. โสเกน เย ธมฺมา. ทุกฺเขน เย ธมฺมา. โทมนสฺเสน เย ธมฺมา. อุปายาเสน เย ธมฺมา. สติปฏฺฐาเนน เย ธมฺมา. สมฺม\-ปฺปธาเนน เย ธมฺมา. อิทฺธิปาเทน เย ธมฺมา. ฌาเนน เย ธมฺมา. อปฺปมญฺญาย เย ธมฺมา. ปญฺจหิ อินฺทฺริเยหิ เย ธมฺมา. ปญฺจหิ พเลหิ เย \csromanfolio{91}ธมฺมา. สตฺตหิ โพชฺฌงฺเคหิ เย ธมฺมา. อริเยน อฏฺฐงฺคิเกน มคฺเคน เย ธมฺมา วิปฺปยุตฺตา. เต ธมฺมา อสงฺขตํ ขนฺธโต ฐเปตฺวา ปญฺจหิ ขนฺเธหิ ทฺวาทสหา\-ยตเนหิ อฏฺฐารสหิ ธาตูหิ สงฺคหิตา. กติหิ อสงฺคหิตา, น เกหิจิ ขนฺเธหิ น เกหิจิ อายตเนหิ น กาหิจิ ธาตูหิ}

\csromanmatikaanchor{mtk.55}
\csromantocmarkref{subsection}{3. ผสฺสาทิสตฺตก}{mtk.55}
\bookmark[dest={mtk.55},level=2]{3. ผสฺสาทิสตฺตก}
\csromanheader{2}{3. ผสฺสา\-ทิ\-สตฺตก}
\proseitem{472}{ผสฺเสน เย ธมฺมา. เวทนาย เย ธมฺมา. สญฺญาย เย ธมฺมา. เจตนาย เย ธมฺมา. จิตฺเตน เย ธมฺมา. มนสิกาเรน เย ธมฺมา วิปฺปยุตฺตา. เต ธมฺมา อสงฺขตํ ขนฺธโต ฐเปตฺวา เอเกน ขนฺเธน เอกาทสหา\-ยตเนหิ เอกาทสหิ ธาตูหิ สงฺคหิตา. กติหิ อสงฺคหิตา, จตูหิ ขนฺเธหิ เอเกนายตเนน สตฺตหิ ธาตูหิ อสงฺคหิตา.}

\proseitem{473}{อธิโมกฺเขน เย ธมฺมา วิปฺปยุตฺตา. เต ธมฺมา อสงฺขตํ ขนฺธโต ฐเปตฺวา ปญฺจหิ ขนฺเธหิ ทฺวาทสหา\-ยตเนหิ สตฺตรสหิ ธาตูหิ สงฺคหิตา. กติหิ อสงฺคหิตา, น เกหิจิ ขนฺเธหิ น เกหิจิ อายตเนหิ เอกาย ธาตุยา อสงฺคหิตา.}

\csromanmatikaanchor{mtk.56}
\csromantocmarkref{subsection}{4. ติก}{mtk.56}
\bookmark[dest={mtk.56},level=2]{4. ติก}
\csromanheader{2}{4. \textbf{ติก}}
\proseitem{474}{กุสเลหิ ธมฺเมหิ เย ธมฺมา. อกุสเลหิ ธมฺเมหิ เย ธมฺมา. สุขาย เวทนาย สมฺปยุตฺเตหิ ธมฺเมหิ เย ธมฺมา. ทุกฺขาย เวทนาย สมฺปยุตฺเตหิ ธมฺเมหิ เย ธมฺมา วิปฺปยุตฺตา. เต ธมฺมา อสงฺขตํ ขนฺธโต ฐเปตฺวา ปญฺจหิ ขนฺเธหิ ทฺวาทสหา\-ยตเนหิ อฏฺฐารสหิ ธาตูหิ สงฺคหิตา. กติหิ อสงฺคหิตา, น เกหิจิ ขนฺเธหิ น เกหิจิ อายตเนหิ น กาหิจิ ธาตูหิ อสงฺคหิตา.}

\proseitem{475}{อพฺยากเตหิ ธมฺเมหิ เย ธมฺมา วิปฺปยุตฺตา. เต ธมฺมา จตูหิ ขนฺเธหิ ทฺวีหายตเนหิ ทฺวีหิ ธาตูหิ สงฺคหิตา. กติหิ อสงฺคหิตา, เอเกน ขนฺเธน ทสหายตเนหิ โสฬสหิ ธาตูหิ อสงฺคหิตา.}

\proseitem{476}{\csromanfolio{92}อทุกฺขม\-สุขาย เวทนาย สมฺปยุตฺเตหิ ธมฺเมหิ เย ธมฺมา. วิปาเกหิ ธมฺเมหิ เย ธมฺมา วิปฺปยุตฺตา. เต ธมฺมา อสงฺขตํ ขนฺธโต ฐเปตฺวา ปญฺจหิ ขนฺเธหิ ทฺวาทสหา\-ยตเนหิ เตรสหิ ธาตูหิ สงฺคหิตา. กติหิ อสงฺคหิตา, น เกหิจิ ขนฺเธหิ น เกหิจิ อายตเนหิ ปญฺจหิ ธาตูหิ อสงฺคหิตา.}

\proseitem{477}{วิปาก\-ธมฺม\-ธมฺเมหิ เย ธมฺมา. สํกิลิฏฺฐ\-สํกิเลสิเกหิ ธมฺเมหิ เย ธมฺมา วิปฺปยุตฺตา. เต ธมฺมา อสงฺขตํ ขนฺธโต ฐเปตฺวา ปญฺจหิ ขนฺเธหิ ทฺวาทสหา\-ยตเนหิ อฏฺฐารสหิ ธาตูหิ สงฺคหิตา. กติหิ อสงฺคหิตา, น เกหิจิ ขนฺเธหิ น เกหิจิ อายตเนหิ น กาหิจิ ธาตูหิ อสงฺคหิตา.}

\proseitem{478}{เนว\-วิปาก\-น\-วิปากธมฺม\-ธมฺเมหิ เย ธมฺมา. อนุปาทินฺนุ\-ปาทานิเยหิ ธมฺเมหิ เย ธมฺมา. อนุปาทินฺน\-อนุปาทานิเยหิ ธมฺเมหิ เย ธมฺมา. อสํกิลิฏฺฐ\-อสํกิเลสิเกหิ ธมฺเมหิ เย ธมฺมา วิปฺปยุตฺตา. เต ธมฺมา จตูหิ ขนฺเธหิ ทฺวีหายตเนหิ อฏฺฐหิ ธาตูหิ สงฺคหิตา. กติหิ อสงฺคหิตา, เอเกน ขนฺเธน ทสหายตเนหิ ทสหิ ธาตูหิ อสงฺคหิตา.}

\proseitem{479}{อุปาทินฺนุ\-ปาทานิเยหิ ธมฺเมหิ เย ธมฺมา วิปฺปยุตฺตา. เต ธมฺมา จตูหิ ขนฺเธหิ ทฺวีหายตเนหิ ตีหิ ธาตูหิ สงฺคหิตา. กติหิ อสงฺคหิตา, เอเกน ขนฺเธน ทสหายตเนหิ ปนฺนรสหิ ธาตูหิ}

\proseitem{480}{อสํกิลิฏฺฐ\-สํกิเลสิเกหิ ธมฺเมหิ เย ธมฺมา วิปฺปยุตฺตา. เต ธมฺมา จตูหิ ขนฺเธหิ ทฺวีหายตเนหิ ทฺวีหิ ธาตูหิ สงฺคหิตา. กติหิ อสงฺคหิตา, เอเกน ขนฺเธน ทสหายตเนหิ โสฬสหิ ธาตูหิ อสงฺคหิตา.}

\proseitem{481}{สวิตกฺก\-สวิจาเรหิ ธมฺเมหิ เย ธมฺมา วิปฺปยุตฺตา. เต ธมฺมา อสงฺขตํ ขนฺธโต ฐเปตฺวา ปญฺจหิ ขนฺเธหิ ทฺวาทสหา\-ยตเนหิ สตฺตรสหิ ธาตูหิ สงฺคหิตา. กติหิ อสงฺคหิตา, น เกหิจิ ขนฺเธหิ น เกหิจิ อายตเนหิ เอกาย ธาตุยา อสงฺคหิตา.}

\proseitem{482}{\csromanfolio{93}อวิตกฺก\-วิจารมตฺเตหิ ธมฺเมหิ เย ธมฺมา. ปีติสหคเตหิ ธมฺเมหิ เย ธมฺมา. สุขสหคเตหิ ธมฺเมหิ เย ธมฺมา วิปฺปยุตฺตา. เต ธมฺมา อสงฺขตํ ขนฺธโต ฐเปตฺวา ปญฺจหิ ขนฺเธหิ ทฺวาทสหา\-ยตเนหิ อฏฺฐารสหิ ธาตูหิ สงฺคหิตา. กติหิ อสงฺคหิตา, น เกหิจิ ขนฺเธหิ น เกหิจิ อายตเนหิ น กาหิจิ ธาตูหิ อสงฺคหิตา.}

\proseitem{483}{อวิตกฺก\-อวิจาเรหิ ธมฺเมหิ เย ธมฺมา วิปฺปยุตฺตา. เต ธมฺมา จตูหิ ขนฺเธหิ ทฺวีหายตเนหิ ตีหิ ธาตูหิ สงฺคหิตา. กติหิ อสงฺคหิตา, เอเกน ขนฺเธน ทสหายตเนหิ ปนฺนรสหิ ธาตูหิ อสงฺคหิตา.}

\proseitem{484}{อุเปกฺขา\-สหคเตหิ ธมฺเมหิ เย ธมฺมา วิปฺปยุตฺตา. เต ธมฺมา อสงฺขตํ ขนฺธโต ฐเปตฺวา ปญฺจหิ ขนฺเธหิ ทฺวาทสหา\-ยตเนหิ เตรสหิ ธาตูหิ สงฺคหิตา. กติหิ อสงฺคหิตา, น เกหิจิ ขนฺเธหิ น เกหิจิ อายตเนหิ ปญฺจหิ ธาตูหิ อสงฺคหิตา.}

\proseitem{485}{ทสฺสเนน ปหาตพฺเพหิ ธมฺเมหิ เย ธมฺมา. ภาวนาย ปหาตพฺเพหิ ธมฺเมหิ เย ธมฺมา. ทสฺสเนน ปหาตพฺพ\-เหตุเกหิ ธมฺเมหิ เย ธมฺมา. ภาวนาย ปหาตพฺพ\-เหตุเกหิ ธมฺเมหิ เย ธมฺมา. อาจยคามีหิ ธมฺเมหิ เย ธมฺมา. อปจยคามีหิ ธมฺเมหิ เย ธมฺมา. เสกฺเขหิ ธมฺเมหิ เย ธมฺมา. อเสกฺเขหิ ธมฺเมหิ เย ธมฺมา. มหคฺคเตหิ ธมฺเมหิ เย ธมฺมา วิปฺปยุตฺตา. เต ธมฺมา อสงฺขตํ ขนฺธโต ฐเปตฺวา ปญฺจหิ ขนฺเธหิ ทฺวาทสหา\-ยตเนหิ อฏฺฐารสหิ ธาตูหิ สงฺคหิตา. กติหิ อสงฺคหิตา, น เกหิจิ ขนฺเธหิ น เกหิจิ อายตเนหิ น กาหิจิ ธาตูหิ อสงฺคหิตา.}

\proseitem{486}{เนว ทสฺสเนน น ภาวนาย ปหาตพฺเพหิ ธมฺเมหิ เย ธมฺมา. เนว ทสฺสเนน น ภาวนาย ปหาตพฺพ\-เหตุเกหิ ธมฺเมหิ เย ธมฺมา. เนวาจยคามินาปจยคามีหิ ธมฺเมหิ เย ธมฺมา. เนว\-เสกฺข\-นาเสกฺเขหิ ธมฺเมหิ เย ธมฺมา. ปริตฺเตหิ ธมฺเมหิ เย ธมฺมา วิปฺปยุตฺตา. เต ธมฺมา จตูหิ ขนฺเธหิ ทฺวีหายตเนหิ ทฺวีหิ ธาตูหิ สงฺคหิตา. กติหิ อสงฺคหิตา, เอเกน ขนฺเธน ทสหายตเนหิ}

\proseitem{487}{\csromanfolio{94}อปฺปมาเณหิ ธมฺเมหิ เย ธมฺมา. ปณีเตหิ ธมฺเมหิ เย ธมฺมา วิปฺปยุตฺตา. เต ธมฺมา จตูหิ ขนฺเธหิ ทฺวีหายตเนหิ อฏฺฐหิ ธาตูหิ สงฺคหิตา. กติหิ อสงฺคหิตา, เอเกน ขนฺเธน ทสหายตเนหิ ทสหิ ธาตูหิ อสงฺคหิตา.}

\proseitem{488}{ปริตฺตา\-รมฺมเณหิ ธมฺเมหิ เย ธมฺมา วิปฺปยุตฺตา. เต ธมฺมา อสงฺขตํ ขนฺธโต ฐเปตฺวา ปญฺจหิ ขนฺเธหิ ทฺวาทสหา\-ยตเนหิ ทฺวาทสหิ ธาตูหิ สงฺคหิตา. กติหิ อสงฺคหิตา, น เกหิจิ ขนฺเธหิ น เกหิจิ อายตเนหิ ฉหิ ธาตูหิ อสงฺคหิตา.}

\proseitem{489}{มหคฺคตา\-รมฺมเณหิ ธมฺเมหิ เย ธมฺมา. อปฺปมาณา\-รมฺมเณหิ ธมฺเมหิ เย ธมฺมา. หีเนหิ ธมฺเมหิ เย ธมฺมา. มิจฺฉตฺต\-นิยเตหิ ธมฺเมหิ เย ธมฺมา. สมฺมตฺต\-นิยเตหิ ธมฺเมหิ เย ธมฺมา. มคฺคารมฺมเณหิ ธมฺเมหิ เย ธมฺมา. มคฺคเหตุเกหิ ธมฺเมหิ เย ธมฺมา. มคฺคาธิปตีหิ ธมฺเมหิ เย ธมฺมา วิปฺปยุตฺตา. เต ธมฺมา อสงฺขตํ ขนฺธโต ฐเปตฺวา ปญฺจหิ ขนฺเธหิ ทฺวาทสหา\-ยตเนหิ อฏฺฐารสหิ ธาตูหิ สงฺคหิตา. กติหิ อสงฺคหิตา, น เกหิจิ ขนฺเธหิ น เกหิจิ อายตเนหิ น กาหิจิ ธาตูหิ}

\proseitem{490}{มชฺฌิเมหิ ธมฺเมหิ เย ธมฺมา. อนิยเตหิ ธมฺเมหิ เย ธมฺมา วิปฺปยุตฺตา. เต ธมฺมา จตูหิ ขนฺเธหิ ทฺวีหายตเนหิ ทฺวีหิ ธาตูหิ สงฺคหิตา. กติหิ อสงฺคหิตา, เอเกน ขนฺเธน ทสหายตเนหิ}

\proseitem{491}{อุปฺปนฺเนหิ ธมฺเมหิ เย ธมฺมา. อนุปฺปนฺเนหิ ธมฺเมหิ เย ธมฺมา. อุปฺปาทีหิ ธมฺเมหิ เย ธมฺมา. อตีเตหิ ธมฺเมหิ เย ธมฺมา. อนาคเตหิ ธมฺเมหิ เย ธมฺมา. ปจฺจุปฺปนฺเนหิ ธมฺเมหิ เย ธมฺมา. อชฺฌตฺเตหิ ธมฺเมหิ เย ธมฺมา. พหิทฺธาหิ ธมฺเมหิ เย ธมฺมา. สนิทสฺสน\-สปฺปฏิเฆหิ ธมฺเมหิ เย ธมฺมา. อนิทสฺสน\-สปฺปฏิเฆหิ ธมฺเมหิ เย ธมฺมา วิปฺปยุตฺตา. เต ธมฺมา จตูหิ ขนฺเธหิ ทฺวีหายตเนหิ อฏฺฐหิ ธาตูหิ สงฺคหิตา. กติหิ อสงฺคหิตา, เอเกน ขนฺเธน ทสหายตเนหิ ทสหิ ธาตูหิ อสงฺคหิตา.}

\proseitem{492}{อตีตารมฺมเณหิ ธมฺเมหิ เย ธมฺมา. อนาคตา\-รมฺมเณหิ ธมฺเมหิ เย ธมฺมา. อชฺฌตฺตา\-รมฺมเณหิ ธมฺเมหิ เย ธมฺมา. พหิทฺธา\-รมฺมเณหิ \csromanfolio{95}ธมฺเมหิ เย ธมฺมา วิปฺปยุตฺตา. เต ธมฺมา อสงฺขตํ ขนฺธโต ฐเปตฺวา ปญฺจหิ ขนฺเธหิ ทฺวาทสหา\-ยตเนหิ อฏฺฐารสหิ ธาตูหิ สงฺคหิตา. กติหิ อสงฺคหิตา, น เกหิจิ ขนฺเธหิ น เกหิจิ อายตเนหิ น กาหิจิ ธาตูหิ อสงฺคหิตา.}

\proseitem{493}{ปจฺจุปฺปนฺนา\-รมฺมเณหิ ธมฺเมหิ เย ธมฺมา. อชฺฌตฺตพหิทฺธา\-รมฺมเณหิ ธมฺเมหิ เย ธมฺมา วิปฺปยุตฺตา. เต ธมฺมา อสงฺขตํ ขนฺธโต ฐเปตฺวา ปญฺจหิ ขนฺเธหิ ทฺวาทสหา\-ยตเนหิ ทฺวาทสหิ ธาตูหิ สงฺคหิตา. กติหิ อสงฺคหิตา, น เกหิจิ ขนฺเธหิ น เกหิจิ อายตเนหิ ฉหิ ธาตูหิ อสงฺคหิตา.}

\csromanmatikaanchor{mtk.57}
\csromantocmarkref{subsection}{5. ทุก}{mtk.57}
\bookmark[dest={mtk.57},level=2]{5. ทุก}
\csromanheader{2}{5. \textbf{ทุก}}
\proseitem{494}{เหตูหิ ธมฺเมหิ เย ธมฺมา. สเหตุเกหิ ธมฺเมหิ เย ธมฺมา. เหตุ\-สมฺปยุตฺเตหิ ธมฺเมหิ เย ธมฺมา. เหตูหิ เจว สเหตุเกหิ จ ธมฺเมหิ เย ธมฺมา. สเหตุเกหิ เจว น จ เหตูหิ ธมฺเมหิ เย ธมฺมา. เหตูหิ เจว เหตุ\-สมฺปยุตฺเตหิ จ ธมฺเมหิ เย ธมฺมา. เหตุ\-สมฺปยุตฺเตหิ เจว น จ เหตูหิ ธมฺเมหิ เย ธมฺมา. น เหตุสเหตุเกหิ ธมฺเมหิ เย ธมฺมา วิปฺปยุตฺตา. เต ธมฺมา อสงฺขตํ ขนฺธโต ฐเปตฺวา ปญฺจหิ ขนฺเธหิ ทฺวาทสหา\-ยตเนหิ อฏฺฐารสหิ ธาตูหิ สงฺคหิตา. กติหิ อสงฺคหิตา, น เกหิจิ ขนฺเธหิ น เกหิจิ อายตเนหิ น กาหิจิ ธาตูหิ อสงฺคหิตา.}

\proseitem{495}{อเหตุเกหิ ธมฺเมหิ เย ธมฺมา. เหตุ\-วิปฺปยุตฺเตหิ ธมฺเมหิ เย ธมฺมา. น เหตุอเหตุเกหิ ธมฺเมหิ เย ธมฺมา วิปฺปยุตฺตา. เต ธมฺมา จตูหิ ขนฺเทหิ ทฺวีหายตเนหิ ทฺวีหิ ธาตูหิ สงฺคหิตา. กติหิ อสงฺคหิตา, เอเกน ขนฺเธน ทสหายตเนหิ โสฬสหิ ธาตูหิ อสงฺคหิตา.}

\proseitem{496}{อปฺปจฺจเยหิ ธมฺเมหิ เย ธมฺมา. อสงฺขเตหิ ธมฺเมหิ เย ธมฺมา. สนิทสฺสเนหิ ธมฺเมหิ เย ธมฺมา. สปฺปฏิเฆหิ ธมฺเมหิ เย ธมฺมา. รูปีหิ ธมฺเมหิ เย ธมฺมา. โลกุตฺตเรหิ ธมฺเมหิ เย ธมฺมา วิปฺปยุตฺตา. เต ธมฺมา จตูหิ ขนฺเธหิ ทฺวีหายตเนหิ อฏฺฐหิ ธาตูหิ สงฺคหิตา. \csromanfolio{96}กติหิ อสงฺคหิตา, เอเกน ขนฺเธน ทสหายตเนหิ ทสหิ ธาตูหิ}

\proseitem{497}{โลกิเยหิ ธมฺเมหิ เย ธมฺมา วิปฺปยุตฺตา. เต ธมฺมา จตูหิ ขนฺเธหิ ทฺวีหายตเนหิ ทฺวีหิ ธาตูหิ สงฺคหิตา. กติหิ อสงฺคหิตา, เอเกน ขนฺเธน ทสหายตเนหิ โสฬสหิ ธาตูหิ อสงฺคหิตา.}

\proseitem{498}{อาสเวหิ ธมฺเมหิ เย ธมฺมา. อาสว\-สมฺปยุตฺเตหิ ธมฺเมหิ เย ธมฺมา. อาสเวหิ เจว สาสเวหิ จ ธมฺเมหิ เย ธมฺมา. อาสเวหิ เจว อาสว\-สมฺปยุตฺเตหิ จ ธมฺเมหิ เย ธมฺมา. อาสว\-สมฺปยุตฺเตหิ เจว โน จ อาสเวหิ ธมฺเมหิ เย ธมฺมา วิปฺปยุตฺตา. เต ธมฺมา อสงฺขตํ ขนฺธโต ฐเปตฺวา ปญฺจหิ ขนฺเธหิ ทฺวาทสหา\-ยตเนหิ อฏฺฐารสหิ ธาตูหิ สงฺคหิตา. กติหิ อสงฺคหิตา, น เกหิจิ ขนฺเธหิ น เกหิจิ อายตเนหิ น กาหิจิ ธาตูหิ อสงฺคหิตา.}

\proseitem{499}{สาสเวหิ ธมฺเมหิ เย ธมฺมา. อาสว\-วิปฺปยุตฺเตหิ ธมฺเมหิ เย ธมฺมา. สาสเวหิ เจว โน จ อาสเวหิ ธมฺเมหิ เย ธมฺมา. อาสว\-วิปฺปยุตฺเตหิ สาสเวหิ ธมฺเมหิ เย ธมฺมา วิปฺปยุตฺตา. เต ธมฺมา จตูหิ ขนฺเธหิ ทฺวีหายตเนหิ ทฺวีหิ ธาตูหิ สงฺคหิตา. กติหิ อสงฺคหิตา, เอเกน ขนฺเธน ทสหายตเนหิ โสฬสหิ ธาตูหิ อสงฺคหิตา.}

\proseitem{500}{อนาสเวหิ ธมฺเมหิ เย ธมฺมา. อาสว\-วิปฺปยุตฺเตหิ อนาสเวหิ ธมฺเมหิ เย ธมฺมา วิปฺปยุตฺตา. เต ธมฺมา จตูหิ ขนฺเธหิ ทฺวีหายตเนหิ อฏฺฐหิ ธาตูหิ สงฺคหิตา. กติหิ อสงฺคหิตา, เอเกน ขนฺเธน ทสหายตเนหิ ทสหิ ธาตูหิ อสงฺคหิตา.}

\proseitem{501}{สญฺโญชเนหิ ธมฺเมหิ เย ธมฺมา. คนฺเถหิ ธมฺเมหิ เย ธมฺมา. โอเฆหิ ธมฺเมหิ เย ธมฺมา. โยเคหิ ธมฺเมหิ เย ธมฺมา. นีวรเณหิ ธมฺเมหิ เย ธมฺมา. ปรามาเสหิ ธมฺเมหิ เย ธมฺมา. ปรามาส\-สมฺปยุตฺเตหิ ธมฺเมหิ เย ธมฺมา. ปรามาเสหิ เจว ปรามฏฺเฐหิ จ ธมฺเมหิ เย ธมฺมา วิปฺปยุตฺตา. เต ธมฺมา อสงฺขตํ ขนฺธโต ฐเปตฺวา ปญฺจหิ ขนฺเธหิ ทฺวาทสหา\-ยตเนหิ อฏฺฐารสหิ ธาตูหิ สงฺคหิตา. กติหิ \csromanfolio{97}อสงฺคหิตา, น เกหิจิ ขนฺเธหิ น เกหิจิ อายตเนหิ น กาหิจิ ธาตูหิ}

\proseitem{502}{ปรามฏฺเฐหิ ธมฺเมหิ เย ธมฺมา. ปรามาส\-วิปฺปยุตฺเตหิ ธมฺเมหิ เย ธมฺมา. ปรามฏฺเฐหิ เจว โน จ ปรามาเสหิ ธมฺเมหิ เย ธมฺมา. ปรามาส\-วิปฺปยุตฺเตหิ ปรามฏฺเฐหิ ธมฺเมหิ เย ธมฺมา วิปฺปยุตฺตา. เต ธมฺมา จตูหิ ขนฺเธหิ ทฺวีหายตเนหิ ทฺวีหิ ธาตูหิ สงฺคหิตา. กติหิ อสงฺคหิตา, เอเกน ขนฺเธน ทสหายตเนหิ โสฬสหิ ธาตูหิ อสงฺคหิตา.}

\proseitem{503}{อปรามฏฺเฐหิ ธมฺเมหิ เย ธมฺมา. ปรามาส\-วิปฺปยุตฺเตหิ อปรามฏฺเฐหิ ธมฺเมหิ เย ธมฺมา วิปฺปยุตฺตา. เต ธมฺมา จตูหิ ขนฺเธหิ ทฺวีหายตเนหิ อฏฺฐหิ ธาตูหิ สงฺคหิตา. กติหิ อสงฺคหิตา, เอเกน ขนฺเธน ทสหายตเนหิ ทสหิ ธาตูหิ อสงฺคหิตา.}

\proseitem{504}{สารมฺมเณหิ ธมฺเมหิ เย ธมฺมา. จิตฺเตหิ ธมฺเมหิ เย ธมฺมา. เจตสิเกหิ ธมฺเมหิ เย ธมฺมา. จิตฺต\-สมฺปยุตฺเตหิ ธมฺเมหิ เย ธมฺมา. จิตฺต\-สํสฏฺเฐหิ ธมฺเมหิ เย ธมฺมา. จิตฺต\-สํสฏฺฐ\-สมุฏฺฐาเนหิ ธมฺเมหิ เย ธมฺมา. จิตฺต\-สํสฏฺฐ\-สมุฏฺฐาน\-สหภูหิ ธมฺเมหิ เย ธมฺมา. จิตฺตสํสฏฺฐสมุฏฺฐานานุปริตฺตีหิ ธมฺเมหิ เย ธมฺมา วิปฺปยุตฺตา. เต ธมฺมา อสงฺขตํ ขนฺธโต ฐเปตฺวา เอเกน ขนฺเธน เอกาทสหา\-ยตเนหิ เอกาทสหิ ธาตูหิ สงฺคหิตา. กติหิ อสงฺคหิตา, จตูหิ ขนฺเธหิ เอเกนายตเนน สตฺตหิ ธาตูหิ อสงฺคหิตา.}

\proseitem{505}{อนารมฺมเณหิ ธมฺเมหิ เย ธมฺมา. จิตฺต\-วิปฺปยุตฺเตหิ ธมฺเมหิ เย ธมฺมา. จิตฺต\-สํสฏฺเฐหิ ธมฺเมหิ เย ธมฺมา. อุปาทาธมฺเมหิ เย ธมฺมา. อนุปาทินฺเนหิ ธมฺเมหิ เย ธมฺมา วิปฺปยุตฺตา. เต ธมฺมา จตูหิ ขนฺเธหิ ทฺวีหายตเนหิ อฏฺฐหิ ธาตูหิ สงฺคหิตา. กติหิ อสงฺคหิตา, เอเกน ขนฺเธน ทสหายตเนหิ ทสหิ ธาตูหิ อสงฺคหิตา.}

\proseitem{506}{อุปาทินฺเนหิ ธมฺเมหิ เย ธมฺมา วิปฺปยุตฺตา. เต ธมฺมา จตูหิ ขนฺเธหิ ทฺวีหายตเนหิ ตีหิ ธาตูหิ สงฺคหิตา. กติหิ อสงฺคหิตา, เอเกน ขนฺเธน ทสหายตเนหิ ปนฺนรสหิ ธาตูหิ อสงฺคหิตา.}

\proseitem{507}{\csromanfolio{98}อุปาทาเนหิ ธมฺเมหิ เย ธมฺมา. กิเลเสหิ ธมฺเมหิ เย ธมฺมา. สํกิลิฏฺเฐหิ ธมฺเมหิ เย ธมฺมา. กิเลส\-สมฺปยุตฺเตหิ ธมฺเมหิ เย ธมฺมา. กิเลเสหิ เจว สํกิเลสิเกหิ จ ธมฺเมหิ เย ธมฺมา. กิเลเสหิ เจว สํกิลิฏฺเฐหิ จ ธมฺเมหิ เย ธมฺมา. สํกิลิฏฺเฐหิ เจว โน จ กิเลเสหิ ธมฺเมหิ เย ธมฺมา. กิเลเสหิ เจว กิเลส\-สมฺปยุตฺเตหิ จ ธมฺเมหิ เย ธมฺมา. กิเลส\-สมฺปยุตฺเตหิ เจว โน จ กิเลเสหิ ธมฺเมหิ เย ธมฺมา วิปฺปยุตฺตา. เต ธมฺมา อสงฺขตํ ขนฺธโต ฐเปตฺวา ปญฺจหิ ขนฺเธหิ ทฺวาทสหา\-ยตเนหิ อฏฺฐารสหิ ธาตูหิ สงฺคหิตา. กติหิ อสงฺคหิตา, น เกหิจิ ขนฺเธหิ น เกหิจิ อายตเนหิ น กาหิจิ ธาตูหิ อสงฺคหิตา.}

\proseitem{508}{สํกิเลสิเกหิ ธมฺเมหิ เย ธมฺมา. อสํกิลิฏฺเฐหิ ธมฺเมหิ เย ธมฺมา. กิเลส\-วิปฺปยุตฺเตหิ ธมฺเมหิ เย ธมฺมา. สํกิเลสิเกหิ เจว โน จ กิเลเสหิ ธมฺเมหิ เย ธมฺมา. กิเลส\-วิปฺปยุตฺเตหิ สํกิเลสิเกหิ ธมฺเมหิ เย ธมฺมา วิปฺปยุตฺตา. เต ธมฺมา จตูหิ ขนฺเธหิ ทฺวีหายตเนหิ ทฺวีหิ ธาตูหิ สงฺคหิตา. กติหิ อสงฺคหิตา, เอเกน ขนฺเธน ทสหายตเนหิ โสฬสหิ ธาตูหิ อสงฺคหิตา.}

\proseitem{509}{อสํกิเลสิเกหิ ธมฺเมหิ เย ธมฺมา. กิเลส\-วิปฺปยุตฺเตหิ อสํกิเลสิเกหิ ธมฺเมหิ เย ธมฺมา วิปฺปยุตฺตา. เต ธมฺมา จตูหิ ขนฺเธหิ ทฺวีหายตเนหิ อฏฺฐหิ ธาตูหิ สงฺคหิตา. กติหิ อสงฺคหิตา, เอเกน ขนฺเธน ทสหายตเนหิ ทสหิ ธาตูหิ อสงฺคหิตา.}

\proseitem{510}{ทสฺสเนน ปหาตพฺเพหิ ธมฺเมหิ เย ธมฺมา. ภาวนาย ปหาตพฺเพหิ ธมฺเมหิ เย ธมฺมา. ทสฺสเนน ปหาตพฺพ\-เหตุเกหิ ธมฺเมหิ เย ธมฺมา. ภาวนาย ปหาตพฺพ\-เหตุเกหิ ธมฺเมหิ เย ธมฺมา วิปฺปยุตฺตา. เต ธมฺมา อสงฺขตํ ขนฺธโต ฐเปตฺวา ปญฺจหิ ขนฺเธหิ ทฺวาทสหา\-ยตเนหิ อฏฺฐารสหิ ธาตูหิ สงฺคหิตา. กติหิ อสงฺคหิตา, น เกหิจิ ขนฺเธหิ น เกหิจิ อายตเนหิ น กาหิจิ ธาตูหิ}

\proseitem{511}{น ทสฺสเนน ปหาตพฺเพหิ ธมฺเมหิ เย ธมฺมา. น ภาวนาย ปหาตพฺเพหิ ธมฺเมหิ เย ธมฺมา. น ทสฺสเนน ปหาตพฺพ\-เหตุเกหิ ธมฺเมหิ เย ธมฺมา. น ภาวนาย ปหาตพฺพ\-เหตุเกหิ ธมฺเมหิ เย \csromanfolio{99}ธมฺมา วิปฺปยุตฺตา. เต ธมฺมา จตูหิ ขนฺเธหิ ทฺวีหายตเนหิ ทฺวีหิ ธาตูหิ สงฺคหิตา. กติหิ อสงฺคหิตา, เอเกน ขนฺเธน ทสหายตเนหิ}

\proseitem{512}{สวิตกฺเกหิ ธมฺเมหิ เย ธมฺมา. สวิจาเรหิ ธมฺเมหิ เย ธมฺมา วิปฺปยุตฺตา. เต ธมฺมา อสงฺขตํ ขนฺธโต ฐเปตฺวา ปญฺจหิ ขนฺเธหิ ทฺวาทสหา\-ยตเนหิ สตฺตรสหิ ธาตูหิ สงฺคหิตา. กติหิ อสงฺคหิตา, น เกหิจิ ขนฺเธหิ น เกหิจิ อายตเนหิ เอกาย ธาตุยา}

\proseitem{513}{สปฺปีติเกหิ ธมฺเมหิ เย ธมฺมา. ปีติสหคเตหิ ธมฺเมหิ เย ธมฺมา. สุขสหคเตหิ ธมฺเมหิ เย ธมฺมา วิปฺปยุตฺตา. เต ธมฺมา อสงฺขตํ ขนฺธโต ฐเปตฺวา ปญฺจหิ ขนฺเธหิ ทฺวาทสหา\-ยตเนหิ อฏฺฐารสหิ ธาตูหิ สงฺคหิตา. กติหิ อสงฺคหิตา, น เกหิจิ ขนฺเธหิ น เกหิจิ อายตเนหิ น กาหิจิ ธาตูหิ อสงฺคหิตา.}

\proseitem{514}{อุเปกฺขา\-สหคเตหิ ธมฺเมหิ เย ธมฺมา วิปฺปยุตฺตา. เต ธมฺมา อสงฺขตํ ขนฺธโต ฐเปตฺวา ปญฺจหิ ขนฺเธหิ ทฺวาทสหา\-ยตเนหิ เตรสหิ ธาตูหิ สงฺคหิตา. กติหิ อสงฺคหิตา, น เกหิจิ ขนฺเธหิ น เกหิจิ อายตเนหิ ปญฺจหิ ธาตูหิ อสงฺคหิตา.}

\proseitem{515}{กามาวจเรหิ ธมฺเมหิ เย ธมฺมา. ปริยาปนฺเนหิ ธมฺเมหิ เย ธมฺมา. สอุตฺตเรหิ ธมฺเมหิ เย ธมฺมา วิปฺปยุตฺตา. เต ธมฺมา จตูหิ ขนฺเธหิ ทฺวีหายตเนหิ ทฺวีหิ ธาตูหิ สงฺคหิตา. กติหิ อสงฺคหิตา, เอเกน ขนฺเธน ทสหายตเนหิ โสฬสหิ ธาตูหิ อสงฺคหิตา.}

\proseitem{516}{น กามาวจเรหิ ธมฺเมหิ เย ธมฺมา. อปริยาปนฺเนหิ ธมฺเมหิ เย ธมฺมา. อนุตฺตเรหิ ธมฺเมหิ เย ธมฺมา วิปฺปยุตฺตา. เต ธมฺมา จตูหิ ขนฺเธหิ ทฺวีหายตเนหิ อฏฺฐหิ ธาตูหิ สงฺคหิตา. กติหิ อสงฺคหิตา, เอเกน ขนฺเธน ทสหายตเนหิ ทสหิ ธาตูหิ อสงฺคหิตา.}

\proseitem{517}{รูปาวจเรหิ ธมฺเมหิ เย ธมฺมา. อรูปาวจเรหิ ธมฺเมหิ เย ธมฺมา. นิยฺยานิเกหิ ธมฺเมหิ เย ธมฺมา. นิยเตหิ ธมฺเมหิ เย ธมฺมา. สรเณหิ ธมฺเมหิ เย ธมฺมา วิปฺปยุตฺตา. เต ธมฺมา อสงฺขตํ ขนฺธโต \csromanfolio{100}ฐเปตฺวา ปญฺจหิ ขนฺเธหิ ทฺวาทสหา\-ยตเนหิ อฏฺฐารสหิ ธาตูหิ สงฺคหิตา. กติหิ อสงฺคหิตา, น เกหิจิ ขนฺเธหิ น เกหิจิ อายตเนหิ น กาหิจิ ธาตูหิ อสงฺคหิตา.}

\proseitem{518}{น รูปาวจเรหิ ธมฺเมหิ เย ธมฺมา. น อรูปาวจเรหิ ธมฺเมหิ เย ธมฺมา. อนิยฺยานิเกหิ ธมฺเมหิ เย ธมฺมา. อนิยเตหิ ธมฺเมหิ เย ธมฺมา. อรเณหิ ธมฺเมหิ เย ธมฺมา วิปฺปยุตฺตา. เต ธมฺมา กติหิ ขนฺเธหิ กติหายตเนหิ กติหิ ธาตูหิ สงฺคหิตา, เต ธมฺมา จตูหิ ขนฺเธหิ ทฺวีหายตเนหิ ทฺวีหิ ธาตูหิ สงฺคหิตา. กติหิ อสงฺคหิตา, เอเกน ขนฺเธน ทสหายตเนหิ โสฬสหิ ธาตูหิ อสงฺคหิตา.}

\setlength{\csromangathapagewidth}{0pt}
\setlength{\csromangathawakwidth}{0pt}
\csromangathameasuregroup{ธมฺมายตนํ ธมฺมธาตุ, \\ สฬายตนํ ชาติชรามตํ, \\ ปฐมนฺตเร สตฺต จ, \\ จุทฺทส ฉ จ มตฺถเก, \\ สมุจฺเฉเท น ลพฺภนฺติ,}{\csromangathabat{ธมฺมายตนํ ธมฺมธาตุ,}{อถ ชีวิตํ นามรูปํ.} \\ \csromangathabat{สฬายตนํ ชาติชรามตํ,}{เทฺว จ ติเก น ลพฺภเร.} \\ \csromangathabat{ปฐมนฺตเร สตฺต จ,}{โคจฺฉเก ทส อปรนฺเต.} \\ \csromangathabat{จุทฺทส ฉ จ มตฺถเก,}{อิจฺเจเต สตฺตจตฺตาลีส ธมฺมา.} \\ \csromangathabat{สมุจฺเฉเท น ลพฺภนฺติ,}{โมฆปุจฺฉเกน จาติ.}}
\csromangathasetleft{ธมฺมายตนํ ธมฺมธาตุ, \\ สฬายตนํ ชาติชรามตํ, \\ ปฐมนฺตเร สตฺต จ, \\ จุทฺทส ฉ จ มตฺถเก, \\ สมุจฺเฉเท น ลพฺภนฺติ,}
\csromangathagroup{\csromangathastanza{\csromangathabat{ธมฺมายตนํ ธมฺมธาตุ,}{อถ ชีวิตํ นามรูปํ.} \\ \csromangathabat{สฬายตนํ ชาติชรามตํ,}{เทฺว จ ติเก น ลพฺภเร.}}\csromangathastanza{\csromangathabat{ปฐมนฺตเร สตฺต จ,}{โคจฺฉเก ทส อปรนฺเต.} \\ \csromangathabat{จุทฺทส ฉ จ มตฺถเก,}{อิจฺเจเต สตฺตจตฺตาลีส ธมฺมา.} \\ \csromangathabat{สมุจฺเฉเท น ลพฺภนฺติ,}{โมฆปุจฺฉเกน จาติ.}}}

\prosewithcloser{วิปฺปยุตฺเตน สงฺคหิตา\-สงฺคหิต\-ปท\-นิทฺเทโส จุทฺทสโม.}{ธาตุกถา\-ปกรณํ นิฏฺฐิตํ.}
\csromanpalirecto
\pitaka{อภิธมฺม\-ปิฏก}
\csromanmatikaanchor{mtk.58}
\csromantocmarknopage{chapter}{ปุคฺคลปญฺญตฺติปาฬิ}{mtk.58}
\bookmark[dest={mtk.58},level=0]{ปุคฺคลปญฺญตฺติปาฬิ}
\gambhira{ปุคฺคล\-ปญฺญตฺติปาฬิ}
\namakkaram{นโม ตสฺส ภควโต อรหโต สมฺมา\-สมฺพุทฺธสฺส.}
\csromanmatikaanchor{mtk.59}
\csromantocmarkref{section}{1. เอกกอุทฺเทส}{mtk.59}
\bookmark[dest={mtk.59},level=1]{1. เอกกอุทฺเทส}
\prose{\csromanfolio{101}1. เอกกอุทฺเทส}

\proseitem{1}{ฉ ปญฺญตฺติโย\csromandash{}ขนฺธ\-ปญฺญตฺติ อายตน\-ปญฺญตฺติ ธาตุปญฺญตฺติ สจฺจปญฺญตฺติ อินฺทฺริย\-ปญฺญตฺติ ปุคฺคลปญฺญตฺตีติ.}

\proseitem{2}{กิตฺตาวตา ขนฺธานํ ขนฺธ\-ปญฺญตฺติ, ยาวตา ปญฺจกฺขนฺธา\csromandash{} รูปกฺขนฺโธ เวทนากฺขนฺโธ สญฺญากฺขนฺโธ สงฺขาร\-กฺขนฺโธ วิญฺญาณ\-กฺขนฺโธ, เอตฺตาวตา ขนฺธานํ ขนฺธ\-ปญฺญตฺติ.}

\proseitem{3}{กิตฺตาวตา อายตนานํ อายตน\-ปญฺญตฺติ, ยาวตา ทฺวาทสา\-ยตนานิ\csromandash{} จกฺขายตนํ รูปายตนํ โสตายตนํ สทฺทายตนํ ฆานายตนํ คนฺธายตนํ ชิวฺหายตนํ รสายตนํ กายายตนํ โผฏฺฐพฺพา\-ยตนํ มนายตนํ ธมฺมายตนํ, เอตฺตาวตา อายตนานํ อายตน\-ปญฺญตฺติ.}

\proseitem{4}{กิตฺตาวตา ธาตูนํ ธาตุปญฺญตฺติ, ยาวตา อฏฺฐารส ธาตุโย\csromandash{} จกฺขุธาตุ รูปธาตุ จกฺขุวิญฺญาณ\-ธาตุ โสตธาตุ สทฺทธาตุ โสตวิญฺญาณ\-ธาตุ ฆานธาตุ คนฺธธาตุ ฆานวิญฺญาณ\-ธาตุ ชิวฺหาธาตุ รสธาตุ ชิวฺหาวิญฺญาณ\-ธาตุ กายธาตุ โผฏฺฐพฺพ\-ธาตุ กายวิญฺญาณ\-ธาตุ มโนธาตุ ธมฺมธาตุ มโนวิญฺญาณ\-ธาตุ, เอตฺตาวตา ธาตูนํ ธาตุปญฺญตฺติ.}

\proseitem{5}{\csromanfolio{102}กิตฺตาวตา สจฺจานํ สจฺจปญฺญตฺติ, ยาวตา จตฺตาริ สจฺจานิ\csromandash{} ทุกฺขสจฺจํ สมุทยสจฺจํ นิโรธสจฺจํ มคฺคสจฺจํ, เอตฺตาวตา สจฺจานํ สจฺจปญฺญตฺติ.}

\proseitem{6}{กิตฺตาวตา อินฺทฺริยานํ อินฺทฺริย\-ปญฺญตฺติ, ยาวตา พาวีสติ\-นฺทฺริยานิ\csromandash{} จกฺขุนฺทฺริยํ โสตินฺทฺริยํ ฆานินฺทฺริยํ ชิวฺหินฺทฺริยํ กายินฺทฺริยํ มนินฺทฺริยํ อิตฺถินฺทฺริยํ ปุริสินฺทฺริยํ ชีวิตินฺทฺริยํ สุขินฺทฺริยํ ทุกฺขินฺทฺริยํ โสมนสฺสิ\-นฺทฺริยํ โทมนสฺสิ\-นฺทฺริยํ อุเปกฺขินฺทฺริยํ สทฺธินฺทฺริยํ วีริยินฺทฺริยํ สตินฺทฺริยํ สมาธินฺทฺริยํ ปญฺญินฺทฺริยํ อนญฺญาต\-ญฺญสฺสามี\-ตินฺทฺริยํ อญฺญินฺทฺริยํ อญฺญาตาวิ\-นฺทฺริยํ, เอตฺตาวตา อินฺทฺริยานํ อินฺทฺริย\-ปญฺญตฺติ.}

\proseitem{7}{กิตฺตาวตา ปุคฺคลานํ ปุคฺคล\-ปญฺญตฺติ,}

\proseitem{1}{สมยวิมุตฺโต,}

\proseitem{2}{อสมยวิมุตฺโต,}

\proseitem{3}{กุปฺปธมฺโม,}

\proseitem{4}{อกุปฺปธมฺโม,}

\proseitem{5}{ปริหานธมฺโม,}

\proseitem{6}{อปริหาน\-ธมฺโม,}

\proseitem{7}{เจตนาภพฺโพ,}

\proseitem{8}{อนุรกฺขณา\-ภพฺโพ,}

\proseitem{9}{ปุถุชฺชโน,}

\proseitem{10}{โคตฺรภู,}

\proseitem{11}{ภยูปรโต,}

\proseitem{12}{อภยูปรโต,}

\proseitem{13}{ภพฺพาคมโน,}

\proseitem{14}{อภพฺพาคมโน,}

\proseitem{15}{นิยโต,}

\proseitem{16}{อนิยโต,}

\proseitem{17}{\csromanfolio{103}ปฏิปนฺนโก,}

\proseitem{18}{ผเลฐิโต,}

\proseitem{19}{สมสีสี,}

\proseitem{20}{ฐิตกปฺปี,}

\proseitem{21}{อริโย,}

\proseitem{22}{อนริโย,}

\proseitem{23}{เสกฺโข,}

\proseitem{24}{อเสกฺโข,}

\proseitem{25}{เนว\-เสกฺข\-นา\-เสกฺโข,}

\proseitem{26}{เตวิชฺโช,}

\proseitem{27}{ฉฬภิญฺโญ,}

\proseitem{28}{สมฺมาสมฺพุทฺโธ,}

\proseitem{29}{ปจฺเจก\-สมฺพุทฺโธ\csromansharedfootnote{0c129bde9ceba893}{ปจฺเจกพุทฺโธ (สี)},}

\proseitem{30}{อุภโตภาค\-วิมุตฺโต,}

\proseitem{31}{ปญฺญาวิมุตฺโต,}

\proseitem{32}{กายสกฺขี,}

\proseitem{33}{ทิฏฺฐิปฺปตฺโต,}

\proseitem{34}{สทฺธาวิมุตฺโต,}

\proseitem{35}{ธมฺมานุสารี,}

\proseitem{36}{สทฺธานุสารี,}

\proseitem{37}{สตฺตกฺขตฺตุ\-ปรโม,}

\proseitem{38}{โกลํโกโล,}

\proseitem{39}{เอกพีชี,}

\csromanmatikaanchor{mtk.60}
\csromantocmarkref{section}{2. ทุกอุทฺเทส}{mtk.60}
\bookmark[dest={mtk.60},level=1]{2. ทุกอุทฺเทส}
\proseitem{40}{\csromanfolio{104}สกทาคามี,}

\proseitem{41}{อนาคามี,}

\proseitem{42}{อนฺตรา\-ปรินิพฺพายี,}

\proseitem{43}{อุปหจฺจ\-ปรินิพฺพายี,}

\proseitem{44}{อสงฺขารปรินิพฺพายี,}

\proseitem{45}{สสงฺขาร\-ปรินิพฺพายี,}

\proseitem{46}{อุทฺธํโสโต\-อกนิฏฺฐคามี,}

\proseitem{47}{โสตาปนฺโน,}

\proseitem{48}{โสตาปตฺติผล\-สจฺฉิกิริยาย ปฏิปนฺโน,}

\proseitem{49}{สกทาคามี,}

\proseitem{50}{สกทาคามิ\-ผล\-สจฺฉิกิริยาย ปฏิปนฺโน,}

\proseitem{51}{อนาคามี,}

\proseitem{52}{อนาคามิ\-ผล\-สจฺฉิกิริยาย ปฏิปนฺโน,}

\proseitem{53}{อรหา,}

\proseitem{54}{อรหตฺตผล\-สจฺฉิกิริยาย\csromansharedfootnote{552f6758ac9cbc73}{อรหตฺตาย (อิ)} ปฏิปนฺโน,}

\vspace{\tipitakaparskip}
\csromancenter{เอกกํ.}
\prose{2. ทุกอุทฺเทส}

\prose{8. \textbf{เทฺว ปุคฺคลา\csromandash{}}}

\proseitem{1}{โกธโน จ, อุปนาหี จ.}

\proseitem{2}{มกฺขี จ, ปฬาสี\csromansharedfootnote{aae52de8288d3969}{ปลาสี (สฺยา, ก)} จ.}

\proseitem{3}{อิสฺสุกี จ, มจฺฉรี จ.}

\proseitem{4}{สโฐ จ, มายาวี จ.}

\proseitem{5}{\csromanfolio{105}อหิริโก จ, อโนตฺตปฺปี จ.}

\proseitem{6}{ทุพฺพโจ จ, ปาปมิตฺโต จ.}

\proseitem{7}{อินฺทฺริเยสุ อคุตฺตทฺวาโร จ, โภชเน อมตฺตญฺญู จ.}

\proseitem{8}{มุฏฺฐสฺสติ จ, อสมฺปชาโน จ.}

\proseitem{9}{สีลวิปนฺโน จ, ทิฏฺฐิวิปนฺโน จ.}

\proseitem{10}{อชฺฌตฺต\-สํโยชโน จ, พหิทฺธา\-สํโยชโน จ.}

\proseitem{11}{อกฺโกธโน จ, อนุปนาหี จ.}

\proseitem{12}{อมกฺขี จ, อปฬาสี จ.}

\proseitem{13}{อนิสฺสุกี จ, อมจฺฉรี จ.}

\proseitem{14}{อสโฐ จ, อมายาวี จ.}

\proseitem{15}{หิริมา จ, โอตฺตปฺปี จ.}

\proseitem{16}{สุวโจ จ, กลฺยาณมิตฺโต จ.}

\proseitem{17}{อินฺทฺริเยสุ คุตฺตทฺวาโร จ, โภชเน มตฺตญฺญู จ.}

\proseitem{18}{อุปฏฺฐิตสฺสติ จ, สมฺปชาโน จ.}

\proseitem{19}{สีลสมฺปนฺโน จ, ทิฏฺฐิสมฺปนฺโน จ.}

\proseitem{20}{เทฺว ปุคฺคลา ทุลฺลภา โลกสฺมิํ.}

\proseitem{21}{เทฺว ปุคฺคลา ทุตฺตปฺปยา.}

\proseitem{22}{เทฺว ปุคฺคลา สุตปฺปยา.}

\proseitem{23}{ทฺวินฺนํ ปุคฺคลานํ อาสวา วฑฺฒนฺติ.}

\proseitem{24}{ทฺวินฺนํ ปุคฺคลานํ อาสวา น วฑฺฒนฺติ.}

\proseitem{25}{หีนาธิมุตฺโต จ, ปณีตาธิมุตฺโต จ.}

\proseitem{26}{ติตฺโต จ, ตปฺเปตา จ.}

\vspace{\tipitakaparskip}
\csromancenter{ทุกํ.}
\csromanpalirecto
\csromanmatikaanchor{mtk.61}
\csromantocmarkref{section}{3. ติกอุทฺเทส}{mtk.61}
\bookmark[dest={mtk.61},level=1]{3. ติกอุทฺเทส}
\csromanmarkh{ปุคฺคล\-ปญฺญตฺติปาฬิ}\csromanmarkhii{3. ติกอุทฺเทส}\csromanheadernomark{1}{ปุคฺคล\-ปญฺญตฺติปาฬิ 3. ติกอุทฺเทส}
\prose{\csromanfolio{106}9. \textbf{ตโย ปุคฺคลา\csromandash{}}}

\proseitem{1}{นิราโส, อาสํโส, วิคตาโส.}

\proseitem{2}{ตโย คิลานูปมา ปุคฺคลา.}

\proseitem{3}{กายสกฺขี, ทิฏฺฐิปฺปตฺโต, สทฺธาวิมุตฺโต.}

\proseitem{4}{คูถภาณี, ปุปฺผภาณี, มธุภาณี.}

\proseitem{5}{อรุกูปมจิตฺโต ปุคฺคโล, วิชฺชูปมจิตฺโต ปุคฺคโล, วชิรูปม\-จิตฺโต ปุคฺคโล.}

\proseitem{6}{อนฺโธ, เอกจกฺขุ, ทฺวิจกฺขุ.}

\proseitem{7}{อวกุชฺชปญฺโญ ปุคฺคโล, อุจฺฉงฺคปญฺโญ\csromansharedfootnote{64abe58e0dd4139f}{อุจฺจงฺคปญฺโญ (สฺยา)} ปุคฺคโล, ปุถุปญฺโญ ปุคฺคโล.}

\proseitem{8}{อตฺเถกจฺโจ ปุคฺคโล กาเมสุ จ ภเวสุ จ อวีตราโค, อตฺเถกจฺโจ ปุคฺคโล กาเมสุ วีตราโค ภเวสุ อวีตราโค, อตฺเถกจฺโจ ปุคฺคโล กาเมสุ จ ภเวสุ จ วีตราโค.}

\proseitem{9}{ปาสาณ\-เลขู\-ปโม ปุคฺคโล, ปถวิ\-เลขู\-ปโม ปุคฺคโล, อุทกเลขูปโม ปุคฺคโล.}

\proseitem{10}{ตโย โปตฺถกูปมา ปุคฺคลา.}

\proseitem{11}{ตโย กาสิก\-วตฺถู\-ปมา ปุคฺคลา.}

\proseitem{12}{สุปฺปเมยฺโย, ทุปฺปเมยฺโย, อปฺปเมยฺโย.}

\proseitem{13}{อตฺเถกจฺโจ ปุคฺคโล น เสวิตพฺโพ น ภชิตพฺโพ น ปยิรุปาสิตพฺโพ, อตฺเถกจฺโจ ปุคฺคโล เสวิตพฺโพ ภชิตพฺโพ ปยิรุปาสิตพฺโพ, อตฺเถกจฺโจ ปุคฺคโล สกฺกตฺวา ครุํ กตฺวา\csromansharedfootnote{7cd8cc2b601a152c}{ครุกตฺวา (สี)} เสวิตพฺโพ ภชิตพฺโพ ปยิรุปสิตพฺโพ.}

\proseitem{14}{\csromanfolio{107}อตฺเถกจฺโจ ปุคฺคโล ชิคุจฺฉิตพฺโพ น เสวิตพฺโพ น ภชิตพฺโพ น ปยิรุปาสิตพฺโพ, อตฺเถกจฺโจ ปุคฺคโล อชฺฌุเปกฺขิตพฺโพ น เสวิตพฺโพ น ภชิตพฺโพ น ปยิรุปาสิตพฺโพ, อตฺเถกจฺโจ ปุคฺคโล เสวิตพฺโพ ภชิตพฺโพ ปยิรุปาสิตพฺโพ.}

\proseitem{15}{อตฺเถกจฺโจ ปุคฺคโล สีเลสุ ปริปูรการี\csromansharedfootnote{f25f5271f7350746}{ปริปูรีการี (สฺยา)} สมาธิสฺมิํ มตฺตโส การี ปญฺญาย มตฺตโส การี, อตฺเถกจฺโจ ปุคฺคโล สีเลสุ จ ปริปูรการี สมาธิสฺมิญฺจ ปริปูรการี ปญฺญาย มตฺตโส การี, อตฺเถกจฺโจ ปุคฺคโล สีเลสุ จ ปริปูรการี สมาธิสฺมิญฺจ ปริปูรการี ปญฺญาย จ ปริปูรการี.}

\proseitem{16}{ตโย สตฺถาโร.}

\proseitem{17}{อปเรปิ ตโย สตฺถาโร.}

\vspace{\tipitakaparskip}
\csromancenter{ติกํ.}
\csromanmatikaanchor{mtk.62}
\csromantocmarkref{section}{4. จตุกฺกอุทฺเทส}{mtk.62}
\bookmark[dest={mtk.62},level=1]{4. จตุกฺกอุทฺเทส}
\csromanheader{1}{4. จตุกฺกอุทฺเทส}
\prose{10. \textbf{จตฺตาโร ปุคฺคลา\csromandash{}}}

\proseitem{1}{อสปฺปุริโส, อสปฺปุริเสน อสปฺปุริสตโร, สปฺปุริโส, สปฺปุริเสน สปฺปุริสตโร.}

\proseitem{2}{ปาโป, ปาเปน ปาปตโร, กลฺยาโณ, กลฺยาเณน กลฺยาณตโร.}

\proseitem{3}{ปาปธมฺโม, ปาปธมฺเมน ปาปธมฺมตโร, กลฺยาณธมฺโม, กลฺยาณ\-ธมฺเมน กลฺยาณ\-ธมฺม\-ตโร.}

\proseitem{4}{สาวชฺโช, วชฺชพหุโล, อปฺปวชฺโช\csromansharedfootnote{e683448f042669e3}{อปฺปสาวชฺโช (สฺยา, ก) อํ 1. 452 ปิฏฺเฐปิ.}, อนวชฺโช.}

\proseitem{5}{อุคฺฆฏิตญฺญู, วิปญฺจิตญฺญู\csromansharedfootnote{ea20b0049a6624c2}{วิปจิตญฺญู (สี) อํ 1. 452 ปิฏฺเฐปิ.}, เนยฺโย, ปทปรโม.}

\proseitem{6}{\csromanfolio{108}ยุตฺต\-ปฺปฏิภาโน โน มุตฺต\-ปฺปฏิภาโน, มุตฺต\-ปฺปฏิภาโน โน ยุตฺต\-ปฺปฏิภาโน, ยุตฺต\-ปฺปฏิภาโน จ มุตฺต\-ปฺปฏิภาโน จ, เนว ยุตฺต\-ปฺปฏิภาโน โน มุตฺต\-ปฺปฏิภาโน.}

\proseitem{7}{จตฺตาโร ธมฺมกถิกา ปุคฺคลา.}

\proseitem{8}{จตฺตาโร วลาหกูปมา ปุคฺคลา.}

\proseitem{9}{จตฺตาโร มูสิกูปมา ปุคฺคลา.}

\proseitem{10}{จตฺตาโร อมฺพูปมา ปุคฺคลา.}

\proseitem{11}{จตฺตาโร กุมฺภูปมา ปุคฺคลา.}

\proseitem{12}{จตฺตาโร อุทกรหทู\-ปมา ปุคฺคลา.}

\proseitem{13}{จตฺตาโร พลีพทฺทูปมา\csromansharedfootnote{55fe36c61c06c1fe}{พลิพทฺทูปมา (สี)} ปุคฺคลา.}

\proseitem{14}{จตฺตาโร อาสีวิสูปมา ปุคฺคลา.}

\proseitem{15}{อตฺเถกจฺโจ ปุคฺคโล อนนุวิจฺจ อปริโยคาเหตฺวา อวณฺณารหสฺส วณฺณํ ภาสิตา โหติ, อตฺเถกจฺโจ ปุคฺคโล อนนุวิจฺจ อปริโยคาเหตฺวา วณฺณารหสฺส อวณฺณํ ภาสิตา โหติ, อตฺเถกจฺโจ ปุคฺคโล อนนุวิจฺจ อปริโยคาเหตฺวา อปฺปสาทนีเย ฐาเน ปสาทํ อุปทํสิตา โหติ, อตฺเถกจฺโจ ปุคฺคโล อนนุวิจฺจ อปริโยคาเหตฺวา ปสาทนีเย ฐาเน อปฺปสาทํ อุปทํสิตา โหติ.}

\proseitem{16}{อตฺเถกจฺโจ ปุคฺคโล อนุวิจฺจ ปริโยคาเหตฺวา อวณฺณารหสฺส อวณฺณํ ภาสิตา โหติ, อตฺเถกจฺโจ ปุคฺคโล อนุวิจฺจ ปริโยคาเหตฺวา วณฺณารหสฺส วณฺณํ ภาสิตา โหติ, อตฺเถกจฺโจ ปุคฺคโล อนุวิจฺจ ปริโยคาเหตฺวา อปฺปสาทนีเย ฐาเน อปฺปสาทํ อุปทํสิตา โหติ, อตฺเถกจฺโจ ปุคฺคโล อนุวิจฺจ ปริโยคาเหตฺวา ปสาทนีเย ฐาเน ปสาทํ อุปทํสิตา โหติ.}

\proseitem{17}{อตฺเถกจฺโจ ปุคฺคโล อวณฺณารหสฺส อวณฺณํ ภาสิตา โหติ ภูตํ ตจฺฉํ กาเลน, โน จ โข วณฺณารหสฺส วณฺณํ ภาสิตา โหติ ภูตํ ตจฺฉํ กาเลน. อตฺเถกจฺโจ ปุคฺคโล วณฺณารหสฺส วณฺณํ}

\prose{\csromanfolio{109}ภาสิตา โหติ ภูตํ ตจฺฉํ กาเลน, โน จ โข อวณฺณารหสฺส อวณฺณํ ภาสิตา โหติ ภูตํ ตจฺฉํ กาเลน. อตฺเถกจฺโจ ปุคฺคโล อวณฺณารหสฺส จ อวณฺณํ ภาสิตา โหติ ภูตํ ตจฺฉํ กาเลน, วณฺณารหสฺส จ วณฺณํ ภาสิตา โหติ ภูตํ ตจฺฉํ กาเลน. อตฺเถกจฺโจ ปุคฺคโล เนว อวณฺณารหสฺส อวณฺณํ ภาสิตา โหติ ภูตํ ตจฺฉํ กาเลน, โน จ วณฺณารหสฺส วณฺณํ ภาสิตา โหติ ภูตํ ตจฺฉํ กาเลน.}

\proseitem{18}{อุฏฺฐาน\-ผลูปชีวี โน ปุญฺญผลู\-ปชีวี, ปุญฺญผลู\-ปชีวี โน อุฏฺฐาน\-ผลูปชีวี, อุฏฺฐาน\-ผลูปชีวี จ ปุญฺญผลู\-ปชีวี จ, เนว อุฏฺฐาน\-ผลูปชีวี โน ปุญฺญผลู\-ปชีวี.}

\proseitem{19}{ตโม ตมปรายโน, ตโม โชติปรายโน, โชติ ตมปรายโน, โชติ โชติปรายโน.}

\proseitem{20}{โอณโตณโต, โอณตุณฺณโต, อุณฺณโตณโต, อุณฺณตุณฺณโต.}

\proseitem{21}{จตฺตาโร รุกฺขูปมา ปุคฺคลา.}

\proseitem{22}{รูปปฺปมาโณ, รูปปฺปสนฺโน, โฆสปฺปมาโณ, โฆสปฺปสนฺโน.}

\proseitem{23}{ลูขปฺปมาโณ, ลูขปฺปสนฺโน, ธมฺมปฺปมาโณ, ธมฺมปฺปสนฺโน.}

\proseitem{24}{อตฺเถกจฺโจ ปุคฺคโล อตฺตหิตาย ปฏิปนฺโน โหติ โน ปรหิตาย, อตฺเถกจฺโจ ปุคฺคโล ปรหิตาย ปฏิปนฺโน โหติ โน อตฺตหิตาย, อตฺเถกจฺโจ ปุคฺคโล อตฺตหิตาย เจว ปฏิปนฺโน โหติ ปรหิตาย จ, อตฺเถกจฺโจ ปุคฺคโล เนว อตฺตหิตาย ปฏิปนฺโน โหติ โน ปรหิตาย.}

\proseitem{25}{อตฺเถกจฺโจ ปุคฺคโล อตฺตนฺตโป โหติ อตฺต\-ปริตาปนา\-นุโยคม\-นุยุตฺโต, อตฺเถกจฺโจ ปุคฺคโล ปรนฺตโป โหติ ปรปริตาปนา\-นุโยคม\-นุยุตฺโต, อตฺเถกจฺโจ ปุคฺคโล อตฺตนฺตโป จ โหติ \csromanfolio{110}อตฺต\-ปริตาปนา\-นุโยคม\-นุยุตฺโต ปรนฺตโป จ ปรปริตาปนา\-นุโยคม\-นุยุตฺโต, อตฺเถกจฺโจ ปุคฺคโล เนว อตฺตนฺตโป โหติ น อตฺต\-ปริตาปนา\-นุโยคม\-นุยุตฺโต น ปรนฺตโป น ปรปริตาปนา\-นุโยคม\-นุยุตฺโต. โส อนตฺตนฺตโป อปรนฺตโป ทิฏฺเฐว ธมฺเม นิจฺฉาโต นิพฺพุโต สีตีภูโต\csromansharedfootnote{37b56aae71fc080c}{สีติภูโต (สี, ก)} สุขปฺปฏิสํเวที พฺรหฺมภูเตน อตฺตนา วิหรติ.}

\proseitem{26}{สราโค, สโทโส, สโมโห, สมาโน.}

\proseitem{27}{อตฺเถกจฺโจ ปุคฺคโล ลาภี โหติ อชฺฌตฺตํ เจโตสมถสฺส น ลาภี อธิปญฺญา\-ธมฺม\-วิปสฺสนาย, อตฺเถกจฺโจ ปุคฺคโล ลาภี โหติ อธิปญฺญา\-ธมฺม\-วิปสฺสนาย น ลาภี อชฺฌตฺตํ เจโตสมถสฺส, อตฺเถกจฺโจ ปุคฺคโล ลาภี เจว โหติ อชฺฌตฺตํ เจโตสมถสฺส ลาภี จ อธิปญฺญา\-ธมฺม\-วิปสฺสนาย, อตฺเถกจฺโจ ปุคฺคโล เนว ลาภี โหติ อชฺฌตฺตํ เจโตสมถสฺส น ลาภี อธิปญฺญา\-ธมฺม\-วิปสฺสนาย.}

\proseitem{28}{อนุโสตคามี ปุคฺคโล, ปฏิโสตคามี ปุคฺคโล, ฐิตตฺโต ปุคฺคโล, ติณฺโณ ปารงฺคโต\csromansharedfootnote{279a8826148dcad3}{ปารคโต (สี,สฺยา)} ถเล ติฏฺฐติ พฺราหฺมโณ.}

\proseitem{29}{อปฺปสฺสุโต สุเตน อนุปปนฺโน, อปฺปสฺสุโต สุเตน อุปปนฺโน, พหุสฺสุโต สุเตน อนุปปนฺโน, พหุสฺสุโต สุเตน อุปปนฺโน.}

\proseitem{30}{สมณมจโล, สมณปทุโม, สมณ\-ปุณฺฑรีโก, สมเณสุ สมณ\-สุขุมาโล.}

\vspace{\tipitakaparskip}
\csromancenter{จตุกฺกํ.}
\csromanmatikaanchor{mtk.63}
\csromantocmarkref{section}{5. ปญฺจกอุทฺเทส}{mtk.63}
\bookmark[dest={mtk.63},level=1]{5. ปญฺจกอุทฺเทส}
\csromanheader{1}{5. ปญฺจกอุทฺเทส}
\prose{11. \textbf{ปญฺจ ปุคฺคลา\csromandash{}}}

\proseitem{25}{อตฺเถกจฺโจ ปุคฺคโล อารภติ\csromansharedfootnote{fdbd7fe7930b6159}{อารมฺภติ (สี, สฺยา)} จ วิปฺปฏิสารี จ โหติ, ตญฺจ เจโตวิมุตฺติํ ปญฺญาวิมุตฺติํ ยถาภูตํ นปฺปชานาติ, ยตฺถสฺส เต อุปฺปนฺนา}

\prose{\csromanfolio{111}ปาปกา อกุสลา ธมฺมา อปริเสสา นิรุชฺฌนฺติ. อตฺเถกจฺโจ ปุคฺคโล อารภติ น วิปฺปฏิสารี จ โหติ, ตญฺจ เจโตวิมุตฺติํ ปญฺญาวิมุตฺติํ ยถาภูตํ นปฺปชานาติ, ยตฺถสฺส เต อุปฺปนฺนา ปาปกา อกุสลา ธมฺมา อปริเสสา นิรุชฺฌนฺติ. อตฺเถกจฺโจ ปุคฺคโล นารภติ วิปฺปฏิสารี จ โหติ, ตญฺจ เจโตวิมุตฺติํ ปญฺญาวิมุตฺติํ ยถาภูตํ นปฺปชานาติ, ยตฺถสฺส เต อุปฺปนฺนา ปาปกา อกุสลา ธมฺมา อปริเสสา นิรุชฺฌนฺติ. อตฺเถกจฺโจ ปุคฺคโล นารภติ น วิปฺปฏิสารี โหติ, ตญฺจ เจโตวิมุตฺติํ ปญฺญาวิมุตฺติํ ยถาภูตํ นปฺปชานาติ, ยตฺถสฺส เต อุปฺปนฺนา ปาปกา อกุสลา ธมฺมา อปริเสสา นิรุชฺฌนฺติ. อตฺเถกจฺโจ ปุคฺคโล นารภติ น วิปฺปฏิสารี โหติ, ตญฺจ เจโตวิมุตฺติํ ปญฺญาวิมุตฺติํ ยถาภูตํ ปชานาติ, ยตฺถสฺส เต อุปฺปนฺนา ปาปกา อกุสลา ธมฺมา อปริเสสา นิรุชฺฌนฺติ.}

\proseitem{2}{ทตฺวา อวชานาติ, สํวาเสน อวชานาติ, อาเธยฺยมุโข โหติ, โลโล โหติ, มนฺโท โมมูโห โหติ.}

\proseitem{3}{ปญฺจ โยธาชีวูปมา ปุคฺคลา.}

\proseitem{4}{ปญฺจ ปิณฺฑปาติกา.}

\proseitem{5}{ปญฺจ ขลุปจฺฉาภตฺติกา.}

\proseitem{6}{ปญฺจ เอกาสนิกา.}

\proseitem{7}{ปญฺจ ปํสุกูลิกา.}

\proseitem{8}{ปญฺจ เตจีวริกา.}

\proseitem{9}{ปญฺจ อารญฺญิกา.}

\proseitem{10}{ปญฺจ รุกฺขมูลิกา.}

\proseitem{11}{ปญฺจ อพฺโภกาสิกา.}

\proseitem{12}{ปญฺจ เนสชฺชิกา.}

\proseitem{13}{ปญฺจ ยถา\-สนฺถติกา.}

\proseitem{14}{ปญฺจ โสสานิกา.}

\vspace{\tipitakaparskip}
\csromancenter{ปญฺจกํ.}
\csromanpalirecto
\csromanmatikaanchor{mtk.64}
\csromantocmarkref{section}{6. ฉกฺกอุทฺเทส}{mtk.64}
\bookmark[dest={mtk.64},level=1]{6. ฉกฺกอุทฺเทส}
\csromanmarkh{ปุคฺคล\-ปญฺญตฺติปาฬิ}\csromanmarkhii{6. ฉกฺกอุทฺเทส}\csromanheadernomark{1}{ปุคฺคล\-ปญฺญตฺติปาฬิ 6. ฉกฺกอุทฺเทส}
\prose{\csromanfolio{112}12. \textbf{ฉ ปุคฺคลา\csromandash{}}}

\proseitem{1}{อตฺเถกจฺโจ ปุคฺคโล ปุพฺเพ อนนุสฺสุเตสุ ธมฺเมสุ สามํ สจฺจานิ อภิสมฺ\-พุชฺฌติ, ตตฺถ จ สพฺพญฺญุตํ ปาปุณาติ พเลสุ\csromansharedfootnote{cd7a9932c7b8c069}{ผเลสุ (อิ)} จ วสีภาวํ. อตฺเถกจฺโจ ปุคฺคโล ปุพฺเพ อนนุสฺสุเตสุ ธมฺเมสุ สามํ สจฺจานิ อภิสมฺ\-พุชฺฌติ, น จ ตตฺถ สพฺพญฺญุตํ ปาปุณาติ น จ พเลสุ วสีภาวํ. อตฺเถกจฺโจ ปุคฺคโล ปุพฺเพ อนนุสฺสุเตสุ ธมฺเมสุ สามํ สจฺจานิ อนภิสมฺพุชฺฌติ, ทิฏฺเฐ เจว ธมฺเม ทุกฺขสฺส\-นฺตกโร โหติ, สาวกปารมิญฺจ ปาปุณาติ. อตฺเถกจฺโจ ปุคฺคโล ปุพฺเพ อนนุสฺสุเตสุ ธมฺเมสุ สามํ สจฺจานิ อนภิสมฺพุชฺฌติ, ทิฏฺเฐว ธมฺเม ทุกฺขสฺส\-นฺตกโร โหติ, น จ สาวกปารมิํ ปาปุณาติ. อตฺเถกจฺโจ ปุคฺคโล ปุพฺเพ อนนุสฺสุเตสุ ธมฺเมสุ สามํ สจฺจานิ อนภิสมฺพุชฺฌติ, น จ ทิฏฺเฐว ธมฺเม ทุกฺขสฺส\-นฺตกโร โหติ, อนาคามี โหติ อนาคนฺตา\csromansharedfootnote{6d67a135ee0569bf}{อนาคนฺตฺวา (สฺยา, ก) อํ 1. 478 ปิฏฺเฐปิ.} อิตฺถตฺตํ. อตฺเถกจฺโจ ปุคฺคโล ปุพฺเพ อนนุสฺสุเตสุ ธมฺเมสุ สามํ สจฺจานิ อนภิสมฺพุชฺฌติ, น จ ทิฏฺเฐว ธมฺเม ทุกฺขสฺส\-นฺตกโร โหติ, อาคามี\csromansharedfootnote{d170ca510c190c75}{โสตาปนฺนสกทาคามี (สฺยา, ก)} โหติ อาคนฺตา อิตฺถตฺตํ.}

\vspace{\tipitakaparskip}
\csromancenter{ฉกฺกํ.}
\csromanmatikaanchor{mtk.65}
\csromantocmarkref{section}{7. สตฺตกอุทฺเทส}{mtk.65}
\bookmark[dest={mtk.65},level=1]{7. สตฺตกอุทฺเทส}
\csromanheader{1}{7. สตฺตกอุทฺเทส}
\prose{13. \textbf{สตฺต ปุคฺคลา\csromandash{}}}

\proseitem{1}{สตฺต อุทกูปมา ปุคฺคลา. สกิํ นิมุคฺโค นิมุคฺโคว โหติ, อุมฺมุชฺชิตฺวา นิมุชฺชติ, อุมฺมุชฺชิตฺวา ฐิโต โหติ, อุมฺมุชฺชิตฺวา วิปสฺสติ วิโลเกติ, อุมฺมุชฺชิตฺวา ปตรติ, อุมฺมุชฺชิตฺวา ปฏิคาธปฺปตฺโต โหติ, อุมฺมุชฺชิตฺวา ติณฺโณ โหติ ปารงฺคโต ถเล ติฏฺฐติ พฺราหฺมโณ.}

\proseitem{2}{\csromanfolio{113}อุภโตภาค\-วิมุตฺโต, ปญฺญาวิมุตฺโต, กายสกฺขี, ทิฏฺฐิปฺปตฺโต, สทฺธาวิมุตฺโต, ธมฺมานุสารี, สทฺธานุสารี.}

\vspace{\tipitakaparskip}
\csromancenter{สตฺตกํ.}
\csromanmatikaanchor{mtk.66}
\csromantocmarkref{section}{8. อฏฺฐกอุทฺเทส}{mtk.66}
\bookmark[dest={mtk.66},level=1]{8. อฏฺฐกอุทฺเทส}
\csromanheader{1}{8. อฏฺฐกอุทฺเทส}
\prose{14. \textbf{อฏฺฐ ปุคฺคลา\csromandash{}}}

\proseitem{1}{จตฺตาโร มคฺคสมงฺคิโน, จตฺตาโร ผลสมงฺคิโน ปุคฺคลา.}

\vspace{\tipitakaparskip}
\csromancenter{อฏฺฐกํ.}
\csromanmatikaanchor{mtk.67}
\csromantocmarkref{section}{9. นวกอุทฺเทส}{mtk.67}
\bookmark[dest={mtk.67},level=1]{9. นวกอุทฺเทส}
\csromanheader{1}{9. นวกอุทฺเทส}
\prose{15. \textbf{นว ปุคฺคลา\csromandash{}}}

\proseitem{1}{สมฺมาสมฺพุทฺโธ, ปจฺเจก\-สมฺพุทฺโธ, อุภโตภาค\-วิมุตฺโต, ปญฺญาวิมุตฺโต, กายสกฺขี, ทิฏฺฐิปฺปตฺโต, สทฺธาวิมุตฺโต, ธมฺมานุสารี, สทฺธานุสารี.}

\vspace{\tipitakaparskip}
\csromancenter{นวกํ.}
\csromanmatikaanchor{mtk.68}
\csromantocmarkref{section}{10. ทสกอุทฺเทส}{mtk.68}
\bookmark[dest={mtk.68},level=1]{10. ทสกอุทฺเทส}
\csromanheader{1}{10. ทสกอุทฺเทส}
\prose{16. \textbf{ทส ปุคฺคลา\csromandash{}}}

\proseitem{1}{ปญฺจนฺนํ อิธ นิฏฺฐา, ปญฺจนฺนํ อิธ วิหาย นิฏฺฐา.}

\vspace{\tipitakaparskip}
\csromancenter{ทสกํ.}
\nitthitam{ปุคฺคลปญฺญตฺติ\-มาติกา นิฏฺฐิตา.}
\csromanmatikaanchor{mtk.69}
\csromantocmarknopage{section}{นิทฺเทส}{mtk.69}
\bookmark[dest={mtk.69},level=1]{นิทฺเทส}
\csromanheader{1}{\textbf{นิทฺเทส}}
\csromanmatikaanchor{mtk.70}
\csromantocmarkref{subsection}{1. เอกกปุคฺคลปญฺญตฺติ}{mtk.70}
\bookmark[dest={mtk.70},level=2]{1. เอกกปุคฺคลปญฺญตฺติ}
\csromanheader{2}{1. เอกก\-ปุคฺคล\-ปญฺญตฺติ}
\proseitem{1}{\csromanfolio{114}กตโม จ ปุคฺคโล สมยวิมุตฺโต, อิเธกจฺโจ ปุคฺคโล กาเลน กาลํ สมเยน สมยํ อฏฺฐ วิโมกฺเข กาเยน ผุสิตฺวา\csromansharedfootnote{e87c423abe96dce4}{ผสฺสิตฺวา (สี, อิ)} วิหรติ, ปญฺญาย จสฺส ทิสฺวา เอกจฺเจ อาสวา ปริกฺขีณา โหนฺติ. อยํ วุจฺจติ ปุคฺคโล สมยวิมุตฺโต.}

\proseitem{2}{กตโม จ ปุคฺคโล อสมยวิมุตฺโต, อิเธกจฺโจ ปุคฺคโล น เหว โข กาเลน กาลํ สมเยน สมยํ อฏฺฐ วิโมกฺเข กาเยน ผุสิตฺวา วิหรติ, ปญฺญาย จสฺส ทิสฺวา อาสวา ปริกฺขีณา โหนฺติ. อยํ วุจฺจติ ปุคฺคโล อสมยวิมุตฺโต. สพฺเพปิ อริยปุคฺคลา อริเย วิโมกฺเข อสมยวิมุตฺตา.}

\proseitem{3}{กตโม จ ปุคฺคโล กุปฺปธมฺโม, อิเธกจฺโจ ปุคฺคโล ลาภี โหติ รูปสหคตานํ วา อรูป\-สหคตานํ วา สมาปตฺตีนํ, โส จ โข น นิกามลาภี โหติ น อกิจฺฉลาภี น อกสิรลาภี, น ยตฺถิจฺฉกํ ยทิจฺฉกํ ยาวติจฺฉกํ สมาปชฺชติปิ วุฏฺฐาติปิ. ฐานํ โข ปเนตํ วิชฺชติ, ยํ ตสฺส ปุคฺคลสฺส ปมาทมาคมฺม ตา สมาปตฺติโย กุปฺเปยฺยุํ. อยํ วุจฺจติ ปุคฺคโล กุปฺปธมฺโม.}

\proseitem{4}{กตโม จ ปุคฺคโล อกุปฺปธมฺโม, อิเธกจฺโจ ปุคฺคโล ลาภี โหติ รูปสหคตานํ วา อรูป\-สหคตานํ วา สมาปตฺตีนํ, โส จ โข นิกามลาภี โหติ อกิจฺฉลาภี อกสิรลาภี, ยตฺถิจฺฉกํ ยทิจฺฉกํ ยาวติจฺฉกํ สมาปชฺชติปิ วุฏฺฐาติปิ. อฏฺฐานเมตํ อนวกาโส, ยํ ตสฺส ปุคฺคลสฺส ปมาทมาคมฺม ตา สมาปตฺติโย กุปฺเปยฺยุํ. อยํ วุจฺจติ ปุคฺคโล อกุปฺปธมฺโม. สพฺเพปิ อริยปุคฺคลา อริเย วิโมกฺเข อกุปฺปธมฺมา.}

\proseitem{5}{กตโม จ ปุคฺคโล ปริหานธมฺโม, อิเธกจฺโจ ปุคฺคโล ลาภี โหติ รูปสหคตานํ วา อรูป\-สหคตานํ วา สมาปตฺตีนํ, โส จ โข น นิกามลาภี โหติ น อกิจฺฉลาภี น อกสิรลาภี,}

\prose{\csromanfolio{115}น ยตฺถิจฺฉกํ ยทิจฺฉกํ ยาวติจฺฉกํ สมาปชฺชติปิ วุฏฺฐาติปิ. ฐานํ โข ปเนตํ วิชฺชติ, ยํ โส ปุคฺคโล ปมาทมาคมฺม ตาหิ สมาปตฺตีหิ ปริหาเยยฺย. อยํ วุจฺจติ ปุคฺคโล ปริหานธมฺโม.}

\proseitem{6}{กตโม จ ปุคฺคโล อปริหาน\-ธมฺโม, อิเธกจฺโจ ปุคฺคโล ลาภี โหติ รูปสหคตานํ วา อรูป\-สหคตานํ วา สมาปฺปตฺตีนํ, โส จ โข นิกามลาภี โหติ อกิจฺฉลาภี อกสิรลาภี, ยตฺถิจฺฉกํ ยทิจฺฉกํ ยาวติจฺฉกํ สมาปชฺชติปิ วุฏฺฐาติปิ. อฏฺฐานเมตํ อนวกาโส, ยํ โส ปุคฺคโล ปมาทมาคมฺม ตาหิ สมาปตฺตีหิ ปริหาเยยฺย. อยํ วุจฺจติ ปุคฺคโล อปริหาน\-ธมฺโม. สพฺเพปิ อริยปุคฺคลา อริเย วิโมกฺเข อปริหาน\-ธมฺมา.}

\proseitem{7}{กตโม จ ปุคฺคโล เจตนาภพฺโพ, อิเธกจฺโจ ปุคฺคโล ลาภี โหติ รูปสหคตานํ วา อรูป\-สหคตานํ วา สมาปตฺตีนํ, โส จ โข น นิกามลาภี โหติ น อกิจฺฉลาภี น อกสิรลาภี, น ยตฺถิจฺฉกํ ยทิจฺฉกํ ยาวติจฺฉกํ สมาปชฺชติปิ วุฏฺฐาติปิ. สเจ อนุสญฺเจเตติ, น ปริหายติ ตาหิ สมาปตฺตีหิ. สเจ น อนุสญฺเจเตติ, ปริหายติ ตาหิ สมาปตฺตีหิ. อยํ วุจฺจติ ปุคฺคโล เจตนาภพฺโพ.}

\proseitem{8}{กตโม จ ปุคฺคโล อนุรกฺขณา\-ภพฺโพ, อิเธกจฺโจ ปุคฺคโล ลาภี โหติ รูปสหคตานํ วา อรูป\-สหคตานํ วา สมาปตฺตีนํ, โส จ โข น นิกามลาภี โหติ น อกิจฺฉลาภี น อกสิรลาภี, น ยตฺถิจฺฉกํ ยทิจฺฉกํ ยาวติจฺฉกํ สมาปชฺชติปิ วุฏฺฐาติปิ. สเจ อนุรกฺขติ, น ปริหายติ ตาหิ สมาปตฺตีหิ. สเจ น อนุรกฺขติ, ปริหายติ ตาหิ สมาปตฺตีหิ. อยํ วุจฺจติ ปุคฺคโล อนุรกฺขณา\-ภพฺโพ.}

\proseitem{9}{กตโม จ ปุคฺคโล ปุถุชฺชโน, ยสฺส ปุคฺคลสฺส ตีณิ สํโยชนานิ อปฺปหีนานิ, น จ เตสํ ธมฺมานํ ปหานาย ปฏิปนฺโน. อยํ วุจฺจติ ปุคฺคโล ปุถุชฺชโน.}

\proseitem{10}{กตโม จ ปุคฺคโล โคตฺรภู, เยสํ ธมฺมานํ สมนนฺตรา อริยธมฺมสฺส อวกฺกนฺติ โหติ เตหิ ธมฺเมหิ สมนฺนาคโต. อยํ วุจฺจติ ปุคฺคโล โคตฺรภู.}

\proseitem{11}{\csromanfolio{116}กตโม จ ปุคฺคโล ภยูปรโต, สตฺต เสกฺขา ภยูปรตา, เย จ ปุถุชฺชนา สีลวนฺโต. อรหา อภยูปรโต.}

\proseitem{12}{กตโม จ ปุคฺคโล อภพฺพาคมโน, เย เต ปุคฺคลา กมฺมาวรเณน สมนฺนาคตา กิเลสาวรเณน สมนฺนาคตา วิปากาวรเณน สมนฺนาคตา อสฺสทฺธา อจฺฉนฺทิกา ทุปฺปญฺญา เอฬา อภพฺพา นิยามํ โอกฺกมิตุํ กุสเลสุ ธมฺเมสุ สมฺมตฺตํ. อิเม วุจฺจนฺติ ปุคฺคลา อภพฺพาคมนา.}

\proseitem{13}{กตโม จ ปุคฺคโล ภพฺพาคมโน, เย เต ปุคฺคลา น กมฺมาวรเณน สมนฺนาคตา น กิเลสาวรเณน สมนฺนาคตา น วิปากาวรเณน สมนฺนาคตา สทฺธา ฉนฺทิกา ปญฺญวนฺโต\csromansharedfootnote{818497e97af3f10c}{ปญฺญวนฺตา (สี)} อเนฬา ภพฺพา นิยามํ โอกฺกมิตุํ กุสเลสุ ธมฺเมสุ สมฺมตฺตํ. อิเม วุจฺจนฺติ ปุคฺคลา ภพฺพาคมนา.}

\proseitem{14}{กตโม จ ปุคฺคโล นิยโต, ปญฺจ ปุคฺคลา อานนฺตริกา, เย จ มิจฺฉาทิฏฺฐิกา นิยตา, อฏฺฐ จ อริยปุคฺคลา นิยตา. อวเสสา ปุคฺคลา อนิยตา.}

\proseitem{15}{กตโม จ ปุคฺคโล ปฏิปนฺนโก, จตฺตาโร มคฺคสมงฺคิโน ปุคฺคลา ปฏิปนฺนกา, จตฺตาโร ผลสมงฺคิโน ปุคฺคลา ผเล ฐิตา.}

\proseitem{16}{กตโม จ ปุคฺคโล สมสีสี, ยสฺส ปุคฺคลสฺส อปุพฺพํ อจริมํ อาสวปริยาทานญฺจ โหติ ชีวิตปริยาทานญฺจ. อยํ วุจฺจติ ปุคฺคโล สมสีสี.}

\proseitem{17}{กตโม จ ปุคฺคโล ฐิตกปฺปี, อยญฺจ ปุคฺคโล โสตาปตฺติผล\-สจฺฉิกิริยาย ปฏิปนฺโน อสฺส, กปฺปสฺส จ อุฑฺฑยฺหนเวลา อสฺส, เนว ตาว กปฺโป อุฑฺฑเยฺหยฺย, ยาวายํ ปุคฺคโล น โสตาปตฺติผลํ สจฺฉิกโรติ. อยํ วุจฺจติ ปุคฺคโล ฐิตกปฺปี. สพฺเพปิ มคฺคสมงฺคิโน ปุคฺคลา ฐิตกปฺปิโน.}

\proseitem{18}{กตโม จ ปุคฺคโล อริโย, อฏฺฐ อริยปุคฺคลา อริยา. อวเสสา ปุคฺคลา อนริยา.}

\proseitem{19}{\csromanfolio{117}กตโม จ ปุคฺคโล เสกฺโข, จตฺตาโร มคฺคสมงฺคิโน ตโย ผลสมงฺคิโน ปุคฺคลา เสกฺขา. อรหา อเสกฺโข. อวเสสา ปุคฺคลา เนว\-เสกฺข\-นาเสกฺขา.}

\proseitem{20}{กตโม จ ปุคฺคโล เตวิชฺโช, ตีหิ วิชฺชาหิ สมนฺนาคโต ปุคฺคโล เตวิชฺโช.}

\proseitem{21}{กตโม จ ปุคฺคโล ฉฬภิญฺโญ, ฉหิ อภิญฺญาหิ สมนฺนาคโต ปุคฺคโล ฉฬภิญฺโญ.}

\proseitem{22}{กตโม จ ปุคฺคโล สมฺมาสมฺพุทฺโธ, อิเธกจฺโจ ปุคฺคโล ปุพฺเพ อนนุสฺสุเตสุ ธมฺเมสุ สามํ สจฺจานิ อภิสมฺ\-พุชฺฌติ, ตตฺถ จ สพฺพญฺญุตํ ปาปุณาติ พเลสุ จ วสีภาวํ. อยํ วุจฺจติ ปุคฺคโล สมฺมาสมฺพุทฺโธ.}

\proseitem{23}{กตโม จ ปุคฺคโล ปจฺเจก\-สมฺพุทฺโธ, อิเธกจฺโจ ปุคฺคโล ปุพฺเพ อนนุสฺสุเตสุ ธมฺเมสุ สามํ สจฺจานิ อภิสมฺ\-พุชฺฌติ, น จ ตตฺถ สพฺพญฺญุตํ ปาปุณาติ น จ พเลสุ วสีภาวํ. อยํ วุจฺจติ ปุคฺคโล ปจฺเจก\-สมฺพุทฺโธ.}

\proseitem{24}{กตโม จ ปุคฺคโล อุภโตภาค\-วิมุตฺโต, อิเธกจฺโจ ปุคฺคโล อฏฺฐ วิโมกฺเข กาเยน ผุสิตฺวา วิหรติ, ปญฺญาย จสฺส ทิสฺวา อาสวา ปริกฺขีณา โหนฺติ. อยํ วุจฺจติ ปุคฺคโล อุภโตภาค\-วิมุตฺโต.}

\proseitem{25}{กตโม จ ปุคฺคโล ปญฺญาวิมุตฺโต, อิเธกจฺโจ ปุคฺคโล น เหว โข อฏฺฐ วิโมกฺเข กาเยน ผุสิตฺวา วิหรติ, ปญฺญาย จสฺส ทิสฺวา อาสวา ปริกฺขีณา โหนฺติ. อยํ วุจฺจติ ปุคฺคโล ปญฺญาวิมุตฺโต.}

\proseitem{26}{กตโม จ ปุคฺคโล กายสกฺขี, อิเธกจฺโจ ปุคฺคโล อฏฺฐ วิโมกฺเข กาเยน ผุสิตฺวา วิหรติ, ปญฺญาย จสฺส ทิสฺวา เอกจฺเจ อาสวา ปริกฺขีณา โหนฺติ. อยํ วุจฺจติ ปุคฺคโล กายสกฺขี.}

\proseitem{27}{กตโม จ ปุคฺคโล ทิฏฺฐิปฺปตฺโต, อิเธกจฺโจ ปุคฺคโล “อิทํ ทุกฺขนฺ”ติ ยถาภูตํ ปชานาติ, “อยํ ทุกฺขสมุทโย”ติ ยถาภูตํ ปชานาติ, “อยํ ทุกฺขนิโรโธ”ติ ยถาภูตํ ปชานาติ, “อยํ ทุกฺขนิโรธ\-คามินี ปฏิปทา”ติ ยถาภูตํ ปชานาติ, ตถาคต\-ปฺปเวทิตา \csromanfolio{118}จสฺส ธมฺมา ปญฺญาย โวทิฏฺฐา โหนฺติ โวจริตา, ปญฺญาย จสฺส ทิสฺวา เอกจฺเจ อาสวา ปริกฺขีณา โหนฺติ. อยํ วุจฺจติ ปุคฺคโล ทิฏฺฐิปฺปตฺโต.}

\proseitem{28}{กตโม จ ปุคฺคโล สทฺธาวิมุตฺโต, อิเธกจฺโจ ปุคฺคโล “อิทํ ทุกฺขนฺ”ติ ยถาภูตํ ปชานาติ, “อยํ ทุกฺขสมุทโย”ติ ยถาภูตํ ปชานาติ, “อยํ ทุกฺขนิโรโธ”ติ ยถาภูตํ ปชานาติ, “อยํ ทุกฺขนิโรธ\-คามินี ปฏิปทา”ติ ยถาภูตํ ปชานาติ, ตถาคต\-ปฺปเวทิตา จสฺส ธมฺมา ปญฺญาย โวทิฏฺฐา โหนฺติ โวจริตา, ปญฺญาย จสฺส ทิสฺวา เอกจฺเจ อาสวา ปริกฺขีณา โหนฺติ, โน จ โข ยถา ทิฏฺฐิ\-ปฺปตฺตสฺส. อยํ วุจฺจติ ปุคฺคโล สทฺธาวิมุตฺโต.}

\proseitem{29}{กตโม จ ปุคฺคโล ธมฺมานุสารี, ยสฺส ปุคฺคลสฺส โสตาปตฺติผล\-สจฺฉิกิริยาย ปฏิปนฺนสฺส ปญฺญินฺทฺริยํ อธิมตฺตํ โหติ, ปญฺญาวาหิํ ปญฺญา\-ปุพฺพงฺคมํ อริยมคฺคํ ภาเวติ. อยํ วุจฺจติ ปุคฺคโล ธมฺมานุสารี. โสตาปตฺติผล\-สจฺฉิกิริยาย ปฏิปนฺโน ปุคฺคโล ธมฺมานุสารี, ผเล ฐิโต ทิฏฺฐิปฺปตฺโต.}

\proseitem{30}{กตโม จ ปุคฺคโล สทฺธานุสารี, ยสฺส ปุคฺคลสฺส โสตาปตฺติผล\-สจฺฉิกิริยาย ปฏิปนฺนสฺส สทฺธินฺทฺริยํ อธิมตฺตํ โหติ, สทฺธาวาหิํ สทฺธา\-ปุพฺพงฺคมํ อริยมคฺคํ ภาเวติ, อยํ วุจฺจติ ปุคฺคโล สทฺธานุสารี. โสตาปตฺติผล\-สจฺฉิกิริยาย ปฏิปนฺโน ปุคฺคโล สทฺธานุสารี, ผเล ฐิโต สทฺธาวิมุตฺโต.}

\proseitem{31}{กตโม จ ปุคฺคโล สตฺตกฺขตฺตุ\-ปรโม อิเธกจฺโจ ปุคฺคโล ติณฺณํ สํโยชนานํ ปริกฺขยา โสตาปนฺโน โหติ อวินิปาต\-ธมฺโม นิยโต สมฺโพธิ\-ปรายโน\csromansharedfootnote{042d9046bc623775}{สมฺโพธิปรายโณ (สี, ก)}, โส สตฺตกฺขตฺตุํ เทเว จ มานุเส จ สนฺธาวิตฺวา สํสริตฺวา ทุกฺขสฺสนฺตํ กโรติ. อยํ วุจฺจติ ปุคฺคโล สตฺตกฺขตฺตุ\-ปรโม.}

\proseitem{32}{กตโม จ ปุคฺคโล โกลํโกโล, อิเธกจฺโจ ปุคฺคโล ติณฺณํ สํโยชนานํ ปริกฺขยา โสตาปนฺโน โหติ อวินิปาต\-ธมฺโม}

\prose{\csromanfolio{119}นิยโต สมฺโพธิ\-ปรายโน, โส เทฺว วา ตีณิ วา กุลานิ สนฺธาวิตฺวา สํสริตฺวา ทุกฺขสฺสนฺตํ กโรติ. อยํ วุจฺจติ ปุคฺคโล โกลํโกโล.}

\proseitem{33}{กตโม จ ปุคฺคโล เอกพีชี, อิเธกจฺโจ ปุคฺคโล ติณฺณํ สํโยชนานํ ปริกฺขยา โสตาปนฺโน โหติ อวินิปาต\-ธมฺโม นิยโต สมฺโพธิ\-ปรายโน, โส เอกํเยว มานุสกํ ภวํ นิพฺพตฺเตตฺวา ทุกฺขสฺสนฺตํ กโรติ. อยํ วุจฺจติ ปุคฺคโล เอกพีชี.}

\proseitem{34}{กตโม จ ปุคฺคโล สกทาคามี, อิเธกจฺโจ ปุคฺคโล ติณฺณํ สํโยชนานํ ปริกฺขยา ราค\-โทส\-โมหานํ ตนุตฺตา สกทาคามี โหติ สกิเทว อิมํ โลกํ อาคนฺตฺวา ทุกฺขสฺสนฺตํ กโรติ. อยํ วุจฺจติ ปุคฺคโล สกทาคามี.}

\proseitem{35}{กตโม จ ปุคฺคโล อนาคามี, อิเธกจฺโจ ปุคฺคโล ปญฺจนฺนํ โอรมฺภาคิยานํ สํโยชนานํ ปริกฺขยา โอปปาติโก โหติ, ตตฺถ ปรินิพฺพายี อนาวตฺติธมฺโม ตสฺมา โลกา. อยํ วุจฺจติ ปุคฺคโล อนาคามี.}

\proseitem{36}{กตโม จ ปุคฺคโล อนฺตรา\-ปรินิพฺพายี, อิเธกจฺโจ ปุคฺคโล ปญฺจนฺนํ โอรมฺภาคิยานํ สํโยชนานํ ปริกฺขยา โอปปาติโก โหติ, ตตฺถ ปรินิพฺพายี อนาวตฺติธมฺโม ตสฺมา โลกา, โส อุปปนฺนํ วา สมนนฺตรา อปฺปตฺตํ วา เวมชฺฌํ อายุปฺปมาณํ อริยมคฺคํ สญฺชเนติ อุปริฏฺฐิมานํ สํโยชนานํ ปหานาย. อยํ วุจฺจติ ปุคฺคโล อนฺตรา\-ปรินิพฺพายี.}

\proseitem{37}{กตโม จ ปุคฺคโล อุปหจฺจ\-ปรินิพฺพายี, อิเธกจฺโจ ปุคฺคโล ปญฺจนฺนํ โอรมฺภาคิยานํ สํโยชนานํ ปริกฺขยา โอปปาติโก โหติ, ตตฺถ ปรินิพฺพายี อนาวตฺติธมฺโม ตสฺมา โลกา, โส อติกฺกมิตฺวา เวมชฺฌํ อายุปฺปมาณํ อุปหจฺจ วา กาลกิริยํ\csromansharedfootnote{c6c73a3cfac16723}{กาลํ กิริยํ (ก)} อริยมคฺคํ สญฺชเนติ อุปริฏฺฐิมานํ สํโยชนานํ ปหานาย. อยํ วุจฺจติ ปุคฺคโล อุปหจฺจ\-ปรินิพฺพายี.}

\proseitem{38}{กตโม จ ปุคฺคโล อสงฺขารปรินิพฺพายี, อิเธกจฺโจ ปุคฺคโล ปญฺจนฺนํ โอรมฺภาคิยานํ สํโยชนานํ ปริกฺขยา โอปปาติโก โหติ, ตตฺถ \csromanfolio{120}ปรินิพฺพายี อนาวตฺติธมฺโม ตสฺมา โลกา, โส อสงฺขาเรน อริยมคฺคํ สญฺชเนติ อุปริฏฺฐิมานํ สํโยชนานํ ปหานาย. อยํ วุจฺจติ ปุคฺคโล อสงฺขารปรินิพฺพายี.}

\proseitem{39}{กตโม จ ปุคฺคโล สสงฺขาร\-ปรินิพฺพายี, อิเธกจฺโจ ปุคฺคโล ปญฺจนฺนํ โอรมฺภาคิยานํ สํโยชนานํ ปริกฺขยา โอปปาติโก โหติ, ตตฺถ ปรินิพฺพายี อนาวตฺติธมฺโม ตสฺมา โลกา, โส สสงฺขาเรน อริยมคฺคํ สญฺชเนติ อุปริฏฺฐิมานํ สํโยชนานํ ปหานาย. อยํ วุจฺจติ ปุคฺคโล สสงฺขาร\-ปรินิพฺพายี.}

\proseitem{40}{กตโม จ ปุคฺคโล อุทฺธํโสโต อกนิฏฺฐคามี, อิเธกจฺโจ ปุคฺคโล ปญฺจนฺนํ โอรมฺภาคิยานํ สํโยชนานํ ปริกฺขยา โอปปาติโก โหติ, ตตฺถ ปรินิพฺพายี อนาวตฺติธมฺโม ตสฺมา โลกา, โส อวิหา จุโต อตปฺปํ คจฺฉติ, อตปฺปา จุโต สุทสฺสํ คจฺฉติ, สุทสฺสา จุโต สุทสฺสิํ คจฺฉติ, สุทสฺสิยา จุโต อกนิฏฺฐํ คจฺฉติ, อกนิฏฺเฐ อริยมคฺคํ สญฺชเนติ อุปริฏฺฐิมานํ สํโยชนานํ ปหานาย. อยํ วุจฺจติ ปุคฺคโล อุทฺธํโสโต อกนิฏฺฐคามี.}

\proseitem{41}{กตโม จ ปุคฺคโล โสตาปนฺโน, โสตาปตฺติผล\-สจฺฉิกิริยาย ปฏิปนฺโน, ติณฺณํ สํโยชนานํ ปหานาย ปฏิปนฺโน ปุคฺคโล โสตาปตฺติผล\-สจฺฉิกิริยาย ปฏิปนฺโน, ยสฺส ปุคฺคลสฺส ตีณิ สํโยชนานิ ปหีนานิ. อยํ วุจฺจติ ปุคฺคโล โสตาปนฺโน.}

\proseitem{42}{กามราค\-พฺยาปาทานํ ตนุภาวาย ปฏิปนฺโน ปุคฺคโล สกทาคามิ\-ผล\-สจฺฉิกิริยาย ปฏิปนฺโน, ยสฺส ปุคฺคลสฺส กามราค\-พฺยาปาทา ตนุภูตา. อยํ วุจฺจติ ปุคฺคโล สกทาคามี.}

\proseitem{43}{กามราค\-พฺยาปาทานํ อนวเสส\-ปฺปหานาย ปฏิปนฺโน ปุคฺคโล อนาคามิ\-ผล\-สจฺฉิกิริยาย ปฏิปนฺโน, ยสฺส ปุคฺคลสฺส กามราค\-พฺยาปาทา อนวเสสา ปหีนา. อยํ วุจฺจติ ปุคฺคโล อนาคามี.}

\proseitem{44}{รูปราคอรูปราคมานอุทฺธจฺจอวิชฺชาย อนวเสส\-ปฺปหานาย ปฏิปนฺโน ปุคฺคโล อรหตฺตผล\-สจฺฉิกิริยาย ปฏิปนฺโน, ยสฺส ปุคฺคลสฺส}

\prose{\csromanfolio{121}รูปราโค อรูปราโค มาโน อุทฺธจฺจํ อวิชฺชา อนวเสสา ปหีนา. อยํ วุจฺจติ ปุคฺคโล อรหา.}

\vspace{\tipitakaparskip}
\csromancenter{เอกกนิทฺเทโส.}
\csromanmatikaanchor{mtk.71}
\csromantocmarkref{subsection}{2. ทุกปุคฺคลปญฺญตฺติ}{mtk.71}
\bookmark[dest={mtk.71},level=2]{2. ทุกปุคฺคลปญฺญตฺติ}
\csromanheader{2}{2. ทุกปุคฺคล\-ปญฺญตฺติ}
\proseitem{45}{กตโม จ ปุคฺคโล โกธโน, ตตฺถ กตโม โกโธ, โย โกโธ กุชฺฌนา กุชฺฌิตตฺตํ โทโส ทุสฺสนา ทุสฺสิตตฺตํ\csromansharedfootnote{deaf2962c9ae776e}{ทูสนา ทูสิตตฺตํ (สฺยา)} พฺยาปตฺติ พฺยาปชฺชนา พฺยาปชฺชิตตฺตํ วิโรโธ ปฏิวิโรโธ จณฺฑิกฺกํ อสุโรโป อนตฺตมนตา จิตฺตสฺส, อยํ วุจฺจติ โกโธ. ยสฺส ปุคฺคลสฺส อยํ โกโธ อปฺปหีโน, อยํ วุจฺจติ ปุคฺคโล โกธโน.}

\proseitem{46}{กตโม จ ปุคฺคโล อุปนาหี, ตตฺถ กตโม อุปนาโห, ปุพฺพกาลํ โกโธ, อปรกาลํ อุปนาโห, โย เอวรูโป อุปนาโห อุปนยฺหนา อุปนยฺหิตตฺตํ อฏฺฐปนา\csromansharedfootnote{42fdcbb5fde9b21d}{อาฐปนา (ก) อภิ 2. 371 ปิฏฺเฐปิ.} ฐปนา สณฺฐปนา อนุสํสนฺทนา อนุปฺปพนฺธนา ทฬฺหีกมฺมํ โกธสฺส, อยํ วุจฺจติ อุปนาโห. ยสฺส ปุคฺคลสฺส อยํ อุปนาโห อปฺปหีโน, อยํ วุจฺจติ ปุคฺคโล อุปนาหี.}

\proseitem{47}{กตโม จ ปุคฺคโล มกฺขี, ตตฺถ กตโม มกฺโข, โย มกฺโข มกฺขายนา มกฺขายิตตฺตํ\csromansharedfootnote{c66dc5098f00bb86}{มกฺขียนา มกฺขียิตตฺตํ (สี), มกฺขิยนา มกฺขิยิตตฺตํ (ก)} นิฏฺฐุริยํ นิฏฺฐุริย\-กมฺมํ, อยํ วุจฺจติ มกฺโข. ยสฺส ปุคฺคลสฺส อยํ มกฺโข อปฺปหีโน, อยํ วุจฺจติ ปุคฺคโล มกฺขี.}

\proseitem{48}{กตโม จ ปุคฺคโล ปฬาสี, ตตฺถ กตโม ปฬาโส, โย ปฬาโส ปฬาสายนา ปฬาสายิตตฺตํ ปฬาสาหาโร วิวาทฏฺฐานํ ยุคคฺคาโห อปฺปฏินิสฺสคฺโค, อยํ วุจฺจติ ปฬาโส. ยสฺส ปุคฺคลสฺส อยํ ปฬาโส อปฺปหีโน, อยํ วุจฺจติ ปุคฺคโล ปฬาสี.}

\proseitem{49}{\csromanfolio{122}กตโม จ ปุคฺคโล อิสฺสุกี, ตตฺถ กตมา อิสฺสา, ยา ปรลาภ\-สกฺการ\-ครุการ\-มานน\-วนฺทน\-ปูชนาสุ อิสฺสา อิสฺสายนา อิสฺสายิตตฺตํ อุสูยา อุสูยนา\csromansharedfootnote{a94fc5360ff7916d}{อุสฺสุยา อุสฺสุยนา (ก) อภิ 2. 372 ปิฏฺเฐปิ.} อุสูยิตตฺตํ, อยํ วุจฺจติ อิสฺสา. ยสฺส ปุคฺคลสฺส อยํ อิสฺสา อปฺปหีนา, อยํ วุจฺจติ ปุคฺคโล อิสฺสุกี.}

\proseitem{50}{กตโม จ ปุคฺคโล มจฺฉรี, ตตฺถ กตมํ มจฺฉริยํ, ปญฺจ มจฺฉริยานิ อาวาส\-มจฺฉริยํ กุลมจฺฉริยํ ลาภ\-มจฺฉริยํ วณฺณ\-มจฺฉริยํ ธมฺม\-มจฺฉริยํ, ยํ เอวรูปํ มจฺเฉรํ มจฺฉรายนา มจฺฉรายิ\-ตตฺตํ เววิจฺฉํ กทริยํ กฏุกญฺจุกตา อคฺคหิตตฺตํ จิตฺตสฺส, อิทํ วุจฺจติ มจฺฉริยํ. ยสฺส ปุคฺคลสฺส อิทํ มจฺฉริยํ อปฺปหีนํ, อยํ วุจฺจติ ปุคฺคโล มจฺฉรี.}

\proseitem{51}{กตโม จ ปุคฺคโล สโฐ, ตตฺถ กตมํ สาเฐยฺยํ, อิเธกจฺโจ สโฐ โหติ ปริสโฐ, ยํ ตตฺถ สฐํ สฐตา สาเฐยฺยํ กกฺกรตา กกฺกริยํ\csromansharedfootnote{f92332c964dcf50e}{กกฺขฬตา กกฺขฬิยํ (สฺยา) เอวํ ขุทฺทกวิภงฺคทุกนิทฺเทเสปิ.} ปริกฺขตฺตตา ปาริกฺขตฺติยํ, อิทํ วุจฺจติ สาเฐยฺยํ. ยสฺส ปุคฺคลสฺส อิทํ สาเฐยฺยํ อปฺปหีนํ, อยํ วุจฺจติ ปุคฺคโล สโฐ.}

\proseitem{52}{กตโม จ ปุคฺคโล มายาวี, ตตฺถ กตมา มายา, อิเธกจฺโจ กาเยน ทุจฺจริตํ จริตฺวา วาจาย ทุจฺจริตํ จริตฺวา มนสา ทุจฺจริตํ จริตฺวา ตสฺส ปฏิจฺฉาทน\-เหตุ ปาปิกํ อิจฺฉํ ปณิทหติ, “มา มํ ชญฺญา”ติ อิจฺฉติ, “มา มํ ชญฺญา”ติ สงฺกปฺปติ, “มา มํ ชญฺญา”ติ วาจํ ภาสติ, “มา มํ ชญฺญา”ติ กาเยน ปรกฺกมติ, ยา เอวรูปา มายา มายาวิตา อจฺจาสรา วญฺจนา นิกติ วิกิรณา ปริหรณา คูหนา ปริคูหนา ฉาทนา ปฏิจฺฉาทนา อนุตฺตานีกมฺมํ อนาวิกมฺมํ โวจฺฉาทนา ปาปกิริยา, อยํ วุจฺจติ มายา. ยสฺส ปุคฺคลสฺส อยํ มายา อปฺปหีนา, อยํ วุจฺจติ ปุคฺคโล มายาวี.}

\proseitem{53}{กตโม จ ปุคฺคโล อหิริโก, ตตฺถ กตมํ อหิริกํ, ยํ น หิรียติ หิริยิตพฺเพน, น หิรียติ ปาปกานํ อกุสลานํ ธมฺมานํ สมาปตฺติยา, อิทํ วุจฺจติ อหิริกํ. อิมินา อหิริเกน สมนฺนาคโต ปุคฺคโล อหิริโก.}

\proseitem{54}{\csromanfolio{123}กตโม จ ปุคฺคโล อโนตฺตปฺปี, ตตฺถ กตมํ อโนตฺตปฺปํ, ยํ น โอตฺตปฺปติ โอตฺตปฺปิตพฺเพน, น โอตฺตปฺปติ ปาปกานํ อกุสลานํ ธมฺมานํ สมาปตฺติยา, อิทํ วุจฺจติ อโนตฺตปฺปํ. อิมินา อโนตฺตปฺเปน สมนฺนาคโต ปุคฺคโล อโนตฺตปฺปี.}

\proseitem{55}{กตโม จ ปุคฺคโล ทุพฺพโจ, ตตฺถ กตมา โทวจสฺสตา, สหธมฺมิเก วุจฺจมาเน โทวจสฺสายํ โทวจสฺสิยํ โทวจสฺสตา วิปฺปฏิกุลคฺคาหิตา วิปจฺจนีก\-สาตตา อนาทริยํ อนาทริยตา อคารวตา อปฺปติสฺสวตา, อยํ วุจฺจติ โทวจสฺสตา. อิมาย โทวจสฺสตาย สมนฺนาคโต ปุคฺคโล ทุพฺพโจ.}

\proseitem{56}{กตโม จ ปุคฺคโล ปาปมิตฺโต, ตตฺถ กตมา ปาปมิตฺตตา, เย เต ปุคฺคลา อสฺสทฺธา ทุสฺสีลา อปฺปสฺสุตา มจฺฉริโน ทุปฺปญฺญา, ยา เตสํ เสวนา นิเสวนา สํเสวนา ภชนา สมฺภชนา ภตฺติ สมฺภตฺติ สมฺปวงฺกตา, อยํ วุจฺจติ ปาปมิตฺตตา. อิมาย ปาปมิตฺตตาย สมนฺนาคโต ปุคฺคโล ปาปมิตฺโต.}

\proseitem{57}{กตโม จ ปุคฺคโล อินฺทฺริเยสุ อคุตฺตทฺวาโร, ตตฺถ กตมา อินฺทฺริเยสุ อคุตฺตทฺวารตา, อิเธกจฺโจ ปุคฺคโล จกฺขุนา รูปํ ทิสฺวา นิมิตฺตคฺคาหี โหติ อนุพฺยญฺชนคฺคาหี, ยตฺวาธิกรณเมนํ จกฺขุนฺทฺริยํ อสํวุตํ วิหรนฺตํ อภิชฺฌา\-โทมนสฺสา ปาปกา อกุสลา ธมฺมา อนฺวาสฺสเวยฺยุํ, ตสฺส สํวราย น ปฏิปชฺชติ, น รกฺขติ จกฺขุนฺทฺริยํ, จกฺขุนฺทฺริเย น สํวรํ อาปชฺชติ. โสเตน สทฺทํ สุตฺวา ฯเปฯ. ฆาเนน คนฺธํ ฆายิตฺวา ฯเปฯ. ชิวฺหาย รสํ สายิตฺวา ฯเปฯ. กาเยน โผฏฺฐพฺพํ ผุสิตฺวา ฯเปฯ. มนสา ธมฺมํ วิญฺญาย นิมิตฺตคฺคาหี โหติ อนุพฺยญฺชนคฺคาหี, ยตฺวาธิกรณเมนํ มนินฺทฺริยํ อสํวุตํ วิหรนฺตํ อภิชฺฌา\-โทมนสฺสา ปาปกา อกุสลา ธมฺมา อนฺวาสฺสเวยฺยุํ, ตสฺส สํวราย น ปฏิปชฺชติ, น รกฺขติ มนินฺทฺริยํ, มนินฺทฺริเย น สํวรํ อาปชฺชติ, ยา อิเมสํ ฉนฺนํ อินฺทฺริยานํ อคุตฺติ อโคปนา อนารกฺโข อสํวโร, อยํ วุจฺจติ อินฺทฺริเยสุ อคุตฺตทฺวารตา. อิมาย อินฺทฺริเยสุ อคุตฺต\-ทฺวารตาย สมนฺนาคโต ปุคฺคโล อินฺทฺริเยสุ อคุตฺตทฺวาโร.}

\proseitem{58}{กตโม จ ปุคฺคโล โภชเน อมตฺตญฺญู, ตตฺถ กตมา โภชเน อมตฺตญฺญุตา, อิเธกจฺโจ ปุคฺคโล อปฺปฏิสงฺขา อโยนิโส \csromanfolio{124}อาหารํ อาหาเรติ ทวาย มทาย มณฺฑนาย วิภูสนาย, ยา ตตฺถ อสนฺตุฏฺฐิตา อมตฺตญฺญุตา อปฺปฏิสงฺขา โภชเน, อยํ วุจฺจติ โภชเน อมตฺตญฺญุตา. อิมาย โภชเน อมตฺตญฺญุตาย สมนฺนาคโต ปุคฺคโล โภชเน อมตฺตญฺญู.}

\proseitem{59}{กตโม จ ปุคฺคโล มุฏฺฐสฺสติ, ตตฺถ กตมํ มุฏฺฐสฺสจฺจํ, ยา อสฺสติ อนนุสฺสติ อปฺปฏิสฺสติ อสฺสติ อสฺสรณตา อธารณตา ปิลาปนตา สมฺมุสนตา, อิทํ วุจฺจติ มุฏฺฐสฺสจฺจํ. อิมินา มุฏฺฐสฺสจฺเจน สมนฺนาคโต ปุคฺคโล มุฏฺฐสฺสติ.}

\proseitem{60}{กตโม จ ปุคฺคโล อสมฺปชาโน, ตตฺถ กตมํ อสมฺปชญฺญํ, ยํ อญฺญาณํ อทสฺสนํ อนภิสมโย อนนุโพโธ อสมฺโพโธ อปฺปฏิเวโธ อสงฺคาหณา อปริโยคาหณา\csromansharedfootnote{ffc00c0827337e00}{อสํคาหนา อปริโยคาหนา (สี, สฺยา, ก)} อสมเปกฺขณา อปจฺจเวกฺขณา อปจฺจกฺขกมฺมํ ทุมฺเมชฺฌํ พาลฺยํ อสมฺปชญฺญํ โมโห ปโมโห สมฺโมโห อวิชฺชา อวิชฺโชโฆ อวิชฺชาโยโค อวิชฺชานุสโย อวิชฺชา\-ปริยุฏฺฐานํ อวิชฺชาลงฺคี โมโห อกุสลมูลํ, อิธํ วุจฺจติ อสมฺปชญฺญํ. อิมินา อสมฺปชญฺเญน สมนฺนาคโต ปุคฺคโล อสมฺปชาโน.}

\proseitem{61}{กตโม จ ปุคฺคโล สีลวิปนฺโน, ตตฺถ กตมา สีลวิปตฺติ, กายิโก วีติกฺกโม วาจสิโก วีติกฺกโม กายิกวาจสิโก วีติกฺกโม, อยํ วุจฺจติ สีลวิปตฺติ. สพฺพมฺปิ ทุสฺสิลฺยํ สีลวิปตฺติ. อิมาย สีลวิปตฺติยา สมนฺนาคโต ปุคฺคโล สีลวิปนฺโน.}

\proseitem{62}{กตโม จ ปุคฺคโล ทิฏฺฐิวิปนฺโน, ตตฺถ กตมา ทิฏฺฐิวิปตฺติ, “นตฺถิ ทินฺนํ, นตฺถิ ยิฏฺฐํ, นตฺถิ หุตํ, นตฺถิ สุกต\-ทุกฺกฏานํ\csromansharedfootnote{5116c917aea66194}{สุกฏทุกฺกฏานํ (สี)} กมฺมานํ ผลํ วิปาโก, นตฺถิ อยํ โลโก, นตฺถิ ปโร โลโก, นตฺถิ มาตา, นตฺถิ ปิตา, นตฺถิ สตฺตา โอปปาติกา, นตฺถิ โลเก สมณพฺราหฺมณา สมฺมคฺคตา\csromansharedfootnote{44029fb9de358ac7}{สมคฺคตา (ก)} สมฺมาปฏิปนฺนา, เย อิมญฺจ โลกํ ปรญฺจ โลกํ สยํ อภิญฺญา สจฺฉิกตฺวา ปเวเทนฺตี”ติ, ยา เอวรูปา ทิฏฺฐิ ทิฏฺฐิคตํ ทิฏฺฐิกนฺตาโร ทิฏฺฐิ\-วิสูกายิกํ ทิฏฺฐิ\-วิปฺผนฺทิตํ ทิฏฺฐิ\-สํโยชนํ คาโห ปฏิคฺคาโห อภินิเวโส ปรามาโส กุมฺมคฺโค มิจฺฉาปโถ มิจฺฉตฺตํ ติตฺถายตนํ วิปริยาส\-คฺคาโห\csromansharedfootnote{030e9d9622d56e71}{วิปริเยสคฺคาโห (สพฺพตฺถ) ปทสิทฺธิ จินฺเตตพฺพา.},}

\prose{\csromanfolio{125}อยํ วุจฺจติ ทิฏฺฐิวิปตฺติ. สพฺพาปิ มิจฺฉาทิฏฺฐิ ทิฏฺฐิวิปตฺติ. อิมาย ทิฏฺฐิ\-วิปตฺติยา สมนฺนาคโต ปุคฺคโล ทิฏฺฐิวิปนฺโน.}

\proseitem{63}{กตโม จ ปุคฺคโล อชฺฌตฺต\-สํโยชโน, ยสฺส ปุคฺคลสฺส ปญฺโจ\-รมฺภาคิยานิ สํโยชนานิ อปฺปหีนานิ, อยํ วุจฺจติ ปุคฺคโล อชฺฌตฺต\-สํโยชโน.}

\proseitem{64}{กตโม จ ปุคฺคโล พหิทฺธา\-สํโยชโน, ยสฺส ปุคฺคลสฺส ปญฺจุ\-ทฺธมฺภาคิยานิ สํโยชนานิ อปฺปหีนานิ, อยํ วุจฺจติ ปุคฺคโล พหิทฺธา\-สํโยชโน.}

\proseitem{65}{กตโม จ ปุคฺคโล อกฺโกธโน, ตตฺถ กตโม โกโธ, โย โกโธ กุชฺฌนา กุชฺฌิตตฺตํ โทโส ทุสฺสนา ทุสฺสิตตฺตํ พฺยาปตฺติ พฺยาปชฺชนา พฺยาปชฺชิตตฺตํ วิโรโธ ปฏิวิโรโธ จณฺฑิกฺกํ อสุโรโป อนตฺตมนตา จิตฺตสฺส, อยํ วุจฺจติ โกโธ. ยสฺส ปุคฺคลสฺส อยํ โกโธ ปหีโน, อยํ วุจฺจติ ปุคฺคโล อกฺโกธโน.}

\proseitem{66}{กตโม จ ปุคฺคโล อนุปนาหี, ตตฺถ กตโม อุปนาโห, ปุพฺพกาลํ โกโธ, อปรกาลํ อุปนาโห, โย เอวรูโป อุปนาโห อุปนยฺหนา อุปนยฺหิตตฺตํ อฏฺฐปนา ฐปนา สณฺฐปนา อนุสํสนฺทนา อนุปฺปพนฺธนา ทฬีกมฺมํ โกธสฺส, อยํ วุจฺจติ อุปนาโห. ยสฺส ปุคฺคลสฺส อยํ อุปนาโห ปหีโน, อยํ วุจฺจติ ปุคฺคโล อนุปนาหี.}

\proseitem{67}{กตโม จ ปุคฺคโล อมกฺขี, ตตฺถ กตโม มกฺโข, โย มกฺโข มกฺขายนา มกฺขายิตตฺตํ นิฏฺฐุริยํ นิฏฺฐุริย\-กมฺมํ, อยํ วุจฺจติ มกฺโข. ยสฺส ปุคฺคลสฺส อยํ มกฺโข ปหีโน, อยํ วุจฺจติ ปุคฺคโล อมกฺขี.}

\proseitem{68}{กตโม จ ปุคฺคโล อปฬาสี, ตตฺถ กตโม ปฬาโส, โย ปฬาโส ปฬาสายนา ปฬาสายิตตฺตํ ปฬาสาหาโร วิวาทฏฺฐานํ ยุคคฺคาโห อปฺปฏินิสฺสคฺโค, อยํ วุจฺจติ ปฬาโส. ยสฺส ปุคฺคลสฺส อยํ ปฬาโส ปหีโน, อยํ วุจฺจติ ปุคฺคโล อปฬาสี.}

\proseitem{69}{\csromanfolio{126}กตโม จ ปุคฺคโล อนิสฺสุกี, ตตฺถ กตมา อิสฺสา, ยา ปรลาภ\-สกฺการ\-ครุการ\-มานน\-วนฺทน\-ปูชนาสุ อิสฺสา อิสฺสายนา อิสฺสายิตตฺตํ อุสูยา อุสูยนา อุสูยิตตฺตํ, อยํ วุจฺจติ อิสฺสา. ยสฺส ปุคฺคลสฺส อยํ อิสฺสา ปหีนา, อยํ วุจฺจติ ปุคฺคโล อนิสฺสุกี.}

\proseitem{70}{กตโม จ ปุคฺคโล อมจฺฉรี, ตตฺถ กตมํ มจฺฉริยํ, ปญฺจ มจฺฉริยานิ อาวาส\-มจฺฉริยํ กุลมจฺฉริยํ ลาภ\-มจฺฉริยํ วณฺณ\-มจฺฉริยํ ธมฺม\-มจฺฉริยํ, ยํ เอวรูปํ มจฺเฉรํ มจฺฉรายนา มจฺฉรายิ\-ตตฺตํ เววิจฺฉํ กทริยํ กฏุกญฺจุกตา อคฺคหิตตฺตํ จิตฺตสฺส, อิทํ วุจฺจติ มจฺฉริยํ. ยสฺส ปุคลสฺส อิทํ มจฺฉริยํ ปหีนํ, อยํ วุจฺจติ ปุคฺคโล อมจฺฉรี.}

\proseitem{71}{กตโม จ ปุคฺคโล อสโฐ, ตตฺถ กตมํ สาเฐยฺยํ, อิเธกจฺโจ สโฐ โหติ ปริสโฐ, ยํ ตตฺถ สฐํ สฐตา สาเฐยฺยํ กกฺกรตา กกฺกริยํ ปริกฺขตฺตตา ปาริกฺขตฺติยํ, อิทํ วุจฺจติ สาเฐยฺยํ. ยสฺส ปุคฺคลสฺส อิทํ สาเฐยฺยํ ปหีนํ, อยํ วุจฺจติ ปุคฺคโล อสโฐ.}

\proseitem{72}{กตโม จ ปุคฺคโล อมายาวี, ตตฺถ กตมา มายา, อิเธกจฺโจ ปุคฺคโล กาเยน ทุจฺจริตํ จริตฺวา วาจาย ทุจฺจริตํ จริตฺวา มนสา ทุจฺจริตํ จริตฺวา ตสฺส ปฏิจฺฉาทน\-เหตุ ปาปิกํ อิจฺฉํ ปณิทหติ, “มา มํ ชญฺญา”ติ อิจฺฉติ, “มา มํ ชญฺญา”ติ สงฺกปฺปติ, “มา มํ ชญฺญา”ติ วาจํ ภาสติ, “มา มํ ชญฺญา”ติ กาเยน ปรกฺกมติ, ยา เอวรูปา มายา มายาวิตา อจฺจาสรา วญฺจนา นิกติ วิกิรณา ปริหรณา คูหนา ปริคูหนา ฉาทนา ปฏิจฺฉาทนา อนุตฺตานีกมฺมํ อนาวิกมฺมํ โวจฺฉาทนา ปาปกิริยา, อยํ วุจฺจติ มายา. ยสฺส ปุคฺคลสฺส อยํ มายา ปหีนา, อยํ วุจฺจติ ปุคฺคโล อมายาวี.}

\proseitem{73}{กตโม จ ปุคฺคโล หิริมา, ตตฺถ กตมา หิรี, ยํ หิรียติ หิริยิตพฺเพน, หิรียติ ปาปกานํ อกุสลานํ ธมฺมานํ สมาปตฺติยา, อยํ วุจฺจติ หิรี. อิมาย หิริยา สมนฺนาคโต ปุคฺคโล หิริมา.}

\proseitem{74}{กตโม จ ปุคฺคโล โอตฺตปฺปี, ตตฺถ กตมํ โอตฺตปฺปํ, ยํ โอตฺตปฺปติ โอตฺตปฺปิตพฺเพน, โอตฺตปฺปติ ปาปกานํ อกุสลานํ ธมฺมานํ สมาปตฺติยา, อิทํ วุจฺจติ โอตฺตปฺปํ. อิมินา โอตฺตปฺเปน สมนฺนาคโต ปุคฺคโล โอตฺตปฺปี.}

\proseitem{75}{\csromanfolio{127}กตโม จ ปุคฺคโล สุวโจ, ตตฺถ กตมา โสวจสฺสตา, สหธมฺมิเก วุจฺจมาเน โสวจสฺสายํ โสวจสฺสิยํ โสวจสฺสตา อวิปฺปฏิกุลคฺคาหิตา อวิปจฺจนีกสาตตา สาทริยํ สาทริยตา สคารวตา สปฺปติสฺสวตา, อยํ วุจฺจติ โสวจสฺสตา. อิมาย โสวจสฺสตาย สมนฺนาคโต ปุคฺคโล สุวโจ.}

\proseitem{76}{กตโม จ ปุคฺคโล กลฺยาณมิตฺโต, ตตฺถ กตมา กลฺยาณมิตฺตตา, เย เต ปุคฺคลา สทฺธา สีลวนฺโต พหุสฺสุตา จาควนฺโต ปญฺญวนฺโต, ยา เตสํ เสวนา นิเสวนา สํเสวนา ภชนา สมฺภชนา ภตฺติ สมฺภตฺติ สมฺปวงฺกตา, อยํ วุจฺจติ กลฺยาณมิตฺตตา. อิมาย กลฺยาณ\-มิตฺตตาย สมนฺนาคโต ปุคฺคโล กลฺยาณมิตฺโต.}

\proseitem{77}{กตโม จ ปุคฺคโล อินฺทฺริเยสุ คุตฺตทฺวาโร, ตตฺถ กตมา อินฺทฺริเยสุ คุตฺตทฺวารตา, อิเธกจฺโจ ปุคฺคโล จกฺขุนา รูปํ ทิสฺวา น นิมิตฺตคฺคาหี โหติ นานุพฺยญฺชน\-คฺคาหี, ยตฺวาธิกรณเมนํ จกฺขุนฺทฺริยํ อสํวุตํ วิหรนฺตํ อภิชฺฌา\-โทมนสฺสา ปาปกา อกุสลา ธมฺมา อนฺวาสฺสเวยฺยุํ, ตสฺส สํวราย ปฏิปชฺชติ, รกฺขติ จกฺขุนฺทฺริยํ, จกฺขุนฺทฺริเย สํวรํ อาปชฺชติ. โสเภน สทฺทํ สุตฺวา ฯเปฯ. ฆาเนน คนฺธํ ฆายิตฺวา ฯเปฯ. ชิวฺหาย รสํ สายิตฺวา ฯเปฯ. กาเยน โผฏฺฐพฺพํ ผุสิตฺวา ฯเปฯ. มนสา ธมฺมํ วิญฺญาย น นิมิตฺตคฺคาหี โหติ นานุพฺยญฺชน\-คฺคาหี, ยตฺวาธิกรณเมนํ มนินฺทฺริยํ อสํวุตํ วิหรนฺตํ อภิชฺฌา\-โทมนสฺสา ปาปกา อกุสลา ธมฺมา อนฺวาสฺสเวยฺยุํ, ตสฺส สํวราย ปฏิปชฺชติ, รกฺขติ มนินฺทฺริยํ, มนินฺทฺริเย สํวรํ อาปชฺชติ. ยา อิเมสํ ฉนฺนํ อินฺทฺริยานํ คุตฺติ โคปนา อารกฺโข สํวโร, อยํ วุจฺจติ อินฺทฺริเยสุ คุตฺตทฺวารตา. อิมาย อินฺทฺริเยสุ คุตฺตทฺวารตาย สมนฺนาคโต ปุคฺคโล อินฺทฺริเยสุ คุตฺตทฺวาโร.}

\proseitem{78}{กตโม จ ปุคฺคโล โภชเน มตฺตญฺญู, ตตฺถ กตมา โภชเน มตฺตญฺญุตา, อิเธกจฺโจ ปุคฺคโล ปฏิสงฺขา โยนิโส อาหารํ อาหาเรติ เนว ทวาย น มทาย น มณฺฑนาย น วิภูสนาย ยาวเทว อิมสฺส กายสฺส ฐิติยา ยาปนาย วิหิํสู\-ปรติยา พฺรหฺมจริยา\-นุคฺคหาย, อิติ ปุราณญฺจ เวทนํ ปฏิหงฺขามิ, นวญฺจ เวทนํ น อุปฺปาเทสฺสามิ, ยาตฺรา จ เม ภวิสฺสติ อนวชฺชตา จ ผาสุวิหาโร จาติ. ยา ตตฺถ สนฺตุฏฺฐิตา \csromanfolio{128}มตฺตญฺญุตา ปฏิสงฺขา โภชเน, อยํ วุจฺจติ โภชเน มตฺตญฺญุตา. อิมาย โภชเน มตฺตญฺญุตาย สมนฺนาคโต ปุคฺคโล โภชเน มตฺตญฺญู.}

\proseitem{79}{กตโม จ ปุคฺคโล อุปฏฺฐิตสฺสติ, ตตฺถ กตมา สติ, ยา สติ อนุสฺสติ ปฏิสฺสติ สติ สรณตา ธารณตา อปิลาปนตา อสมฺมุสนตา สติ สตินฺทฺริยํ สติพลํ สมฺมาสติ, อยํ วุจฺจติ สติ. อิมาย สติยา สมนฺนาคโต ปุคฺคโล อุปฏฺฐิตสฺสติ.}

\proseitem{80}{กตโม จ ปุคฺคโล สมฺปชาโน, ตตฺถ กตมํ สมฺปชญฺญํ, ยา ปญฺญา ปชานนา วิจโย ปวิจโย ธมฺมวิจโย สลฺลกฺขณา อุปลกฺขณา ปจฺจุ\-ปลกฺขณา ปณฺฑิจฺจํ โกสลฺลํ เนปุญฺญํ เวภพฺยา จินฺตา อุปปริกฺขา ภูรี เมธา ปริณายิกา วิปสฺสนา สมฺปชญฺญํ ปโตโท ปญฺญา ปญฺญินฺทฺริยํ ปญฺญาพลํ ปญฺญาสตฺถํ ปญฺญาปาสาโท ปญฺญาอาโลโก ปญฺญา- โอภาโส ปญฺญาปชฺโชโต ปญฺญารตนํ อโมโห ธมฺมวิจโย สมฺมาทิฏฺฐิ, อิทํ วุจฺจติ สมฺปชญฺญํ. อิมินา สมฺปชญฺเญน สมนฺนาคโต ปุคฺคโล สมฺปชาโน.}

\proseitem{81}{กตโม จ ปุคฺคโล สีลสมฺปนฺโน, ตตฺถ กตมา สีลสมฺปทา, กายิโก อวีติกฺกโม วาจสิโก อวีติกฺกโม กายิกวาจสิโก อวีติกฺกโม, อยํ วุจฺจติ สีลสมฺปทา, สพฺโพปิ สีลสํวโร สีลสมฺปทา. อิมาย สีลสมฺปทาย สมนฺนาคโต ปุคฺคโล สีลสมฺปนฺโน.}

\proseitem{82}{กตโม จ ปุคฺคโล ทิฏฺฐิสมฺปนฺโน. ตตฺถ กตมา ทิฏฺฐิสมฺปทา, “อตฺถิ ทินฺนํ, อตฺถิ ยิฏฺฐํ, อตฺถิ หุตํ, อตฺถิ สุกต\-ทุกฺกฏานํ กมฺมานํ ผลํ วิปาโก, อตฺถิ อยํ โลโก, อตฺถิ ปโร โลโก, อตฺถิ มาตา, อตฺถิ ปิตา, อตฺถิ สตฺตา โอปปาติกา, อตฺถิ โลเก สมณพฺราหฺมณา สมฺมคฺคตา สมฺมาปฏิปนฺนา, เย อิมญฺจ โลกํ ปรญฺจ โลกํ สยํ อภิญฺญา สจฺฉิกตฺวา ปเวเทนฺตี”ติ, ยา เอวรูปา ปญฺญา ปชานนา ฯเปฯ อโมโห ธมฺมวิจโย สมฺมาทิฏฺฐิ, อยํ วุจฺจติ ทิฏฺฐิสมฺปทา, สพฺพาปิ สมฺมาทิฏฺฐิ ทิฏฺฐิสมฺปทา. อิมาย ทิฏฺฐิ\-สมฺปทาย สมนฺนาคโต ปุคฺคโล ทิฏฺฐิสมฺปนฺโน.}

\proseitem{83}{กตเม เทฺว ปุคฺคลา ทุลฺลภา โลกสฺมิํ, โย จ ปุพฺพการี, โย จ กตญฺญู กตเวที. อิเม เทฺว ปุคฺคลา ทุลฺลภา โลกสฺมิํ.}

\proseitem{84}{\csromanfolio{129}กตเม เทฺว ปุคฺคลา ทุตฺตปฺปยา, โย จ ลทฺธํ ลทฺธํ นิกฺขิปติ, โย จ ลทฺธํ ลทฺธํ วิสฺสชฺเชติ. อิเม เทฺว ปุคฺคลา ทุตฺตปฺปยา.}

\proseitem{85}{กตเม เทฺว ปุคฺคลา สุตปฺปยา, โย จ ลทฺธํ ลทฺธํ น นิกฺขิปติ, โย จ ลทฺธํ ลทฺธํ น วิสฺสชฺเชติ. อิเม เทฺว ปุคฺคลา สุตปฺปยา.}

\proseitem{86}{กตเมสํ ทฺวินฺนํ ปุคฺคลานํ อาสวา วฑฺฒนฺติ, โย จ น กุกฺกุจฺจายิ\-ตพฺพํ กุกฺกุจฺจายติ, โย จ กุกฺกุจฺจายิ\-ตพฺพํ น กุกฺกุจฺจายติ. อิเมสํ ทฺวินฺนํ ปุคฺคลานํ อาสวา วฑฺฒนฺติ.}

\proseitem{87}{กตเมสํ ทฺวินฺนํ ปุคฺคลานํ อาสวา น วฑฺฒนฺติ, โย จ น กุกฺกุจฺจายิ\-ตพฺพํ น กุกฺกุจฺจายติ, โย จ กุกฺกุจฺจายิ\-ตพฺพํ กุกฺกุจฺจายติ. อิเมสํ ทฺวินฺนํ ปุคฺคลานํ อาสวา น วฑฺฒนฺติ.}

\proseitem{88}{กตโม จ ปุคฺคโล หีนาธิมุตฺโต, อิเธกจฺโจ ปุคฺคโล ทุสฺสีโล โหติ ปาปธมฺโม, โส อญฺญํ ทุสฺสีลํ ปาปธมฺมํ เสวติ ภชติ ปยิรุปาสติ, อยํ วุจฺจติ ปุคฺคโล หีนาธิมุตฺโต.}

\proseitem{89}{กตโม จ ปุคฺคโล ปณีตาธิมุตฺโต, อิเธกจฺโจ ปุคฺคโล สีลวา โหติ กลฺยาณธมฺโม, โส อญฺญํ สีลวนฺตํ กลฺยาณธมฺมํ เสวติ ภชติ ปยิรุปาสติ, อยํ วุจฺจติ ปุคฺคโล ปณีตาธิมุตฺโต.}

\proseitem{90}{กตโม จ ปุคฺคโล ติตฺโต, ปจฺเจก\-สมฺพุทฺธา\csromansharedfootnote{0c129bde9ceba893}{ปจฺเจกพุทฺโธ (สี)} เย จ ตถาคตสฺส สาวกา อรหนฺโต ติตฺตา. สมฺมาสมฺพุทฺโธ ติตฺโต จ ตปฺเปตา จ\csromansharedfootnote{70a3fd7afa068419}{ตปฺเปตา จ, อยํ วุจฺจติ ปุคฺคโล ติตฺโต (สี)}.}

\vspace{\tipitakaparskip}
\csromancenter{ทุกนิทฺเทโส.}
\csromanmatikaanchor{mtk.72}
\csromantocmarkref{subsection}{3. ติกปุคฺคลปญฺญตฺติ}{mtk.72}
\bookmark[dest={mtk.72},level=2]{3. ติกปุคฺคลปญฺญตฺติ}
\csromanheader{2}{3. ติกปุคฺคล\-ปญฺญตฺติ}
\proseitem{91}{กตโม จ ปุคฺคโล นิราโส, อิเธกจฺโจ ปุคฺคโล ทุสฺสีโล โหติ ปาปธมฺโม อสุจิ สงฺกสฺสร\-สมาจาโร ปฏิจฺฉนฺน\-กมฺมนฺโต \csromanfolio{130}อสฺสมโณ สมณปฏิญฺโญ อพฺรหฺมจารี พฺรหฺมจาริ\-ปฏิญฺโญ อนฺโตปูติ อวสฺสุโต กสมฺพุชาโต, โส สุณาติ “อิตฺถนฺนาโม กิร ภิกฺขุ อาสวานํ ขยา อนาสวํ เจโตวิมุตฺติํ ปญฺญาวิมุตฺติํ ทิฏฺเฐว ธมฺเม สยํ อภิญฺญา สจฺฉิกตฺวา อุปสมฺปชฺช วิหรตี”ติ, ตสฺส น เอวํ โหติ “กุทาสฺสุ นามาหมฺปิ อาสวานํ ขยา อนาสวํ เจโตวิมุตฺติํ ปญฺญาวิมุตฺติํ ทิฏฺเฐว ธมฺเม สยํ อภิญฺญา สจฺฉิกตฺวา อุปสมฺปชฺช วิหริสฺสามี”ติ. อยํ วุจฺจติ ปุคฺคโล นิราโส.}

\proseitem{92}{กตโม จ ปุคฺคโล อาสํโส, อิเธกจฺโจ ปุคฺคโล สีลวา โหติ กลฺยาณธมฺโม, โส สุณาติ “อิตฺถนฺนาโม กิร ภิกฺขุ อาสวานํ ขยา อนาสวํ เจโตวิมุตฺติํ ปญฺญาวิมุตฺติํ ทิฏฺเฐว ธมฺเม สยํ อภิญฺญา สจฺฉิกตฺวา อุปสมฺปชฺช วิหรตี”ติ, ตสฺส เอวํ โหติ “กุทาสฺสุ นามาหมฺปิ อาสวานํ ขยา อนาสวํ เจโตวิมุตฺติํ ปญฺญาวิมุตฺติํ ทิฏฺเฐว ธมฺเม สยํ อภิญฺญา สจฺฉิกตฺวา อุปสมฺปชฺช วิหริสฺสามี”ติ. อยํ วุจฺจติ ปุคฺคโล อาสํโส.}

\proseitem{93}{กตโม จ ปุคฺคโล วิคตาโส, อิเธกจฺโจ ปุคฺคโล อาสวานํ ขยา อนาสวํ เจโตวิมุตฺติํ ปญฺญาวิมุตฺติํ ทิฏฺเฐว ธมฺเม สยํ อภิญฺญา สจฺฉิกตฺวา อุปสมฺปชฺช วิหรติ, โส สุณาติ “อิตฺถนฺนาโม กิร ภิกฺขุ อาสวานํ ขยา อนาสวํ เจโตวิมุตฺติํ ปญฺญาวิมุตฺติํ ทิฏฺเฐว ธมฺเม สยํ อภิญฺญา สจฺฉิกตฺวา อุปสมฺปชฺช วิหรตี”ติ, ตสฺส น เอวํ โหติ “กุทาสฺสุ นามาหมฺปิ อาสวานํ ขยา อนาสวํ เจโตวิมุตฺติํ ปญฺญาวิมุตฺติํ ทิฏฺเฐว ธมฺเม สยํ อภิญฺญา สจฺฉิกตฺวา อุปสมฺปชฺช วิหริสฺสามี”ติ. ตํ กิสฺส เหตุ, ยา หิสฺส ปุพฺเพ อวิมุตฺตสฺส วิมุตฺตาสา, สา ปฏิปฺปสฺสทฺธา, อยํ วุจฺจติ ปุคฺคโล วิคตาโส.}

\proseitem{94}{ตตฺถ กตเม ตโย คิลานูปมา ปุคฺคลา. ตโย คิลานา\csromandash{}อิเธกจฺโจ คิลาโน ลภนฺโต วา สปฺปายานิ โภชนานิ อลภนฺโต วา สปฺปายานิ โภชนานิ ลภนฺโต วา สปฺปายานิ เภสชฺชานิ อลภนฺโต วา สปฺปายานิ เภสชฺชานิ ลภนฺโต วา ปติรูปํ อุปฏฺฐากํ อลภนฺโต วา ปติรูปํ อุปฏฺฐากํ เนว วุฏฺฐาติ ตมฺหา อาพาธา. (1)}

\prose{อิธ ปเนกจฺโจ คิลาโน ลภนฺโต วา สปฺปายานิ โภชนานิ อลภนฺโต วา สปฺปายานิ โภชนานิ ลภนฺโต วา สปฺปายานิ}

\prose{\csromanfolio{131}เภสชฺชานิ อลภนฺโต วา สปฺปายานิ เภสชฺชานิ ลภนฺโต วา ปติรูปํ อุปฏฺฐากํ อลภนฺโต วา ปติรูปํ อุปฏฺฐากํ วุฏฺฐาติ ตมฺหา อาพาธา. (2)}

\prose{อิธ ปเนกจฺโจ คิลาโน ลภนฺโต สปฺปายานิ โภชนานิ โน อลภนฺโต, ลภนฺโต สปฺปายานิ เภสชฺชานิ โน อลภนฺโต, ลภนฺโต ปติรูปํ อุปฏฺฐากํ โน อลภนฺโต วุฏฺฐาติ ตมฺหา อาพาธา. (3)}

\prose{ตตฺร ยฺวายํ คิลาโน ลภนฺโต สปฺปายานิ โภชนานิ โน อลภนฺโต, ลภนฺโต สปฺปายานิ เภสชฺชานิ โน อลภนฺโต, ลภนฺโต ปติรูปํ อุปฏฺฐากํ โน อลภนฺโต วุฏฺฐาติ ตมฺหา อาพาธา, อิมํ คิลานํ ปฏิจฺจ ภควตา คิลานภตฺตํ อนุญฺญาตํ, คิลาน\-เภสชฺชํ อนุญฺญาตํ, คิลานุปฏฺฐาโก อนุญฺญาโต, อิมญฺจ ปน คิลานํ ปฏิจฺจ อญฺเญปิ คิลานา อุปฏฺฐาตพฺพา.}

\prose{เอวเมวํ\csromansharedfootnote{b06969151201b6ae}{เอวเมว (สี)} ตโย คิลานูปมา ปุคฺคลา สนฺโต สํวิชฺชมานา โลกสฺมิํ. กตเม ตโย, อิเธกจฺโจ ปุคฺคโล ลภนฺโต วา ตถาคตํ ทสฺสนาย อลภนฺโต วา ตถาคตํ ทสฺสนาย, ลภนฺโต วา ตถาคต\-ปฺปเวทิตํ ธมฺมวินยํ สวณาย อลภนฺโต วา ตถาคต\-ปฺปเวทิตํ ธมฺมวินยํ สวณาย เนว โอกฺกมติ นิยามํ กุสเลสุ ธมฺเมสุ สมฺมตฺตํ. (1)}

\prose{อิธ ปเนกจฺโจ ปุคฺคโล ลภนฺโต วา ตถาคตํ ทสฺสนาย อลภนฺโต วา ตถาคตํ ทสฺสนาย, ลภนฺโต วา ตถาคต\-ปฺปเวทิตํ ธมฺมวินยํ สวณาย อลภนฺโต วา ตถาคต\-ปฺปเวทิตํ ธมฺมวินยํ สวณาย โอกฺกมติ นิยามํ กุสเลสุ ธมฺเมสุ สมฺมตฺตํ. (2)}

\prose{อิธ ปเนกจฺโจ ปุคฺคโล ลภนฺโต ตถาคตํ ทสฺสนาย โน อลภนฺโต, ลภนฺโต ตถาคต\-ปฺปเวทิตํ ธมฺมวินยํ สวณาย โน อลภนฺโต โอกฺกมติ นิยามํ กุสเลสุ ธมฺเมสุ สมฺมตฺตํ. (3)}

\prose{ตตฺร ยฺวายํ ปุคฺคโล ลภนฺโต ตถาคตํ ทสฺสนาย โน อลภนฺโต, ลภนฺโต ตถาคต\-ปฺปเวทิตํ ธมฺมวินยํ สวณาย โน}

\prose{\csromanfolio{132}อลภนฺโต โอกฺกมติ นิยามํ กุสเลสุ ธมฺเมสุ สมฺมตฺตํ, อิมํ ปุคฺคลํ ปฏิจฺจ ภควตา ธมฺมเทสนา อนุญฺญาตา, อิมญฺจ ปุคฺคลํ ปฏิจฺจ อญฺเญสมฺปิ ธมฺโม เทเสตพฺโพ. อิเม ตโย คิลานูปมา ปุคฺคลา สนฺโต สํวิชฺชมานา โลกสฺมิํ.}

\proseitem{95}{กตโม จ ปุคฺคโล กายสกฺขี, อิเธกจฺโจ ปุคฺคโล อฏฺฐ วิโมกฺเข กาเยน ผุสิตฺวา วิหรติ, ปญฺญาย จสฺส ทิสฺวา เอกจฺเจ อาสวา ปริกฺขีณา โหนฺติ. อยํ วุจฺจติ ปุคฺคโล กายสกฺขี.}

\proseitem{96}{กตโม จ ปุคฺคโล ทิฏฺฐิปฺปตฺโต, อิเธกจฺโจ ปุคฺคโล “อิทํ ทุกฺขนฺ”ติ ยถาภูตํ ปชานาติ, “อยํ ทุกฺขสมุทโย”ติ ยถาภูตํ ปชานาติ, “อยํ ทุกฺขนิโรโธ”ติ ยถาภูตํ ปชานาติ, “อยํ ทุกฺขนิโรธ\-คามินี ปฏิปทา”ติ ยถาภูตํ ปชานาติ, ตถาคต\-ปฺปเวทิตา จสฺส ธมฺมา ปญฺญาย โวทิฏฺฐา โหนฺติ โวจริตา, ปญฺญาย จสฺส ทิสฺวา เอกจฺเจ อาสวา ปริกฺขีณา โหนฺติ. อยํ วุจฺจติ ปุคฺคโล ทิฏฺฐิปฺปตฺโต.}

\proseitem{97}{กตโม จ ปุคฺคโล สทฺธาวิมุตฺโต, อิเธกจฺโจ ปุคฺคโล “อิทํ ทุกฺขนฺ”ติ ยถาภูตํ ปชานาติ ฯเปฯ ตถาคต\-ปฺปเวทิตา จสฺส ธมฺมา ปญฺญาย โวทิฏฺฐา โหนฺติ โวจริตา, ปญฺญาย จสฺส ทิสฺวา เอกจฺเจ อาสวา ปริกฺขีณา โหนฺติ, โน จ โข ยถา ทิฏฺฐิ\-ปฺปตฺตสฺส. อยํ วุจฺจติ ปุคฺคโล สทฺธาวิมุตฺโต.}

\proseitem{98}{กตโม จ ปุคฺคโล คูถภาณี, อิเธกจฺโจ ปุคฺคโล มุสาวาที โหติ สภคฺคโต วา ปริสคฺคโต วา ญาติมชฺฌคโต วา ปูคมชฺฌคโต วา ราชกุล\-มชฺฌคโต วา อภินีโต สกฺขิปุฏฺโฐ “เอหมฺโภ\csromansharedfootnote{42547c9fcdd6a660}{เอหิ โภ (สฺยา, ก) ม 1. 355; อํ 1. 125 ปิฏฺเฐสุปิ.} ปุริส ยํ ชานาสิ, ตํ วเทหี”ติ, โส อชานํ วา อาห “ชานามี”ติ, ชานํ วา อาห “น ชานามี”ติ, อปสฺสํ วา อาห “ปสฺสามี”ติ, ปสฺสํ วา อาห “น ปสฺสามี”ติ. อิติ อตฺตเหตุ วา ปรเหตุ วา อามิส\-กิญฺจิกฺข\-เหตุ วา สมฺปชานมุสา ภาสิตา โหติ. อยํ วุจฺจติ ปุคฺคโล คูถภาณี.}

\proseitem{99}{\csromanfolio{133}กตโม จ ปุคฺคโล ปุปฺผภาณี, อิเธกจฺโจ ปุคฺคโล มุสาวาทํ ปหาย มุสาวาทา ปฏิวิรโต โหติ สภคฺคโต วา ปริสคฺคโต วา ญาติมชฺฌคโต วา ปูคมชฺฌคโต วา ราชกุล\-มชฺฌคโต วา อภินีโต สกฺขิปุฏฺโฐ “เอหมฺโภ ปุริส ยํ ชานาสิ, ตํ วเทหี”ติ, โส อชานํ วา อาห “น ชานามี”ติ, ชานํ วา อาห “ชานามี”ติ, อปสฺสํ วา อาห “น ปสฺสามี”ติ, ปสฺสํ วา อาห “ปสฺสามี”ติ. อิติ อตฺตเหตุ วา ปรเหตุ วา อามิส\-กิญฺจิกฺข\-เหตุ วา น สมฺปชานมุสา ภาสิตา โหติ. อยํ วุจฺจติ ปุคฺคโล ปุปฺผภาณี.}

\proseitem{100}{กตโม จ ปุคฺคโล มธุภาณี, อิเธกจฺโจ ปุคฺคโล ยา สา วาจา เนลา กณฺณสุขา เปมนิยา หทยงฺคมา โปรี พหุชนกนฺตา พหุชนมนาปา, ตถารูปิํ วาจํ ภาสิตา โหติ. อยํ วุจฺจติ ปุคฺคโล มธุภาณี.}

\proseitem{101}{กตโม จ ปุคฺคโล อรุกูปมจิตฺโต, อิเธกจฺโจ ปุคฺคโล โกธโน โหติ อุปายาสพหุโล, อปฺปมฺปิ วุตฺโต สมาโน อภิสชฺชติ กุปฺปติ พฺยาปชฺชติ ปติตฺถียติ\csromansharedfootnote{3e34ca7488820369}{ปติฏฺฐียติ (สฺยา, ก) อํ 1. 122 ปิฏฺเฐปิ.}, โกปญฺจ โทสญฺจ อปฺปจฺจยญฺจ ปาตุกโรติ. เสยฺยถาปิ นาม ทุฏฺฐารุโก กฏฺเฐน วา กฐลาย\csromansharedfootnote{7d61a4a8e7b2bdd5}{กถลาย (ก), กถเลน (อฏฺฐกถา) อํ 1. 122 ปิฏฺเฐปิ.} วา ฆฏฺฏิโต ภิยฺโยโส มตฺตาย อาสวํ เทติ\csromansharedfootnote{240e0831a04ac08b}{อสฺสวโนติ (สี)}, เอวเมวํ อิเธกจฺโจ ปุคฺคโล โกธโน โหติ อุปายาสพหุโล, อปฺปมฺปิ วุตฺโต สมาโน อภิสชฺชติ กุปฺปติ พฺยาปชฺชติ ปติตฺถียติ, โกปญฺจ โทสญฺจ อปฺปจฺจยญฺจ ปาตุกโรติ. อยํ วุจฺจติ ปุคฺคโล อรุกูปมจิตฺโต.}

\proseitem{102}{กตโม จ ปุคฺคโล วิชฺชูปมจิตฺโต, อิเธกจฺโจ ปุคฺคโล “อิทํ ทุกฺขนฺ”ติ ยถาภูตํ ปชานาติ, “อยํ ทุกฺขสมุทโย”ติ ยถาภูตํ ปชานาติ, “อยํ ทุกฺขนิโรโธ”ติ ยถาภูตํ ปชานาติ, “อยํ ทุกฺขนิโรธ\-คามินี ปฏิปทา”ติ ยถาภูตํ ปชานาติ. เสยฺยถาปิ นาม จกฺขุมา ปุริโส รตฺต\-นฺธการ\-ติมิสาย วิชฺชนฺตริกาย รูปานิ ปสฺเสยฺย, เอวเมวํ อิเธกจฺโจ ปุคฺคโล “อิทํ ทุกฺขนฺ”ติ ยถาภูตํ ปชานาติ, “อยํ ทุกฺขสมุทโย”ติ ยถาภูตํ ปชานาติ, “อยํ ทุกฺขนิโรโธ”ติ ยถาภูตํ \csromanfolio{134}ปชานาติ, “อยํ ทุกฺขนิโรธ\-คามินี ปฏิปทา”ติ ยถาภูตํ ปชานาติ. อยํ วุจฺจติ ปุคฺคโล วิชฺชูปมจิตฺโต.}

\proseitem{103}{กตโม จ ปุคฺคโล วชิรูปม\-จิตฺโต, อิเธกจฺโจ ปุคฺคโล อาสวานํ ขยา อนาสวํ เจโตวิมุตฺติํ ปญฺญาวิมุตฺติํ ทิฏฺเฐว ธมฺเม สยํ อภิญฺญา สจฺฉิกตฺวา อุปสมฺปชฺช วิหรติ. เสยฺยถาปิ นาม วชิรสฺส นตฺถิ กิญฺจิ อเภชฺชํ มณิ วา ปาสาโณ วา, เอวเมวํ อิเธกจฺโจ ปุคฺคโล อาสวานํ ขยา อนาสวํ เจโตวิมุตฺติํ ปญฺญาวิมุตฺติํ ทิฏฺเฐว ธมฺเม สยํ อภิญฺญา สจฺฉิกตฺวา อุปสมฺปชฺช วิหรติ. อยํ วุจฺจติ ปุคฺคโล วชิรูปม\-จิตฺโต.}

\proseitem{104}{กตโม จ ปุคฺคโล อนฺโธ, อิเธกจฺจสฺส ปุคฺคลสฺส ตถารูปํ จกฺขุ น โหติ, ยถารูเปน จกฺขุนา อนธิคตํ วา โภคํ อธิคจฺเฉยฺย, อธิคตํ วา โภคํ ผาติํ กเรยฺย. ตถารูปมฺปิสฺส จกฺขุ น โหติ, ยถารูเปน จกฺขุนา กุสลากุสเล ธมฺเม ชาเนยฺย, สาวชฺชานวชฺเช ธมฺเม ชาเนยฺย, หีนปฺปณีเต ธมฺเม ชาเนยฺย, กณฺห\-สุกฺก\-สปฺปฏิภาเค ธมฺเม ชาเนยฺย. อยํ วุจฺจติ ปุคฺคโล อนฺโธ.}

\proseitem{105}{กตโม จ ปุคฺคโล เอกจกฺขุ, อิเธกจฺจสฺส ปุคฺคลสฺส ตถารูปํ จกฺขุ โหติ, ยถารูเปน จกฺขุนา อนธิคตํ วา โภคํ อธิคจฺเฉยฺย, อธิคตํ วา โภคํ ผาติํ กเรยฺย. ตถารูปมฺปิสฺส จกฺขุ น โหติ, ยถารูเปน จกฺกุนา กุสลากุสเล ธมฺเม ชาเนยฺย, สาวชฺชานวชฺเช ธมฺเม ชาเนยฺย, หีนปฺปณีเต ธมฺเม ชาเนยฺย, กณฺห\-สุกฺก\-สปฺปฏิภาเค ธมฺเม ชาเนยฺย. อยํ วุจฺจติ ปุคฺคโล เอกจกฺขุ.}

\proseitem{106}{กตโม จ ปุคฺคโล ทฺวิจกฺขุ, อิเธกจฺจสฺส ปุคฺคลสฺส ตถารูปํ จกฺขุ โหติ, ยถารูเปน จกฺขุนา อนธิคตํ วา โภคํ อธิคจฺเฉยฺย, อธิคตํ วา โภคํ ผาติํ กเรยฺย. ตถารูปมฺปิสฺส จกฺขุ โหติ, ยถารูเปน จกฺขุนา กุสลากุสเล ธมฺเม ชาเนยฺย, สาวชฺชานวชฺเช ธมฺเม ชาเนยฺย, หีนปฺปณีเต ธมฺเม ชาเนยฺย, กณฺห\-สุกฺก\-สปฺปฏิภาเค ธมฺเม ชาเนยฺย. อยํ วุจฺจติ ปุคฺคโล ทฺวิจกฺขุ.}

\proseitem{107}{\csromanfolio{135}กตโม จ ปุคฺคโล อวกุชฺชปญฺโญ, อิเธกจฺโจ ปุคฺคโล อารามํ คนฺตา\csromansharedfootnote{fb49894a4c820838}{คโต (สี), คนฺตฺวา (สฺยา)} โหติ อภิกฺขณํ ภิกฺขูนํ สนฺติเก ธมฺมสฺส\-วณาย\csromansharedfootnote{ed434fc0eb16b2f4}{ธมฺมสฺสวนาย (สฺยา)}, ตสฺส ภิกฺขู ธมฺมํ เทเสนฺติ อาทิกลฺยาณํ มชฺเฌกลฺยาณํ ปริโยสาน\-กลฺยาณํ สาตฺถํ สพฺยญฺชนํ เกวล\-ปริปุณฺณํ ปริสุทฺธํ พฺรหฺมจริยํ ปกาเสนฺติ. โส ตสฺมิํ อาสเน นิสินฺโน ตสฺสา กถาย เนว อาทิํ มนสิ กโรติ, น มชฺฌํ มนสิ กโรติ, น ปริโยสานํ มนสิ กโรติ, วุฏฺฐิโตปิ ตมฺหา อาสนา ตสฺสา กถาย เนว อาทิํ มนสิ กโรติ, น มชฺฌํ มนสิ กโรติ, น ปริโยสานํ มนสิ กโรติ. เสยฺยถาปิ นาม กุมฺโภ นิกฺกุชฺโช\csromansharedfootnote{ec957e1f9b63aeca}{นิกุชฺโช (สฺยา) อํ 1. 128 ปิฏฺเฐปิ.} ตตฺร อุทกํ อาสิตฺตํ วิวฏฺฏติ โน สณฺฐาติ. เอวเมวํ อิเธกจฺโจ ปุคฺคโล อารามํ คนฺตา โหติ อภิกฺขณํ ภิกฺขูนํ สนฺติเก ธมฺมสฺส\-วณาย, ตสฺส ภิกฺขู ธมฺมํ เทเสนฺติ อาทิกลฺยาณํ มชฺเฌกลฺยาณํ ปริโยสาน\-กลฺยาณํ สาตฺถํ สพฺยญฺชนํ เกวล\-ปริปุณฺณํ ปริสุทฺธํ พฺรหฺมจริยํ ปกาเสนฺติ. โส ตสฺมิํ อาสเน นิสินฺโน ตสฺสา กถาย เนว อาทิํ มนสิ กโรติ, น มชฺฌํ มนสิ กโรติ, น ปริโยสานํ มนสิ กโรติ, วุฏฺฐิโตปิ ตมฺหา อาสนา ตสฺสา กถาย เนว อาทิํ มนสิ กโรติ, น มชฺฌํ มนสิ กโรติ, น ปริโยสานํ มนสิ กโรติ. อยํ วุจฺจติ ปุคฺคโล อวกุชฺชปญฺโญ.}

\proseitem{108}{กตโม จ ปุคฺคโล อุจฺฉงฺคปญฺโญ, อิเธกจฺโจ ปุคฺคโล อารามํ คนฺตา โหติ อภิกฺขณํ ภิกฺขูนํ สนฺติเก ธมฺมสฺส\-วณาย, ตสฺส ภิกฺขู ธมฺมํ เทเสนฺติ อาทิกลฺยาณํ มชฺเฌกลฺยาณํ ปริโยสาน\-กลฺยาณํ สาตฺถํ สพฺยญฺชนํ เกวล\-ปริปุณฺณํ ปริสุทฺธํ พฺรหฺมจริยํ ปกาเสนฺติ. โส ตสฺมิํ อาสเน นิสินฺโน ตสฺสา กถาย อาทิมฺปิ มนสิ กโรติ, มชฺฌมฺปิ มนสิ กโรติ, ปริโยสานมฺปิ มนสิ กโรติ. วุฏฺฐิโต จ โข ตมฺหา อาสนา ตสฺสา กถาย เนว อาทิํ มนสิ กโรติ, น มชฺฌํ มนสิ กโรติ, น ปริโยสานํ มนสิ กโรติ. เสยฺยถาปิ นาม ปุริสสฺส อุจฺฉงฺเค นานาขชฺชกานิ อากิณฺณานิ ติลา ตณฺฑุลา\csromansharedfootnote{94a3a29ddf2cbc77}{ติลตณฺฑุลา (ก) อํ 1. 128 ปิฏฺเฐปิ.} โมทกา พทรา. โส ตมฺหา อาสนา วุฏฺฐหนฺโต สติสมฺโมสา ปกิเรยฺย. เอวเมวํ อิเธกจฺโจ ปุคฺคโล อารามํ คนฺตา โหติ อภิกฺขณํ ภิกฺขูนํ สนฺติเก ธมฺมสฺส\-วณาย, ตสฺส ภิกฺขู ธมฺมํ เทเสนฺติ อาทิกลฺยาณํ มชฺเฌกลฺยาณํ \csromanfolio{136}ปริโยสาน\-กลฺยาณํ สาตฺถํ สพฺยญฺชนํ เกวล\-ปริปุณฺณํ ปริสุทฺธํ พฺรหฺมจริยํ ปกาเสนฺติ. โส ตสฺมิํ อาสเน นิสินฺโน ตสฺสา กถาย อาทิมฺปิ มนสิ กโรติ, มชฺฌมฺปิ มนสิ กโรติ, ปริโยสานมฺปิ มนสิ กโรติ, วุฏฺฐิโต จ โข ตมฺหา อาสนา ตสฺสา กถาย เนว อาทิมฺปิ มนสิ กโรติ, น มชฺฌมฺปิ มนสิ กโรติ, น ปริโยสานมฺปิ มนสิ กโรติ. อยํ วุจฺจติ ปุคฺคโล อุจฺฉงฺคปญฺโญ.}

\proseitem{109}{กตโม จ ปุคฺคโล ปุถุปญฺโญ, อิเธกจฺโจ ปุคฺคโล อารามํ คนฺตา โหติ อภิกฺขณํ ภิกฺขูนํ สนฺติเก ธมฺมสฺส\-วณาย, ตสฺส ภิกฺขู ธมฺมํ เทเสนฺติ อาทิกลฺยาณํ มชฺเฌกลฺยาณํ ปริโยสาน\-กลฺยาณํ สาตฺถํ สพฺยญฺชนํ เกวล\-ปริปุณฺณํ ปริสุทฺธํ พฺรหฺมจริยํ ปกาเสนฺติ. โส ตสฺมิํ อาสเน นิสินฺโน ตสฺสา กถาย อาทิมฺปิ มนสิ กโรติ, มชฺฌมฺปิ มนสิ กโรติ, ปริโยสานมฺปิ มนสิ กโรติ, วุฏฺฐิโตปิ ตมฺหา อาสนา ตสฺสา กถาย อาทิมฺปิ มนสิ กโรติ, มชฺฌมฺปิ มนสิ กโรติ, ปริโยสานมฺปิ มนสิ กโรติ. เสยฺยถาปิ นาม กุมฺโภ อุกฺกุชฺโช ตตฺร อุทกํ อาสิตฺตํ สณฺฐาติ โน วิวฏฺฏติ. เอวเมวํ อิเธกจฺโจ ปุคฺคโล อารามํ คนฺตา โหติ อภิกฺขณํ ภิกฺขูณํ สนฺติเก ธมฺมสฺส\-วณาย, ตสฺส ภิกฺขู ธมฺมํ เทเสนฺติ อาทิกลฺยาณํ มชฺเฌกลฺยาณํ ปริโยสาน\-กลฺยาณํ สาตฺถํ สพฺยญฺชนํ เกวล\-ปริปุณฺณํ ปริสุทฺธํ พฺรหฺมจริยํ ปกาเสนฺติ. โส ตสฺมิํ อาสเน นิสินฺโน ตสฺสา กถาย อาทิมฺปิ มนสิ กโรติ, มชฺฌมฺปิ มนสิ กโรติ, ปริโยสานมฺปิ มนสิ กโรติ, วุฏฺฐิโตปิ ตมฺหา อาสนา ตสฺสา กถาย อาทิมฺปิ มนสิ กโรติ, มชฺฌมฺปิ มนสิ กโรติ, ปริโยสานมฺปิ มนสิ กโรติ. อยํ วุจฺจติ ปุคฺคโล ปุถุปญฺโญ.}

\proseitem{110}{กตโม จ ปุคฺคโล กาเมสุ จ ภเวสุ จ อวีตราโค, โสตาปนฺน\-สกทาคามิโน. อิเม วุจฺจนฺติ ปุคฺคลา กาเมสุ จ ภเวสุ จ อวีตราคา.}

\proseitem{111}{กตโม จ ปุคฺคโล กาเมสุ วีตราโค ภเวสุ อวีตราโค, อนาคามี. อยํ วุจฺจติ ปุคฺคโล กาเมสุ วีตราโค ภเวสุ อวีตราโค.}

\proseitem{112}{กตโม จ ปุคฺคโล กาเมสุ จ ภเวสุ จ วีตราโค, อรหา. อยํ วุจฺจติ ปุคฺคโล กาเมสุ จ ภเวสุ จ วีตราโค.}

\proseitem{113}{\csromanfolio{137}กตโม จ ปุคฺคโล ปาสาณ\-เลขู\-ปโม, อิเธกจฺโจ ปุคฺคโล อภิณฺหํ กุชฺฌติ, โส จ ขฺวสฺส โกโธ จิรํ ทีฆรตฺตํ อนุเสติ. เสยฺยถาปิ นาม ปาสาเณ เลขา น ขิปฺปํ ลุชฺชติ วาเตน วา อุทเกน วา, จิรฏฺฐิติกา โหติ. เอวเมวํ อิเธกจฺโจ ปุคฺคโล อภิณฺหํ กุชฺฌติ, โส จ ขฺวสฺส โกโธ จิรํ ทีฆรตฺตํ อนุเสติ. อยํ วุจฺจติ ปุคฺคโล ปาสาณ\-เลขู\-ปโม.}

\proseitem{114}{กตโม จ ปุคฺคโล ปถวิ\-เลขู\-ปโม, อิเธกจฺโจ ปุคฺคโล อภิณฺหํ กุชฺฌติ, โส จ ขฺวสฺส โกโธ น จิรํ ทีฆรตฺตํ อนุเสติ. เสยฺยถาปิ นาม ปถวิยา\csromansharedfootnote{bf5ff9d0e4932223}{ปฐวิยา (สี, สฺยา)} เลขา ขิปฺปํ ลุชฺชติ วาเตน วา อุทเกน วา, น จิรฏฺฐิติกา โหติ. เอวเมวํ อิเธกจฺโจ ปุคฺคโล อภิณฺหํ กุชฺฌติ, โส จ ขฺวสฺส โกโธ น จิรํ ทีฆรตฺตํ อนุเสติ. อยํ วุจฺจติ ปุคฺคโล ปถวิ\-เลขู\-ปโม.}

\proseitem{115}{กตโม จ ปุคฺคโล อุทกเลขูปโม, อิเธกจฺโจ ปุคฺคโล อาคาเฬฺหนปิ วุจฺจมาโน ผรุเสนปิ วุจฺจมาโน อมนาเปนปิ วุจฺจมาโน สํสนฺทติเมว\csromansharedfootnote{8353efe94ed60d58}{...เจว (สฺยา) อํ 1. 287 ปิฏฺเฐปิ.} สนฺธิยติเมว\csromansharedfootnote{8353efe94ed60d58}{...เจว (สฺยา) อํ 1. 287 ปิฏฺเฐปิ.} สมฺโมทติเมว\csromansharedfootnote{8353efe94ed60d58}{...เจว (สฺยา) อํ 1. 287 ปิฏฺเฐปิ.}. เสยฺยถาปิ นาม อุทเก เลขา ขิปฺปํ ลุชฺชติ, น จิรฏฺฐิติกา โหติ. เอวเมวํ อิเธกจฺโจ ปุคฺคโล อาคาเฬฺหนปิ วุจฺจมาโน ผรุเสนปิ วุจฺจมาโน อมนาเปนปิ วุจฺจมาโน สํสนฺทติเมว สนฺธิยติเมว สมฺโมทติเมว. อยํ วุจฺจติ ปุคฺคโล อุทกเลขูปโม.}

\proseitem{116}{ตตฺถ กตเม ตโย โปตฺถกูปมา ปุคฺคลา. ตโย โปตฺถกา นโวปิ โปตฺถโก ทุพฺพณฺโณ เจว โหติ ทุกฺข\-สมฺผสฺโส จ อปฺปคฺโฆ จ, มชฺฌิโมปิ โปตฺถโก ทุพฺพณฺโณ เจว โหติ ทุกฺข\-สมฺผสฺโส จ อปฺปคฺโฆ จ, ชิณฺโณปิ โปตฺถโก ทุพฺพณฺโณ เจว โหติ ทุกฺข\-สมฺผสฺโส จ อปฺปคฺโฆ จ, ชิณฺณมฺปิ โปตฺถกํ อุกฺขลิ\-ปริมชฺชนํ วา กโรนฺติ, สงฺการกูเฏ วา นํ ฉฑฺเฑนฺติ.}

\prose{เอวเมวํ ตโยเม โปตฺถกูปมา ปุคฺคลา สนฺโต สํวิชฺชมานา ภิกฺขูสุ. กตเม ตโย, นโว เจปิ ภิกฺขุ โหติ ทุสฺสีโล ปาปธมฺโม, อิมสฺส ทุพฺพณฺณตาย. เสยฺยถาปิ โส โปตฺถโก ทุพฺพณฺโณ, \csromanfolio{138}ตถูปโม อยํ ปุคฺคโล. เย โข ปนสฺส เสวนฺติ ภชนฺติ ปยิรุปาสนฺติ ทิฏฺฐานุคติํ อาปชฺชนฺติ, เตสํ ตํ โหติ ทีฆรตฺตํ อหิตาย ทุกฺขาย, อิทมสฺส ทุกฺข\-สมฺผสฺสตาย. เสยฺยถาปิ โส โปตฺถโก ทุกฺข\-สมฺผสฺโส, ตถูปโม อยํ ปุคฺคโล. เยสํ โข ปน โส\csromansharedfootnote{f6f8679f1b60b661}{เยสํ โข ปน (สพฺพตฺถ) อํ 1. 249 ปิฏฺเฐปิ.} ปฏิคฺคณฺหาติ\csromansharedfootnote{12a0af67e46bb649}{ปติคณฺหาติ (สี) รูปสิทฺธิฏีกาย ปน สเมติ.} จีวร\-ปิณฺฑปาต\-เสนาสน\-คิลานปจฺจย\-เภสชฺช\-ปริกฺขารํ, เตสํ ตํ น มหปฺผลํ โหติ น มหานิสํสํ, อิทมสฺส อปฺปคฺฆตาย. เสยฺยถาปิ โส โปตฺถโก อปฺปคฺโฆ, ตถูปโม อยํ ปุคฺคโล.}

\prose{มชฺฌิโม เจปิ ภิกฺขุ โหติ ฯเปฯ. เถโร เจปิ ภิกฺขุ โหติ ทุสฺสีโล ปาปธมฺโม, อิทมสฺส ทุพฺพณฺณตาย. เสยฺยถาปิ โส โปตฺถโก ทุพฺพณฺโณ, ตถูปโม อยํ ปุคฺคโล. เย โข ปนสฺส เสวนฺติ ภชนฺติ ปยิรุปาสนฺติ ทิฏฺฐานุคติํ อาปฺปชฺชนฺติ, เตสํ ตํ โหติ ทีฆรตฺตํ อหิตาย ทุกฺขาย, อิทมสฺส ทุกฺข\-สมฺผสฺสตาย. เสยฺยถาปิ โส โปตฺถโก ทุกฺข\-สมฺผสฺโส, ตถูปโม อยํ ปุคฺคโล. เยสํ โข ปน โส ปฏิคฺคณฺหาติ จีวร ปิณฺฑปาต เสนาสน คิลานปจฺจย เภสชฺช\-ปริกฺขารํ, เตสํ ตํ น มหปฺผลํ โหติ น มหานิสํสํ, อิทมสฺส อปฺปคฺฆตาย. เสยฺยถาปิ โส โปตฺถโก อปฺปคฺโฆ, ตถูปโม อยํ ปุคฺคโล.}

\prose{เอวรูโป เจ เถโร ภิกฺขุ สํฆมชฺเฌ ภณติ, ตเมนํ ภิกฺขู เอวมาหํสุ “กิํ นุ โข ตุยฺหํ พาลสฺส อพฺยตฺตสฺส ภณิเตน, ตฺวมฺปิ นาม ภณิตพฺพํ มญฺญสี”ติ. โส กุปิโต อนตฺตมโน ตถารูปิํ วาจํ นจฺฉาเรติ, ยถารูปาย วาจาย สํโฆ ตํ อุกฺขิปติ สงฺการกูเฏว นํ โปตฺถกํ. อิเม ตโย โปตฺถกูปมา ปุคฺคลา สนฺโต สํวิชฺชมานา ภิกฺขูสุ.}

\proseitem{117}{ตตฺถ กตเม ตโย กาสิก\-วตฺถู\-ปมา ปุคฺคลา. ตีณิ กาสิกวตฺถานิ นวมฺปิ กาสิกวตฺถํ วณฺณวนฺต\-ญฺเจว โหติ สุขสมฺผสฺสญฺจ มหคฺฆญฺจ, มชฺฌิมมฺปิ กาสิกวตฺถํ วณฺณวนฺต\-ญฺเจว โหติ สุขสมฺผสฺสญฺจ มหคฺฆญฺจ, ชิณฺณมฺปิ กาสิกวตฺถํ วณฺณวนฺต\-ญฺเจว โหติ สุขสมฺผสฺสญฺจ มหคฺฆญฺจ, ชิณฺณมฺปิ กาสิกวตฺถํ รตน\-ปลิเวฐนํ วา กโรนฺติ คนฺธ\-กรณฺฑเก วา นํ นิกฺขิปนฺติ.}

\prose{\csromanfolio{139}เอวเมวํ ตโยเม กาสิก\-วตฺถู\-ปมา ปุคฺคลา สนฺโต สํวิชฺชมานา ภิกฺขูสุ. กตเม ตโย, นโว เจปิ ภิกฺขุ โหติ สีลวา กลฺยาณธมฺโม, อิทมสฺส สุวณฺณตาย. เสยฺยถาปิ ตํ กาสิกวตฺถํ วณฺณวนฺตํ ตถูปโม อยํ ปุคฺคโล. เย โข ปนสฺส เสวนฺติ ภชนฺติ ปยิรุปาสนฺติ ทิฏฺฐานุคติํ อาปชฺชนฺติ, เตสํ ตํ โหติ ทีฆรตฺตํ หิตาย สุขาย อิทมสฺส สุขสมฺผสฺสตาย. เสยฺยถาปิ ตํ กาสิกวตฺถํ สุขสมฺผสฺสํ, ตถูปโม อยํ ปุคฺคโล. เยสํ โข ปน โส ปฏิคฺคณฺหาติ จีวร\-ปิณฺฑปาต\-เสนาสน\-คิลานปจฺจย\-เภสชฺช\-ปริกฺขารํ. เตสํ ตํ มหปฺผลํ โหติ มหานิสํสํ, อิทมสฺส มหคฺฆตาย. เสยฺยถาปิ ตํ กาสิกวตฺถํ มหคฺฆํ, ตถูปโม อยํ ปุคฺคโล.}

\prose{มชฺฌิโม เจปิ ภิกฺขุ ฯเปฯ. เถโร เจปิ ภิกฺขุ โหติ สีลวา กลฺยาณธมฺโม, อิทมสฺส สุวณฺณตาย. เสยฺยถาปิ ตํ กาสิกวตฺถํ วณฺณวนฺตํ, ตถูปโม อยํ ปุคฺคโล. เย โข ปนสฺส เสวนฺติ ภชนฺติ ปยิรุปาสนฺติ ทิฏฺฐานุคติํ อาปชฺชนฺติ, เตสํ ตํ โหติ ทีฆรตฺตํ หิตาย สุขาย, อิทมสฺส สุขสมฺผสฺสตาย. เสยฺยถาปิ ตํ กาสิกวตฺถํ สุขสมฺผสฺสํ, ตถูปโม อยํ ปุคฺคโล. เยสํ โข ปน โส ปฏิคฺคณฺหาติ จีวร ปิณฺฑปาต เสนาสน คิลานปจฺจย เภสชฺช ปริกฺขารํ, เตสํ ตํ มหปฺผลํ โหติ มหานิสํสํ, อิทมสฺส มหคฺฆตาย. เสยฺยถาปิ ตํ กาสิกวตฺถํ มหคฺฆํ, ตถูปโม อยํ ปุคฺคโล.}

\prose{เอวรูโป เจ เถโร ภิกฺขุ สํฆมชฺเฌ ภณติ, ตเมนํ ภิกฺขู เอวมาหํสุ “อปฺปสทฺทา อายสฺมนฺโต โหถ, เถโร ภิกฺขุ ธมฺมญฺจ วินยญฺจ ภณตี”ติ. ตสฺส ตํ วจนํ อาเธยฺยํ คจฺฉติ, คนฺธ\-กรณฺฑเกว นํ กาสิกวตฺถํ. อิเม ตโย กาสิก\-วตฺถู\-ปมา ปุคฺคลา สนฺโต สํวิชฺชมานา ภิกฺขูสุ.}

\proseitem{118}{กตโม จ ปุคฺคโล สุปฺปเมยฺโย, อิเธกจฺโจ ปุคฺคโล อุทฺธโต โหติ อุนฺนโฬ จปโล มุขโร วิกิณฺณวาโจ มุฏฺฐสฺสติ อสมฺปชาโน อสมาหิโต วิพฺภนฺตจิตฺโต ปากฏินฺทฺริโย. อยํ วุจฺจติ ปุคฺคโล สุปฺปเมยฺโย.}

\proseitem{119}{กตโม จ ปุคฺคโล ทุปฺปเมยฺโย, อิเธกจฺโจ ปุคฺคโล อนุทฺธโต โหติ อนุนฺนโฬ อจปโล อมุขโร อวิกิณฺณวาโจ \csromanfolio{140}อุปฏฺฐิตสฺสติ สมฺปชาโน สมาหิโต เอกคฺคจิตฺโต สํวุตินฺทฺริโย. อยํ วุจฺจติ ปุคฺคโล ทุปฺปเมยฺโย.}

\proseitem{120}{กตโม จ ปุคฺคโล อปฺปเมยฺโย, อิเธกจฺโจ ปุคฺคโล อาสวานํ ขยา อนาสวํ เจโตวิมุตฺติํ ปญฺญาวิมุตฺติํ ทิฏฺเฐว ธมฺเม สยํ อภิญฺญา สจฺฉิกตฺวา อุปสมฺปชฺช วิหรติ. อยํ วุจฺจติ ปุคฺคโล อปฺปเมยฺโย.}

\proseitem{121}{กตโม จ ปุคฺคโล น เสวิตพฺโพ น ภชิตพฺโพ น ปยิรุปาสิตพฺโพ, อิเธกจฺโจ ปุคฺคโล หีโน โหติ สีเลน สมาธินา ปญฺญาย, เอวรูโป ปุคฺคโล น เสวิตพฺโพ น ภชิตพฺโพ น ปยิรุปาสิตพฺโพ อญฺญตฺร อนุทฺทยา อญฺญตฺร อนุกมฺปา.}

\proseitem{122}{กตโม จ ปุคฺคโล เสวิตพฺโพ ภชิตพฺโพ ปยิรุปาสิตพฺโพ, อิเธกจฺโจ ปุคฺคโล สทิโส โหติ สีเลน สมาธินา ปญฺญาย, เอวรูโป ปุคฺคโล เสวิตพฺโพ ภชิตพฺโพ ปยิรุปาสิตพฺโพ. ตํ กิสฺส เหตุ, สีลสามญฺญคตานํ สตํ สีลกถา จ โน ภวิสฺสติ, สา จ โน ผาสุ ภวิสฺสติ, สา จ โน ปวตฺตินี\csromansharedfootnote{0357e7ca1f4a23c7}{ปวตฺตนี (สี) อํ 1. 123 ปิฏฺเฐปิ.} ภวิสฺสติ. สมาธิ\-สามญฺญคตานํ สตํ สมาธิกถา จ โน ภวิสฺสติ, สา จ โน ผาสุ ภวิสฺสติ, สา จ โน ปวตฺตินี ภวิสฺสติ. ปญฺญา\-สามญฺญคตานํ สตํ ปญฺญากถา จ โน ภวิสฺสติ, สา จ โน ผาสุ ภวิสฺสติ, สา จ โน ปวตฺตินี ภวิสฺสตีติ. ตสฺมา เอวรูโป ปุคฺคโล เสวิตพฺโพ ภชิตพฺโพ ปยิรุปาสิตพฺโพ.}

\proseitem{123}{กตโม จ ปุคฺคโล สกฺกตฺวา ครุํ กตฺวา เสวิตพฺโพ ภชิตพฺโพ ปยิรุปาสิตพฺโพ, อิเธกจฺโจ ปุคฺคโล อธิโก โหติ สีเลน สมาธินา ปญฺญาย, เอวรูโป ปุคฺคโล สกฺกตฺวา ครุํ กตฺวา เสวิตพฺโพ ภชิตพฺโพ ปยิรุปาสิตพฺโพ. ตํ กิสฺส เหตุ, อปริปูรํ วา สีลกฺขนฺธํ ปริปูเรสฺสามิ, ปริปูรํ วา สีลกฺขนฺธํ ตตฺถ ตตฺถ ปญฺญาย อนุคฺคเหสฺสามิ. อปริปูรํ วา สมาธิ\-กฺขนฺธํ ปริปูเรสฺสามิ, ปริปูรํ วา สมาธิ\-กฺขนฺธํ ตตฺถ ตตฺถ ปญฺญาย อนุคฺคเหสฺสามิ. อปริปูรํ วา ปญฺญากฺขนฺธํ ปริปูเรสฺสามิ, ปริปูรํ วา ปญฺญากฺขนฺธํ ตตฺถ ตตฺถ ปญฺญาย อนุคฺคเหสฺสามีติ. ตสฺมา เอวรูโป ปุคฺคโล สกฺกตฺวา ครุํ กตฺวา เสวิตพฺโพ ภชิตพฺโพ ปยิรุปาสิตพฺโพ.}

\proseitem{124}{\csromanfolio{141}กตโม จ ปุคฺคโล ชิคุจฺฉิตพฺโพ น เสวิตพฺโพ น ภชิตพฺโพ น ปยิรุปาสิตพฺโพ, อิเธกจฺโจ ปุคฺคโล ทุสฺสีโล โหติ ปาปธมฺโม อสุจิ สงฺกสฺสร\-สมาจาโร ปฏิจฺฉนฺน\-กมฺมนฺโต อสฺสมโณ สมณปฏิญฺโญ อพฺรหฺมจารี พฺรหฺมจาริ\-ปฏิญฺโญ อนฺโตปูติ อวสฺสุโต กสมฺพุชาโต. เอวรูโป ปุคฺคโล ชิคุจฺฉิตพฺโพ น เสวิตพฺโพ น ภชิตพฺโพ น ปยิรุปาสิตพฺโพ. ตํ กิสฺส เหตุ, กิญฺจาปิ เอวรูปสฺส ปุคฺคลสฺส น ทิฏฺฐานุคติํ อาปชฺชติ, อถ โข นํ ปาปโก กิตฺติสทฺโท อพฺภุคฺคจฺฉติ “ปาปมิตฺโต ปุริสปุคฺคโล ปาปสหาโย ปาปสมฺปวงฺโก”ติ. เสยฺยถาปิ นาม อหิ คูถคโต กิญฺจาปิ น ฑํสติ, อถ โข นํ มกฺเขติ. เอวเมวํ กิญฺจาปิ เอวรูปสฺส ปุคฺคลสฺส น ทิฏฺฐานุคติํ อาปชฺชติ, อถ โข นํ ปาปโก กิตฺติสทฺโท อพฺภุคฺคจฺฉติ “ปาปมิตฺโต ปุริสปุคฺคโล ปาปสหาโย ปาปสมฺปวงฺโก”ติ. ตสฺมา เอวรูโป ปุคฺคโล ชิคุจฺฉิตพฺโพ น เสวิตพฺโพ น ภชิตพฺโพ น ปยิรุปาโสตพฺโพ.}

\proseitem{125}{กตโม จ ปุคฺคโล อชฺฌุเปกฺขิตพฺโพ น เสวิตพฺโพ น ภชิตพฺโพ น ปยิรุปาสิตพฺโพ, อิเธกจฺโจ ปุคฺคโล โกธโน โหติ อุปายาสพหุโล, อปฺปมฺปิ วุตฺโต สมาโน อภิสชฺชติ กุปฺปติ พฺยาปชฺชติ ปติตฺถียติ, โกปญฺจ โทสญฺจ อปฺปจฺจยญฺจ ปาตุกโรติ. เสยฺยถาปิ นาม ทุฏฺฐารุโก กฏฺเฐน วา กฐลาย วา ฆฏฺฏิโต ภิยฺโยโส มตฺตาย อาสวํ เทติ. เอวเมวํ อิเธกจฺโจ ปุคฺคโล โกธโน โหติ อุปายาสพหุโล, อปฺปมฺปิ วุตฺโต สมาโน อภิสชฺชติ กุปฺปติ พฺยาปชฺชติ ปติตฺถียติ, โกปญฺจ โทสญฺจ อปฺปจฺจยญฺจ ปาตุกโรติ. เสยฺยถาปิ นาม ตินฺทุกาลาตํ กฏฺเฐน วา กฐลาย วา ฆฏฺฏิตํ ภิยฺโยโส มตฺตาย จิจฺจิฏายติ จิฏิจิฏายติ. เอวเมวํ อิเธกจฺโจ ปุคฺคโล โกธโน โหติ อุปายาสพหุโล, อปฺปมฺปิ วุตฺโต สมาโน อภิสชฺชติ กุปฺปติ พฺยาปชฺชติ ปติตฺถียติ, โกปญฺจ โทสญฺจ อปฺปจฺจยญฺจ ปาตุกโรติ. เสยฺยถาปิ นาม คูถกูโป กฏฺเฐน วา กฐลาย วา ฆฏฺฏิโต ภิยฺโยโส มตฺตาย ทุคฺคนฺโธ โหติ. เอวเมวํ อิเธกจฺโจ ปุคฺคโล โกธโน โหติ อุปายาสพหุโล, อปฺปมฺปิ วุตฺโต สมาโน อภิสชฺชติ กุปฺปติ พฺยาปชฺชติ ปติตฺถียติ, โกปญฺจ โทสญฺจ อปฺปจฺจยญฺจ ปาตุกโรติ, เอวรูโป ปุคฺคโล อชฺฌุเปกฺขิตพฺโพ น เสวิตพฺโพ น ภชิตพฺโพ น ปยิรุปาสิตพฺโพ. ตํ กิสฺส เหตุ, อกฺโกเสยฺยปิ มํ, ปริภาเสยฺยปิ มํ, \csromanfolio{142}อนตฺถมฺปิ เม กเรยฺยาติ. ตสฺมา เอวรูโป ปุคฺคโล อชฺฌุเปกฺขิตพฺโพ น เสวิตพฺโพ น ภชิตพฺโพ น ปยิรุปาสิตพฺโพ.}

\proseitem{126}{กตโม จ ปุคฺคโล เสวิตพฺโพ ภชิตพฺโพ ปยิรุปาสิตพฺโพ, อิเธกจฺโจ ปุคฺคโล สีลวา โหติ กลฺยาณธมฺโม, เอวรูโป ปุคฺคโล เสวิตพฺโพ ภชิตพฺโพ ปยิรุปาสิตพฺโพ. ตํ กิสฺส เหตุ, กิญฺจาปิ เอวรูปสฺส ปุคฺคลสฺส น ทิฏฺฐานุคติํ อาปชฺชติ, อถ โข นํ กลฺยาโณ กิตฺติสทฺโท อพฺภุคฺคจฺฉติ “กลฺยาณมิตฺโต ปุริสปุคฺคโล กลฺยาณสหาโย กลฺยาณ\-สมฺปวงฺโก”ติ. ตสฺมา เอวรูโป ปุคฺคโล เสวิตพฺโพ ภชิตพฺโพ ปยิรุปาสิตพฺโพ.}

\proseitem{127}{กตโม จ ปุคฺคโล สีเลสุ ปริปูรการี สมาธิสฺมิํ มตฺตโส การี ปญฺญาย มตฺตโส การี, โสตาปนฺน\-สกทาคามิโน. อิเม วุจฺจนฺติ ปุคฺคลา สีเลสุ ปริปูรการิโน สมาธิสฺมิํ มตฺตโส การิโน ปญฺญาย มตฺตโส การิโน.}

\proseitem{128}{กตโม จ ปุคฺคโล สีเลสุ จ ปริปูรการี สมา ธิสฺมิญฺจ ปริปูรการี ปญฺญาย มตฺตโส การี, อนาคามี. อยํ วุจฺจติ ปุคฺคโล สีเลสุ จ ปริปูรการี สมาธิสฺมิญฺจ ปริปูรการี ปญฺญาย มตฺตโส การี.}

\proseitem{129}{กตโม จ ปุคฺคโล สีเลสุ จ ปริปูรการี สมาธิสฺมิญฺจ ปริปูรการี ปญฺญาย จ ปริปูรการี, อรหา. อยํ วุจฺจติ ปุคฺคโล สีเลสุ จ ปริปูรการี สมาธิสฺมิญฺจ ปริปูรการี ปญฺญาย จ ปริปูรการี.}

\proseitem{130}{ตตฺถ กตเม ตโย สตฺถาโร, อิเธกจฺโจ สตฺถา กามานํ ปริญฺญํ ปญฺญเปติ\csromansharedfootnote{41d129e9bea9532b}{ปญฺญาเปติ (สี, สฺยา)}, น รูปานํ ปริญฺญํ ปญฺญเปติ, น เวทนานํ ปริญฺญํ ปญฺญเปติ. อิธ ปเนกจฺโจ สตฺถา กามานญฺจ ปริญฺญํ ปญฺญเปติ, รูปานญฺจ ปริญฺญํ ปญฺญเปติ, น เวทนานํ ปริญฺญํ ปญฺญเปติ. อิธ ปเนกจฺโจ สตฺถา กามานญฺจ ปริญฺญํ ปญฺญเปติ, รูปานญฺจ ปริญฺญํ ปญฺญเปติ, เวทนานญฺจ ปริญฺญํ ปญฺญเปติ.}

\prose{ตตฺร ยฺวายํ สตฺถา กามานํ ปริญฺญํ ปญฺญเปติ, น รูปานํ ปริญฺญํ ปญฺญเปติ, น เวทนานํ ปริญฺญํ ปญฺญเปติ, รูปาวจร\-สมาปตฺติยา ลาภี สตฺถา เตน ทฏฺฐพฺโพ. ตตฺร ยฺวายํ สตฺถา กามานญฺจ ปริญฺญํ ปญฺญเปติ,}

\prose{\csromanfolio{143}รูปานญฺจ ปริญฺญํ ปญฺญเปติ, น เวทนานํ ปริญฺญํ ปญฺญเปติ, อรูปาวจร\-สมาปตฺติยา ลาภี สตฺถา เตน ทฏฺฐพฺโพ. ตตฺร ยฺวายํ สตฺถา กามานญฺจ ปริญฺญํ ปญฺญเปติ, รูปานญฺจ ปริญฺญํ ปญฺญเปติ, เวทนานญฺจ ปริญฺญํ ปญฺญเปติ, สมฺมาสมฺพุทฺโธ สตฺถา เตน ทฏฺฐพฺโพ. อิเม ตโย สตฺถาโร.}

\proseitem{131}{ตตฺถ กตเม อปเรปิ ตโย สตฺถาโร, อิเธกจฺโจ สตฺถา ทิฏฺเฐ เจว ธมฺเม อตฺตานํ สจฺจโต เถตโต ปญฺญเปติ, อภิสมฺปรายญฺจ อตฺตานํ สจฺจโต เถตโต ปญฺญเปติ. อิธ ปเนกจฺโจ สตฺถา ทิฏฺเฐ เจว ธมฺเม อตฺตานํ สจฺจโต เถตโต ปญฺญเปติ, โน จ โข อภิสมฺปรายํ อตฺตานํ สจฺจโต เถตโต ปญฺญเปติ. อิธ ปเนกจฺโจ สตฺถา ทิฏฺเฐ เจว ธมฺเม อตฺตานํ สจฺจโต เถตโต น ปญฺญเปติ, อภิสมฺปรายญฺจ อตฺตานํ สจฺจโต เถตโต น ปญฺญเปติ.}

\prose{ตตฺร ยฺวายํ สตฺถา ทิฏฺเฐ เจว ธมฺเม อตฺตานํ สจฺจโต เถตโต ปญฺญเปติ, อภิสมฺปรายญฺจ อตฺตานํ สจฺจโต เถตโต ปญฺญเปติ, สสฺสตวาโท สตฺถา เตน ทฏฺฐพฺโพ. ตตฺร ยฺวายํ สตฺถา ทิฏฺเฐ เจว ธมฺเม อตฺตานํ สจฺจโต เถตโต ปญฺญเปติ, โน จ โข อภิสมฺปรายํ อตฺตานํ สจฺจโต เถตโต ปญฺญเปติ, อุจฺเฉทวาโท สตฺถา เตน ทฏฺฐพฺโพ. ตตฺร ยฺวายํ สตฺถา ทิฏฺเฐ เจว ธมฺเม อตฺตานํ สจฺจโต เถตโต น ปญฺญเปติ, อภิสมฺปรายญฺจ อตฺตานํ สจฺจโต เถตโต น ปญฺญเปติ, สมฺมาสมฺพุทฺโธ สตฺถา เตน ทฏฺฐพฺโพ. อิเม อปเรปิ ตโย สตฺถาโร.}

\vspace{\tipitakaparskip}
\csromancenter{ติกนิทฺเทโส.}
\csromanmatikaanchor{mtk.73}
\csromantocmarkref{subsection}{4. จตุกฺกปุคฺคลปญฺญตฺติ}{mtk.73}
\bookmark[dest={mtk.73},level=2]{4. จตุกฺกปุคฺคลปญฺญตฺติ}
\csromanheader{2}{4. จตุกฺก\-ปุคฺคล\-ปญฺญตฺติ}
\proseitem{132}{กตโม จ ปุคฺคโล อสปฺปุริโส, อิเธกจฺโจ ปุคฺคโล ปาณาติปาตี โหติ, อทินฺนาทายี โหติ, กาเมสุ\-มิจฺฉาจารี โหติ, มุสาวาที โหติ, สุรา\-เมรย\-มชฺช\-ปมาท\-ฏฺฐายี โหติ. อยํ วุจฺจติ ปุคฺคโล อสปฺปุริโส.}

\proseitem{133}{กตโม จ ปุคฺคโล อสปฺปุริเสน อสปฺปุริสตโร, อิเธกจฺโจ ปุคฺคโล อตฺตนา จ ปาณาติปาตี โหติ, ปรญฺจ ปาณาติปาเต \csromanfolio{144}สมาทเปติ. อตฺตนา จ อทินฺนาทายี โหติ, ปรญฺจ อทินฺนาทาเน สมาทเปติ. อตฺตนา จ กาเมสุ\-มิจฺฉาจารี โหติ, ปรญฺจ กาเมสุ\-มิจฺฉาจาเร สมาทเปติ. อตฺตนา จ มุสาวาที โหติ, ปรญฺจ มุสาวาเท สมาทเปติ. อตฺตนา จ สุรา\-เมรย\-มชฺช\-ปมาท\-ฏฺฐายี โหติ, ปรญฺจ สุรา\-เมรย\-มชฺช\-ปมาทฏฺฐาเน สมาทเปติ. อยํ วุจฺจติ ปุคฺคโล อสปฺปุริเสน อสปฺปุริสตโร.}

\proseitem{134}{กตโม จ ปุคฺคโล สปฺปุริโส, อิเธกจฺโจ ปุคฺคโล ปาณาติปาตา ปฏิวิรโต โหติ, อทินฺนาทานา ปฏิวิรโต โหติ, กาเมสุ\-มิจฺฉาจารา ปฏิวิรโต โหติ, มุสาวาทา ปฏิวิรโต โหติ, สุรา\-เมรย\-มชฺช\-ปมาทฏฺฐานา ปฏิวิรโต โหติ. อยํ วุจฺจติ ปุคฺคโล สปฺปุริโส.}

\proseitem{135}{กตโม จ ปุคฺคโล สปฺปุริเสน สปฺปุริสตโร, อิเธกจฺโจ ปุคฺคโล อตฺตนา จ ปาณาติปาตา ปฏิวิรโต โหติ, ปรญฺจ ปาณาติปาตา เวรมณิยา สมาทเปติ. อตฺตนา จ อทินฺนาทานา ปฏิวิรโต โหติ, ปรญฺจ อทินฺนาทานา เวรมณิยา สมาทเปติ. อตฺตนา จ กาเมสุ\-มิจฺฉาจารา ปฏิวิรโต โหติ, ปรญฺจ กาเมสุ\-มิจฺฉาจารา เวรมณิยา สมาทเปติ. อตฺตนา จ มุสาวาทา ปฏิวิรโต โหติ, ปรญฺจ มุสาวาทา เวรมณิยา สมาทเปติ. อตฺตนา จ สุรา\-เมรย\-มชฺช\-ปมาทฏฺฐานา ปฏิวิรโต โหติ, ปรญฺจ สุราเมรยมชฺชปมาทฏฺฐนา เวรมณิยา สมาทเปติ. อยํ วุจฺจติ ปุคฺคโล สปฺปุริเสน สปฺปุริสตโร.}

\proseitem{136}{กตโม จ ปุคฺคโล ปาโป, อิเธกจฺโจ ปุคฺคโล ปาณาติปาตี โหติ, อทินฺนาทายี โหติ, กาเมสุ\-มิจฺฉาจารี โหติ, มุสาวาที โหติ, ปิสุณวาโจ\csromansharedfootnote{f14a8426ac926009}{ปิสุณาวาโจ} โหติ, ผรุสวาโจ\csromansharedfootnote{f8a3e15cee38302b}{ผรุสาวาโจ (สี) ที 3. 67 ปิฏฺเฐปิ.} โหติ, สมฺผปฺปลาปี โหติ, อภิชฺฌาลุ โหติ, พฺยาปนฺนจิตฺโต โหติ, มิจฺฉาทิฏฺฐิ\csromansharedfootnote{1ba141f937734156}{มิจฺฉาทิฏฺฐี (ก)} โหติ. อยํ วุจฺจติ ปุคฺคโล ปาโป.}

\proseitem{137}{กตโม จ ปุคฺคโล ปาเปน ปาปตโร, อิเธกจฺโจ ปุคฺคโล อตฺตนา จ ปาณาติปาตี โหติ, ปรญฺจ ปาณาติปาเต สมาทเปติ.}

\prose{\csromanfolio{145}อตฺตนา จ อทินฺนาทายี โหติ, ปรญฺจ อทินฺนาทาเน สมาทเปติ. อตฺตนา จ กาเมสุ\-มิจฺฉาจารี โหติ, ปรญฺจ กาเมสุ\-มิจฺฉาจาเร สมาทเปติ. อตฺตนา จ มุสาวาที โหติ, ปรญฺจ มุสาวาเท สมาทเปติ. อตฺตนา จ ปิสุณาวาโจ โหติ, ปรญฺจ ปิสุณาย วาจาย สมาทเปติ. อตฺตนา จ ผรุสวาโจ โหติ, ปรญฺจ ผรุสาย วาจาย สมาทเปติ. อตฺตนา จ สมฺผปฺปลาปี โหติ, ปรญฺจ สมฺผปฺปลาเป สมาทเปติ. อตฺตนา จ อภิชฺฌาลุ โหติ, ปรญฺจ อภิชฺฌาย สมาทเปติ. อตฺตนา จ พฺยาปนฺนจิตฺโต โหติ, ปรญฺจ พฺยาปาเท สมาทเปติ. อตฺตนา จ มิจฺฉาทิฏฺฐิ โหติ, ปรญฺจ มิจฺฉาทิฏฺฐิยา สมาทเปติ. อยํ วุจฺจติ ปุคฺคโล ปาเปน ปาปตโร.}

\proseitem{138}{กตโม จ ปุคฺคโล กลฺยาโณ, อิเธกจฺโจ ปุคฺคโล, ปาณาติปาตา ปฏิวิรโต โหติ, อทินฺนาทานา ปฏิวิรโต โหติ, กาเมสุ\-มิจฺฉาจารา ปฏิวิรโต โหติ, มุสาวาทา ปฏิวิรโต โหติ, ปิสุณาย วาจาย ปฏิวิรโต โหติ, ผรุสาย วาจาย ปฏิวิรโต โหติ, สมฺผปฺปลาปา ปฏิวิรโต โหติ, อนภิชฺฌาลุ โหติ, อพฺยาปนฺนจิตฺโต โหติ, สมฺมาทิฏฺฐิ\csromansharedfootnote{3f61ac8498a45edc}{สมฺมาทิฏฺฐี (ก)} โหติ. อยํ วุจฺจติ ปุคฺคโล กลฺยาโณ.}

\proseitem{139}{กตโม จ ปุคฺคโล กลฺยาเณน กลฺยาณตโร, อิเธกจฺโจ ปุคฺคโล อตฺตนา จ ปาณาติปาตา ปฏิวิรโต โหติ, ปรญฺจ ปาณาติปาตา เวรมณิยา สมาทเปติ. อตฺตนา จ อทินฺนาทานา ปฏิวิรโต โหติ, ปรญฺจ อทินฺนาทานา เวรมณิยา สมาทเปติ. อตฺตนา จ กาเมสุ\-มิจฺฉาจารา ปฏิวิรโต โหติ, ปรญฺจ กาเมสุ\-มิจฺฉาจารา เวรมณิยา สมาทเปติ. อตฺตนา จ มุสาวาทา ปฏิวิรโต โหติ, ปรญฺจ มุสาวาทา เวรมณิยา สมาทเปติ. อตฺตนา จ ปิสุณาย วาจาย ปฏิวิรโต โหติ, ปรญฺจ ปิสุณาย วาจาย เวรมณิยา สมาทเปติ. อตฺตนา จ ผรุสาย วาจาย ปฏิวิรโต โหติ, ปรญฺจ ผรุสาย วาจาย เวรมณิยา สมาทเปติ. อตฺตนา จ สมฺผปฺปลาปา ปฏิวิรโต โหติ, ปรญฺจ สมฺผปฺปลาปา เวรมณิยา สมาทเปติ. อตฺตนา จ อนภิชฺฌาลุ โหติ, ปรญฺจ อนภิชฺฌาย สมาทเปติ. อตฺตนา จ อพฺยาปนฺนจิตฺโต โหติ, ปรญฺจ อพฺยาปาเท สมาทเปติ. อตฺตนา จ สมฺมาทิฏฺฐิ โหติ, \csromanfolio{146}ปรญฺจ สมฺมาทิฏฺฐิยา สมาทเปติ. อยํ วุจฺจติ ปุคฺคโล กลฺยาเณน กลฺยาณตโร.}

\proseitem{140}{กตโม จ ปุคฺคโล ปาปธมฺโม, อิเธกจฺโจ ปุคฺคโล ปาณาติปาตี โหติ, อทินฺนาทายี โหติ ฯเปฯ มิจฺฉาทิฏฺฐิ โหติ. อยํ วุจฺจติ ปุคฺคโล ปาปธมฺโม.}

\proseitem{141}{กตโม จ ปุคฺคโล ปาปธมฺเมน ปาปธมฺมตโร, อิเธกจฺโจ ปุคฺคโล อตฺตนา จ ปาณาติปาตี โหติ, ปรญฺจ ปาณาติปาเต สมาทเปติ. อตฺตนา จ อทินฺนาทายี โหติ, ปรญฺจ อทินฺนาทาเน สมาทเปติ ฯเปฯ. อตฺตนา จ มิจฺฉาทิฏฺฐิ โหติ, ปรญฺจ มิจฺฉาทิฏฺฐิยา สมาทเปติ. อยํ วุจฺจติ ปุคฺคโล ปาปธมฺเมน ปาปธมฺมตโร.}

\proseitem{142}{กตโม จ ปุคฺคโล กลฺยาณธมฺโม, อิเธกจฺโจ ปุคฺคโล ปาณาติปาตา ปฏิวิรโต โหติ, อทินฺนาทานา ปฏิวิรโต โหติ ฯเปฯ สมฺมาทิฏฺฐิ โหติ. อยํ วุจฺจติ ปุคฺคโล กลฺยาณธมฺโม.}

\proseitem{143}{กตโม จ ปุคฺคโล กลฺยาณ\-ธมฺเมน กลฺยาณ\-ธมฺม\-ตโร, อิเธกจฺโจ ปุคฺคโล อตฺตนา จ ปาณาติปาตา ปฏิวิรโต โหติ, ปรญฺจ ปาณาติปาตา เวรมณิยา สมาทเปติ ฯเปฯ. อตฺตนา จ สมฺมาทิฏฺฐิ โหติ, ปรญฺจ สมฺมาทิฏฺฐิยา สมาทเปติ. อยํ วุจฺจติ ปุคฺคโล กลฺยาณ\-ธมฺเมน กลฺยาณ\-ธมฺม\-ตโร.}

\proseitem{144}{กตโม จ ปุคฺคโล สาวชฺโช, อิเธกจฺโจ ปุคฺคโล สาวชฺเชน กายกมฺเมน สมนฺนาคโต โหติ, สาวชฺเชน วจีกมฺเมน สมนฺนาคโต โหติ, สาวชฺเชน มโนกมฺเมน สมนฺนาคโต โหติ. อยํ วุจฺจติ ปุคฺคโล สาวชฺโช.}

\proseitem{145}{กตโม จ ปุคฺคโล วชฺชพหุโล, อิเธกจฺโจ ปุคฺคโล สาวชฺเชน พหุลํ กายกมฺเมน สมนฺนาคโต โหติ อปฺปํ อนวชฺเชน, สาวชฺเชน พหุลํ วจีกมฺเมน สมนฺนาคโต โหติ อปฺปํ อนวชฺเชน, สาวชฺเชน พหุลํ มโนกมฺเมน สมนฺนาคโต โหติ อปฺปํ อนวชฺเชน. อยํ วุจฺจติ ปุคฺคโล วชฺชพหุโล.}

\proseitem{146}{\csromanfolio{147}กตโม จ ปุคฺคโล อปฺปวชฺโช, อิเธกจฺโจ ปุคฺคโล อนวชฺเชน พหุลํ กายกมฺเมน สมนฺนาคโต โหติ อปฺปํ สาวชฺเชน, อนวชฺเชน พหุลํ วจีกมฺเมน สมนฺนาคโต โหติ อปฺปํ สาวชฺเชน, อนวชฺเชน พหุลํ มโนกมฺเมน สมนฺนาคโต โหติ อปฺปํ สาวชฺเชน. อยํ วุจฺจติ ปุคฺคโล อปฺปวชฺโช.}

\proseitem{147}{กตโม จ ปุคฺคโล อนวชฺโช, อิเธกจฺโจ ปุคฺคโล อนวชฺเชน กายกมฺเมน สมนฺนาคโต โหติ, อนวชฺเชน วจีกมฺเมน สมนฺนาคโต โหติ, อนวชฺเชน มโนกมฺเมน สมนฺนาคโต โหติ. อยํ วุจฺจติ ปุคฺคโล อนวชฺโช.}

\proseitem{148}{กตโม จ ปุคฺคโล อุคฺฆฏิตญฺญู, ยสฺส ปุคฺคลสฺส สห อุทาหฏเวลาย ธมฺมา\-ภิสมโย โหติ. อยํ วุจฺจติ ปุคฺคโล อุคฺฆฏิตญฺญู.}

\proseitem{149}{กตโม จ ปุคฺคโล วิปญฺจิตญฺญู, ยสฺส ปุคฺคลสฺส สํขิตฺเตน ภาสิตสฺส วิตฺถาเรน อตฺเถ วิภชิยมาเน ธมฺมา\-ภิสมโย โหติ. อยํ วุจฺจติ ปุคฺคโล วิปญฺจิตญฺญู.}

\proseitem{150}{กตโม จ ปุคฺคโล เนยฺโย, ยสฺส ปุคฺคลสฺส อุทฺเทสโต ปริปุจฺฉโต โยนิโส มนสิกโรโต กลฺยาณมิตฺเต เสวโต ภชโต ปยิรุปาสโต เอวํ อนุปุพฺเพน ธมฺมา\-ภิสมโย โหติ. อยํ วุจฺจติ ปุคฺคโล เนยฺโย.}

\proseitem{151}{กตโม จ ปุคฺคโล ปทปรโม, ยสฺส ปุคฺคลสฺส พหุมฺปิ สุณโต พหุมฺปิ ภณโต พหุมฺปิ ธารยโต พหุมฺปิ วาจยโต น ตาย ชาติยา ธมฺมา\-ภิสมโย โหติ. อยํ วุจฺจติ ปุคฺคโล ปทปรโม.}

\proseitem{152}{กตโม จ ปุคฺคโล ยุตฺต\-ปฺปฏิภาโน\csromansharedfootnote{abcde61138da3c77}{ยุตฺตปฏิภาโณ (สฺยา) อํ 1. 452 ปิฏฺเฐปิ.} โน มุตฺต\-ปฺปฏิภาโน, อิเธกจฺโจ ปุคฺคโล ปญฺหํ ปุฏฺโฐ สมาโน ยุตฺตํ วทติ โน สีฆํ. อยํ วุจฺจติ ปุคฺคโล ยุตฺต\-ปฺปฏิภาโน โน มุตฺต\-ปฺปฏิภาโน.}

\proseitem{153}{\csromanfolio{148}กตโม จ ปุคฺคโล มุตฺต\-ปฺปฏิภาโน โน ยุตฺต\-ปฺปฏิภาโน, อิเธกจฺโจ ปุคฺคโล ปญฺหํ ปุฏฺโฐ สมาโน สีฆํ วทติ โน ยุตฺตํ. อยํ วุจฺจติ ปุคฺคโล มุตฺต\-ปฺปฏิภาโน โน ยุตฺต\-ปฺปฏิภาโน.}

\proseitem{154}{กตโม จ ปุคฺคโล ยุตฺต\-ปฺปฏิภาโน จ มุตฺตปฺปฏภาโน จ, อิเธกจฺโจ ปุคฺคโล ปญฺหํ ปุฏฺโฐ สมาโน ยุตฺตญฺจ วทติ สีฆญฺจ. อยํ วุจฺจติ ปุคฺคโล ยุตฺต\-ปฺปฏิภาโน จ มุตฺต\-ปฺปฏิภาโน จ.}

\proseitem{155}{กตโม จ ปุคฺคโล เนว ยุตฺต\-ปฺปฏิภาโน โน มุตฺต\-ปฺปฏิภาโน, อิเธกจฺโจ ปุคฺคโล ปญฺหํ ปุฏฺโฐ สมาโน เนว ยุตฺตํ วทติ โน สีฆํ. อยํ วุจฺจติ ปุคฺคโล เนว ยุตฺต\-ปฺปฏิภาโน โน มุตฺต\-ปฺปฏิภาโน.}

\proseitem{156}{กตฺถ กตเม จตฺตาโร ธมฺมกถิกา ปุคฺคลา, อิเธกจฺโจ ธมฺมกถิโก อปฺปญฺจ ภาสติ อสหิตญฺจ, ปริสา จสฺส น กุสลา โหติ สหิตา\-สหิตสฺส, เอวรูโป ธมฺมกถิโก เอวรูปาย ปริสาย ธมฺมกถิโกเตฺวว สงฺขํ คจฺฉติ.}

\prose{อิธ ปเนกจฺโจ ธมฺมกถิโก อปฺปญฺจ ภาสติ สหิตญฺจ, ปริสา จสฺส กุสลา โหติ สหิตา\-สหิตสฺส, เอวรูโป ธมฺมกถิโก เอวรูปาย ปริสาย ธมฺมกถิโกเตฺวว สงฺขํ คจฺฉติ.}

\prose{อิธ ปเนกจฺโจ ธมฺมกถิโก พหุญฺจ ภาสติ อสหิตญฺจ, ปริสา จสฺส น กุสลา โหติ สหิตา\-สหิตสฺส, เอวรูโป ธมฺมกถิโก เอวรูปาย ปริสาย ธมฺมกถิโกเตฺวว สงฺขํ คจฺฉติ.}

\prose{อิธ ปเนกจฺโจ ธมฺมกถิโก พหุญฺจ ภาสติ สหิตญฺจ, ปริสา จสฺส กุสลา โหติ สหิตา\-สหิตสฺส, เอวรูโป ธมฺมกถิโก เอวรูปาย ปริสาย ธมฺมกถิโกเตฺวว สงฺขํ คจฺฉติ. อิเม จตฺตาโร ธมฺมกถิกา ปุคฺคลา.}

\proseitem{157}{ตตฺถ กตเม จตฺตาโร วลาหกูปมา ปุคฺคลา. จตฺตาโร วลาหกา, คชฺชิตา โน วสฺสิตา, วสฺสิตา โน คชฺชิตา, คชฺชิตา จ วสฺสิตา จ, เนว คชฺชิตา โน วสฺสิตา. เอวเมวํ จตฺตาโรเม วลาหกูปมา ปุคฺคลา สนฺโต สํวิชฺชมานา โลกสฺมิํ. กตเม}

\prose{\csromanfolio{149}จตฺตาโร, คชฺชิตา โน วสฺสิตา, วสฺสิตา โน คชฺชิตา, คชฺชิตา จ วสฺสิตา จ, เนว คชฺชิตา โน วสฺสิตา.}

\prose{กถญฺจ ปุคฺคโล คชฺชิตา โหติ โน วสฺสิตา, อิเธกจฺโจ ปุคฺคโล ภาสิตา โหติ โน กตฺตา, เอวํ ปุคฺคโล คชฺชิตา โหติ โน วสฺสิตา. เสยฺยถาปิ โส วลาหโก คชฺชิตา โน วสฺสิตา, ตถูปโม อยํ ปุคฺคโล. (1)}

\prose{กถญฺจ ปุคฺคโล วสฺสิตา โหติ โน คชฺชิตา, อิเธกจฺโจ ปุคฺคโล กตฺตา โหติ โน ภาสิตา, เอวํ ปุคฺคโล วสฺสิตา โหติ โน คชฺชิตา. เสยฺยถาปิ โส วลาหโก วสฺสิตา โน คชฺชิตา, ตถูปโม อยํ ปุคฺคโล. (2)}

\prose{กถญฺจ ปุคฺคโล คชฺชิตา จ โหติ วสฺสิตา จ, อิเธกจฺโจ ปุคฺคโล ภาสิตา จ โหติ กตฺตา จ, เอวํ ปุคฺคโล คชฺชิตา จ โหติ วสฺสิตา จ. เสยฺยถาปิ โส วลาหโก คชฺชิตา จ วสฺสิตา จ, ตถูปโม อยํ ปุคฺคโล. (3)}

\prose{กถญฺจ ปุคฺคโล เนว คชฺชิตา โหติ โน วสฺสิตา, อิเธกจฺโจ ปุคฺคโล เนว ภาสิตา โหติ โน กตฺตา, เอวํ ปุคฺคโล เนว คชฺชิตา โหติ โน วสฺสิตา. เสยฺยถาปิ โส วลาหโก เนว คชฺชิตา โน วสฺสิตา, ตถูปโม อยํ ปุคฺคโล. (4)}

\prose{อิเม จตฺตาโร วลาหกูปมา ปุคฺคลา สนฺโต สํวิชฺชมานา โลกสฺมิํ.}

\proseitem{158}{ตตฺถ กตเม จตฺตาโร มูสิกูปมา ปุคฺคลา. จตสฺโส มูสิกา, คาธํ กตฺตา โน วสิตา, วสิตา โน คาธํ กตฺตา, คาธํ กตฺตา จ วสิตา จ, เนว คาธํ กตฺตา โน วสิตา. เอวเมวํ จตฺตาโรเม มูสิกูปมา ปุคฺคลา สนฺโต สํวิชฺชมานา โลกสฺมิํ. กตเม จตฺตาโร, คาธํ กตฺตา โน วสิตา, วสิตา โน คาธํ กตฺตา, คาธํ กตฺตา จ วสิตา จ, เนว คาธํ กตฺตา โน วสิตา.}

\prose{กถญฺจ ปุคฺคโล คาธํ กตฺตา โหติ โน วสิตา, อิเธกจฺโจ ปุคฺคโล ธมฺมํ ปริยาปุณาติ สุตฺตํ เคยฺยํ เวยฺยากรณํ คาถํ อุทานํ อิติวุตฺตกํ ชาตกํ อพฺภุตธมฺมํ เวทลฺลํ, โส “อิทํ ทุกฺขนฺ”ติ ยถาภูตํ นปฺปชานาติ, “อยํ ทุกฺขสมุทโย”ติ ยถาภูตํ นปฺปชานาติ, “อยํ \csromanfolio{150}ทุกฺขนิโรโธ”ติ ยถาภูตํ นปฺปชานาติ, “อยํ ทุกฺขนิโรธ\-คามินี ปฏิปทา”ติ ยถาภูตํ นปฺปชานาติ, เอวํ ปุคฺคโล คาธํ กตฺตา โหติ โน วสิตา. เสยฺยถาปิ สา มูสิกา คาธํ กตฺตา โน วสิตา, ตถูปโม อยํ ปุคฺคโล. (1)}

\prose{กถญฺจ ปุคฺคโล วสิตา โหติ โน คาธํ กตฺตา, อิเธกจฺโจ ปุคฺคโล ธมฺมํ น ปริยาปุณาติ สุตฺตํ เคยฺยํ เวยฺยากรณํ คาถํ อุทานํ อิติวุตฺตกํ ชาตกํ อพฺภุตธมฺมํ เวทลฺลํ, โส “อิทํ ทุกฺขนฺ”ติ ยถาภูตํ ปชานาติ, “อยํ ทุกฺขสมุทโย”ติ ยถาภูตํ ปชานาติ, “อยํ ทุกฺขนิโรโธ”ติ ยถาภูตํ ปชานาติ, “อยํ ทุกฺขนิโรธ\-คามินี ปฏิปทา”ติ ยถาภูตํ ปชานาติ, เอวํ ปุคฺคโล วสิตา โหติ โน คาธํ กตฺตา. เสยฺยถาปิ สา มูสิกา วสิตา โน คาธํ กตฺตา, ตถูปโม อยํ ปุคฺคโล. (2)}

\prose{กถญฺจ ปุคฺคโล คาธํ กตฺตา จ โหติ วสิตา จ, อิเธกจฺโจ ปุคฺคโล ธมฺมํ ปริยาปุณาติ สุตฺตํ เคยฺยํ เวยฺยากรณํ คาถํ อุทานํ อิติวุตฺตกํ ชาตกํ อพฺภุตธมฺมํ เวทลฺลํ, โส “อิทํ ทุกฺขนฺ”ติ ยถาภูตํ ปชานาติ, “อยํ ทุกฺขสมุทโย”ติ ยถาภูตํ ปชานาติ, “อยํ ทุกฺขนิโรโธ”ติ ยถาภูตํ ปชานาติ, “อยํ ทุกฺขนิโรธ\-คามินี ปฏิปทา”ติ ยถาภูตํ ปชานาติ, เอวํ ปุคฺคโล คาธํ กตฺตา จ โหติ วสิตา จ. เสยฺยถาปิ สา มูสิกา คาธํ กตฺตา จ วสิตา จ, ตถูปโม อยํ ปุคฺคโล. (3)}

\prose{กถญฺจ ปุคฺคโล เนว คาธํ กตฺตา โหติ โน วสิตา, อิเธกจฺโจ ปุคฺคโล ธมฺมํ น ปริยาปุณาติ สุตฺตํ เคยฺยํ เวยฺยากรณํ คาถํ อุทานํ อิติวุตฺตกํ ชาตกํ อพฺภุตธมฺมํ เวทลฺลํ, โส “อิทํ ทุกฺขนฺ”ติ ยถาภูตํ นปฺปชานาติ, “อยํ ทุกฺขสมุทโย”ติ ยถาภูตํ นปฺปชานาติ, “อยํ ทุกฺขนิโรโธ”ติ ยถาภูตํ นปฺปชานาติ, “อยํ ทุกฺขนิโรธ\-คามินี ปฏิปทา”ติ ยถาภูตํ นปฺปชานาติ, เอวํ ปุคฺคโล เนว คาธํ กตฺตา โหติ โน วสิตา. เสยฺยถาปิ สา มูสิกา เนว คาธํ กตฺตา โน วสิตา, ตถูปโม อยํ ปุคฺคโล. (4)}

\prose{อิเม จตฺตาโร มูสิกูปมา ปุคฺคลา สนฺโต สํวิชฺชมานา โลกสฺมิํ.}

\proseitem{159}{\csromanfolio{151}ตตฺถ กตเม จตฺตาโร อมฺพูปมา ปุคฺคลา. จตฺตาริ อมฺพานิ, อามํ ปกฺกวณฺณิ\csromansharedfootnote{797ee7fb49228b7e}{ปกฺกวณฺณี}, ปกฺกํ อามวณฺณิ\csromansharedfootnote{5018aaa5197b2cfc}{อามวณฺณี (สฺยา, ก) อํ 1. 421 ปิฏฺเฐปิ.}, อามํ อามวณฺณิ, ปกฺกํ ปกฺกวณฺณิ. เอวเมวํ จตฺตาโรเม อมฺพูปมา ปุคฺคลา สนฺโต สํวิชฺชมานา โลกสฺมิํ. กตเม จตฺตาโร, อาโม ปกฺกวณฺณี, ปกฺโก อามวณฺณี, อาโม อามวณฺณี, ปกฺโก ปกฺกวณฺณี.}

\prose{กถญฺจ ปุคฺคโล อาโม โหติ ปกฺกวณฺณี, อิเธกจฺจสฺส ปุคฺคลสฺส ปาสาทิกํ โหติ อภิกฺกนฺตํ ปฏิกฺกนฺตํ อาโลกิตํ วิโลกิตํ สมิญฺชิตํ\csromansharedfootnote{6a6ed8b24ec6283c}{สมฺมิญฺชิตํ (สี, สฺยา)} ปสาริตํ สงฺฆาฏิ\-ปตฺต\-จีวร\-ธารณํ, โส “อิทํ ทุกฺขนฺ”ติ ยถาภูตํ นปฺปชานาติ, “อยํ ทุกฺขสมุทโย”ติ ยถาภูตํ นปฺปชานาติ, “อยํ ทุกฺขนิโรโธ”ติ ยถาภูตํ นปฺปชานาติ, “อยํ ทุกฺขนิโรธ\-คามินี ปฏิปทา”ติ ยถาภูตํ นปฺปชานาติ. เอวํ ปุคฺคโล อาโม โหติ ปกฺกวณฺณี. เสยฺยถาปิ ตํ อมฺพํ อามํ ปกฺกวณฺณิ, ตถูปโม อยํ ปุคฺคโล. (1)}

\prose{กถญฺจ ปุคฺคโล ปกฺโก โหติ อามวณฺณี, อิเธกจฺจสฺส ปุคฺคลสฺส น ปาสาทิกํ โหติ อภิกฺกนฺตํ ปฏิกฺกนฺตํ อาโลกิตํ วิโลกิตํ สมิญฺชิตํ ปสาริตํ สงฺฆาฏิ\-ปตฺต\-จีวร\-ธารณํ, โส “อิทํ ทุกฺขนฺ”ติ ยถาภูตํ ปชานาติ, “อยํ ทุกฺขสมุทโย”ติ ยถาภูตํ ปชานาติ, “อยํ ทุกฺขนิโรโธ”ติ ยถาภูตํ ปชานาติ, “อยํ ทุกฺขนิโรธ\-คามินี ปฏิปทา”ติ ยถาภูตํ ปชานาติ. เอวํ ปุคฺคโล ปกฺโก โหติ อามวณฺณี เสยฺยถาปิ ตํ อมฺพํ ปกฺกํ อามวณฺณิ, ตถูปโม อยํ ปุคฺคโล. (2)}

\prose{กถญฺจ ปุคฺคโล อาโม โหติ อามวณฺณี, อิเธกจฺจสฺส ปุคฺคลสฺส น ปาสาทิกํ โหติ อภิกฺกนฺตํ ปฏิกฺกนฺตํ อาโลกิตํ วิโลกิตํ สมิญฺชิตํ ปสาริตํ สงฺฆาฏิ\-ปตฺต\-จีวร\-ธารณํ, โส “อิทํ ทุกฺขนฺ”ติ ยถาภูตํ นปฺปชานาติ ฯเปฯ “อยํ ทุกฺขนิโรธ\-คามินี ปฏิปทา”ติ ยถาภูตํ นปฺปชานาติ. เอวํ ปุคฺคโล อาโม โหติ อามวณฺณี. เสยฺยถาปิ ตํ อมฺพํ อามํ อามวณฺณิ, ตถูปโม อยํ ปุคฺคโล. (3)}

\prose{กถญฺจ ปุคฺคโล ปกฺโก โหติ ปกฺกวณฺณี, อิเธกจฺจสฺส ปุคฺคลสฺส ปาสาทิกํ โหติ อภิกฺกนฺตํ ปฏิกฺกนฺตํ อาโลกิตํ วิโลกิตํ สมิญฺชิตํ ปสาริตํ สงฺฆาฏิ\-ปตฺต\-จีวร\-ธารณํ, โส “อิทํ ทุกฺขนฺ”ติ ยถาภูตํ \csromanfolio{152}ปชานาติ ฯเปฯ “อยํ ทุกฺขนิโรธ\-คามินี ปฏิปทา”ติ ยถาภูตํ ปชานาติ. เอวํ ปุคฺคโล ปกฺโก โหติ ปกฺกวณฺณี. เสยฺยถาปิ ตํ อมฺพํ ปกฺกํ ปกฺกวณฺณิ, ตถูปโม อยํ ปุคฺคโล. (4)}

\prose{อิเม จตฺตาโร อมฺพูปมา ปุคฺคลา สนฺโต สํวิชฺชมานา โลกสฺมิํ.}

\proseitem{160}{ตตฺถ กตเม จตฺตาโร กุมฺภูปมา ปุคฺคลา. จตฺตาโร กุมฺภา, ตุจฺโฉ ปิหิโต, ปูโร วิวโฏ, ตุจฺโฉ วิวโฏ, ปูโร ปิหิโต. เอวเมวํ จตฺตาโรเม กุมฺภูปมา ปุคฺคลา สนฺโต สํวิชฺชมานา โลกสฺมิํ. กตเม จตฺตาโร, ตุจฺโฉ ปิหิโต, ปูโร วิวโฏ, ตุจฺโฉ วิวโฏ, ปูโร ปิหิโต.}

\prose{กถญฺจ ปุคฺคโล ตุจฺโฉ โหติ ปิหิโต, อิเธกจฺจสฺส ปุคฺคลสฺส ปาสาทิกํ โหติ อภิกฺกนฺตํ ปฏิกฺกนฺตํ อาโลกิตํ วิโลกิตํ สมิญฺชิตํ ปสาริตํ สงฺฆาฏิ\-ปตฺต\-จีวร\-ธารณํ, โส “อิทํ ทุกฺขนฺ”ติ ยถาภูตํ นปฺปชานาติ ฯเปฯ “อยํ ทุกฺขนิโรธ\-คามินี ปฏิปทา”ติ ยถาภูตํ นปฺปชานาติ. เอวํ ปุคฺคโล ตุจฺโฉ โหติ ปิหิโต. เสยฺยถาปิ โส กุมฺโภ ตุจฺโฉ ปิหิโต, ตถูปโม อยํ ปุคฺคโล.}

\prose{กถญฺจ ปุคฺคโล ปูโร โหติ วิวโฏ, อิเธกจฺจสฺส ปุคฺคลสฺส น ปาสาทิกํ โหติ อภิกฺกนฺตํ ปฏิกฺกนฺตํ อาโลกิตํ วิโลกิตํ สมิญฺชิตํ ปสาริตํ สงฺฆาฏิ\-ปตฺต\-จีวร\-ธารณํ, โส “อิทํ ทุกฺขนฺ”ติ ยถาภูตํ ปชานาติ ฯเปฯ “อยํ ทุกฺขนิโรธ\-คามินี ปฏิปทา”ติ ยถาภูตํ ปชานาติ. เอวํ ปุคฺคโล ปูโร โหติ วิวโฏ. เสยฺยถาปิ โส กุมฺโภ ปูโร วิวโฏ, ตถูปโม อยํ ปุคฺคโล.}

\prose{กถญฺจ ปุคฺคโล ตุจฺโฉ โหติ วิวโฏ, อิเธกจฺจสฺส ปุคฺคลสฺส น ปาสาทิกํ โหติ อภิกฺกนฺตํ ปฏิกฺกนฺตํ อาโลกิตํ วิโลกิตํ สมิญฺชิตํ ปสาริตํ สงฺฆาฏิ\-ปตฺต\-จีวร\-ธารณํ, โส “อิทํ ทุกฺขนฺ”ติ ยถาภูตํ นปฺปชานาติ ฯเปฯ “อยํ ทุกฺขนิโรธ\-คามินี ปฏิปทา”ติ ยถาภูตํ นปฺปชานาติ. เอวํ ปุคฺคโล ตุจฺโฉ โหติ วิวโฏ. เสยฺยถาปิ โส กุมฺโภ ตุจฺโฉ วิวโฏ, ตถูปโม อยํ ปุคฺคโล.}

\prose{กถญฺจ ปุคฺคโล ปูโร โหติ ปิหิโต, อิเธกจฺจสฺส ปุคฺคลสฺส ปาสาทิกํ โหติ อภิกฺกนฺตํ ปฏิกฺกนฺตํ อาโลกิตํ วิโลกิตํ}

\prose{\csromanfolio{153}สมิญฺชิตํ ปสาริตํ สงฺฆาฏิ\-ปตฺต\-จีวร\-ธารณํ, โส “อิทํ ทุกฺขนฺ”ติ ยถาภูตํ ปชานาติ ฯเปฯ “อยํ ทุกฺขนิโรธ\-คามินี ปฏิปทา”ติ ยถาภูตํ ปชานาติ. เอวํ ปุคฺคโล ปูโร โหติ ปิหิโต. เสยฺยถาปิ โส กุมฺโภ ปูโร ปิหิโต, ตถูปโม อยํ ปุคฺคโล. อิเม จตฺตาโร กุมฺภูปมา ปุคฺคลา สนฺโต สํวิชฺชมานา โลกสฺมิํ.}

\proseitem{161}{ตตฺถ กตเม จตฺตาโร อุทกรหทู\-ปมา ปุคฺคลา. จตฺตาโร อุทกรหทา\csromandash{}อุตฺตาโน คมฺภีโรภาโส, คมฺภีโร อุตฺตาโนภาโส, อุตฺตาโน อุตฺตาโนภาโส, คมฺภีโร คมฺภีโรภาโส. เอวเมวํ จตฺตาโรเม อุทกรหทู\-ปมา ปุคฺคลา สนฺโต สํวิชฺชมานา โลกสฺมิํ. กตเม จตฺตาโร, อุตฺตาโน คมฺภีโรภาโส, คมฺภีโร อุตฺตาโนภาโส, อุตฺตาโน อุตฺตาโนภาโส, คมฺภีโร คมฺภีโรภาโส.}

\prose{กถญฺจ ปุคฺคโล อุตฺตาโน โหติ คมฺภีโรภาโส, อิเธกจฺจสฺส ปุคฺคลสฺส ปาสาทิกํ โหติ อภิกฺกนฺตํ ปฏิกฺกนฺตํ อาโลกิตํ วิโลกิตํ สมิญฺชิตํ ปสาริตํ สงฺฆาฏิ\-ปตฺต\-จีวร\-ธารณํ, โส “อิทํ ทุกฺขนฺ”ติ ยถาภูตํ นปฺปชานาติ ฯเปฯ “อยํ ทุกฺขนิโรธ\-คามินี ปฏิปทา”ติ ยถาภูตํ นปฺปชานาติ. เอวํ ปุคฺคโล อุตฺตาโน โหติ คมฺภีโรภาโส. เสยฺยถาปิ โส อุทกรหโท อุตฺตาโน คมฺภีโรภาโส, ตถูปโม อยํ ปุคฺคโล.}

\prose{กถญฺจ ปุคฺคโล คมฺภีโร โหติ อุตฺตาโนภาโส, อิเธกจฺจสฺส ปุคฺคลสฺส น ปาสาทิกํ โหติ อภิกฺกนฺตํ ปฏิกฺกนฺตํ อาโลกิตํ วิโลกิตํ สมิญฺชิตํ ปสาริตํ สงฺฆาฏิ\-ปตฺต\-จีวร\-ธารณํ, โส “อิทํ ทุกฺขนฺ”ติ ยถาภูตํ ปชานาติ ฯเปฯ “อยํ ทุกฺขนิโรธ\-คามินี ปฏิปทา”ติ ยถาภูตํ ปชานาติ. เอวํ ปุคฺคโล คมฺภีโร โหติ อุตฺตาโนภาโส. เสยฺยถาปิ โส อุทกรหโท คมฺภีโร อุตฺตาโนภาโส, ตถูปโม อยํ ปุคฺคโล.}

\prose{กถญฺจ ปุคฺคโล อุตฺตาโน โหติ อุตฺตาโนภาโส, อิเธกจฺจสฺส ปุคฺคลสฺส น ปาสาทิกํ โหติ อภิกฺกนฺตํ ปฏิกฺกนฺตํ อาโลกิตํ วิโลกิตํ สมิญฺชิตํ ปสาริตํ สงฺฆาฏิ\-ปตฺต\-จีวร\-ธารณํ, โส “อิทํ ทุกฺขนฺ”ติ ยถาภูตํ นปฺปชานาติ ฯเปฯ “อยํ ทุกฺขนิโรธ\-คามินี ปฏิปทา”ติ \csromanfolio{154}ยถาภูตํ นปฺปชาติ. เอวํ ปุคฺคโล อุตฺตาโน โหติ อุตฺตาโนภาโส. เสยฺยถาปิ โส อุทกรหโท อุตฺตาโน อุตฺตาโนภาโส, ตถูปโม อยํ ปุคฺคโล.}

\prose{กถญฺจ ปุคฺคโล คมฺภีโร โหติ คมฺภีโรภาโส, อิเธกจฺจสฺส ปุคฺคลสฺส ปาสาทิกํ โหติ อภิกฺกนฺตํ ปฏิกฺกนฺตํ อาโลกิตํ วิโลกิตํ สมิญฺชิตํ ปสาริตํ สงฺฆาฏิ\-ปตฺต\-จีวร\-ธารณํ, โส “อิทํ ทุกฺขนฺ”ติ ยถาภูตํ ปชานาติ ฯเปฯ “อยํ ทุกฺขนิโรธ\-คามินี ปฏิปทา”ติ ยถาภูตํ ปชานาติ. เอวํ ปุคฺคโล คมฺภีโร โหติ คมฺภีโรภาโส. เสยฺยถาปิ โส อุทกรหโท คมฺภีโร คมฺภีโรภาโส, ตถูปโม อยํ ปุคฺคโล. อิเม จตฺตาโร อุทกรหทู\-ปมา ปุคฺคลา สนฺโต สํวิชฺชมานา โลกสฺมิํ.}

\proseitem{162}{ตตฺถ กตเม จตฺตาโร พลีพทฺทูปมา ปุคฺคลา. จตฺตาโร พลีพทฺทา\csromansharedfootnote{42bf05725d868dfd}{พลิพทฺธา (สฺยา)}\csromandash{}สกควจณฺโฑ\csromansharedfootnote{53a46b5a646bd770}{สควจณฺโฑ (กสี) อํ 1. 423 ปิฏฺเฐปิ.} โน ปรควจณฺโฑ, ปรควจณฺโฑ โน สกควจณฺโฑ, สกควจณฺโฑ จ ปรควจณฺโฑ จ, เนว สกควจณฺโฑ โน ปรควจณฺโฑ. เอวเมวํ จตฺตาโรเม พลีพทฺทูปมา ปุคฺคลา สนฺโต สํวิชฺชมานา โลกสฺมิํ. กตเม จตฺตาโร, สกควจณฺโฑ โน ปรควจณฺโฑ, ปรควจณฺโฑ โน สกควจณฺโฑ, สกควจณฺโฑ จ ปรควจณฺโฑ จ, เนว สกควจณฺโฑ โน ปรควจณฺโฑ.}

\prose{กถญฺจ ปุคฺคโล สกควจณฺโฑ โหติ โน ปรควจณฺโฑ, อิเธกจฺโจ ปุคฺคโล สกปริสํ อุพฺเพชิตา โหติ โน ปรปริสํ. เอวํ ปุคฺคโล สกควจณฺโฑ โหติ โน ปรควจณฺโฑ. เสยฺยถาปิ โส พลีพทฺโท สกควจณฺโฑ โน ปรควจณฺโฑ, ตถูปโม อยํ ปุคฺคโล.}

\prose{กถญฺจ ปุคฺคโล ปรควจณฺโฑ โหติ โน สกควจณฺโฑ, อิเธกจฺโจ ปุคฺคโล ปรปริสํ อุพฺเพชิตา โหติ โน สกปริสํ. เอวํ ปุคฺคโล ปรควจณฺโฑ โหติ โน สกควจณฺโฑ. เสยฺยถาปิ โส พลีพทฺโท ปรควจณฺโฑ โน สกควจณฺโฑ, ตถูปโม อยํ ปุคฺคโล.}

\prose{กถญฺจ ปุคฺคโล สกควจณฺโฑ จ โหติ ปรควจณฺโฑ จ, อิเธกจฺโจ ปุคฺคโล สกปริสญฺจ อุพฺเพชิตา โหติ ปรปริสญฺจ. เอวํ}

\prose{\csromanfolio{155}ปุคฺคโล สกควจณฺโฑ จ โหติ ปรควจณฺโฑ จ. เสยฺยถาปิ โส พลีพทฺโท สกควจณฺโฑ จ ปรควจณฺโฑ จ, ตถูปโม อยํ ปุคฺคโล.}

\prose{กถญฺจ ปุคฺคโล เนว สกควจณฺโฑ โหติ โน ปรควจณฺโฑ, อิเธกจฺโจ ปุคฺคโล เนว สกปริสํ อุพฺเพชิตา โหติ โน ปรปริสํ. เอวํ ปุคฺคโล เนว สกควจณฺโฑ โหติ โน ปรควจณฺโฑ. เสยฺยถาปิ โส พลีพทฺโท เนว สกควจณฺโฑ โน ปรควจณฺโฑ, ตถูปโม อยํ ปุคฺคโล. อิเม จตฺตาโร พลีพทฺทูปมา ปุคฺคลา สนฺโต สํวิชฺชมานา โลกสฺมิํ.}

\proseitem{163}{ตตฺถ กตเม จตฺตาโร อาสีวิสูปมา ปุคฺคลา. จตฺตาโร อาสีวิสา\csromansharedfootnote{f64d1c07819c5aaa}{อาสิวิสา (สฺยา)}\csromandash{} อาคตวิโส โน โฆรวิโส, โฆรวิโส โน อาคตวิโส, อาคตวิโส จ โฆรวิโส จ, เนว อาคตวิโส โน โฆรวิโส. เอวเมวํ จตฺตาโรเม อาสีวิสูปมา ปุคฺคลา สนฺโต สํวิชฺชมานา โลกสฺมิํ. กตเม จตฺตาโร, อาคตวิโส โน โฆรวิโส, โฆรวิโส โน อาคตวิโส, อาคตวิโส จ โฆรวิโส จ, เนว อาคตวิโส โน โฆรวิโส.}

\prose{กถญฺจ ปุคฺคโล อาคตวิโส โหติ โน โฆรวิโส, อิเธกจฺโจ ปุคฺคโล อภิณฺหํ กุชฺฌติ, โส จ ขฺวสฺส โกโธ น จิรํ ทีฆรตฺตํ อนุเสติ. เอวํ ปุคฺคโล อาคตวิโส โหติ โน โฆรวิโส. เสยฺยถาปิ โส อาสีวิโส อาคตวิโส โน โฆรวิโส, ตถูปโม อยํ ปุคฺคโล.}

\prose{กถญฺจ ปุคฺคโล โฆรวิโส โหติ โน อาคตวิโส, อิเธกจฺโจ ปุคฺคโล นเหว โข\csromansharedfootnote{0f1e3b79200f8cbb}{เนว โข (สี) อํ 1.426 ปิฏฺเฐปิ.} อภิณฺหํ กุชฺฌติ, โส จ ขฺวสฺส โกโธ จิรํ ทีฆรตฺตํ อนุเสติ. เอวํ ปุคฺคโล โฆรวิโส โหติ โน อาคตวิโส. เสยฺยถาปิ โส อาสีวิโส โฆรวิโส โน อาคตวิโส, ตถูปโม อยํ ปุคฺคโล.}

\prose{กถญฺจ ปุคฺคโล อาคตวิโส จ โหติ โฆรวิโส จ, อิเธกจฺโจ ปุคฺคโล อภิณฺหํ กุชฺฌติ, โส จ ขฺวสฺส โกโธ จิรํ ทีฆรตฺตํ อนุเสติ. เอวํ ปุคฺคโล อาคตวิโส จ โหติ โฆรวิโส จ. \csromanfolio{156}เสยฺยถาปิ โส อาสีวิโส อาคตวิโส จ โฆรวิโส จ, ตถูปโม อยํ ปุคฺคโล.}

\prose{กถญฺจ ปุคฺคโล เนว อาคตวิโส โหติ โน โฆรวิโส, อิเธกจฺโจ ปุคฺคโล นเหว โข อภิณฺหํ กุชฺฌติ, โส จ ขฺวสฺส โกโธ น จิรํ ทีฆรตฺตํ อนุเสติ. เอวํ ปุคฺคโล เนว อาคตวิโส โหติ โน โฆรวิโส. เสยฺยถาปิ โส อาสีวิโส เนว อาคตวิโส โน โฆรวิโส, ตถูปโม อยํ ปุคฺคโล. อิเม จตฺตาโร อาสีวิสูปมา ปุคฺคลา สนฺโต สํวิชฺชมานา โลกสฺมิํ.}

\proseitem{164}{กถญฺจ ปุคฺคโล อนนุวิจฺจ อปริโยคาเหตฺวา อวณฺณารหสฺส วณฺณํ ภาสิตา โหติ, อิเธกจฺโจ ปุคฺคโล ทุปฺปฏิปนฺนานํ มิจฺฉา\-ปฏิปนฺนานํ ติตฺถิยานํ ติตฺถิย\-สาวกานํ วณฺณํ ภาสติ “สุปฺปฏิปนฺนา อิติปิ ‘สมฺมา ปฏิปนฺนา’ อิติปี”ติ. เอวํ ปุคฺคโล อนนุวิจฺจ อปริโยคาเหตฺวา อวณฺณารหสฺส วณฺณํ ภาสิตา โหติ. (1)}

\prose{กถญฺจ ปุคฺคโล อนนุวิจฺจ อปริโยคาเหตฺวา วณฺณารหสฺส อวณฺณํ ภาสิตา โหติ, อิเธกจฺโจ ปุคฺคโล สุปฺปฏิปนฺนานํ สมฺมา\-ปฏิปนฺนานํ พุทฺธานํ พุทฺธ\-สาวกานํ อวณฺณํ ภาสติ “ทุปฺปฏิปนฺนา อิติปิ ‘มิจฺฉา\-ปฏิปนฺนา’ อิติปี”ติ. เอวํ ปุคฺคโล อนนุวิจฺจ อปริโยคาเหตฺวา วณฺณารหสฺส อวณฺณํ ภาสิตา โหติ. (2)}

\prose{กถญฺจ ปุคฺคโล อนนุวิจฺจ อปริโยคาเหตฺวา อปฺปสาทนีเย ฐาเน ปสาทํ อุปทํสิตา โหติ, อิเธกจฺโจ ปุคฺคโล ทุปฺปฏิปทาย มิจฺฉา\-ปฏิปทาย ปสาทํ ชเนติ “สุปฺปฏิปทา อิติปิ ‘สมฺมาปฏิปทา’ อิติปี”ติ. เอวํ ปุคฺคโล อนนุวิจฺจ อปริโยคาเหตฺวา อปฺปสาทนีเย ฐาเน ปสาทํ อุปทํสิตา โหติ. (3)}

\prose{กถญฺจ ปุคฺคโล อนนุวิจฺจ อปริโยคาเหตฺวา ปสาทนีเย ฐาเน อปฺปสาทํ อุปทํสิตา โหติ, อิเธกจฺโจ ปุคฺคโล สุปฺปฏิปทาย สมฺมา\-ปฏิปทาย อปฺปสาทํ ชเนติ “ทุปฺปฏิปทา อิติปิ ‘มิจฺฉาปฏิปทา’ อิติปี”ติ. เอวํ ปุคฺคโล อนนุวิจฺจ อปริโยคาเหตฺวา ปสาทนีเย ฐาเน อปฺปสาทํ อุปทํสิตา โหติ. (4)}

\proseitem{165}{\csromanfolio{157}กถญฺจ ปุคฺคโล อนุวิจฺจ ปริโยคาเหตฺวา อวณฺณารหสฺส อวณฺณํ ภาสิตา โหติ, อิเธกจฺโจ ปุคฺคโล ทุปฺปฏิปนฺนานํ มิจฺฉา\-ปฏิปนฺนานํ ติตฺถิยานํ ติตฺถิย\-สาวกานํ อวณฺณํ ภาสติ “ทุปฺปฏิปนฺนา อิติปิ ‘มิจฺฉา\-ปฏิปนฺนา’ อิติปี”ติ. เอวํ ปุคฺคโล อนุวิจฺจ ปริโยคาเหตฺวา อวณฺณารหสฺส อวณฺณํ ภาสิตา โหติ. (1)}

\prose{กถญฺจ ปุคฺคโล อนุวิจฺจ ปริโยคาเหตฺวา วณฺณารหสฺส วณฺณํ ภาสิตา โหติ, อิเธกจฺโจ ปุคฺคโล สุปฺปฏิปนฺนานํ สมฺมา\-ปฏิปนฺนานํ พุทฺธานํ พุทฺธ\-สาวกานํ วณฺณํ ภาสติ “สุปฺปฏิปนฺนา อิติปิ ‘สมฺมาปฏิปนฺนา’ อิติปี”ติ. เอวํ ปุคฺคโล อนุวิจฺจ ปริโยคาเหตฺวา วณฺณารหสฺส วณฺณํ ภาสิตา โหติ. (2)}

\prose{กถญฺจ ปุคฺคโล อนุวิจฺจ ปริโยคาเหตฺวา อปฺปสาทนีเย ฐาเน อปฺปสาทํ อุปทํสิตา โหติ, อิเธกจฺโจ ปุคฺคโล ทุปฺปฏิปทาย มิจฺฉา\-ปฏิปทาย อปฺปสาทํ ชเนติ “ทุปฺปฏิปทา อิติปิ ‘มิจฺฉาปฏิปทา’ อิติปี”ติ. เอวํ ปุคฺคโล อนุวิจฺจ ปริโยคาเหตฺวา อปฺปสาทนีเย ฐาเน อปฺปสาทํ อุปทํสิตา โหติ. (3)}

\prose{กถญฺจ ปุคฺคโล อนุวิจฺจ ปริโยคาเหตฺวา ปสาทนีเย ฐาเน ปสาทํ อุปทํสิตา โหติ, อิเธกจฺโจ ปุคฺคโล สุปฺปฏิปทาย สมฺมา\-ปฏิปทาย ปสาทํ ชเนติ “สุปฺปฏิปทา อิติปิ ‘สมฺมาปฏิปทา’ อิติปี”ติ. เอวํ ปุคฺคโล อนุวิจฺจ ปริโยคาเหตฺวา ปสาทนีเย ฐาเน ปสาทํ อุปทํสิตา โหติ. (4)}

\proseitem{166}{กถญฺจ ปุคฺคโล อวณฺณารหสฺส อวณฺณํ ภาสิตา โหติ ภูตํ ตจฺฉํ กาเลน, โน จ โข วณฺณารหสฺส วณฺณํ ภาสิตา โหติ ภูตํ ตจฺฉํ กาเลน, อิเธกจฺโจ ปุคฺคโล วณฺโณปิ สํวิชฺชติ, อวณฺโณปิ สํวิชฺชติ. โย ตตฺถ อวณฺโณ, ตํ ภณติ ภูตํ ตจฺฉํ กาเลน. โย ตตฺถ วณฺโณ, ตํ น ภณติ ภูตํ ตจฺฉํ กาเลน. เอวํ ปุคฺคโล อวณฺณารหสฺส อวณฺณํ ภาสิตา โหติ ภูตํ ตจฺฉํ กาเลน, โน จ โข วณฺณารหสฺส วณฺณํ ภาสิตา โหติ ภูตํ ตจฺฉํ กาเลน. (1)}

\prose{กถญฺจ ปุคฺคโล วณฺณารหสฺส วณฺณํ ภาสิตา โหติ ภูตํ ตจฺฉํ กาเลน, โน จ โข อวณฺณารหสฺส อวณฺณํ ภาสิตา โหติ ภูตํ ตจฺฉํ กาเลน, อิเธกจฺโจ ปุคฺคโล วณฺโณปิ สํวิชฺชติ, อวณฺโณปิ \csromanfolio{158}สํวิชฺชติ. โย ตตฺถ วณฺโณ, ตํ ภณติ ภูตํ ตจฺฉํ กาเลน. โย ตตฺถ อวณฺโณ, ตํ น ภณติ ภูตํ ตจฺฉํ กาเลน. เอวํ ปุคฺคโล วณฺณารหสฺส วณฺณํ ภาสิตา โหติ ภูตํ ตจฺฉํ กาเลน, โน จ โข อวณฺณารหสฺส อวณฺณํ ภาสิตา โหติ ภูตํ ตจฺฉํ กาเลน. (2)}

\prose{กถญฺจ ปุคฺคโล อวณฺณารหสฺส จ อวณฺณํ ภาสิตา โหติ ภูตํ ตจฺฉํ กาเลน, วณฺณารหสฺส จ วณฺณํ ภาสิตา โหติ ภูตํ ตจฺฉํ กาเลน, อิเธกจฺโจ ปุคฺคโล วณฺโณปิ สํวิชฺชติ, อวณฺโณปิ สํวิชฺชติ. โย ตตฺถ อวณฺโณ, ตํ ภณติ ภูตํ ตจฺฉํ กาเลน. โยปิ ตตฺถ วณฺโณ, ตมฺปิ ภณติ ภูตํ ตจฺฉํ กาเลน. ตตฺร กาลญฺญู โหติ ตสฺส ปญฺหสฺส เวยฺยากรณาย. เอวํ ปุคฺคโล อวณฺณารหสฺส จ อวณฺณํ ภาสิตา โหติ ภูตํ ตจฺฉํ กาเลน, วณฺณารหสฺส จ วณฺณํ ภาสิตา โหติ ภูตํ ตจฺฉํ กาเลน. (3)}

\prose{กถญฺจ ปุคฺคโล เนว อวณฺณารหสฺส อวณฺณํ ภาสิตา โหติ ภูตํ ตจฺฉํ กาเลน, โนปิ วณฺณารหสฺส วณฺณํ ภาสิตา โหติ ภูตํ ตจฺฉํ กาเลน, อิเธกจฺโจ ปุคฺคโล วณฺโณปิ สํวิชฺชติ, อวณฺโณปิ สํวิชฺชติ. โย ตตฺถ อวณฺโณ, ตํ น ภณติ ภูตํ ตจฺฉํ กาเลน. โยปิ ตตฺถ วณฺโณ, ตมฺปิ น ภณติ ภูตํ ตจฺฉํ กาเลน. อุเปกฺขโก วิหรติ สโต สมฺปชาโน. เอวํ ปุคฺคโล เนว อวณฺณารหสฺส อวณฺณํ ภาสิตา โหติ ภูตํ ตจฺฉํ กาเลน, โนปิ วณฺณารหสฺส วณฺณํ ภาสิตา โหติ ภูตํ ตจฺฉํ กาเลน. (4)}

\proseitem{167}{กตโม จ ปุคฺคโล อุฏฺฐาน\-ผลูปชีวี โน ปุญฺญผลู\-ปชีวี, ยสฺส ปุคฺคลสฺส อุฏฺฐหโต ฆฏโต วายมโต อาชีโว อภินิพฺพตฺตติ โน ปุญฺญโต. อยํ วุจฺจติ ปุคฺคโล อุฏฺฐาน\-ผลูปชีวี โน ปุญฺญผลู\-ปชีวี.}

\prose{กตโม จ ปุคฺคโล ปุญฺญผลู\-ปชีวี โน อุฏฺฐาน\-ผลูปชีวี, ปรนิมฺมิต\-วส\-วตฺตี เทเว\csromansharedfootnote{98cd814497a2437c}{ปรนิมฺมิตวสวตฺติเทเว (สี, สฺยา)} อุปาทาย ตตูปริ เทวา ปุญฺญผลู\-ปชีวิโน น อุฏฺฐาน\-ผลูปชีวิโน.}

\prose{\csromanfolio{159}กตโม จ ปุคฺคโล อุฏฺฐาน\-ผลูปชีวี จ ปุญฺญผลู\-ปชีวี จ, ยสฺส ปุคฺคลสฺส อุฏฺฐหโต ฆฏโต วายมโต อาชีโว อภินิพฺพตฺตติ ปุญฺญโต จ. อยํ วุจฺจติ ปุคฺคโล อุฏฺฐาน\-ผลูปชีวี จ ปุญฺญผลู\-ปชีวี จ.}

\prose{กตโม จ ปุคฺคโล เนว อุฏฺฐาน\-ผลูปชีวี โน ปุญฺญผลู\-ปชีวี เนรยิกา เนว อุฏฺฐาน\-ผลูปชีวิโน โน ปุญฺญผลู\-ปชีวิโน.}

\proseitem{168}{กถญฺจ ปุคฺคโล ตโม โหติ ตมปรายโน, อิเธกจฺโจ ปุคฺคโล นีเจ กุเล ปจฺจาชาโต โหติ จณฺฑาลกุเล วา เนสาทกุเล วา เวนกุเล\csromansharedfootnote{66d4fa4e1648d9f6}{เวณกุเล (สี, สฺยา)} วา รถการกุเล วา ปุกฺกุสกุเล วา ทลิทฺเท\csromansharedfootnote{ce85339a021fe0b0}{ทฬิทฺเท (สี) อํ 1. 497 ปิฏฺเฐปิ.} อปฺป\-นฺน\-ปาน\-โภชเน กสิรวุตฺติเก, ยตฺถ กสิเรน ฆาสจฺฉาโท ลพฺภติ, โส จ โหติ ทุพฺพณฺโณ ทุทฺทสิโก โอโกฏิมโก พหฺวาพาโธ กาโณ วา กุณี วา ขญฺโช วา ปกฺขหโต วา, น ลาภี อนฺนสฺส ปานสฺส วตฺถสฺส ยานสฺส มาลา\-คนฺธ\-วิเลปนสฺส เสยฺยา\-วสถ\-ปทีเปยฺยสฺส. โส กาเยน ทุจฺจริตํ จรติ, วาจาย ทุจฺจริตํ จรติ, มนสา ทุจฺจริตํ จรติ. โส กาเยน ทุจฺจริตํ จริตฺวา วาจาย ทุจฺจริตํ จริตฺวา มนสา ทุจฺจริตํ จริตฺวา กายสฺส เภทา ปรํ มรณา อปายํ ทุคฺคติํ วินิปาตํ นิรยํ อุปปชฺชติ. เอวํ ปุคฺคโล ตโม โหติ ตมปรายโน. (1)}

\prose{กถญฺจ ปุคฺคโล ตโม โหติ โชติปรายโน, อิเธกจฺโจ ปุคฺคโล นีเจ กุเล ปจฺจาชาโต โหติ จณฺฑาลกุเล วา เนสาทกุเล วา เวนกุเล วา รถการกุเล วา ปุกฺกุสกุเล วา ทลิทฺเท อปฺป\-นฺน\-ปาน\-โภชเน กสิรวุตฺติเก, ยตฺถ กสิเรน ฆาสจฺฉาโท ลพฺภติ, โส จ โหติ ทุพฺพณฺโณ ทุทฺทสิโก โอโกฏิมโก พหฺวาพาโธ กาโณ วา กุณี วา ขญฺโช วา ปกฺขหโต วา, น ลาภี อนฺนสฺส ปานสฺส วตฺถสฺส ยานสฺส มาลา\-คนฺธ\-วิเลปนสฺส เสยฺยา\-วสถ\-ปทีเปยฺยสฺส. โส กาเยน สุจริตํ จรติ, วาจาย สุจริตํ จรติ, มนสา สุจริตํ จรติ. โส กาเยน สุจริตํ จริตฺวา วาจาย สุจริตํ จริตฺวา มนสา สุจริตํ จริตฺวา กายสฺส เภทา ปรํ มรณา สุคติํ สคฺคํ โลกํ อุปปชฺชติ. เอวํ ปุคฺคโล ตโม โหติ โชติปรายโน. (2)}

\prose{กถญฺจ ปุคฺคโล โชติ โหติ ตมปรายโน, อิเธกจฺโจ ปุคฺคโล อุจฺเจ กุเล ปจฺจาชาโต โหติ ขตฺติย\-มหาสาล\-กุเล วา \csromanfolio{160}พฺราหฺมณ\-มหาสาล\-กุเล วา คหปติ\-มหาสาล\-กุเล วา อฑฺเฒ มหทฺธเน มหาโภเค ปหูต\-ชาตรูป\-รชเต ปหูต\-วิตฺตู\-ปกรเณ ปหูต\-ธน\-ธญฺเญ, โส จ โหติ อภิรูโป ทสฺสนีโย ปาสาทิโก ปรมาย วณฺณ\-โปกฺขร\-ตาย สมนฺนาคโต ลาภี อนฺนสา ปานสฺส วตฺถสฺส ยานสฺส มาลา\-คนฺธ\-วิเลปนสฺส เสยฺยา\-วสถ\-ปทีเปยฺยสฺส. โส กาเยน ทุจฺจริตํ จรติ, วาจาย ทุจฺจริตํ จรติ, มนสา ทุจฺจริตํ จรติ. โส กาเยน ทุจฺจริตํ จริตฺวา วาจาย ทุจฺจริตํ จริตฺวา มนสา ทุจฺจริตํ จริตฺวา กายสฺส เภทา ปรํ มรณา อปายํ ทุคฺคติํ วินิปาตํ นิรยํ อุปปชฺชติ. เอวํ ปุคฺคโล โชติ โหติ ตมปรายโน. (3)}

\prose{กถญฺจ ปุคฺคโล โชติ โหติ โชติปรายโน, อิเธกจฺโจ ปุคฺคโล อุจฺเจ กุเล ปจฺจาชาโต โหติ ขตฺติย\-มหาสาล\-กุเล วา พฺราหฺมณ\-มหาสาล\-กุเล วา คหปติ\-มหาสาล\-กุเล วา อฑฺเฒ มหทฺธเน มหาโภเค ปหูต\-ชาตรูป\-รชเต ปหูต\-วิตฺตู\-ปกรเณ ปหูต\-ธน\-ธญฺเญ, โส จ โหติ อภิรูโป ทสฺสนีโย ปาสาทิโก ปรมาย วณฺณ\-โปกฺขร\-ตาย สมนฺนาคโต ลาภี อนฺนสฺส ปานสฺส วตฺถสฺส ยานสฺส มาลา\-คนฺธ\-วิเลปนสฺส เสยฺยา\-วสถ\-ปทีเปยฺยสฺส. โส กาเยน สุจริตํ จรติ, วาจาย สุจริตํ จรติ, มนสา สุจริตํ จรติ. โส กาเยน สุจริตํ จริตฺวา วาจาย สุจริตํ จริตฺวา มนสา สุจริตํ จริตฺวา กายสฺส เภทา ปรํ มรณา สุคติํ สคฺคํ โลกํ อุปปชฺชติ. เอวํ ปุคฺคโล โชติ โหติ โชติปรายโน. (4)}

\proseitem{169}{กถญฺจ ปุคฺคโล โอณโตณโต โหติ ฯเปฯ. เอวํ ปุคฺคโล โอณโตณโต โหติ.}

\prose{กถญฺจ ปุคฺคโล โอณตุณฺณโต โหติ ฯเปฯ. เอวํ ปุคฺคโล โอณตุณฺณโต โหติ.}

\prose{กถญฺจ ปุคฺคโล อุณฺณโตณโต โหติ ฯเปฯ. เอวํ ปุคฺคโล อุณฺณโตณโต โหติ.}

\prose{กถญฺจ ปุคฺคโล อุณฺณตุณฺณโต โหติ ฯเปฯ. เอวํ ปุคฺคโล อุณฺณตุณฺณโต โหติ.}

\proseitem{170}{\csromanfolio{161}ตตฺถ กตเม จตฺตาโร รุกฺขูปมา ปุคฺคลา. จตฺตาโร รุกฺขา, เผคฺคุ สารปริวาโร, สาโร เผคฺคุปริวาโร, เผคฺคุ เผคฺคุปริวาโร, สาโร สารปริวาโร. เอวเมวํ จตฺตาโรเม รุกฺขูปมา ปุคฺคลา สนฺโต สํวิชฺชมานา โลกสฺมิํ. กตเม จตฺตาโร เผคฺคุ สารปริวาโร, สาโร เผคฺคุปริวาโร, เผคฺคุ เผคฺคุปริวาโร, สาโร สารปริวาโร.}

\prose{กถญฺจ ปุคฺคโล เผคฺคุ โหติ สารปริวาโร, อิเธกจฺโจ ปุคฺคโล ทุสฺสีโล โหติ ปาปธมฺโม, ปริสา จ ขฺวสฺส โหติ สีลวตี กลฺยาณธมฺมา. เอวํ ปุคฺคโล เผคฺคุ โหติ สารปริวาโร. เสยฺยถาปิ โส รุกฺโข เผคฺคุ สารปริวาโร, ตถูปโม อยํ ปุคฺคโล.}

\prose{กถญฺจ ปุคฺคโล สาโร โหติ เผคฺคุปริวาโร, อิเธกจฺโจ ปุคฺคโล สีลวา โหติ กลฺยาณธมฺโม, ปริสา จ ขฺวสฺส โหติ ทุสฺสีลา ปาปธมฺมา. เอวํ ปุคฺคโล สาโร โหติ เผคฺคุปริวาโร. เสยฺยถาปิ โส รุกฺโข สาโร เผคฺคุปริวาโร, ตถูปโม อยํ ปุคฺคโล.}

\prose{กถญฺจ ปุคฺคโล เผคฺคุ โหติ เผคฺคุปริวาโร, อิเธกจฺโจ ปุคฺคโล ทุสฺสีโล โหติ ปาปธมฺโม, ปริสาปิสฺส โหติ ทุสฺสีลา ปาปธมฺมา. เอวํ ปุคฺคโล เผคฺคุ โหติ เผคฺคุปริวาโร. เสยฺยถาปิ โส รุกฺโข เผคฺคุ เผคฺคุปริวาโร, ตถูปโม อยํ ปุคฺคโล.}

\prose{กถญฺจ ปุคฺคโล สาโร โหติ สารปริวาโร, อิเธกจฺโจ ปุคฺคโล สีลวา โหติ กลฺยาณธมฺโม, ปริสาปิสฺส โหติ สีลวตี กลฺยาณธมฺมา. เอวํ ปุคฺคโล สาโร โหติ สารปริวาโร. เสยฺยถาปิ โส รุกฺโข สาโร สารปริวาโร, ตถูปโม อยํ ปุคฺคโล. อิเม จตฺตาโร รุกฺขูปมา ปุคฺคลา สนฺโต สํวิชฺชมานา โลกสฺมิํ.}

\proseitem{171}{กตโม จ ปุคฺคโล รูปปฺปมาโณ รูปปฺปสนฺโน, อิเธกจฺโจ ปุคฺคโล อาโรหํ วา ปสฺสิตฺวา ปริณาหํ วา ปสฺสิตฺวา สณฺฐานํ วา ปสฺสิตฺวา ปาริปูริํ วา ปสฺสิตฺวา ตตฺถ ปมาณํ คเหตฺวา ปสาทํ ชเนติ, อยํ วุจฺจติ ปุคฺคโล รูปปฺปมาโณ รูปปฺปสนฺโน.}

\prose{\csromanfolio{162}กตโม จ ปุคฺคโล โฆสปฺปมาโณ โฆสปฺปสนฺโน, อิเธกจฺโจ ปุคฺคโล ปรวณฺณนาย ปรโถมนาย ปรปสํสนาย ปรวณฺณ\-หาริกาย\csromansharedfootnote{57843ad66bfc5326}{ปรวณฺณหาริยา (สี)} ตตฺถ ปมาณํ คเหตฺวา ปสาทํ ชเนติ, อยํ วุจฺจติ ปุคฺคโล โฆสปฺปมาโณ โฆสปฺปสนฺโน.}

\proseitem{172}{กตโม จ ปุคฺคโล ลูขปฺปมาโณ ลูขปฺปสนฺโน, อิเธกจฺโจ ปุคฺคโล จีวรลูขํ วา ปสฺสิตฺวา ปตฺตลูขํ วา ปสฺสิตฺวา เสนาสนลูขํ วา ปสฺสิตฺวา วิวิธํ วา ทุกฺกรการิกํ ปสฺสิตฺวา ตตฺถ ปมาณํ คเหตฺวา ปสาทํ ชเนติ, อยํ วุจฺจติ ปุคฺคโล ลูขปฺปมาโณ ลูขปฺปสนฺโน.}

\prose{กตโม จ ปุคฺคโล ธมฺมปฺปมาโณ ธมฺมปฺปสนฺโน, อิเธกจฺโจ ปุคฺคโล สีลํ วา ปสฺสิตฺวา สมาธิํ วา ปสฺสิตฺวา ปญฺญํ วา ปสฺสิตฺวา ตตฺถ ปมาณํ คเหตฺวา ปสาทํ ชเนติ, อยํ วุจฺจติ ปุคฺคโล ธมฺมปฺปมาโณ ธมฺมปฺปสนฺโน.}

\proseitem{173}{กถญฺจ ปุคฺคโล อตฺตหิตาย ปฏิปนฺโน โหติ โน ปรหิตาย, อิเธกจฺโจ ปุคฺคโล อตฺตนา สีลสมฺปนฺโน โหติ, โน ปรํ สีลสมฺปทาย สมาทเปติ. อตฺตนา สมาธิ\-สมฺปนฺโน โหติ, โน ปรํ สมาธิ\-สมฺปทาย สมาทเปติ. อตฺตนา ปญฺญาสมฺปนฺโน โหติ, โน ปรํ ปญฺญาสมฺปทาย สมาทเปติ. อตฺตนา วิมุตฺติ\-สมฺปนฺโน โหติ, โน ปรํ วิมุตฺติ\-สมฺปทาย สมาทเปติ. อตฺตนา วิมุตฺติ\-ญาณ\-ทสฺสน\-สมฺปนฺโน โหติ, โน ปรํ วิมุตฺติ\-ญาณ\-ทสฺสน\-สมฺปทาย สมาทเปติ. เอวํ ปุคฺคโล อตฺตหิตาย ปฏิปนฺโน โหติ โน ปรหิตาย.}

\prose{กถญฺจ ปุคฺคโล ปรหิตาย ปฏิปนฺโน โหติ โน อตฺตหิตาย, อิเธกจฺโจ ปุคฺคโล อตฺตนา น สีลสมฺปนฺโน โหติ, ปรํ สีลสมฺปทาย สมาทเปติ. อตฺตนา น สมาธิ\-สมฺปนฺโน โหติ, ปรํ สมาธิ\-สมฺปทาย สมาทเปติ. อตฺตนา น ปญฺญาสมฺปนฺโน โหติ, ปรํ ปญฺญาสมฺปทาย สมาทเปติ. อตฺตนา น วิมุตฺติ\-สมฺปนฺโน โหติ, ปรํ วิมุตฺติ\-สมฺปทาย สมาทเปติ. อตฺตนา น วิมุตฺติ\-ญาณ\-ทสฺสน\-สมฺปนฺโน โหติ, ปรํ วิมุตฺติ\-ญาณ\-ทสฺสน\-สมฺปทาย สมาทเปติ. เอวํ ปุคฺคโล ปรหิตาย ปฏิปนฺโน โหติ โน อตฺตหิตาย.}

\prose{\csromanfolio{163}กถญฺจ ปุคฺคโล อตฺตหิตาย เจว ปฏิปนฺโน โหติ ปรหิตาย จ, อิเธกจฺโจ ปุคฺคโล อตฺตนา จ สีลสมฺปนฺโน โหติ, ปรญฺจ สีลสมฺปทาย สมาทเปติ. อตฺตนา จ สมาธิ\-สมฺปนฺโน โหติ, ปรญฺจ สมาธิ\-สมฺปทาย สมาทเปติ. อตฺตนา จ ปญฺญาสมฺปนฺโน โหติ, ปรญฺจ ปญฺญาสมฺปทาย สมาทเปติ, อตฺตนา จ วิมุตฺติ\-สมฺปนฺโน โหติ, ปรญฺจ วิมุตฺติ\-สมฺปทาย สมาทเปติ. อตฺตนา จ วิมุตฺติ\-ญาณ\-ทสฺสน\-สมฺปนฺโน โหติ, ปรญฺจ วิมุตฺติ\-ญาณ\-ทสฺสน\-สมฺปทาย สมาทเปติ. เอวํ ปุคฺคโล อตฺตหิตาย เจว ปฏิปนฺโน โหติ ปรหิตาย จ.}

\prose{กถญฺจ ปุคฺคโล เนว อตฺตหิตาย ปฏิปนฺโน โหติ โน ปรหิตาย, อิเธกจฺโจ ปุคฺคโล อตฺตนา น สีลสมฺปนฺโน โหติ, โน ปรํ สีลสมฺปทาย สมาทเปติ. อตฺตนา น สมาธิ\-สมฺปนฺโน โหติ, โน ปรํ สมาธิ\-สมฺปทาย สมาทเปติ. อตฺตนา น ปญฺญาสมฺปนฺโน โหติ, โน ปรํ ปญฺญาสมฺปทาย สมาทเปติ. อตฺตนา น วิมุตฺติ\-สมฺปนฺโน โหติ, โน ปรํ วิมุตฺติ\-สมฺปทาย สมาทเปติ. อตฺตนา น วิมุตฺติ\-ญาณ\-ทสฺสน\-สมฺปนฺโน โหติ, โน ปรํ วิมุตฺติ\-ญาณ\-ทสฺสน\-สมฺปทาย สมาทเปติ. เอวํ ปุคฺคโล เนว อตฺตหิตาย ปฏิปนฺโน โหติ โน ปรหิตาย.}

\proseitem{174}{กถญฺจ ปุคฺคโล อตฺตนฺตโป โหติ อตฺต\-ปริตาปนา\-นุโยคม\-นุยุตฺโต, อิเธกจฺโจ ปุคฺคโล อเจลโก โหติ มุตฺตาจาโร, หตฺถา\-ปเลขโน\csromansharedfootnote{429ddba78d8d1da6}{หตฺถาวเลขโน (สฺยา)}, น- เอหิภทฺทนฺติโก, นติฏฺฐ\-ภทฺทนฺติโก, นาภิหฏํ, น อุทฺทิสฺสกตํ, น นิมนฺตนํ สาทิยติ, โส น กุมฺภิมุขา ปฏิคฺคณฺหาติ, น กโฬปิมุขา\csromansharedfootnote{1fca7e984041f476}{กโลปิมุขา (สี, สฺยา) ม 2. 5 ปิฏฺเฐปิ.} ปฏิคฺคณฺหาติ, น เอฬกมนฺตรํ, น ทณฺฑมนฺตรํ, น มุสลมนฺตรํ, น ทฺวินฺนํ ภุญฺชมานานํ, น คพฺภินิยา, น ปายมานาย, น ปุริส\-นฺตร\-คตาย, น สงฺกิตฺตีสุ, น ยตฺถ สา อุปฏฺฐิโต โหติ, น ยตฺถ มกฺขิกา สณฺฑ\-สณฺฑ\-จารินี, น มจฺฉํ, น มํสํ, น สุรํ, น เมรยํ, น ถุโสทกํ ปิวติ, โส เอกาคาริโก วา โหติ เอกาโลปิโก, ทฺวาคาริโก วา โหติ ทฺวาโลปิโก ฯเปฯ สตฺตาคาริโก วา โหติ สตฺตาโลปิโก, เอกิสฺสาปิ ทตฺติยา ยาเปติ, ทฺวีหิปิ ทตฺตีหิ ยาเปติ ฯเปฯ สตฺตหิปิ ทตฺตีหิ ยาเปติ, เอกาหิกมฺปิ อาหารํ อาหาเรติ, ทฺวีหิกมฺปิ\csromansharedfootnote{3e840b4295cf6330}{ทฺวาหิกมฺปิ (สี)} อาหารํ อาหาเรติ ฯเปฯ สตฺตาหิกมฺปิ อาหารํ อาหาเรหิ, อิติ \csromanfolio{164}เอวรูปํ อฑฺฒมาสิกมฺปิ ปริยาย\-ภตฺตโภชนานุโยคมนุยุตฺโต วิหรติ. โส สากภกฺโข วา โหติ, สามากภกฺโข วา โหติ, นีวารภกฺโข วา โหติ, ททฺทุลภกฺโข วา โหติ, หฏภกฺโข วา โหติ, กณภกฺโข วา โหติ, อาจามภกฺโข วา โหติ, ปิญฺญากภกฺโข วา โหติ, ติณภกฺโข วา โหติ, โคมยภกฺโข วา โหติ, วนมูล\-ผลา\-หาโร ยาเปติ ปวตฺต\-ผล\-โภชี. โส สาณานิปิ ธาเรติ, มสาณานิปิ ธาเรติ, ฉวทุสฺสานิปิ ธาเรติ, ปํสุกูลานิปิ ธาเรติ, ติรีฏานิปิ ธาเรติ, อชินมฺปิ ธาเรติ, อชินกฺขิปมฺปิ ธาเรติ, กุสจีรมฺปิ ธาเรติ, วากจีรมฺปิ ธาเรติ, ผลกจีรมฺปิ ธาเรติ, เกสกมฺพลมฺปิ ธาเรติ, วาฬกมฺพลมฺปิ ธาเรติ, อุลูกปกฺขมฺปิ\csromansharedfootnote{d84edafe5a03b46a}{อุลุกปกฺขมฺปิ (สี, สฺยา)} ธาเรติ, เกสมสฺสุ\-โลจโกปิ โหติ เกสมสฺสุ\-โลจนา\-นุโยคม\-นุยุตฺโต, อุพฺภฏฺฐโกปิ โหติ อาสน\-ปฏิกฺขิตฺโต, อุกฺกุฏิโกปิ โหติ อุกฺกุฏิก\-ปฺปธานมนุยุตฺโต, กณฺฏกา\-ปสฺสยิ\-โกปิ โหติ กณฺฏกาปสฺสเย เสยฺยํ กปฺเปติ, สาย\-ตติยกมฺปิ\csromansharedfootnote{9f7425f33f903dae}{สายํตติยกมฺปิ (สฺยา, ก) ม 2. 6 ปิฏฺเฐปิ.} อุทโก\-โรหนา\-นุโยคม\-นุยุตฺโต วิหรติ, อิติ เอวรูปํ อเนกวิหิตํ กายสฺส อาตาปน\-ปริตาปนา\-นุโยคม\-นุยุตฺโต วิหรติ. เอวํ ปุคฺคโล อตฺตนฺตโป โหติ อตฺต\-ปริตาปนา\-นุโยคม\-นุยุตฺโต.}

\proseitem{175}{กถญฺจ ปุคฺคโล ปรนฺตโป โหติ ปรปริตาปนา\-นุโยคม\-นุยุตฺโต, อิเธกจฺโจ ปุคฺคโล โอรพฺภิโก โหติ สูกริโก สากุณิโก มาควิโก ลุทฺโท มจฺฉฆาตโก โจโร โจรฆาตโก โคฆาตโก พนฺธนาคาริโก, เย วา ปนญฺเญปิ เกจิ กุรูรกมฺมนฺตา. เอวํ ปุคฺคโล ปรนฺตโป โหติ ปรปริตาปนา\-นุโยคม\-นุยุตฺโต.}

\proseitem{176}{กถญฺจ ปุคฺคโล อตฺตนฺตโป จ โหติ อตฺต\-ปริตาปนา\-นุโยคม\-นุยุตฺโต ปรนฺตโป จ ปรปริตาปนา\-นุโยคม\-นุยุตฺโต, อิเธกจฺโจ ปุคฺคโล ราชา วา โหติ ขตฺติโย มุทฺธาวสิตฺโต\csromansharedfootnote{0a23b209103c8ced}{มุทฺธาภิสิตฺโต (สฺยา, ก)} พฺราหฺมโณ วา มหาสาโล, โส ปุรตฺถิเมน นครสฺส นวํ สนฺธาคารํ\csromansharedfootnote{d3a2a2f7512b22fc}{สนฺถาคารํ (สฺยา), ยญฺญาคารํ (สี)} การาเปตฺวา เกสมสฺสุํ โอหาเรตฺวา ขราชินํ\csromansharedfootnote{caec3bb479d500f0}{ขุราชินํ (สฺยา, ก)} นิวาเสตฺวา สปฺปิเตเลน กายํ}

\prose{\csromanfolio{165}อพฺภญฺชิตฺวา มิควิสาเณน ปิฏฺฐิํ กณฺฑุวมาโน\csromansharedfootnote{dafac62e193e1e32}{กณฺฑูยมาโน (สี)} สนฺธาคารํ ปวิสติ สทฺธิํ มเหสิยา พฺราหฺมเณน จ ปุโรหิเตน. โส ตตฺถ อนนฺตรหิตาย ภูมิยา หริตุ\-ปลิตฺตาย เสยฺยํ กปฺเปติ. เอกิสฺสา คาวิยา สรูปวจฺฉาย ยํ เอกสฺมิํ ถเน ขีรํ โหติ, เตน ราชา ยาเปติ. ยํ ทุติยสฺมิํ ถเน ขีรํ โหติ, เตน มเหสี ยาเปติ. ยํ ตติยสฺมิํ ถเน ขีรํ โหติ, เตน พฺราหฺมโณ ปุโรหิโต ยาเปติ. ยํ จตุตฺถสฺมิํ ถเน ขีรํ โหติ, เตน อคฺคิํ ชุหติ. อวเสเสน วจฺฉโก ยาเปติ. โส เอวมาห “เอตฺตกา อุสภา หญฺญนฺตุ ยญฺญตฺถาย, เอตฺตกา วจฺฉตรา หญฺญนฺตุ ยญฺญตฺถาย, เอตฺตกา วจฺฉตริโย หญฺญนฺตุ ยญฺญตฺถาย, เอตฺตกา อชา หญฺญนฺตุ ยญฺญตฺถาย, เอตฺตกา อุรพฺภา หญฺญนฺตุ ยญฺญตฺถาย, (เอตฺตกา อสฺสา หญฺญนฺตุ ยญฺญตฺถาย)\csromansharedfootnote{9920403ae6d8d2e9}{( ) นตฺถิ สีหฬโปตฺถเก. มชฺฌิมนิกาเย กนฺทรกสุตฺเตปิ เอวเมว.} เอตฺตกา รุกฺขา ฉิชฺชนฺตุ ยูปตฺถาย, เอตฺตกา ทพฺภา ลูยนฺตุ พริหิสตฺถายา”ติ\csromansharedfootnote{a54f7d05f04035f6}{ปริหิํสตฺถายาติ (สี, สฺยา, ก) ม 2. 7 ปิฏฺเฐปิ.}. เยปิสฺส เต โหนฺติ ทาสาติ วา เปสฺสาติ วา กมฺมกราติ วา, เตปิ ทณฺฑตชฺชิตา ภยตชฺชิตา อสฺสุมุขา รุทมานา ปริกมฺมานิ กโรนฺติ. เอวํ ปุคฺคโล อตฺตนฺตโป จ โหติ อตฺต\-ปริตาปนา\-นุโยคม\-นุยุตฺโต ปรนฺตโป จ ปรปริตาปนา\-นุโยคม\-นุยุตฺโต.}

\proseitem{177}{กถญฺจ ปุคฺคโล เนว อตฺตนฺตโป จ โหติ น อตฺต\-ปริตาปนา\-นุโยคม\-นุยุตฺโต น ปรนฺตโป น ปรปริตาปนา\-นุโยคม\-นุยุตฺโต โส อนตฺตนฺตโป อปรนฺตโป ทิฏฺเฐว ธมฺเม นิจฺฉาโต นิพฺพุโต สีตีภูโต สุขปฺปฏิสํเวที พฺรหฺมภูเตน อตฺตนา วิหรติ.}

\prose{อิธ ตถาคโต โลเก อุปฺปชฺชติ อรหํ สมฺมาสมฺพุทฺโธ วิชฺชา\-จรณ\-สมฺปนฺโน สุคโต โลกวิทู อนุตฺตโร ปุริส\-ทมฺม\-สารถิ สตฺถา เทวมนุสฺสานํ พุทฺโธ ภควา. โส อิมํ โลกํ สเทวกํ สมารกํ สพฺรหฺมกํ สสฺสมณ\-พฺราหฺมณิํ ปชํ สเทวมนุสฺสํ สยํ อภิญฺญา สจฺฉิกตฺวา ปเวเทติ, โส ธมฺมํ เทเสติ อาทิกลฺยาณํ มชฺเฌกลฺยาณํ ปริโยสาน\-กลฺยาณํ สาตฺถํ สพฺยญฺชนํ เกวล\-ปริปุณฺณํ ปริสุทฺธํ พฺรหฺมจริยํ ปกาเสติ, ตํ ธมฺมํ สุณาติ คหปติ วา คหปติปุตฺโต วา อญฺญตรสฺมิํ วา กุเล \csromanfolio{166}ปจฺจาชาโต, โส ตํ ธมฺมํ สุตฺวา ตถาคเต สทฺธํ ปฏิลภติ, โส เตน สทฺธา\-ปฏิลาเภน สมนฺนาคโต อิติ ปฏิสญฺจิกฺขติ “สมฺพาโธ ฆราวาโส รชาปโถ, อพฺโภกาโส ปพฺพชฺชา, นยิทํ สุกรํ อคารํ อชฺฌาวสตา เอกนฺต\-ปริปุณฺณํ เอกนฺต\-ปริสุทฺธํ สงฺขลิขิตํ พฺรหฺมจริยํ จริตุํ, ยํนูนาหํ เกสมสฺสุํ โอหาเรตฺวา กาสายานิ วตฺถานิ อจฺฉาเทตฺวา อคารสฺมา อนคาริยํ ปพฺพเชยฺยนฺ”ติ. โส อปเรน สมเยน อปฺปํ วา โภคกฺขนฺธํ ปหาย มหนฺตํ วา โภคกฺขนฺธํ ปหาย อปฺปํ วา ญาติปริวฏฺฏํ ปหาย มหนฺตํ วา ญาติปริวฏฺฏํ ปหาย เกสมสฺสุํ โอหาเรตฺวา กาสายานิ วตฺถานิ อจฺฉาเทตฺวา อคารสฺมา อนคาริยํ ปพฺพชติ.}

\proseitem{178}{โส เอวํ ปพฺพชิโต สมาโน ภิกฺขูนํ สิกฺขา\-สาชีว\-สมาปนฺโน ปาณาติปาตํ ปหาย ปาณาติปาตา ปฏิวิรโต โหติ, นิหิตทณฺโฑ นิหิตสตฺโถ ลชฺชี ทยาปนฺโน สพฺพปาณภูต\-หิตา\-นุกมฺปี วิหรติ.}

\prose{อทินฺนาทานํ ปหาย อทินฺนาทานา ปฏิวิรโต โหติ, ทินฺนาทายี ทินฺน\-ปาฏิกงฺขี อเถเนน สุจิภูเตน อตฺตนา วิหรติ.}

\prose{อพฺรหฺมจริยํ ปหาย พฺรหฺมจารี โหติ อาราจารี\csromansharedfootnote{0999306d47cbf569}{อนาจารี (ก)} ปฏิวิรโต เมถุนา คามธมฺมา.}

\prose{มุสาวาทํ ปหาย มุสาวาทา ปฏิวิรโต โหติ สจฺจวาที สจฺจสนฺโธ เถโต ปจฺจยิโก อวิสํวาทโก โลกสฺส.}

\prose{ปิสุณํ วาจํ ปหาย ปิสุณาย วาจาย ปฏิวิรโต โหติ, อิโต สุตฺวา น อมุตฺร อกฺขาตา อิเมสํ เภทาย, อมุตฺร วา สุตฺวา น อิเมสํ อกฺขาตา อมูสํ เภทาย, อิติ ภินฺนานํ วา สนฺธาตา สหิตานํ วา อนุปฺปทาตา สมคฺคาราโม สมคฺครโต สมคฺคนนฺที สมคฺคกรณิํ วาจํ ภาสิตา โหติ.}

\prose{ผรุสํ วาจํ ปหาย ผรุสาย วาจาย ปฏิวิรโต โหติ, ยา สา วาจา เนลา กณฺณสุขา เปมนียา หทยงฺคมา โปรี พหุชนกนฺตา พหุชนมนาปา ตถารูปิํ วาจํ ภาสิตา โหติ.}

\prose{\csromanfolio{167}สมฺผปฺปลาปํ ปหาย สมฺผปฺปลาปา ปฏิวิรโต โหติ กาลวาที ภูตวาที อตฺถวาที ธมฺมวาที วินยวาที นิธานวติํ วาจํ ภาสิตา กาเลน สาปเทสํ ปริยนฺตวติํ อตฺถสํหิตํ.}

\proseitem{179}{โส ภีชคามภูตคามสมารมฺภา ปฏิวิรโต โหติ. เอกภตฺติโก โหติ รตฺตูปรโต, วิรโต วิกาลโภชนา. นจฺจ\-คีต\-วาทิต\-วิสูก\-ทสฺสนา ปฏิวิรโต โหติ. มาลา\-คนฺธ\-วิเลปน\-ธารณ\-มณฺฑน\-วิภูสน\-ฏฺฐานา ปฏิวิรโต โหติ. อุจฺจาสยน\-มหาสยนา ปฏิวิรโต โหติ. ชาตรูป รชต ปฏิคฺคหณา ปฏิวิรโต โหติ.}

\prose{อามก\-ธญฺญ\-ปฏิคฺคหณา ปฏิวิรโต โหติ, อามก\-มํส\-ปฏิคฺคหณา ปฏิวิรโต โหติ, อิตฺถิกุมาริกาปฏิคฺคหณา ปฏิวิรโต โหติ, ทาสิ\-ทาส\-ปฏิคฺคหณา ปฏิวิรโต โหติ, อเช\-ฬก\-ปฏิคฺคหณา ปฏิวิรโต โหติ, กุกฺกุฏ\-สูกร\-ปฏิคฺคหณา ปฏิวิรโต โหติ, หตฺถิ\-ควา\-สฺส\-วฬว\-ปฏิคฺคหณา ปฏิวิรโต โหติ, เขตฺตวตฺถุ\-ปฏิคฺคหณา ปฏิวิรโต โหติ, ทูเตยฺย\-ปหิณ\-คมนา\-นุโยคา ปฏิวิรโต โหติ, กยวิกฺกยา ปฏิวิรโต โหติ, ตุลา\-กูฏ\-กํส\-กูฏ\-มาน\-กูฏา ปฏิวิรโต โหติ, อุกฺโกฏน\-วญฺจน\-นิกติ\-สาจิโยคา\csromansharedfootnote{ae1d07b3c9adcc76}{...สาวิโยคา (สฺยา, ก)} ปฏิวิรโต โหติ, เฉทน\-วธ\-พนฺธน\-วิปราโมส\-อาโลป\-สหสาการา ปฏิวิรโต โหติ.}

\proseitem{180}{โส สนฺตุฏฺโฐ โหติ กาย\-ปริหาริเกน จีวเรน กุจฺฉิ\-ปริหาริเกน ปิณฺฑปาเตน, โส เยน เยเนว ปกฺกมติ, สมาทาเยว ปกฺกมติ. เสยฺยถาปิ นาม ปกฺขี สกุโณ เยน เยเนว เฑติ, สปตฺตภาโรว เฑติ. เอวเมวํ ภิกฺขุ สนฺตุฏฺโฐ โหติ กาย\-ปริหาริเกน จีวเรน กุจฺฉิ\-ปริหาริเกน ปิณฺฑปาเตน, โส เยน เยเนว ปกฺกมติ, สมาทาเยว ปกฺกมติ. โส อิมินา อริเยน สีลกฺขนฺเธน สมนฺนาคโต อชฺฌตฺตํ อนวชฺชสุขํ ปฏิสํเวเทติ.}

\proseitem{181}{โส จกฺขุนา รูปํ ทิสฺวา น นิมิตฺตคฺคาหี โหติ นานุพฺยญฺชน\-คฺคาหี, ยตฺวาธิกรณเมนํ จกฺขุนฺทฺริยํ อสํวุตํ วิหรนฺตํ อภิชฺฌา\-โทมนสฺสา ปาปกา อกุสลา ธมฺมา อนฺวาสฺสเวยฺยุํ, ตสฺส สํวราย ปฏิปชฺชติ, รกฺขติ จกฺขุนฺทฺริยํ, จกฺขุนฺทฺริเย สํวรํ อาปชฺชติ. โสเตน สทฺทํ สุตฺวา ฯเปฯ. \csromanfolio{168}ฆาเนน คนฺธํ ฆายิตฺวา ฯเปฯ. ชิวฺหาย รสํ สายิตฺวา ฯเปฯ. กาเยน โผฏฺฐพฺพํ ผุสิตฺวา ฯเปฯ. มนสา ธมฺมํ วิญฺญาย น นิมิตฺตคฺคาหี โหติ นานุพฺยญฺชน\-คฺคาหี, ยตฺวาธิกรณเมนํ มนินฺทฺริยํ อสํวุตํ วิหรนฺตํ อภิชฺฌา\-โทมนสฺสา ปาปกา อกุสลา ธมฺมา อนฺวาสฺสเวยฺยุํ, ตสฺส สํวราย ปฏิปชฺชติ, รกฺขติ มนินฺทฺริยํ, มนินฺทฺริเย สํวรํ อาปชฺชติ. โส อิมินา อริเยน อินฺทฺริย\-สํวเรน สมนฺนาคโต อชฺฌตฺตํ อพฺยาเสกสุขํ ปฏิสํเวเทติ.}

\proseitem{182}{โส อภิกฺกนฺเต ปฏิกฺกนฺเต สมฺปชานการี โหติ, อาโลกิเต วิโลกิเต สมฺปชานการี โหติ, สมิญฺชิเต ปสาริเต สมฺปชานการี โหติ, สงฺฆาฏิ\-ปตฺต\-จีวร\-ธารเณ สมฺปชานการี โหติ, อสิเต ปีเต ขายิเต สายิเต สมฺปชานการี โหติ, อุจฺจาร\-ปสฺสาว\-กมฺเม สมฺปชานการี โหติ, คเต ฐิเต นิสินฺเน สุตฺเต ชาคริเต ภาสิเต ตุณฺหีภาเว สมฺปชานการี โหติ.}

\prose{โส อิมินา จ อริเยน สีลกฺขนฺเธน สมนฺนาคโต อิมินา จ อริเยน อินฺทฺริย\-สํวเรน สมนฺนาคโต อิมินา จ อริเยน สติสมฺปชญฺเญน สมนฺนาคโต อิมาย จ อริยาย สนฺตุฏฺฐิยา สมนฺนาคโต\csromansharedfootnote{1ec179c55ee0c54c}{ปสฺส ม 2. กนฺทรกสุตฺเต อํ 1. อตฺตนฺตปสุตฺเต จ.} วิวิตฺตํ เสนาสนํ ภชติ อรญฺญํ รุกฺขมูลํ ปพฺพตํ กนฺทรํ คิริคุหํ สุสานํ วนปตฺถํ อพฺโภกาสํ ปลาลปุญฺชํ. โส ปจฺฉาภตฺตํ ปิณฺฑปาต\-ปฏิกฺกนฺโต นิสีทติ ปลฺลงฺกํ อาภุชิตฺวา อุชุํ กายํ ปณิธาย ปริมุขํ สติํ อุปฏฺฐเปตฺวา, โส อภิชฺฌํ โลเก ปหาย วิคตา\-ภิชฺเฌน เจตสา วิหรติ, อภิชฺฌาย จิตฺตํ ปริโสเธติ. พฺยาปาทปโทสํ ปหาย อพฺยาปนฺนจิตฺโต วิหรติ สพฺพปาณภูต\-หิตา\-นุกมฺปี, พฺยาปาทปโทสา จิตฺตํ ปริโสเธติ. ถินมิทฺธํ\csromansharedfootnote{2b81a700dc7a1cda}{ถีนมิทฺธํ (สี, สฺยา)} ปหาย วิคต\-ถินมิทฺโธ วิหรติ อาโลกสญฺญี สโต สมฺปชาโน, ถินมิทฺธา จิตฺตํ ปริโสเธติ. อุทฺธจฺจ\-กุกฺกุจฺจํ ปหาย อนุทฺธโต วิหรติ อชฺฌตฺตํ วูปสนฺตจิตฺโต, อุทฺธจฺจ\-กุกฺกุจฺจา จิตฺตํ ปริโสเธติ. วิจิกิจฺฉํ ปหาย ติณฺณ\-วิจิกิจฺโฉ วิหรติ อกถํกถี กุสเลสุ ธมฺเมสุ, วิจิกิจฺฉาย จิตฺตํ ปริโสเธติ.}

\proseitem{183}{โส อิเม ปญฺจ นีวรเณ ปหาย เจตโส อุปกฺกิเลเส ปญฺญาย ทุพฺพลีกรเณ วิวิจฺเจว กาเมหิ วิวิจฺจ อกุสเลหิ ธมฺเมหิ}

\prose{\csromanfolio{169}สวิตกฺกํ สวิจารํ วิเวกชํ ปีติสุขํ ปฐมํ ฌานํ อุปสมฺปชฺช วิหรติ. วิตกฺก\-วิจารานํ วูปสมา อชฺฌตฺตํ สมฺปสาทนํ เจตโส เอโกทิภาวํ อวิตกฺกํ อวิจารํ สมาธิชํ ปีติสุขํ ทุติยํ ฌานํ อุปสมฺปชฺช วิหรติ. ปีติยา จ วิราคา อุเปกฺขโก จ วิหรติ สโต จ สมฺปชาโน สุขญฺจ กาเยน ปฏิสํเวเทติ, ยํ ตํ อริยา อาจิกฺขนฺติ “อุเปกฺขโก สติมา สุขวิหารี”ติ, ตติยํ ฌานํ อุปสมฺปชฺช วิหรติ. สุขสฺส จ ปหานา ทุกฺขสฺส จ ปหานา ปุพฺเพว โสมนสฺส\-โทมนสฺสานํ อตฺถงฺคมา อทุกฺขมสุขํ อุเปกฺขา\-สติ\-ปาริสุทฺธิํ จตุตฺถํ ฌานํ อุปสมฺปชฺช วิหรติ.}

\prose{โส เอวํ สมาหิเต จิตฺเต ปริสุทฺเธ ปริโยทาเต อนงฺคเณ วิคตู\-ปกฺกิเลเส มุทุภูเต กมฺมนิเย ฐิเต อาเนญฺชปฺปตฺเต ปุพฺเพ\-นิวาสา\-นุสฺสติ\-ญาณาย จิตฺตํ อภินินฺนาเมติ, โส อเนกวิหิตํ ปุพฺเพนิวาสํ อนุสฺสรติ. เสยฺยถิทํ, เอกมฺปิ ชาติํ เทฺวปิ ชาติโย ติสฺโสปิ ชาติโย จตสฺโสปิ ชาติโย ปญฺจปิ ชาติโย ทสปิ ชาติโย วีสมฺปิ ชาติโย ติํสมฺปิ ชาติโย จตฺตาลีสมฺปิ ชาติโย ปญฺญาสมฺปิ ชาติโย ชาติสตมฺปิ ชาติสหสฺสมฺปิ ชาติ\-สต\-สหสฺสมฺปิ อเนเกปิ สํวฏฺฏกปฺเป อเนเกปิ วิวฏฺฏกปฺเป อเนเกปิ สํวฏฺฏ\-วิวฏฺฏ\-กปฺเป, “อมุตฺราสิํ เอวํนาโม เอวํโคตฺโต เอวํวณฺโณ เอวมาหาโร เอวํ\-สุขทุกฺข\-ปฏิสํเวที เอวมา\-ยุปริยนฺโต, โส ตโต จุโต อมุตฺร อุทปาทิํ, ตตฺราปาสิํ เอวํนาโม เอวํโคตฺโต เอวํวณฺโณ เอวมาหาโร เอวํ\-สุขทุกฺข\-ปฏิสํเวที เอวมา\-ยุปริยนฺโต, โส ตโต จุโต อิธูปปนฺโน”ติ, อิติ สาการํ สอุทฺเทสํ อเนกวิหิตํ ปุพฺเพนิวาสํ อนุสฺสรติ.}

\proseitem{184}{โส เอวํ สมาหิเต จิตฺเต ปริสุทฺเธ ปริโยทาเต อนงฺคเณ วิคตู\-ปกฺกิเลเส มุทุภูเต กมฺมนิเย ฐิเต อาเนญฺชปฺปตฺเต สตฺตานํ จุตูปปาต\-ญาณาย จิตฺตํ อภินินฺนาเมติ, โส ทิพฺเพน จกฺขุนา วิสุทฺเธน อติกฺกนฺต\-มานุสเกน สตฺเต ปสฺสติ จวมาเน อุปปชฺชมาเน หีเน ปณีเต สุวณฺเณ ทุพฺพณฺเณ สุคเต ทุคฺคเต ยถากมฺมูปเค สตฺเต ปชานาติ “อิเม วต โภนฺโต สตฺตา กาย\-ทุจฺจริเตน สมนฺนาคตา วจีทุจฺจริเตน สมนฺนาคตา มโน\-ทุจฺจริเตน สมนฺนาคตา, อริยานํ อุปวาทกา มิจฺฉาทิฏฺฐิกา มิจฺฉาทิฏฺฐิ\-กมฺม\-สมาทานา, เต กายสฺส เภทา ปรํ มรณา อปายํ ทุคฺคติํ วินิปาตํ นิรยํ อุปปนฺนา. อิเม วา ปน โภนฺโต สตฺตา \csromanfolio{170}กายสุจริเตน สมนฺนาคตา วจีสุจริเตน สมนฺนาคตา มโนสุจริเตน สมนฺนาคตา, อริยานํ อนุปวาทกา สมฺมาทิฏฺฐิกา สมฺมา\-ทิฏฺฐิ\-กมฺม\-สมาทานา, เต กายสฺส เภทา ปรํ มรณา สุคติํ สคฺคํ โลกํ อุปปนฺนา”ติ. โส อิติ ทิพฺเพน จกฺขุนา วิสุทฺเธน อติกฺกนฺต\-มานุสเกน สตฺเต ปสฺสติ จวมาเน อุปปชฺชมาเน หีเน ปณีเต สุวณฺเณ ทุพฺพณฺเณ สุคเต ทุคฺคเต ยถากมฺมูปเค สตฺเต ปชานาติ.}

\proseitem{185}{โส เอวํ สมาหิเต จิตฺเต ปริสุทฺเธ ปริโยทาเต อนงฺคเณ วิคตู\-ปกฺกิเลเส มุทุภูเต กมฺมนิเย ฐิเต อาเนญฺชปฺปตฺเต อาสวานํ ขยญาณาย จิตฺตํ อภินินฺนาเมติ, โส “อิทํ ทุกฺขนฺ”ติ ยถาภูตํ ปชานาติ, “อยํ ทุกฺขสมุทโย”ติ ยถาภูตํ ปชานาติ, “อยํ ทุกฺขนิโรโธ”ติ ยถาภูตํ ปชานาติ, “อยํ ทุกฺขนิโรธ\-คามินี ปฏิปทา”ติ ยถาภูตํ ปชานาติ, “อิเม อาสวา”ติ ยถาภูตํ ปชานาติ, “อยํ อาสวสมุทโย”ติ ยถาภูตํ ปชานาติ, “อยํ อาสวนิโรโธ”ติ ยถาภูตํ ปชานาติ, “อยํ อาสว\-นิโรธ\-คามินี ปฏิปทา”ติ ยถาภูตํ ปชานาติ. ตสฺส เอวํ ชานโต เอวํ ปสฺสโต กามาสวาปิ จิตฺตํ วิมุจฺจติ, ภวาสวาปิ จิตฺตํ วิมุจฺจติ, อวิชฺชาสวาปิ จิตฺตํ วิมุจฺจติ, วิมุตฺตสฺมิํ “วิมุตฺตมฺ”อิติ ญาณํ โหติ, “ขีณา ชาติ, วุสิตํ พฺรหฺมจริยํ, กตํ กรณียํ, นาปรํ อิตฺถตฺตายา”ติ ปชานาติ. เอวํ ปุคฺคโล เนว อตฺตนฺตโป จ โหติ น อตฺตปริตาปนานุโยคมนุตฺโต น ปรนฺตโป น ปรปริตาปนา\-นุโยคม\-นุยุตฺโต, โส อนตฺตนฺตโป อปรนฺตโป ทิฏฺเฐว ธมฺเม นิจฺฉาโต นิพฺพุโต สีตีภูโต สุขปฺปฏิสํเวที พฺรหฺมภูเตน อตฺตนา วิหรติ.}

\proseitem{186}{กตโม จ ปุคฺคโล สราโค, ยสฺส ปุคฺคลสฺส ราโค อปฺปหีโน, อยํ วุจฺจติ ปุคฺคโล สราโค.}

\prose{กตโม จ ปุคฺคโล สโทโส, ยสฺส ปุคฺคลสฺส โทโส อปฺปหีโน, อยํ วุจฺจติ ปุคฺคโล สโทโส.}

\prose{กตโม จ ปุคฺคโล สโมโห, ยสฺส ปุคฺคลสฺส โมโห อปฺปหีโน, อยํ วุจฺจติ ปุคฺคโล สโมโห.}

\prose{กตโม จ ปุคฺคโล สมาโน, ยสฺส ปุคฺคลสฺส มาโน อปฺปหีโน, อยํ วุจฺจติ ปุคฺคโล สมาโน.}

\proseitem{187}{\csromanfolio{171}กถญฺจ ปุคฺคโล ลาภี โหติ อชฺฌตฺตํ เจโตสมถสฺส น ลาภี อธิปญฺญา\-ธมฺม\-วิปสฺสนาย, อิเธกจฺโจ ปุคฺคโล ลาภี โหติ รูปสหคตานํ วา อรูป\-สหคตานํ วา สมาปตฺตีนํ, น ลาภี โลกุตฺตร\-มคฺคสฺส วา ผลสฺส วา, เอวํ ปุคฺคโล ลาภี โหติ อชฺฌตฺตํ เจโตสมถสฺส น ลาภี อธิปญฺญา\-ธมฺม\-วิปสฺสนาย.}

\prose{กถญฺจ ปุคฺคโล ลาภี โหติ อธิปญฺญา\-ธมฺม\-วิปสฺสนาย น ลาภี อชฺฌตฺตํ เจโตสมถสฺส, อิเธกจฺโจ ปุคฺคโล ลาภี โหติ โลกุตฺตร\-มคฺคสฺส วา ผลสฺส วา, น ลาภี รูปสหคตานํ วา อรูป\-สหคตานํ วา สมาปตฺตีนํ, เอวํ ปุคฺคโล ลาภี โหติ อธิปญฺญา\-ธมฺม\-วิปสฺสนาย น ลาภี อชฺฌตฺตํ เจโตสมถสฺส.}

\prose{กถญฺจ ปุคฺคโล ลาภี เจว โหติ อชฺฌตฺตํ เจโตสมถสฺส ลาภี จ อธิปญฺญา\-ธมฺม\-วิปสฺสนาย, อิเธกจฺโจ ปุคฺคโล ลาภี โหติ รูปสหคตานํ วา อรูป\-สหคตานํ วา สมาปตฺตีนํ, ลาภี โลกุตฺตร\-มคฺคสฺส วา ผลสฺส วา, เอวํ ปุคฺคโล ลาภี เจว โหติ อชฺฌตฺตํ เจโตสมถสฺส ลาภี จ อธิปญฺญา\-ธมฺม\-วิปสฺสนาย.}

\prose{กถญฺจ ปุคฺคโล เนว ลาภี โหติ อชฺฌตฺตํ เจโตสมถสฺส น ลาภี อธิปญฺญา\-ธมฺม\-วิปสฺสนาย, อิเธกจฺโจ ปุคฺคโล เนว ลาภี โหติ รูปสหคตานํ วา อรูป\-สหคตานํ วา สมาปตฺตีนํ, น ลาภี โลกุตฺตร\-มคฺคสฺส วา ผลสฺส วา, เอวํ ปุคฺคโล เนว ลาภี โหติ อชฺฌตฺตํ เจโตสมถสฺส น ลาภี อธิปญฺญา\-ธมฺม\-วิปสฺสนาย.}

\proseitem{188}{กตโม จ ปุคฺคโล อนุโสตคามี อิเธกจฺโจ ปุคฺคโล กาเม จ ปฏิเสวติ, ปาปญฺจ กมฺมํ กโรติ, อยํ วุจฺจติ ปุคฺคโล อนุโสตคามี.}

\prose{กตโม จ ปุคฺคโล ปฏิโสตคามี, อิเธกจฺโจ ปุคฺคโล กาเม จ น ปฏิเสวติ ปาปญฺจ กมฺมํ น กโรติ, โส สหาปิ ทุกฺเขน สหาปิ โทมนสฺเสน อสฺสุมุเขนปิ รุทมาโน ปริปุณฺณํ ปริสุทฺธํ พฺรหฺมจริยํ จรติ, อยํ วุจฺจติ ปุคฺคโล ปฏิโสตคามี.}

\prose{\csromanfolio{172}กตโม จ ปุคฺคโล ฐิตตฺโต, อิเธกจฺโจ ปุคฺคโล ปญฺจนฺนํ โอรมฺภาคิยานํ สํโยชนานํ ปริกฺขยา โอปปาติโก โหติ ตตฺถ ปรินิพฺพายี อนาวตฺติธมฺโม ตสฺมา โลกา. อยํ วุจฺจติ ปุคฺคโล ฐิตตฺโต.}

\prose{กตโม จ ปุคฺคโล ติณฺโณ ปารงฺคโต ถเล ติฏฺฐติ พฺราหฺมโณ, อิเธกจฺโจ ปุคฺคโล อาสวานํ ขยา อนาสวํ เจโตวิมุตฺติํ ปญฺญาวิมุตฺติํ ทิฏฺเฐว ธมฺเม สยํ อภิญฺญา สจฺฉิกตฺวา อุปสมฺปชฺช วิหรติ, อยํ วุจฺจติ ปุคฺคโล ติณฺโณ ปารงฺคโต ถเล ติฏฺฐติ พฺราหฺมโณ.}

\proseitem{189}{กถญฺจ ปุคฺคโล อปฺปสฺสุโต โหติ สุเตน อนุปปนฺโน, อิเธกจฺจสฺส ปุคฺคลสฺส อปฺปกํ สุตํ โหติ สุตฺตํ เคยฺยํ เวยฺยากรณํ คาถํ อุทานํ อิติวุตฺตกํ ชาตกํ อพฺภุตธมฺมํ เวทลฺลํ, โส ตสฺส อปฺปกสฺส สุตสฺส น อตฺถมญฺญาย ธมฺมมญฺญาย ธมฺมา\-นุธมฺม\-ปฺปฏิปนฺโน\csromansharedfootnote{f76dad9b8277910c}{น ธมฺมมญฺญาย น ธมฺมานุธมฺมปฏิปนฺโน (สฺยา), น ธมฺมมญฺญาย ธมฺมานุธมฺมปฺปฏิปนฺโน (ก)} โหติ. เอวํ ปุคฺคโล อปฺปสฺสุโต โหติ สุเตน อนุปปนฺโน.}

\prose{กถญฺจ ปุคฺคโล อปฺปสฺสุโต โหติ สุเตน อุปปนฺโน, อิเธกจฺจสฺส ปุคฺคลสฺส อปฺปกํ สุตํ โหติ สุตฺตํ เคยฺยํ เวยฺยากรณํ คาถํ อุทานํ อิติวุตฺตกํ ชาตกํ อพฺภุตธมฺมํ เวทลฺลํ, โส ตสฺส อปฺปกสฺส สุตสฺส อตฺถมญฺญาย ธมฺมมญฺญาย ธมฺมา\-นุธมฺม\-ปฺปฏิปนฺโน โหติ, เอวํ ปุคฺคโล อปฺปสฺสุโต โหติ สุเตน อุปปนฺโน.}

\prose{กถญฺจ ปุคฺคโล พหุสฺสุโต โหติ สุเตน อนุปปนฺโน, อิเธกจฺจสฺส ปุคฺคลสฺส พหุกํ สุตํ โหติ สุตฺตํ เคยฺยํ เวยฺยากรณํ คาถํ อุทานํ อิติวุตฺตกํ ชาตกํ อพฺภุตธมฺมํ เวทลฺลํ, โส ตสฺส พหุกสฺส สุตสฺส น อตฺถมญฺญาย ธมฺมมญฺญาย ธมฺมา\-นุธมฺม\-ปฺปฏิปนฺโน โหติ, เอวํ ปุคฺคโล พหุสฺสุโต โหติ สุเตน อนุปปนฺโน.}

\prose{กถญฺจ ปุคฺคโล พหุสฺสุโต โหติ สุเตน อุปปนฺโน, อิเธกจฺจสฺส ปุคฺคลสฺส พหุกํ สุตํ โหติ สุตฺตํ เคยฺยํ เวยฺยากรณํ คาถํ อุทานํ อิติวุตฺตกํ ชาตกํ อพฺภุตธมฺมํ เวทลฺลํ, โส ตสฺส พหุกสฺส สุตสฺส อตฺถมญฺญาย ธมฺมมญฺญาย ธมฺมา\-นุธมฺม\-ปฺปฏิปนฺโน โหติ, เอวํ ปุคฺคโล พหุสฺสุโต โหติ สุเตน อุปปนฺโน.}

\csromanmatikaanchor{mtk.74}
\csromantocmarkref{subsection}{5. ปญฺจกปุคฺคลปญฺญตฺติ}{mtk.74}
\bookmark[dest={mtk.74},level=2]{5. ปญฺจกปุคฺคลปญฺญตฺติ}
\proseitem{190}{\csromanfolio{173}กตโม จ ปุคฺคโล สมณมจโล, อิเธกจฺโจ ปุคฺคโล ติณฺณํ สํโยชนานํ ปริกฺขยา โสตาปนฺโน โหติ อวินิปาต\-ธมฺโม นิยโต สมฺโพธิ\-ปรายโน. อยํ วุจฺจติ ปุคฺคโล สมณมจโล.}

\prose{กตโม จ ปุคฺคโล สมณปทุโม, อิเธกจฺโจ ปุคฺคโล ติณฺณํ สํโยชนานํ ปริกฺขยา ราค\-โทส\-โมหานํ ตนุตฺตา สกทาคามี โหติ, สกิเทว อิมํ โลกํ อาคนฺตฺวา ทุกฺขสฺสนฺตํ กโรติ. อยํ วุจฺจติ ปุคฺคโล สมณปทุโม.}

\prose{กตโม จ ปุคฺคโล สมณ\-ปุณฺฑรีโก, อิเธกจฺโจ ปุคฺคโล ปญฺจนฺนํ โอรมฺภาคิยานํ สํโยชนานํ ปริกฺขยา โอปปาติโก โหติ, ตตฺถ ปรินิพฺพายี อนาวตฺติธมฺโม ตสฺมา โลกา. อยํ วุจฺจติ ปุคฺคโล สมณ\-ปุณฺฑรีโก.}

\prose{กตโม จ ปุคฺคโล สมเณสุ สมณ\-สุขุมาโล, อิเธกจฺโจ ปุคฺคโล อาสวานํ ขยา อนาสวํ เจโตวิมุตฺติํ ปญฺญาวิมุตฺติํ ทิฏฺเฐว ธมฺเม สยํ อภิญฺญา สจฺฉิกตฺวา อุปสมฺปชฺช วิหรติ. อยํ วุจฺจติ ปุคฺคโล สมเณสุ สมณ\-สุขุมาโลติ.}

\vspace{\tipitakaparskip}
\csromancenter{จตุกฺกนิทฺเทโส.}
\csromanheader{2}{5. \textbf{ปญฺจกปคฺคลปญฺญตฺติ}}
\proseitem{191}{ตตฺร ยฺวายํ ปุคฺคโล อารภติ จ วิปฺปฏิสารี จ โหติ, ตญฺจ เจโตวิมุตฺติํ ปญฺญาวิมุตฺติํ ยถาภูตํ นปฺปชานาติ, ยตฺถสฺส เต อุปฺปนฺนา ปาปกา อกุสลา ธมฺมา อปริเสสา นิรุชฺฌนฺติ, โส เอวมสฺส วจนีโย “อายสฺมโต โข อารมฺภชา\csromansharedfootnote{90a8f0f309e5f23a}{อารพฺภชา (ก) อํ 2. 147 ปิฏฺเฐปิ.} อาสวา สํวิชฺชนฺติ, วิปฺปฏิสารชา อาสวา ปวฑฺฒนฺติ, สาธุ วตายสฺมา อารมฺภเช อาสเว ปหาย วิปฺปฏิสารเช อาสเว ปฏิวิโนเทตฺวา จิตฺตํ ปญฺญญฺจ ภาเวตุ, เอวมายสฺมา อมุนา ปญฺจเมน ปุคฺคเลน สมสโม ภวิสฺสตี”ติ.}

\prose{ตตฺร ยฺวายํ ปุคฺคโล อารภติ น วิปฺปฏิสารี โหติ, ตญฺจ เจโตวิมุตฺติํ ปญฺญาวิมุตฺติํ ยถาภูตํ นปฺปชานาติ, ยตฺถสฺส เต อุปฺปนฺนา \csromanfolio{174}ปาปกา อกุสลา ธมฺมา อปริเสสา นิรุชฺฌนฺติ, โส เอวมสฺส วจนีโย “อายสฺมโต โข อารมฺภชา อาสวา สํวิชฺชนฺติ, วิปฺปฏิสารชา อาสวา นปฺปวฑฺฒนฺติ, สาธุ วตายสฺมา อารมฺภเช อาสเว ปหาย จิตฺตํ ปญฺญญฺจ ภาเวตุ, เอวมายสฺมา อมุนา ปญฺจเมน ปุคฺคเลน สมสโม ภวิสฺสตี”ติ.}

\prose{ตตฺร ยฺวายํ ปุคฺคโล น อารภติ วิปฺปฏิสารี โหติ, ตญฺจ เจโตวิมุตฺติํ ปญฺญาวิมุตฺติํ ยถาภูตํ นปฺปชานาติ, ยตฺถสฺส เต อุปฺปนฺนา ปาปกา อกุสลา ธมฺมา อปริเสสา นิรุชฺฌนฺติ, โส เอวมสฺส วจนีโย “อายสฺมโต โข อารมฺภชา อาสวา น สํวิชฺชนฺติ, วิปฺปฏิสารชา อาสวา ปวฑฺฒนฺติ, สาธุ วตายสฺมา วิปฺปฏิสารเช อาสเว ปฏิวิโนเทตฺวา จิตฺตํ ปญฺญญฺจ ภาเวตุ, เอวมายสฺมา อมุนา ปญฺจเมน ปุคฺคเลน สมสโม ภวิสฺสตี”ติ.}

\prose{ตตฺร ยฺวายํ ปุคฺคโล น อารภติ น วิปฺปฏิสารี โหติ, ตญฺจ เจโตวิมุตฺติํ ปญฺญาวิมุตฺติํ ยถาภูตํ นปฺปชานาติ, ยตฺถสฺส เต ปาปกา อกุสลา ธมฺมา อปริเสสา นิรุชฺฌนฺติ, โส เอวมสฺส วจนีโย “อายสฺมโต โข อารมฺภชา อาสวา น สํวิชฺชนฺติ, วิปฺปฏิสารชา อาสวา นปฺปวฑฺฒนฺติ, สาธุ วตายสฺมา จิตฺตํ ปญฺญญฺจ ภาเวตุ, เอวมายสฺมา อมุนา ปญฺจเมน ปุคฺคเลน สมสโม ภวิสฺสตี”ติ. อิเม จตฺตาโร ปุคฺคลา อมุนา ปญฺจเมน ปุคฺคเลน เอวํ โอวทิยมานา เอวํ อนุสาสิยมานา อนุปุพฺเพน อาสวานํ ขยํ ปาปุณนฺติ.}

\proseitem{192}{กถญฺจ ปุคฺคโล ทตฺวา อวชานาติ, อิเธกจฺโจ ปุคฺคโล ยสฺส ปุคฺคลสฺส เทติ จีวร\-ปิณฺฑปาต\-เสนาสน\-คิลานปจฺจย\-เภสชฺช\-ปริกฺขารํ, ตสฺส เอวํ โหติ “อหํ ทมฺมิ, อยํ\csromansharedfootnote{94029671090718a9}{อยํ ปน (สฺยา, ก) อํ 2. 145 ปิฏฺเฐปิ.} ปฏิคฺคณฺหาตี”ติ, ตเมนํ ทตฺวา อวชานาติ. เอวํ ปุคฺคโล ทตฺวา อวชานาติ.}

\prose{กถญฺจ ปุคฺคโล สํวาเสน อวชานาติ, อิเธกจฺโจ ปุคฺคโล เยน ปุคฺคเลน สทฺธิํ สํวสติ เทฺว วา ตีณิ วา วสฺสานิ, ตเมนํ สํวาเสน อวชานาติ. เอวํ ปุคฺคโล สํวาเสน อวชานาติ.}

\prose{\csromanfolio{175}กถญฺจ ปุคฺคโล อาเธยฺยมุโข โหติ, อิเธกจฺโจ ปุคฺคโล ปรสฺส วณฺเณ วา อวณฺเณ วา ภาสิยมาเน ขิปฺปญฺเญว อธิมุจฺจิตา โหติ. เอวํ ปุคฺคโล อาเธยฺยมุโข โหติ.}

\prose{กถญฺจ ปุคฺคโล โลโล โหติ, อิเธกจฺโจ ปุคฺคโล อิตฺตรสทฺโธ โหติ อิตฺตรภตฺตี อิตฺตรเปโม อิตฺตรปฺปสาโท. เอวํ ปุคฺคโล โลโล โหติ.}

\prose{กถญฺจ ปุคฺคโล มนฺโท โมมูโห โหติ, อิเธกจฺโจ ปุคฺคโล กุสลากุสเล ธมฺเม น ชานาติ, สาวชฺชานวชฺเช ธมฺเม น ชานาติ, หีนปฺปณีเต ธมฺเม น ชานาติ, กณฺห\-สุกฺก\-สปฺปฏิภาเค ธมฺเม น ชานาติ. เอวํ ปุคฺคโล มนฺโท โมมูโห โหติ.}

\proseitem{193}{ตตฺถ กตเม ปญฺจ โยธาชีวูปมา ปุคฺคลา. ปญฺจ โยธาชีวา, อิเธกจฺโจ โยธาชีโว รชคฺคญฺเญว ทิสฺวา สํสีทติ วิสีทติ น สนฺถมฺภติ\csromansharedfootnote{eef3faa0963842ca}{สตฺถมฺภติ (สี) อํ 2. 78 ปิฏฺเฐปิ.} น สกฺโกติ สงฺคามํ โอตริตุํ, เอวรูโปปิ อิเธกจฺโจ โยธาชีโว โหติ. อยํ ปฐโม โยธาชีโว สนฺโต สํวิชฺชมาโน โลกสฺมิํ.}

\prose{ปุน จปรํ อิเธกจฺโจ โยธาชีโว สหติ รชคฺคํ, อปิ จ โข ธชคฺคญฺเญว ทิสฺวา สํสีทติ วิสีทติ น สนฺถมฺภติ น สกฺโกติ สงฺคามํ โอตริตุํ, เอวรูโปปิ อิเธกจฺโจ โยธาชีโว โหติ. อยํ ทุติโย โยธาชีโว สนฺโต สํวิชฺชมาโน โลกสฺมิํ.}

\prose{ปุน จปรํ อิเธกจฺโจ โยธาชีโว สหติ รชคฺคํ, สหติ ธชคฺคํ, อปิ จ โข อุสฺสารณญฺเญว\csromansharedfootnote{e65db24f1063ee27}{อุสฺสาทนํเยว (สี), อุสฺสาทนญฺเญว (สฺยา, ก) อํ 2. 78 ปิฏฺเฐปิ.} สุตฺวา สํสีทติ วิสีทติ น สนฺถมฺภติ น สกฺโกติ สงฺคามํ โอตริตุํ, เอวรูโปปิ อิเธกจฺโจ โยธาชีโว โหติ. อยํ ตติโย โยธาชีโว สนฺโต สํวิชฺชมาโน โลกสฺมิํ.}

\prose{ปุน จปรํ อิเธกจฺโจ โยธาชีโว สหติ รชคฺคํ, สหติ ธชคฺคํ, สหติ อุสฺสารณํ, อปิ จ โข สมฺปหาเร หญฺญติ พฺยาปชฺชติ, เอวรูโปปิ อิเธกจฺโจ โยธาชีโว โหติ. อยํ จตุตฺโถ โยธาชีโว สนฺโต สํวิชฺชมาโน โลกสฺมิํ.}

\prose{\csromanfolio{176}ปุน จปรํ อิเธกจฺโจ โยธาชีโว สหติ รชคฺคํ, สหติ ธชคฺคํ, สหติ อุสฺสารณํ, สหติ สมฺปหารํ, โส ตํ สงฺคามํ อภิวิชินิตฺวา วิชิตสงฺคาโม ตเมว สงฺคามสีสํ อชฺฌาวสติ, เอวรูโปปิ อิเธกจฺโจ โยธาชีโว โหติ. อยํ ปญฺจโม โยธาชีโว สนฺโต สํวิชฺชมาโน โลกสฺมิํ. อิเม ปญฺจ โยธาชีวา สนฺโต สํวิชฺชมานา โลกสฺมิํ.}

\proseitem{194}{เอวเมวํ ปญฺจิเม โยธาชีวูปมา ปุคฺคลา สนฺโต สํวิชฺชมานา ภิกฺขูสุ. กตเม ปญฺจ, อิเธกจฺโจ ภิกฺขุ รชคฺคญฺเญว ทิสฺวา สํสีทติ วิสีทติ น สนฺถมฺภติ น สกฺโกติ พฺรหฺมจริยํ สนฺธาเรตุํ\csromansharedfootnote{e0ea0421807f1133}{สนฺตาเนตุํ (สี, สฺยา) อํ 2. 79 ปิฏฺเฐปิ.}, สิกฺขา\-ทุพฺพลฺยํ อาวิกตฺวา\csromansharedfootnote{c46cd67850c4ca0a}{อาวีกตฺวา (สี)} สิกฺขํ ปจฺจกฺขาย หีนายาวตฺตติ. กิมสฺส รชคฺคสฺมิํ, อิธ ภิกฺขุ สุณาติ “อสุกสฺมิํ นาม คาเม วา นิคเม วา อิตฺถี วา กุมารี วา อภิรูปา ทสฺสนียา ปาสาทิกา ปรมาย วณฺณ\-โปกฺขร\-ตาย สมนฺนาคตา”ติ, โส ตํ สุตฺวา สํสีทติ วิสีทติ น สนฺถมฺภติ น สกฺโกติ พฺรหฺมจริยํ สนฺธาเรตุํ, สิกฺขา\-ทุพฺพลฺยํ อาวิกตฺวา สิกฺขํ ปจฺจกฺขาย หีนายาวตฺตติ. อิทมสฺส รชคฺคสฺมิํ.}

\prose{เสยฺยถาปิ โส โยธาชีโว รชคฺคญฺเญว ทิสฺวา สํสีทติ วิสีทติ น สนฺถมฺภติ น สกฺโกติ สงฺคามํ โอตริตุํ, ตถูปโม อยํ ปุคฺคโล. เอวรูโปปิ อิเธกจฺโจ ปุคฺคโล โหติ. อยํ ปฐโม โยธาชีวูปโม ปุคฺคโล สนฺโต สํวิชฺชมาโน ภิกฺขูสุ.}

\proseitem{195}{ปุน จปรํ อิเธกจฺโจ ภิกฺขุ สหติ รชคฺคํ, อปิ จ โข ธชคฺคญฺเญว ทิสฺวา สํสีทติ วิสีทติ น สนฺถมฺภติ น สกฺโกติ พฺรหฺมจริยํ สนฺธาเรตุํ, สิกฺขา\-ทุพฺพลฺยํ อาวิกตฺวา สิกฺขํ ปจฺจกฺขาย หีนายาวตฺตติ. กิมสฺส ธชคฺคสฺมิํ, อิธ ภิกฺขุ น เหว โข สุณาติ “อสุกสฺมิํ นาม คาเม วา นิคเม วา อิตฺถี วา กุมารี วา อภิรูปา ทสฺสนียา ปาสาทิกา ปรมาย วณฺณ\-โปกฺขร\-ตาย สมนฺนาคตา”ติ. อปิ จ โข สามํ\csromansharedfootnote{21ceb31ff6d7b0ee}{สามํเยว (สี)} ปสฺสติ อิตฺถิํ วา กุมาริํ วา อภิรูปํ ทสฺสนียํ ปาสาทิกํ ปรมาย วณฺณ\-โปกฺขร\-ตาย สมนฺนาคตํ, โส ตํ ทิสฺวา สํสีทติ วิสีทติ น สนฺถมฺภติ น สกฺโกติ พฺรหฺมจริยํ สนฺธาเรตุํ, สิกฺขา\-ทุพฺพลฺยํ อาวิกตฺวา สิกฺขํ ปจฺจกฺขาย หีนายาวตฺตติ. อิทมสฺส ธชคฺคสฺมิํ.}

\prose{\csromanfolio{177}เสยฺยถาปิ โส โยธาชีโว สหติ รชคฺคํ, อปิ จ โข ธชคฺคญฺเญว ทิสฺวา สํสีทติ วิสีทติ น สนฺถมฺภติ น สกฺโกติ สงฺคามํ โอตริตุํ, ตถูปโม อยํ ปุคฺคโล. เอวรูโปปิ อิเธกจฺโจ ปุคฺคโล โหติ. อยํ ทุติโย โยธาชีปูปโม ปุคฺคโล สนฺโต สํวิชฺชมาโน ภิกฺขูสุ.}

\proseitem{196}{ปุน จปรํ อิเธกจฺโจ ภิกฺขุ สหติ รชคฺคํ, สหติ ธชคฺคํ, อปิ จ โข อุสฺสารณญฺเญว สุตฺวา สํสีทติ วิสีทติ น สนฺถมฺภติ น สกฺโกติ พฺรหฺมจริยํ สนฺธาเรตุํ, สิกฺขา\-ทุพฺพลฺยํ อาวิกตฺวา สิกฺขํ ปจฺจกฺขาย หีนายาวตฺตติ, กิมสฺส อุสฺสารณาย. อิธ ภิกฺขุํ อรญฺญคตํ วา รุกฺขมูลคตํ วา สุญฺญาคารคตํ วา มาตุคาโม อุปสงฺกมิตฺวา อูหสติ\csromansharedfootnote{f21aa10707ceafc6}{อุหสติ (อฏฺฐกถา) อํ 2. 80 ปิฏฺเฐปิ.} อุลฺลปติ อุชฺชคฺฆติ อุปฺปณฺเฑติ, โส มาตุคาเมน อูหสิยมาโน อุลฺลปิยมาโน อุชฺชคฺฆิยมาโน อุปฺปณฺฑิยมาโน สํสีทติ วิสีทติ น สนฺถมฺภติ น สกฺโกติ พฺรหฺมจริยํ สนฺธาเรตุํ, สิกฺขา\-ทุพฺพลฺยํ อาวิกตฺวา สิกฺขํ ปจฺจกฺขาย หีนายาวตฺตติ. อิทมสฺส อุสฺสารณาย.}

\prose{เสยฺยถาปิ โส โยธาชีโว สหติ รชคฺคํ, สหติ ธชคฺคํ, อปิ จ โข อุสฺสารณญฺเญว สุตฺวา สํสีทติ วิสีทติ น สนฺถมฺภติ น สกฺโกติ สงฺคามํ โอตริตุํ, ตถูปโม อยํ ปุคฺคโล. เอวรูโปปิ อิเธกจฺโจ ปุคฺคโล โหติ. อยํ ตติโย โยธาชีวูปโม ปุคฺคโล สนฺโต สํวิชฺชมาโน ภิกฺขูสุ.}

\proseitem{197}{ปุน จปรํ อิเธกจฺโจ ภิกฺขุ สหติ รชคฺคํ, สหติ ธชคฺคํ, สหติ อุสฺสารณํ, อปิ จ โข สมฺปหาเร หญฺญติ พฺยาปชฺชติ. กิมสฺส สมฺปหารสฺมิํ, อิธ ภิกฺขุํ อรญฺญคตํ วา รุกฺขมูลคตํ วา สุญฺญาคารคตํ วา มาตุคาโม อุปสงฺกมิตฺวา อภินิสีทติ อภินิปชฺชติ อชฺโฌตฺถรติ. โส มาตุคาเมน อภินิ\-สีทิย\-มาโน, อภินิ\-ปชฺชิย\-มาโน อชฺโฌตฺถริยมาโน สิกฺขํ อปฺปจฺจกฺขาย ทุพฺพลฺยํ อนาวิกตฺวา เมถุนํ ธมฺมํ ปฏิเสวติ. อิทมสฺส สมฺปหารสฺมิํ.}

\prose{เสยฺยถาปิ โส โยธาชีโว สหติ รชคฺคํ, สหติ ธชคฺคํ, สหติ อุสฺสารณํ, อปิ จ โข สมฺปหาเร หญฺญติ พฺยาปชฺชติ, \csromanfolio{178}ตถูปโม อยํ ปุคฺคโล. เอวรูโปปิ อิเธกจฺโจ ปุคฺคโล โหติ. อยํ จตุตฺโถ โยธาชีวูปโม ปุคฺคโล สนฺโต สํวิชฺชมาโน ภิกฺขูสุ.}

\proseitem{198}{ปุน จปรํ อิเธกจฺโจ ภิกฺขุ สหติ รชคฺคํ, สหติ ธชคฺคํ, สหติ อุสฺสารณํ, สหติ สมฺปหารํ. โส ตํ สงฺคามํ อภิวิชินิตฺวา วิชิตสงฺคาโม ตเมว สงฺคามสีสํ อชฺฌาวสติ. กิมสฺส สงฺคาม\-วิชยสฺมิํ, อิธ ภิกฺขุํ อรญฺญคตํ วา รุกฺขมูลคตํ วา สุญฺญาคารคตํ วา มาตุคาโม อุปสงฺกมิตฺวา อภินิสีทติ อภินิปชฺชติ อชฺโฌตฺถรติ. โส มาตุคาเมน อภินิ\-สีทิย\-มาโน อภินิ\-ปชฺชิย\-มาโน อชฺโฌตฺถริยมาโน วินิเวเฐตฺวา วินิโมเจตฺวา เยน กามํ ปกฺกมติ.}

\prose{โส วิวิตฺตํ เสนาสนํ ภชติ อรญฺญํ รุกฺขมูลํ ปพฺพตํ กนฺทรํ คิริคุหํ สุสานํ วนปตฺถํ อพฺโภกาสํ ปลาลปุญฺชํ. โส อรญฺญคโต วา รุกฺขมูลคโต วา สุญฺญาคารคโต วา นิสีทติ ปลฺลงฺกํ อาภุชิตฺวา อุชุํ กายํ ปณิธาย ปริมุขํ สติํ อุปฏฺฐเปตฺวา, โส อภิชฺฌํ โลเก ปหาย วิคตา\-ภิชฺเฌน เจตสา วิหรติ, อภิชฺฌาย จิตฺตํ ปริโสเธติ. พฺยาปาทปโทสํ ปหาย อพฺยาปนฺนจิตฺโต วิหรติ สพฺพปาณภูต\-หิตา\-นุกมฺปี, พฺยาปาทปโทสา จิตฺตํ ปริโสเธติ. ถินมิทฺธํ ปหาย วิคต\-ถินมิทฺโธ วิหรติ อาโลกสญฺญี สโต สมฺปชาโน, ถินมิทฺธา จิตฺตํ ปริโสเธติ. อุทฺธจฺจ\-กุกฺกุจฺจํ ปหาย อนุทฺธโต วิหรติ อชฺฌตฺตํ วูปสนฺตจิตฺโต, อุทฺธจฺจ\-กุกฺกุจฺจา จิตฺตํ ปริโสเธติ. วิจิกิจฺฉํ ปหาย ติณฺณ\-วิจิกิจฺโฉ วิหรติ อกถํกถี กุสเลสุ ธมฺเมสุ, วิจิกิจฺฉาย จิตฺตํ ปริโสเธติ.}

\prose{โส อิเม ปญฺจ นีวรเณ ปหาย เจตโส อุปกฺกิเลเส ปญฺญาย ทุพฺพลีกรเณ วิวิจฺเจว กาเมหิ วิวิจฺจ อกุสเลหิ ธมฺเมหิ สวิตกฺกํ สวิจารํ วิเวกชํ ปีติสุขํ ปฐมํ ฌานํ อุปสมฺปชฺช วิหรติ, วิตกฺก\-วิจารานํ วูปสมา ทุติยํ ฌานํ ฯเปฯ ตติยํ ฌานํ ฯเปฯ จตุตฺถํ ฌานํ อุปสมฺปชฺช วิหรติ.}

\prose{โส เอวํ สมาหิเต จิตฺเต ปริสุทฺเธ ปริโยทาเต อนงฺคเณ วิคตู\-ปกฺกิเลเส มุทุภูเต กมฺมนิเย ฐิเต อาเนญฺชปฺปตฺเต อาสวานํ ขยญาณาย จิตฺตํ อภินินฺนาเมติ, โส “อิทํ ทุกฺขนฺ”ติ ยถาภูตํ ปชานาติ, “อยํ ทุกฺขสมุทโย”ติ ยถาภูตํ ปชานาติ, “อยํ ทุกฺขนิโรโธ”ติ ยถาภูตํ ปชานาติ “อยํ ทุกฺขนิโรธ\-คามินี ปฏิปทา”ติ ยถาภูตํ ปชานาติ, “อิเม อาสวา”ติ ยถาภูตํ ปชานาติ, “อยํ}

\prose{\csromanfolio{179}อาสวสมุทโย”ติ ยถาภูตํ ปชานาติ, “อยํ อาสวนิโรโธ”ติ ยถาภูตํ ปชานาติ, “อยํ อาสว\-นิโรธ\-คามินี ปฏิปทา”ติ, ยถาภูตํ ปชานาติ.}

\prose{ตสฺส เอวํ ชานโต เอวํ ปสฺสโต กามาสวาปิ จิตฺตํ วิมุจฺจติ, ภวาสวาปิ จิตฺตํ วิมุจฺจติ, อวิชฺชาสวาปิ จิตฺตํ วิมุจฺจติ, วิมุตฺตสฺมิํ “วิมุตฺตมฺ”อิติ ญาณํ โหติ, “ขีณา ชาติ, วุสิตํ พฺรหฺมจริยํ, กตํ กรณียํ, นาปรํ อิตฺถตฺตายา”ติ ปชานาติ. อิทมสฺส สงฺคาม\-วิชยสฺมิํ. เสยฺยถาปิ โส โยธาชีโว สหติ รชคฺคํ, สหติ ธชคฺคํ, สหติ อุสฺสารณํ, สหติ สมฺปหารํ. โส ตํ สงฺคามํ อภิวิชินิตฺวา วิชิตสงฺคาโม ตเมว สงฺคามสีสํ อชฺฌาวสติ, ตถูปโม อยํ ปุคฺคโล. เอวรูโปปิ อิเธกจฺโจ ปุคฺคโล โหติ. อยํ ปญฺจโม โยธาชีวูปโม ปุคฺคโล สนฺโต สํวิชฺชมาโน ภิกฺขูสุ. อิเม ปญฺจ โยธาชีวูปมา ปุคฺคลา สนฺโต สํวิชฺชมานา ภิกฺขูสุ.}

\proseitem{199}{ตตฺถ กตเม ปญฺจ ปิณฺฑปาติกา, มนฺทตฺตา โมมูหตฺตา ปิณฺฑปาติโก โหติ, ปาปิจฺโฉ อิจฺฉาปกโต ปิณฺฑปาติโก โหติ, อุมฺมาทา จิตฺตวิกฺเขปา ปิณฺฑปาติโก โหติ, วณฺณิตํ พุทฺเธหิ พุทฺธสาวเกหีติ ปิณฺฑปาติโก โหติ, อปิ จ อปฺปิ\-จฺฉ\-ตํเยว\csromansharedfootnote{44d35dd4760b0b07}{อปฺปิจฺฉํเยว (สฺยา) อํ 2. 193 ปิฏฺเฐปิ.} นิสฺสาย สนฺตุฏฺฐิํ เยว นิสฺสาย สลฺเลขํเยว นิสฺสาย อิทมตฺถิตํเยว\csromansharedfootnote{fbcfb0ecefa26ca7}{อิทมฏฺฐิตํเยว (สี)} นิสฺสาย ปิณฺฑปาติโก โหติ. ตตฺร ยฺวายํ ปิณฺฑปาติโก อปฺปิ\-จฺฉ\-ตํเยว นิสฺสาย สนฺตุฏฺฐิํเยว นิสฺสาย สลฺเลขํเยว นิสฺสาย อิทมตฺถิตํเยว นิสฺสาย ปิณฺฑปาติโก, อยํ อิเมสํ ปญฺจนฺนํ ปิณฺฑปาติกานํ อคฺโค จ เสฏฺโฐ จ ปาโมกฺโข\csromansharedfootnote{60b98e8f34384703}{โมกฺโข (สี)} จ อุตฺตโม จ ปวโร จ.}

\prose{เสยฺยถาปิ นาม ควา ขีรํ, ขีรมฺหา ทธิ, ทธิมฺหา นวนีตํ\csromansharedfootnote{3d672f0f7dd5684e}{โนนีตํ (สี)}, นวนีตมฺหา สปฺปิ, สปฺปิมฺหา สปฺปิมณฺโฑ, สปฺปิมณฺฑํ ตตฺถ อคฺคมกฺขายติ. เอวเมวํ ยฺวายํ ปิณฺฑปาติโก อปฺปิ\-จฺฉ\-ตํเยว นิสฺสาย สนฺตุฏฺฐิํเยว นิสฺสาย สลฺเลขํเยว นิสฺสาย อิทมตฺถิตํเยว นิสฺสาย ปิณฺฑปาติโก โหติ, อยํ อิเมสํ ปญฺจนฺนํ ปิณฺฑปาติกานํ อคฺโค จ เสฏฺโฐ จ ปาโมกฺโข จ อุตฺตโม จ ปวโร จ. อิเม ปญฺจ ปิณฺฑปาติกา.}

\proseitem{200}{\csromanfolio{180}ตตฺถ กตเม ปญฺจ ขลุปจฺฉาภตฺติกา ฯเปฯ. ปญฺจ เอกาสนิกา ฯเปฯ. ปญฺจ ปํสุกูลิกา ฯเปฯ. ปญฺจ เตจีวริกา ฯเปฯ. ปญฺจ อารญฺญิกา ฯเปฯ. ปญฺจ รุกฺขมูลิกา ฯเปฯ. ปญฺจ อพฺโภกาสิกา ฯเปฯ. ปญฺจ เนสชฺชิกา ฯเปฯ. ปญฺจ ยถา\-สนฺถติกา ฯเปฯ.}

\proseitem{201}{ตตฺถ กตเม ปญฺจ โสสานิกา, มนฺทตฺตา โมมูหตฺตา โสสานิโก โหติ, ปาปิจฺโฉ อิจฺฉาปกโต โสสานิโก โหติ, อุมฺมาทา จิตฺตวิกฺเขปา โสสานิโก โหติ, วณฺณิตํ พุทฺเธหิ พุทฺธสาวเกหีติ โสสานิโก โหติ, อปิ จ อปฺปิ\-จฺฉ\-ตํเยว นิสฺสาย สนฺตุฏฺฐิํเยว นิสฺสาย สลฺเลขํเยว นิสฺสาย อิทมตฺถิตํเยว นิสฺสาย โสสานิโก โหติ. ตตฺร ยฺวายํ โสสานิโก อปฺปิ\-จฺฉ\-ตํเยว นิสฺสาย สนฺตุฏฺฐิํเยว นิสฺสาย สลฺเลขํเยว นิสฺสาย อิทมตฺถิตํเยว นิสฺสาย โสสานิโก, อยํ อิเมสํ ปญฺจนฺนํ โสสานิกานํ อคฺโค จ เสฏฺโฐ จ ปาโมกฺโข จ อุตฺตโม จ ปวโร จ.}

\prose{เสยฺยถาปิ นาม ควา ขีรํ, ขีรมฺหา ทธิ, ทธิมฺหา นวนีตํ, นวนีตมฺหา สปฺปิ, สปฺปิมฺหา สปฺปิมณฺโฑ, สปฺปิมณฺฑํ ตตฺถ อคฺคมกฺขายติ. เอวเมวํ ยฺวายํ โสสานิโก อปฺปิ\-จฺฉ\-ตํเยว นิสฺสาย สนฺตุฏฺฐิํเยว นิสฺสาย สลฺเลขํเยว นิสฺสาย อิทมตฺถิตํเยว นิสฺสาย โสสานิโก โหติ, อยํ อิเมสํ ปญฺจนฺนํ โสสานิกานํ อคฺโค จ เสฏฺโฐ จ ปาโมกฺโข จ อุตฺตโม จ ปวโร จ. อิเม ปญฺจ โสสานิกา.}

\vspace{\tipitakaparskip}
\csromancenter{ปญฺจกนิทฺเทโส.}
\csromanmatikaanchor{mtk.75}
\csromantocmarkref{subsection}{6. ฉกฺกปุคฺคลปญฺญตฺติ}{mtk.75}
\bookmark[dest={mtk.75},level=2]{6. ฉกฺกปุคฺคลปญฺญตฺติ}
\csromanheader{2}{6. ฉกฺก\-ปุคฺคล\-ปญฺญตฺติ}
\proseitem{202}{ตตฺร ยฺวายํ ปุคฺคโล ปุพฺเพ อนนุสฺสุเตสุ ธมฺเมสุ สามํ สจฺจานิ อภิสมฺ\-พุชฺฌติ, ตตฺถ จ สพฺพญฺญุตํ ปาปุณาติ พเลสุ จ วสีภาวํ, สมฺมาสมฺพุทฺโธ เตน ทฏฺฐพฺโพ.}

\prose{ตตฺร ยฺวายํ ปุคฺคโล ปุพฺเพ อนนุสฺสุเตสุ ธมฺเมสุ สามํ สจฺจานิ อภิสมฺ\-พุชฺฌติ, น จ ตตฺถ สพฺพญฺญุตํ ปาปุณาติ น จ พเลสุ วสีภาวํ ปจฺเจก\-สมฺพุทฺโธ เตน ทฏฺฐพฺโพ.}

\prose{\csromanfolio{181}ตตฺร ยฺวายํ ปุคฺคโล ปุพฺเพ อนนุสฺสุเตสุ ธมฺเมสุ สามํ สจฺจานิ อนภิสมฺพุชฺฌติ, ทิฏฺเฐว ธมฺเม ทุกฺขสฺส\-นฺตกโร โหติ\csromansharedfootnote{9c7dba504037352d}{ทุกฺขสฺสนฺตํ กโรติ (สี) เอวมุปริปิ.} สาวกปารมิญฺจ ปาปุณาติ, สาริปุตฺต\-โมคฺคลฺลานา เตน ทฏฺฐพฺพา.}

\prose{ตตฺร ยฺวายํ ปุคฺคโล ปุพฺเพ อนนุสฺสุเตสุ ธมฺเมสุ สามํ สจฺจานิ อนภิสมฺพุชฺฌติ, ทิฏฺเฐว ธมฺเม ทุกฺขสฺส\-นฺตกโร โหติ, น จ สาวกปารมิํ ปาปุณาติ, อวเสสา อรหนฺตา เตน ทฏฺฐพฺพา.}

\prose{ตตฺร ยฺวายํ ปุคฺคโล ปุพฺเพ อนนุสฺสุเตสุ ธมฺเมสุ สามํ สจฺจานิ อนภิสมฺพุชฺฌติ, น จ ทิฏฺเฐว ธมฺเม ทุกฺขสฺส\-นฺตกโร โหติ, อนาคามี โหติ อนาคนฺตา อิตฺถตฺตํ, อนาคามี เตน ทฏฺฐพฺโพ.}

\prose{ตตฺร ยฺวายํ ปุคฺคโล ปุพฺเพ อนนุสฺสุเตสุ ธมฺเมสุ สามํ สจฺจานิ อนภิสมฺพุชฺฌติ, น จ ทิฏฺเฐว ธมฺเม ทุกฺขสฺส\-นฺตกโร โหติ, อาคนฺตา อิตฺถตฺตํ, โสตาปนฺน\-สกทาคามิโน เตน ทฏฺฐพฺพา.}

\vspace{\tipitakaparskip}
\csromancenter{ฉกฺกนิทฺเทโส.}
\csromanmatikaanchor{mtk.76}
\csromantocmarkref{subsection}{7. สตฺตกปุคฺคลปญฺญตฺติ}{mtk.76}
\bookmark[dest={mtk.76},level=2]{7. สตฺตกปุคฺคลปญฺญตฺติ}
\csromanheader{2}{7. สตฺตก\-ปุคฺคล\-ปญฺญตฺติ}
\proseitem{203}{กถญฺจ ปุคฺคโล สกิํ นิมุคฺโค นิมุคฺโคว โหติ, อิเธกจฺโจ ปุคฺคโล สมนฺนาคโต โหติ เอกนฺตกาฬเกหิ อกุสเลหิ ธมฺเมหิ. เอวํ ปุคฺคโล สกิํ นิมุคฺโค นิมุคฺโคว โหติ.}

\prose{กถญฺจ ปุคฺคโล อุมฺมุชฺชิตฺวา นิมุชฺชติ, อิเธกจฺโจ ปุคฺคโล อุมฺมุชฺชติ “สาหุ สทฺธา กุสเลสุ ธมฺเมสุ, สาธุ\csromansharedfootnote{dd12a048b05902ce}{สาหุ (สี, สฺยา) เอวํ ตีสุ ฏฺฐาเนสุปิ.} หิรี กุสเลสุ ธมฺเมสุ, สาธุ โอตฺตปฺปํ กุสเลสุ ธมฺเมสุ, สาธุ วีริยํ\csromansharedfootnote{29fffb8d80c6a398}{วิริยํ (สี, สฺยา)} กุสเลสุ ธมฺเมสุ, สาธุ ปญฺญา กุสเลสุ ธมฺเมสู”ติ. ตสฺส สา สทฺธา เนว ติฏฺฐติ โน วฑฺฒติ หายติเยว, ตสฺส สา หิรี เนว ติฏฺฐติ โน วฑฺฒติ หายติเยว, ตสฺส ตํ โอตฺตปฺปํ เนว ติฏฺฐติ โน วฑฺฒติ หายติเยว, ตสฺส ตํ วีริยํ เนว ติฏฺฐติ โน วฑฺฒติ หายติเยว, ตสฺส สา ปญฺญา เนว ติฏฺฐติ โน วฑฺฒติ หายติเยว. เอวํ ปุคฺคโล อุมฺมุชฺชิตฺวา นิมุชฺชติ.}

\prose{\csromanfolio{182}กถญฺจ ปุคฺคโล อุมฺมุชฺชิตฺวา ฐิโต โหติ, อิเธกจฺโจ ปุคฺคโล อุมฺมุชฺชติ “สาหุ สทฺธา กุสเลสุ ธมฺเมสุ, สาธุ หิรี กุสเลสุ ธมฺเมสุ, สาธุ โอตฺตปฺปํ กุสเลสุ ธมฺเมสุ, สาธุ วีริยํ กุสเลสุ ธมฺเมสุ, สาธุ ปญฺญา กุสเลสุ ธมฺเมสู”ติ. ตสฺส สา สทฺธา เนว หายติ โน วฑฺฒติ ฐิตา โหติ, ตสฺส สา หิรี เนว หายติ โน วฑฺฒติ ฐิตา โหติ, ตสฺส ตํ โอตฺตปฺปํ เนว หายติ โน วฑฺฒติ ฐิตํ โหติ, ตสฺส ตํ วีริยํ เนว หายติ โน วฑฺฒติ ฐิตํ โหติ, ตสฺส สา ปญฺญา เนว หายติ โน วฑฺฒติ ฐิตา โหติ, เอวํ ปุคฺคโล อุมฺมุชฺชิตฺวา ฐิโต โหติ.}

\prose{กถญฺจ ปุคฺคโล อุมฺมุชฺชิตฺวา วิปสฺสติ วิโลเกติ, อิเธกจฺโจ ปุคฺคโล อุมฺมุชฺชติ “สาหุ สทฺธา กุสเลสุ ธมฺเมสุ, สาธุ หิรี กุสเลสุ ธมฺเมสุ, สาธุ โอตฺตปฺปํ กุสเลสุ ธมฺเมสุ, สาธุ วีริยํ กุสเลสุ ธมฺเมสุ, สาธุ ปญฺญา กุสเลสุ ธมฺเมสู”ติ. โส ติณฺณํ สํโยชนานํ ปริกฺขยา โสตาปนฺโน โหติ อวินิปาต\-ธมฺโม นิยโต สมฺโพธิ\-ปรายโน. เอวํ ปุคฺคโล อุมฺมุชฺชิตฺวา วิปสฺสติ วิโลเกติ.}

\prose{กถญฺจ ปุคฺคโล อุมฺมุชฺชิตฺวา ปตรติ, อิเธกจฺโจ ปุคฺคโล อุมฺมุชฺชติ “สาหุ สทฺธา กุสเลสุ ธมฺเมสุ, สาธุ หิรี กุสเลสุ ธมฺเมสุ, สาธุ โอตฺตปฺปํ กุสเลสุ ธมฺเมสุ, สาธุ วีริยํ กุสเลสุ ธมฺเมสุ, สาธุ ปญฺญา กุสเลสุ ธมฺเมสู”ติ. โส ติณฺณํ สํโยชนานํ ปริกฺขยา ราค\-โทส\-โมหานํ ตนุตฺตา สกทาคามี โหติ สกิเทว อิมํ โลกํ อาคนฺตฺวา ทุกฺขสฺส\-นฺตกโร โหติ. เอวํ ปุคฺคโล อุมฺมุชฺชิตฺวา ปตรติ.}

\prose{กถญฺจ ปุคฺคโล อุมฺมุชฺชิตฺวา ปติคาธ\-ปฺปตฺโต โหติ, อิเธกจฺโจ ปุคฺคโล อุมฺมุชฺชติ “สาหุ สทฺธา กุสเลสุ ธมฺเมสุ, สาธุ หิรี กุสเลสุ ธมฺเมสุ, สาธุ โอตฺตปฺปํ กุสเลสุ ธมฺเมสุ, สาธุ วีริยํ กุสเลสุ ธมฺเมสุ, สาธุ ปญฺญา กุสเลสุ ธมฺเมสู”ติ. โส ปญฺจนฺนํ โอรมฺภาคิยานํ สํโยชนานํ ปริกฺขยา โอปปาติโก โหติ ตตฺถ ปรินิพฺพายี อนาวตฺติธมฺโม ตสฺมา โลกา. เอวํ ปุคฺคโล อุมฺมุชฺชิตฺวา ปติคาธ\-ปฺปตฺโต โหติ.}

\prose{กถญฺจ ปุคฺคโล อุมฺมุชฺชิตฺวา ติณฺโณ โหติ ปารงฺคโต ถเล ติฏฺฐติ พฺราหฺมโณ, อิเธกจฺโจ ปุคฺคโล อุมฺมุชฺชติ “สาหุ สทฺธา กุสเลสุ}

\prose{\csromanfolio{183}ธมฺเมสุ, สาธุ หิรี กุสเลสุ ธมฺเมสุ, สาธุ โอตฺตปฺปํ กุสเลสุ ธมฺเมสุ, สาธุ วีริยํ กุสเลสุ ธมฺเมสุ, สาธุ ปญฺญา กุสเลสุ ธมฺเมสู”ติ. โส อาสวานํ ขยา อนาสวํ เจโตวิมุตฺติํ ปญฺญาวิมุตฺติํ ทิฏฺเฐว ธมฺเม สยํ อภิญฺญา สจฺฉิกตฺวา อุปสมฺปชฺช วิหรติ. เอวํ ปุคฺคโล อุมฺมุชฺชิตฺวา ติณฺโณ โหติ ปารงฺคโต ถเล ติฏฺฐติ พฺราหฺมโณ.}

\proseitem{204}{กตโม จ ปุคฺคโล อุภโตภาค\-วิมุตฺโต, อิเธกจฺโจ ปุคฺคโล อฏฺฐ วิโมกฺเข กาเยน ผุสิตฺวา วิหรติ, ปญฺญาย จสฺส ทิสฺวา อาสวา ปริกฺขีณา โหนฺติ. อยํ วุจฺจติ ปุคฺคโล อุภโตภาค\-วิมุตฺโต.}

\proseitem{205}{กตโม จ ปุคฺคโล ปญฺญาวิมุตฺโต ฯเปฯ. กายสกฺขี. ทิฏฺฐิปฺปตฺโต. สทฺธาวิมุตฺโต. ธมฺมานุสารี.}

\proseitem{206}{กตโม จ ปุคฺคโล สทฺธานุสารี, ยสฺส ปุคฺคลสฺส โสตาปตฺติผล\-สจฺฉิกิริยาย ปฏิปนฺนสฺส สทฺธินฺทฺริยํ อธิมตฺตํ โหติ, สทฺธาวาหิํ สทฺธา\-ปุพฺพงฺคมํ อริยมคฺคํ ภาเวติ. อยํ วุจฺจติ ปุคฺคโล สทฺธานุสารี, โสตาปตฺติผล\-สจฺฉิกิริยาย ปฏิปนฺโน ปุคฺคโล สทฺธานุสารี, ผเล ฐิโต สทฺธา\-วิมุตฺโตติ.}

\vspace{\tipitakaparskip}
\csromancenter{สตฺตกนิทฺเทโส.}
\csromanmatikaanchor{mtk.77}
\csromantocmarkref{subsection}{8. อฏฺฐกปุคฺคลปญฺญตฺติ}{mtk.77}
\bookmark[dest={mtk.77},level=2]{8. อฏฺฐกปุคฺคลปญฺญตฺติ}
\csromanheader{2}{8. อฏฺฐก\-ปุคฺคล\-ปญฺญตฺติ}
\proseitem{207}{ตตฺถ กตเม จตฺตาโร มคฺคสมงฺคิโน จตฺตาโร ผลสมงฺคิโน ปุคฺคลา. โสตาปนฺโน, โสตาปตฺติผล\-สจฺฉิกิริยาย ปฏิปนฺโน. สกทาคามี, สกทาคามิ\-ผล\-สจฺฉิกิริยาย ปฏิปนฺโน. อนาคามี, อนาคามิ\-ผล\-สจฺฉิกิริยาย ปฏิปนฺโน. อรหา, อรหตฺตผล\-สจฺฉิกิริยาย\csromansharedfootnote{b7198701db8c741d}{อรหตฺตาย (สฺยา, ก) อํ 3. 115 ปิฏฺเฐปิ.} ปฏิปนฺโน. อิเม จตฺตาโร มคฺคสมงฺคิโน, อิเม จตฺตาโร ผลสมงฺคิโน ปุคฺคลา.}

\vspace{\tipitakaparskip}
\csromancenter{อฏฺฐกนิทฺเทโส.}
\csromanmatikaanchor{mtk.78}
\csromantocmarkref{subsection}{9. นวกปุคฺคลปญฺญตฺติ}{mtk.78}
\bookmark[dest={mtk.78},level=2]{9. นวกปุคฺคลปญฺญตฺติ}
\csromanmarkh{ปุคฺคล\-ปญฺญตฺติปาฬิ}\csromanmarkhii{9. นวก\-ปุคฺคล\-ปญฺญตฺติ}\csromanheadernomark{2}{ปุคฺคล\-ปญฺญตฺติปาฬิ 9. นวก\-ปุคฺคล\-ปญฺญตฺติ}
\proseitem{208}{\csromanfolio{184}กตโม จ ปุคฺคโล สมฺมาสมฺพุทฺโธ, อิเธกจฺโจ ปุคฺคโล ปุพฺเพ อนนุสฺสุเตสุ ธมฺเมสุ สามํ สจฺจานิ อภิสมฺ\-พุชฺฌติ, ตตฺถ จ สพฺพญฺญุตํ ปาปุณาติ พเลสุ จ วสีภาวํ. อยํ วุจฺจติ ปุคฺคโล สมฺมาสมฺพุทฺโธ.}

\prose{กตโม จ ปุคฺคโล ปจฺเจก\-สมฺพุทฺโธ, อิเธกจฺโจ ปุคฺคโล ปุพฺเพ อนนุสฺสุเตสุ ธมฺเมสุ สามํ สจฺจานิ อภิสมฺ\-พุชฺฌติ, น จ ตตฺถ สพฺพญฺญุตํ ปาปุณาติ น จ พเลสุ วสีภาวํ. อยํ วุจฺจติ ปุคฺคโล ปจฺเจก\-สมฺพุทฺโธ.}

\prose{กตโม จ ปุคฺคโล อุภโตภาค\-วิมุตฺโต, อิเธกจฺโจ ปุคฺคโล อฏฺฐ วิโมกฺเข กาเยน ผุสิตฺวา วิหรติ, ปญฺญาย จสฺส ทิสฺวา อาสวา ปริกฺขีณา โหนฺติ. อยํ วุจฺจติ ปุคฺคโล อุภโตภาค\-วิมุตฺโต.}

\prose{กตโม จ ปุคฺคโล ปญฺญาวิมุตฺโต, อิเธกจฺโจ ปุคฺคโล น เหว โข อฏฺฐ วิโมกฺเข กาเยน ผุสิตฺวา วิหรติ, ปญฺญาย จสฺส ทิสฺวา อาสวา ปริกฺขีณา โหนฺติ. อยํ วุจฺจติ ปุคฺคโล ปญฺญาวิมุตฺโต.}

\prose{กตโม จ ปุคฺคโล กายสกฺขี, อิเธกจฺโจ ปุคฺคโล อฏฺฐ วิโมกฺเข กาเยน ผุสิตฺวา วิหรติ, ปญฺญาย จสฺส ทิสฺวา เอกจฺเจ อาสวา ปริกฺขีณา โหนฺติ. อยํ วุจฺจติ ปุคฺคโล กายสกฺขี.}

\prose{กตโม จ ปุคฺคโล ทิฏฺฐิปฺปตฺโต, อิเธกจฺโจ ปุคฺคโล “อิทํ ทุกฺขนฺ”ติ ยถาภูตํ ปชานาติ ฯเปฯ “อยํ ทุกฺขนิโรธ\-คามินี ปฏิปทา”ติ ยถาภูตํ ปชานาติ, ตถาคต\-ปฺปเวทิตา จสฺส ธมฺมา ปญฺญาย โวทิฏฺฐา โหนฺติ โวจริตา, ปญฺญาย จสฺส ทิสฺวา เอกจฺเจ อาสวา ปริกฺขีณา โหนฺติ. อยํ วุจฺจติ ปุคฺคโล ทิฏฺฐิปฺปตฺโต.}

\prose{กตโม จ ปุคฺคโล สทฺธาวิมุตฺโต, อิเธกจฺโจ ปุคฺคโล “อิทํ ทุกฺขนฺ”ติ ยถาภูตํ ปชานาติ ฯเปฯ “อยํ ทุกฺขนิโรธ\-คามินี ปฏิปทา”ติ ยถาภูตํ ปชานาติ, ตถาคต\-ปฺปเวทิตา จสฺส ธมฺมา ปญฺญาย}

\prose{\csromanfolio{185}โวทิฏฺฐา โหนฺติ โวจริตา, ปญฺญาย จสฺส ทิสฺวา เอกจฺเจ อาสวา ปริกฺขีณา โหนฺติ, โน จ โข ยถา ทิฏฺฐิ\-ปฺปตฺตสฺส. อยํ วุจฺจติ ปุคฺคโล สทฺธาวิมุตฺโต.}

\prose{กตโม จ ปุคฺคโล ธมฺมานุสารี, ยสฺส ปุคฺคลสฺส โสตาปตฺติผล\-สจฺฉิกิริยาย ปฏิปนฺนสฺส ปญฺญินฺทฺริยํ อธิมตฺตํ โหติ, ปญฺญาวาหิํ ปญฺญา\-ปุพฺพงฺคมํ อริยมคฺคํ ภาเวติ. อยํ วุจฺจติ ปุคฺคโล ธมฺมานุสารี. โสตาปตฺติผล\-สจฺฉิกิริยาย ปฏิปนฺโน ปุคฺคโล ธมฺมานุสารี. ผเล ฐิโต ทิฏฺฐิปฺปตฺโต.}

\prose{กตโม จ ปุคฺคโล สทฺธานุสารี, ยสฺส ปุคฺคลสฺส โสตาปตฺติผล\-สจฺฉิกิริยาย ปฏิปนฺนสฺส สทฺธินฺทฺริยํ อธิมตฺตํ โหติ, สทฺธาวาหิํ สทฺธา\-ปุพฺพงฺคมํ อริยมคฺคํ ภาเวติ. อยํ วุจฺจติ ปุคฺคโล สทฺธานุสารี. โสตาปตฺติผล\-สจฺฉิกิริยาย ปฏิปนฺโน ปุคฺคโล สทฺธานุสารี. ผเล ฐิโต สทฺธา\-วิมุตฺโตติ.}

\vspace{\tipitakaparskip}
\csromancenter{นวกนิทฺเทโส.}
\csromanmatikaanchor{mtk.79}
\csromantocmarkref{subsection}{10. ทสกปุคฺคลปญฺญตฺติ}{mtk.79}
\bookmark[dest={mtk.79},level=2]{10. ทสกปุคฺคลปญฺญตฺติ}
\csromanheader{2}{10. ทสก\-ปุคฺคล\-ปญฺญตฺติ}
\proseitem{209}{กตเมสํ ปญฺจนฺนํ อิธ นิฏฺฐา, สตฺตกฺขตฺตุ\-ปรมสฺส, โกลํโกลสฺส, เอกพีชิสฺส, สกทาคามิสฺส, โย จ ทิฏฺเฐว ธมฺเม อรหา, อิเมสํ ปญฺจนฺนํ อิธ นิฏฺฐา.}

\prose{กตเมสํ ปญฺจนฺนํ อิธ วิหาย นิฏฺฐา, อนฺตรา\-ปรินิพฺพายิสฺส, อุปหจฺจ ปรินิพฺพายิสฺส, อสงฺขารปรินิพฺพายิสฺส, สสงฺขาร\-ปรินิพฺพายิสฺส, อุทฺธํโสตสฺส, อกนิฏฺฐ\-คามิโน, อิเมสํ ปญฺจนฺนํ อิธ วิหาย นิฏฺฐาติ. เอตฺตาวตา ปุคฺคลานํ ปุคฺคลปญฺญตฺตีติ.}

\vspace{\tipitakaparskip}
\csromancenter{ทสกนิทฺเทโส.}
\nitthitam{ปุคฺคลปญฺญตฺติ\-ปกรณํ นิฏฺฐิตํ.}
